\section{Theta lifts and a common period lattice}
\label{sec:periods}

Fix one translate of the family chosen in Section~\ref{sec:assembly},
with CM extension \(E/F\). A word \(w\in W\) is a sequence of characters
with associated column types and levels. For each word we compare its
top theta periods with Bessel periods at the bottom of the towers.
All words at this translate must use the same bottom row coordinates
to compute the dimension of the sum of the period images.

\subsection{The towers and the top pairing}

\teamanchor{proof.normalized-towers}
Take even \(k,m\) with \(k\geq m+3\) and \(m\geq2\), and odd \(t\geq r+5\),
where \(r=k+m-1\), as in the
\releasecite{Formalism/LocalRingCardinality/Params.lean}{6}{selected rank data}.
The spaces \(X_j\)
are definite of rank \(k+j(t+1)\). The spaces \(Y_j\) have rank \(r+jt\),
and negative index \(\lfloor(m-j)/2\rfloor\)
at the real places. For \(j<m\), the construction embeds \(X_j\) in
\(Y_j\). At the top \(Y_m\) is isometric to a hyperplane
\(X'\subset X_m\) perpendicular to a prescribed anisotropic line.
At a real place the entry \(-e_i\) of the character used at stage
\(i<j\) contributes the string \(-e_i+(t-1)/2-a\), \(0\leq a<t\), to
the row Harish--Chandra parameter at \(j\). The corresponding column
string is \(-e_i+t/2-a\), \(0\leq a\leq t\). The row parameter also
contains the signed entries from stage zero, with the labels prescribed
at stage \(j\). A common shift by the half-integral splitting parameter
is added to the row entries when \(j\) is odd. These are the
\releasecite{Formalism/LocalRingCardinality/DataTheta.lean}{535}{parameters of the two towers}.

Write \(H_w^{\rm full}\) for the full space at the bottom column level.
For each designated residue-unit group, averaging its action gives
the invariant projector. The ramified column space \(\mathcal M_w\)
is the intersection of the kernels of these projectors. For a tuple
\(\boldsymbol\mu\) of row characters, write
\(\Theta^{[m]}_{w,\boldsymbol\mu}\) for the composite of the \(m\)
upward lifts on the \(X\)-tower. For the unitary pair on a Hermitian
space \(A\) and a skew-Hermitian space \(B\) over \(E\), fix the additive
character used in its Weil representation. Kudla's construction splits
the pullback of the local metaplectic extension over \(U(A)\times U(B)\)
using \(\chi_B\) on \(U(A)\) and \(\chi_A\) on \(U(B)\). The restrictions
of \(\chi_A,\chi_B\) to the ground field's multiplicative group are
\(\eta_{E/F}^{\dim A},\eta_{E/F}^{\dim B}\), with the local quadratic
character in each local construction. With Hecke characters having
these restrictions on \(\mathbb A_F^\times\), the local splittings
assemble to a splitting over the adelic product. A Hermitian space
can be given the skew-Hermitian convention by scaling its form by a
nonzero trace-zero element of \(E\).
\cite[\S\S3.1.1 and 4.1.1]{p2borade-et-al}
Each step uses the prescribed additive and splitting characters and
finite and real tests. In the definite class model the step \(j\to j+1\)
is the sum
\[
 (T_{j,w,\boldsymbol\mu}\phi)(c')
 =\sum_c \frac{K_{j,w,\boldsymbol\mu}(c,c')\phi(c)}
                 {|\operatorname{Stab}_j(c)|}.
\]
Here \(K_{j,w,\boldsymbol\mu}(c,c')\) is a map of coefficient spaces. It
incorporates the determinant characters, the unfolded Schwartz test,
and the determinant factors in the input and output identifications of
the coefficient spaces. The determinant characters and Schwartz tests
both depend on the tuple. The identification at stage zero realizes the
same input \(\phi\in\mathcal M_w\) in the column model for that tuple.

Write \([G]=G(F)\backslash G(\mathbb A_F)\). Let \(\iota_m\) denote the
top inclusion of unitary groups and \(\operatorname{res}_m\) the
coefficient restriction to \(X'\). If \(f_m\) is the automorphic form
representing the top row \(F_m\), the top period in bilinear conventions is
\begin{equation}
\label{eq:period-top}
 b_{w,\boldsymbol\mu}(\phi,F_m)
 =m_Y(\boldsymbol\mu,m)\int_{[U(Y_m)]}
     (\operatorname{res}_m\Theta^{[m]}_{w,\boldsymbol\mu}\phi)(h)
     f_m(h)\,dh .
\end{equation}
The product denotes contraction in the coefficient model, with the
lift restricted along \(h\mapsto\iota_m(h)\). For a Petersson pairing
conjugate-linear in its first variable, the first input is the
conjugate of the lifted column. The factor \(m_Y(\boldsymbol\mu,j)\)
is the product of the prescribed local parahoric indices at stage
\(j\). It equals \(1\) at all stages for the base tuple
\(\boldsymbol\mu_0\) and at \(j=0\) for every tuple.

In the opposite direction, pairing the downward theta integral against
the row multiplicity space at \(j\) gives the row map from \(j+1\) to
\(j\). This integral uses the factor \(m_Y(\boldsymbol\mu,j+1)\).
Composing the row maps gives \(\iota_{\boldsymbol\mu}\) into the common
bottom row space \(E_{\rm row}\). We use its restriction to the image
\(\mathcal V^0_{\boldsymbol\mu}\) of the chosen top row projector. The
\releasecite{Formalism/LocalRingCardinality/DataTheta.lean}{727}{definitions of the row measures and pairings}
specify the same coefficient and measure conventions in both directions.

\subsection{Changing the path of the local tests}

\teamanchor{proof.matched-square-tests}
For two dual pairs \((G,H')\) and \((H,G')\) in the same symplectic
group, suppose \(H\subset G\), \(H'\subset G'\), and the restricted
Weil kernels agree on \(H\times H'\) via the chosen splittings. In
Kudla's seesaw identity, for cusp forms \(f\) on \(H\) and \(f'\) on
\(H'\), with compatible quotient measures and an absolutely integrable
product of the common kernel with \(f(h)f'(h')\), one has
\[
 \int_{[H']}(\Theta_{H\to G'}f)|_{H'}(h')\,f'(h')\,dh'
 =
 \int_{[H]} f(h)\,(\Theta_{H'\to G}f')|_H(h)\,dh .
\]
Here each \(\Theta\) integrates the corresponding kernel against the
input cusp form. The bilinear convention places the Petersson
conjugation in the choice of dual inputs. Both sides integrate the
same function on \([H]\times[H']\).
\cite[\S1, Definition~1.3 and (1.14)]{seesaw}
The two paths below use compatible restrictions of the Weil splittings.

The integral row tests use an auxiliary grid with rank
\(r+jt+i(r+6)\) at \((i,j)\). A rightward step increases this rank
by \(r+6\), and a downward step decreases it by \(t\). The definite
column lifts defining the top pairing follow the \(X_j\) tower of
rank \(k+j(t+1)\). The
\releasecite{Development/AbelianVarietyTheory/Definitions/Divisors/ExistsGivenStageEqOne.lean}{138}{local grid realization}
uses the horizontal increment specified in the
\releasecite{Development/AbelianVarietyTheory/Definitions/Divisors/LeastSplitAuxiliaryPrimeThetaChainSpec.lean}{121}{auxiliary rank data}.
At a fixed tuple the square with lower left corner \((i,j)\)
compares the routes
\[
 \begin{aligned}
 (i,j+1)&\longrightarrow(i,j)\longrightarrow(i+1,j)
             &&\text{(down, then right)},\\
 (i,j+1)&\longrightarrow(i+1,j+1)\longrightarrow(i+1,j)
             &&\text{(right, then down)}.
 \end{aligned}
\]

\noindent\leavevmode\teamanchor{candidate.transfer}%
\textbf{AI-proposed candidate (level-preserving transfer of the split finite tests).}
At a prescribed split finite place \(v\), work over \(F_v\) with
uniformizer \(\varpi\) and the chosen additive character of conductor
\(\mathcal O_v\). Take the prescribed depth \(e\), the maximum of
one and the coset exponent at \(v\). The two determinant characters
\(\alpha,\beta\) of the sections have the prescribed ramified unitary
ratio \(\nu=\alpha/\beta\), with conductor exponent at most \(e\).
The vertex levels are open
subgroups of the integral general linear groups. Take an admissible
test on the old route, supported on integral matrices and invariant
modulo \(\varpi^e\) there, satisfying the old level and character
actions. Hold the variables outside the square fixed. Inside it
concatenate the two test matrices in column blocks, with the prescribed
transpose on one matrix. The numbers \(a,a'\) of rows on the new
and old routes are the ranks at \((i+1,j+1)\) and \((i,j)\).
The number \(n\) of columns is the sum of the other two ranks, so
\(n=a+a'\).

The local transfer applies the finite Fourier sum in the first
column block \(P\), with the factor \(q_v^{-ea'|P|}\) for residue
cardinality \(q_v\). It gives a function \(\Psi\) on \(a'\)-by-\(n\)
matrices. The transfer then uses the primitive-frame transform
below and the inverse Fourier sum in the complementary block.
Let \(\Lambda'\) be the row lattice after the first Fourier sum,
with \(\varpi^{-e}\mathcal O_v\) on the transformed coordinates and
\(\mathcal O_v\) on the others. For matrices with \(\ell\) rows write
\[
 I_\ell(\Phi,X)=\int_{\mathrm{GL}_\ell(F_v)}
        \Phi(gX)\,\nu(\det g)|\det g|^{n/2}\,dg .
\]
The measure for \(\ell=a'\) gives the inverse-transpose of the
prescribed lower level volume one. For \(\ell=a\) the subgroup
\(\mathrm{GL}_a(\mathcal O_v)\) has volume one. Write
\(X_{\rm st}=(I_a\ 0)\), \(Y_{\rm st}=(0\ I_{a'})\) for the
complementary coordinate frames. For \(h\in\mathrm{GL}_n(F_v)\),
the full row rank of \(Y_{\rm st}h\) confines the support of
\(\Psi(gY_{\rm st}h)\) to a compact set of matrices \(g\). Relative to
additive matrix measure, the absolute determinant weight in
\(I_{a'}\) is proportional to \(|\det g|^{(a-a')/2}\), which is
bounded on that set since \(a-a'=t+r+6>0\).
Thus \(I_{a'}(\Psi,Y_{\rm st}h)\) converges absolutely.
For \(X=X_{\rm st}h^{-\mathsf T}\) primitive in \(\varpi^{-e}(\Lambda')^\vee\),
the transform sets
\(\phi_A(X)=\nu(\det h)|\det h|^{n/2} I_{a'}(\Psi,Y_{\rm st}h)\).
It is zero off the primitive frames. Changing \(h\) with \(X\) fixed
changes \(Y_{\rm st}h\) by a left basis change. The inverse
determinant factor from the integral cancels the change in the
external weight \(\nu(\det h)|\det h|^{n/2}\).

The resulting zeta sections satisfy, for every \(h\),
\begin{equation}
\label{eq:period-complementary-sections}
 \begin{aligned}
 &\beta(\det h)|\det h|^{-a/2}I_a(\phi_A,X_{\rm st}h^{-\mathsf T})\\
 &\hspace{2em}=(\nu\beta)(\det h)|\det h|^{a'/2}I_{a'}(\Psi,Y_{\rm st}h).
 \end{aligned}
\end{equation}
For the first integral one can sum over Cartan shells. For fixed
\(h\), the basis changes making \(X_{\rm st}h^{-\mathsf T}\) a primitive
frame form one coset of \(\mathrm{GL}_a(\mathcal O_v)\). Its volume is
one and it lies in a single shell. Replacing \(h\) by
\(\operatorname{diag}(g^{-\mathsf T},I)h\) shows that the determinant
weights for \(g\) cancel on the coset. This gives the
\releasecite{Development/AbelianVarietyTheory/Proofs/DivisorProofs/PrimitiveFrameDatumZetaIntegralEqSmaller.lean}{143}{complementary-frame identity}.
To check depth, lift a change of a primitive frame modulo
\(\varpi^e\) to a principal-congruence column change. The conductor
bound makes its determinant character trivial. This calculation
and finite Fourier inversion give the
\releasecite{Development/AbelianVarietyTheory/Proofs/DivisorProofs/IsNatDatumSquareFlip.lean}{40}{prescribed support and depth}
on the new route. Here depth \(e\) means constancy modulo
\(\varpi^e\) times the prescribed matrix lattice. Use
\(\mathrm{GL}_a(\mathcal O_v)\) at the new intermediate vertex and
the original levels at the other vertices. This is the prescribed
level at the selected interior vertices. The
\releasecite{Development/AbelianVarietyTheory/Proofs/TwistedOperators/SplitThetaGridIsAdmissibleFlipTensor.lean}{180}{comparison of the level actions}
preserves admissibility with the required character on the new path.
Eischen--Harris--Li--Skinner use a partial Fourier transform in matrix
blocks to construct split local doubling sections with specified
support and congruence characters. \cite[\S4.3, (58) and
Lemma~4.3.2]{p2eischen-doubling}
In the square the transform acts on the prescribed depth space of
admissible tests, and the section identity relates the two complementary
ranks in the indicated measures.

\teamanchor{proof.principal-global-flip}
\subsubsection{The adelic path comparison}
Let \(\mathcal L_\gamma(\Xi)\) be the iterated normalized theta integral
on a path \(\gamma\) with test \(\Xi\). The arithmetic calculation uses
the following adelic matching on its prescribed tests. At each square,
choose real Fock tests on the new route in the same local quotient as
the old ones, and use the specified nonsplit tests and the split
transfer above. The comparison required after integrating the kernels
at the two intermediate vertices is a single nonzero scalar
\(c_{ij,\boldsymbol\mu}\) in the prescribed quotient measures. At each prefix of the chosen sequence of squares the identity on the
same cusp forms, with the same exterior factors, is
\begin{equation}
\label{eq:period-global-square}
 \mathcal L_{\rm right,down}(\Xi')
   =c_{ij,\boldsymbol\mu}\mathcal L_{\rm down,right}(\Xi).
\end{equation}
The scalar is independent of the cusp form and of the chosen tensor
expansion of the test for that square.
The adelic comparison follows by matching the real and finite models
of the integrated kernels, applying the regularized theta identity
and producing a nonzero leading integral for each square.
\releasecite{Development/AutomorphicArithmetic/ThetaCorrespondence/ThetaPairLifts/LocalTypeDictionary/TwistedSquareThetaIntegralsEq.lean}{17}{Integrated kernel comparison}
\releasecite{Development/AutomorphicArithmetic/ThetaCorrespondence/ThetaPairLifts/LocalTypeDictionary/ConstK.lean}{15}{Nonvanishing}
\releasecite{Development/AbelianVarietyTheory/Proofs/TwistedOperators/GridSquareFlipSmallRouteEqInvSmulBigRoute.lean}{314}{Application on a path}.
The split-place equality \eqref{eq:period-complementary-sections}
gives the local transferred section. The real and inert
normalized quotient maps, the placed nonzero test, and the
regularized comparison are proved in the square calculation
below. Its use in \eqref{eq:period-global-square} is on the
admissible placed path tests. The convergence argument below verifies the incoming cusp type and
the replacement tests at every prefix.

Given \eqref{eq:period-global-square}, expand the tested Schwartz
tensor into a finite sum of pure tensors and apply the square equality
with the exterior arrows fixed in each summand. The two quotient
integrations can be interchanged when the joined kernel times the cusp
forms is absolutely integrable, as in the seesaw above. Swapping each
downward arrow past each rightward arrow takes \(m^2\) swaps, so
induction gives the inverse factor
\[
 (c^\square_{\boldsymbol\mu})^{-1},
 \qquad c^\square_{\boldsymbol\mu}=\prod_{0\le i,j<m}c_{ij,\boldsymbol\mu}.
\]
If the original expression is divided by a character value
\(\gamma_{\boldsymbol\mu}\), the opposite path thus has denominator
\(\gamma_{\boldsymbol\mu}c^\square_{\boldsymbol\mu}\).

\teamanchor{proof.period-grid-prefixes}
\subsubsection{Convergence through the chosen squares}

Fix an installed tuple, its permitted finite row test and the real
polynomial tensor on the initial path. First consider the finite
factors. On the boundary \(i=0\) they are the placed rectangular
type tests, the tame parahoric tests at the tuple primes, and the
lattice vacuum elsewhere among the selected split places. The last
boundary edge uses the frame of \(X'\). Each symbol is supported
on the integral matrix lattice and constant modulo the residue
prime. For a type test this follows because the condition on its
triangular blocks and the tame determinant characters depend
only on the reductions. For the parahoric and lattice tests the
same holds by their support conditions. At the other boundary
\(j=0\), the coset tests are indicators of an integral framed
lattice coset at the prescribed depth. Their translations are
integral because the marked rational maps take the source
lattice to the target lattice. The local path function, the
product of these edge tests, is therefore supported on the
integral matrices and constant modulo the installed depth on
each of its edges.

In the finite matrix action at an interior vertex the absolute
values of integral unit determinants are one. The extra parity
character is unramified and is also one there. Along \(i=0\)
the two consecutive edge laws contribute the intervening type
character and its inverse. At \((0,0)\) the endpoint character
on the column side is trivial. Along \(j=0\) the vertex group
preserves the indicated coset and the edge determinant characters
are trivial on its congruence subgroup. Consequently the initial
path is admissible at the prescribed levels.
\releasecite{Development/AbelianVarietyTheory/Proofs/DivisorProofs/UnderMemAndPlaceAboveEqOfMem.lean}{432}{Support and depth},
\releasecite{Development/AbelianVarietyTheory/Proofs/DivisorProofs/UnderMemAndPlaceAboveEqOfMem.lean}{467}{the horizontal and downward level laws} and
\releasecite{Development/AbelianVarietyTheory/Proofs/DivisorProofs/UnderMemAndPlaceAboveEqOfMem.lean}{516}{their product on the chosen path}.

The ratio at a cell is the character \(\mu_{j+1}\xi_H\) divided
by the horizontal character indexed by \(m-i\).
At the deep auxiliary place the numerator is unramified and
the denominator is ramified. At the other fixed split grid
places the denominator is nontrivial on the first principal
congruence subgroup, where the two numerator characters are
trivial. At the extra tuple places the horizontal and the
reference vertical characters are unramified, while \(\xi_H\)
is ramified. The level data give the stated upper conductor
bound in each case. The parity character in the two splitting
characters is unramified. The ordinary additive character has
conductor \(\mathcal O_v\) at these places. This verifies the
hypotheses of the split transfer for each cell at the same
chosen split half and depth. The opposite middle vertex is new
to a monotone path through the cell. It receives the integral
level in the transfer, and the exterior characters and levels
remain those already used for the path. Induction on the
ordered squares now proves admissibility at every prefix.
\releasecite{Development/AbelianVarietyTheory/Proofs/DivisorProofs/ValuationLeOneOfMem.lean}{299}{The three conductor cases}
\releasecite{Development/AbelianVarietyTheory/Proofs/TwistedOperators/LocalFlipStageIsTensorPlaceDepth.lean}{25}{and the prescribed order and depths}.

The same frames give the finite factors of the actual boundary
lifts. Along the downward boundary, the shallow residue symbol is
the specified local lift vector in its two framed polarizations.
Outside the selected split halves it is the normalized vacuum.
The local changes of position leave only values of the relative
finite character. In the last term this uses the determinant
of the fixed frame on \(X'\) and its conjugate half.
Along the rightward boundary the factors are the prescribed
coset vectors. Multiplying the inverses of the downward
character factors gives the initial scalar
\(\gamma_{\boldsymbol\mu}\) in the representation of the path
tensor. If the downward factor at step \(j\) is \(c_j\) times
the section in the split chart, this scalar is \(\prod_jc_j^{-1}\).
The boundary polynomials use the very same real vectors as the
row chain. Replacing a fundamental domain for a rational
quotient has no effect since the integrand descends and the Haar
normalization is fixed by the vertex level.
\releasecite{Development/AbelianVarietyTheory/Proofs/TwistedOperators/GivenFiniteRepresentationChain.lean}{20}{The placed boundary vectors and scalar}
\releasecite{Development/AbelianVarietyTheory/Proofs/TwistedOperators/GivenPathRepresentationFinite.lean}{101}{and the initialized analytic path}.

For the cuspidality test at any intermediate vertex, work at
one real place and suppress that index. Use
\(R_{ij}=r+jt+i(r+6)\), let \(s=r+6\), and put
\(\delta_{ij}=-1/2\) for odd \(i+j\) and \(0\) otherwise.
Before taking the dual, the signed parameter at \((i,j)\) is
the following union translated by \(\delta_{ij}\):
the entries \(v_2,\ldots,v_k\) with positive sign,
the entries \(-b_\ell\) with positive sign when
\(\ell\le j+1\) or \(\ell-j-1\) is even and negative sign
otherwise, the positive \(t\)-strings centered at
\(-e_1,\ldots,-e_j\), and the positive \(s\)-strings centered
at \(c_{m-i+1},\ldots,c_m\). A string of length \(d\) centered
at \(c\) has entries \(c+(d-1)/2,\ldots,c-(d-1)/2\).
Here each \(c_\ell\) is the prescribed integral midpoint
between the successive upper nodes. Both \(t\) and \(s\) are
odd. Thus \(R_{ij}\) is odd exactly when \(i+j\) is even.
The splitting exponent for the parameter of the smaller
member, after the determinant-cover correction, is the
unshifted center plus \(\delta_{ij}\). The exponent for the
larger member is the same center plus the shift of that
member, by the opposite rank parity.

For a rightward step take \(d=s\) and center \(c_{m-i}\).
The defining upper gap exceeds \(2N_{\rm amb}^{3}+2s+2\),
where \(N_{\rm amb}\) bounds the grid ranks. The centers are
in descending order, with rounding error at most one.
If \(i>0\), the largest old upper string ends at
\(c_{m-i+1}+(s-1)/2\). The cloud and the lower strings are
below the old upper strings and below the first upper center. Every old entry is therefore at
distance at least \((s+1)/2\) below the new center. In the
negated larger Hermitian form the old labels reverse and the
new string has negative sign, as in the real theta rule.
For a downward step compare the pair in increasing rank,
with \(d=t\) and center \(-e_{j+1}\). The two inequalities
\[
 e_\ell-b_\ell\ \ge\ (t+3)/2,\qquad
 b_{\ell+1}-e_\ell\ \ge\ (t+3)/2
\]
show that the old entries below this center are exactly
the \(-b_\ell\) with \(\ell>j+1\). All old entries are at
distance at least \((t+3)/2\). Increasing \(j\) reverses
exactly these labels and inserts the positive \(t\)-string.
These assertions include the boundary cases \(i=0\), \(j=0\)
and \(j+1=m\), where the descriptions use the transported
\(X'\) parameter. Paul’s discrete-series theta rule gives
the specified partners with these signed lists, including
their duals in the bilinear theta integral.
\cite[Theorem~2.7(b) and Proposition~3.4(ii)]{p2paul-real-theta}
\releasecite{Development/AbelianVarietyTheory/Definitions/Foundations/StrictInterlaceIffDescEnum.lean}{808}{The signed lists},
\releasecite{Development/AbelianVarietyTheory/Proofs/TwistedOperators/ContragredientThetaParamsAddBlockArrow.lean}{17}{the \(s\)-step}
\releasecite{Development/AbelianVarietyTheory/Proofs/TwistedOperators/ContragredientThetaParamsRemoveBlockArrow.lean}{407}{and the \(t\)-step}.

Here is the lower-tower check in both directions. With the
splitting center removed, the smaller parameter is integral
with no zero, so for the odd rank difference \(d\) the integer
\(k_\pi\) of Ichino's criterion is \(-1\).
The signature of the larger partner has \(l=0\) and
\(t_{\rm Ichino}=(d-1)/2\), after exchanging the two signatures
when the inserted string is negative. The two nearby signed
sets in the nonvanishing criterion are empty by the preceding
distances. A lower member in that Witt tower has negative \(l\)
and the lift is zero. In the downward direction the larger
input has the complete centered \(d\)-string and no longer one.
The signs of its component characters alternate along the
string, because its signature labels there are constant.
Hence \(k_\pi=d\). The small partner has rank \(n-d\),
where \(n=R_{i,j+1}\) is the input rank. This is the first
permissible rank. Indeed the criterion for a discrete-series input with
\(k_\pi=d\) requires \(l\ge0\) in the balanced signature case
and \(l\ge d\) otherwise. The only exception at \(l=-1\)
requires an even-multiplicity weight, absent in a regular
discrete series. Thus no member of smaller rank in this
tower can occur.
\cite[\S6.1, Theorem~6.1 and Remark~6.2]{p2ichino-real-theta}

For any of the finite tensors just described, the theta kernel
and its derivatives have polynomial height growth on the
quotient in the integration variable. The Heisenberg--Poisson
bound used below for the triangle tests gives this estimate
before integration. Rapid decay of a smooth cuspidal input
therefore gives absolute integrability, locally uniform in the
target and polynomially bounded in its height. Real smooth
Howe duality on the specified pair gives the stated real type
in the image. The tower constant-term formula for a theta
lift of a cusp form expresses the proper-parabolic terms
through lower theta lifts in that target tower. The real
vanishing just checked kills each such term. Thus the next
function is again a smooth cusp form in the specified type
and satisfies the same rapid decay hypothesis for the next
integration.
\cite[Theorems~A and~2.23--2.24]{p2miller-schmid-decay}
\cite[Proposition~3.1]{p2wu-unitary-tower}
\releasecite{Development/AbelianVarietyTheory/Proofs/TwistedOperators/CuspAlongOfPairLifts.lean}{419}{The type on every edge}
\releasecite{Development/AbelianVarietyTheory/Proofs/TwistedOperators/PathConvZipFock.lean}{17}{and the path induction}.

At an active split place, an admissible path function is on
the finite product of its residue matrix spaces. Expanding in
products of coset indicators and applying the finite vertex
averages expresses it as a finite sum of allowed pure
Schwartz tensors. Exterior factors can be held fixed while
partitioning the two edges of a cell. Real polynomial tensors
admit a finite expansion for the same reason of finite
dimensionality in each degree, and outside the active finite
places the vacua are fixed. Apply the preceding cusp argument
to each summand. The absolute bounds for the two integrated
kernels of the square, multiplied by the rapidly decaying
incoming cusp form, justify Fubini with the fixed exterior
variables. This compares the two paths in the same vertex
measures. Together with the finite admissibility induction
it proves convergence and the stated equality after every
flip, for each installed tuple and every allowed test on the
original cusp input space.
\releasecite{Development/AbelianVarietyTheory/Proofs/TwistedOperators/GridSquareFlipSmallRouteEqInvSmulBigRoute.lean}{224}{Finite decomposition with exterior factors}
\releasecite{Development/AbelianVarietyTheory/Proofs/TwistedOperators/PathConvZipFock.lean}{40}{and convergence for the resulting sums}.

\subsection{The principal-level row test}
\teamanchor{proof.principal-row-tests}
Work in the prescribed principal-level model before imposing the
auxiliary bottom level character. Let \(j_{\boldsymbol\mu}^{\rm top}x\)
be the automorphic realization of a rational top coordinate \(x\) in
the good isotypic part, and let \(\Theta_{\boldsymbol\mu}^{\downarrow}\)
be the downward row chain in this model, with its fixed real vectors.
The functionals \(\psi_a\) on the bottom multiplicity are indexed by
the prescribed grid data \(\Phi_a\) and column tests \(\Psi_a\).
Each value is the Bessel integral of the column test against the
further grid lift of the bottom row, in the prescribed measures.
Write \(p_{\boldsymbol\mu,a}(x)=
\psi_a(\Theta_{\boldsymbol\mu}^{\downarrow}j_{\boldsymbol\mu}^{\rm top}x)\).

\teamanchor{proof.principal-level-projection}%
\subsubsection{Changing the bottom level}
The level-change group of the row chain is
\(\prod_{v\in P_g}\mathrm{GL}_r(k_v)\), where \(P_g\) is the set of
chosen split components at the two type primes and \(k_v\) their
residue fields. In the frames at those places the flag blocks have
sizes \(r,t,\ldots,t\), with one extra \(t\)-block on the larger
vertex. To represent a residue element \(g_0\), use a block
diagonal lift \(\operatorname{diag}(\widetilde g_0,I,\ldots,I)\)
in both vertices and the identity away from \(P_g\). The split
unitary group identifies with the general linear group at its
chosen component. An integral lift there gives the inverse dual
action at its paired component, so it preserves the self-dual
framed lattice. The lifts exist because a matrix with invertible
reduction has unit determinant. Both have flag character one.
Two lifts with the same reduction differ by a principal-level
element, whose character is canceled in the normalized
right-translation action.

At \(v\in P_g\) the row-step test in the Schr\"odinger model is
supported on integral matrices whose reduction is \(\binom{A}{0}\),
with \(A\) block upper triangular, invertible and with initial
\(r\)-block \(I_r\). For the step \(j+1\to j\) its coefficient
there is
\[
   \prod_{h=1}^j\chi_{h,v}(\det A_h)\,
          \chi_{j+1,v}(\det A)^{-1},
\]
where \(A_h\) are its \(t\)-blocks. The same formula applies
at the last step with \(X'\) as the larger vertex. For lifts
\(g,h\) on the small and large vertex use the matrix substitution
\(Y\mapsto h^{-1}Yg\). The preferred lifts preserve this support,
all the block determinants and the determinant of \(A\). The
split Weil factor evaluates the two splitting characters on
\((\det g)^{-1}\) and \(\det h\), with unit absolute
determinants. The two characters agree with
\(\chi_{j+1,v}^{-1}\) on the units, including in the last step
where the extra unramified character is trivial on units.
The two values cancel. This proves simultaneous invariance of
each step test under the preferred lifts. There is no action
at the nonsplit places, at the twist places, or on the real
Fock factor for these lifts.
\releasecite{Development/AbelianVarietyTheory/Proofs/TwistedOperators/TypePlacesLocalDatumBlockTriangularLaw.lean}{124}{The split row-step test}
\releasecite{Development/AbelianVarietyTheory/Proofs/TwistedOperators/LastStepVectorLiftPairFlagLevelCharacters.lean}{160}{The last-step model}.

In the theta integral against a right translate on the large
vertex, change variables by that translate on its rational
quotient. Simultaneous invariance of the Schwartz factor changes
the translated kernel to a right translate on the small vertex.
This proves the intertwining identity for each step on the
convergent automorphic integrals, and hence for their composition.
Apply the group-algebra idempotents on both sides by linearity.
The bottom good projector fixes the lift of a good top
coordinate. The subgroup character projector, when imposed at
both ends, restricts the same downward map to that level.
If the uncut image is the full permitted bottom space, its image
after the cut is the corresponding bottom projector image: apply
the top projector to any preimage of a vector in this image.

Denote by \(Q_{\boldsymbol\mu,a}(x)\) the corresponding Bessel
integral on the opposite path after the swaps. The remaining
matching required for this same \(\Phi_a,\Psi_a,x\) is the normalized
expansion
\begin{equation}
\label{eq:period-principal-unfolding}
\begin{aligned}
 Q_{\boldsymbol\mu,a}(x)
   &= (\gamma'_{\boldsymbol\mu})^{-1}
      s_{\rm AA}\sum_{\beta}
      \left(\prod_v d_{\beta v}\right)A_{\boldsymbol\mu,a,\beta}(x),\\
 A_{\boldsymbol\mu,a,\beta}(x)
   &=\kappa J_{\boldsymbol\mu,a,\beta}(x).
\end{aligned}
\end{equation}
Here \(\beta\) runs over the supported local coset families of
\(\Phi_a\), and \(d_{\beta v}\) is the rational measure weight of its
cell at \(v\). The factor \(s_{\rm AA}\) is the product of the
prescribed vacuum factors at the other finite places. The term
\(A_{\boldsymbol\mu,a,\beta}(x)\) is the sum of definite periods of
the two ascending lifts with the real Fock tensors from unfolding
and the finite tests from an expansion of that cell.
The arithmetic term \(J_{\boldsymbol\mu,a,\beta}(x)\) is the
corresponding class sum of the normalized jet pairings. Its
summands at each class have the reciprocal of the product of the
row and column stabilizer orders as weights. The measures in
\(A_{\boldsymbol\mu,a,\beta}\) are the measures for these particular
term levels. The second identity requires a single \(\kappa\) for
the real Fock tensor used in the unfolding, independent of
\(\boldsymbol\mu,a,\beta,x\) and the choice of finite expansion.
The two character values \(\gamma_{\boldsymbol\mu},
\gamma'_{\boldsymbol\mu}\) belong to the subgroup generated by the
local values of the prescribed finite-order character
\(\xi_{H,\boldsymbol\mu}\).

\teamanchor{proof.principal-opposite-cells}
\subsubsection{Finite row orbits and their measures}
Only the vertices \(I=\{1,\ldots,m-2\}\) of the last downward
column are integrated over a row-group quotient. Write \(R_j\)
for the alternating Bessel and Fourier--Jacobi row groups, \(G_j\)
for their ambient groups, and \(K_j\) for the full finite lattice
stabilizers. Keep the fixed initial inverse staircase and the neutral
endpoints. At an interior vertex the base finite point is the finite
part of the fixed flag tilt \(\delta_j\). Its translates used here
are
\[
 b_{\beta,j}=\delta_{j,f}\prod_{v\in S}
                s_{v,j}g_{\beta,v,j}.
\]
At a tested split place \(v\) in \(S\), identify \(G_j(F_v)\)
with \(\mathrm{GL}_{n_j+c_j}(F_v)\), where \(n_j\) is the
definite column rank and \(c_j=m-1-j\). The sign \(s_{v,j}\)
is diagonal and integral. It acts trivially on the middle block
and changes only the designated trailing flag entries at odd \(j\).
The matrices \(g_{\beta,v,j}\) represent the local double cosets
for \(R_j(F_v)\) and \(K_{j,v}=\mathrm{GL}_{n_j+c_j}(\mathcal O_v)\).
The local row group is the product of the middle
\(\mathrm{GL}_{n_j}(F_v)\) and the unipotent group with entries in
positions of strictly increasing flag weight. The middle block has
weight \(c_j+1\); the successive one-dimensional flag weights are
\(2,4,\ldots,2c_j\). Conjugation by the tilt and by the signs
preserves this unipotent group and acts trivially on the middle
Levi. Thus the labels are the double cosets in the stated frames,
with no further identification. To see surjectivity and injectivity
of this description on the ordinary-place support, lift each local
left factor on one split half by its Hermitian inverse transpose
on the other half, and take the identity elsewhere. Lift the right
integral factors in the same way. Self-duality gives membership in
\(K_j\) on both halves. The restricted-product factorization then
gives the required global left and right factors. Conversely each
global equality of double cosets gives its local equality by
projection. The real component, which is taken separately in the
quotient integral, does not affect a label.
\releasecite{Development/AbelianVarietyTheory/Definitions/TwistedInnerProducts/IsRealFam.lean}{147}{The fixed tilt, base points and representatives}
\releasecite{Development/AbelianVarietyTheory/Proofs/TwistedOperators/CellVolumeMulCoef.lean}{26}{The split row subgroup}
\releasecite{Development/AbelianVarietyTheory/Proofs/TwistedOperators/TiltedPointsCoveredByRepresentativeCells.lean}{28}{Local representatives in the finite row quotient}.

Here the ordinary-place support means the single orbit of the
base lattice at every finite place outside \(S\). In the actual
unfolded product, nonvanishing forces each two-arrow contraction
to be nonzero. Its two flat factors are then both nonzero.
The ordinary-place support calculation can be made one arrow
at a time. With the base moves already applied, use the
self-dual lattice vacuum at a place outside \(S\). The
Bessel-to-Fourier--Jacobi condition says that a nonzero flat
factor with source in its base orbit forces the next
lattice point into the Fourier--Jacobi base orbit. The
Fourier--Jacobi-to-Bessel condition forces the following
point into the Bessel base orbit. Starting from the fixed
initial point, induction gives the ordinary-place support at
every interior vertex. At the fixed auxiliary split places
within \(S\), only the reference staircase is allowed. There the
local input is the one-staircase support of the transferred
tensor at the exact test depth. It is applied to the *flattened*
tensor, rather than to each entry of the original path test.
These two support conditions refer to the placed flat
contractions described below.
\releasecite{Development/AbelianVarietyTheory/Proofs/TwistedOperators/SingleCosetOffTwistPlacesLargeDepth.lean}{189}{Induction with the two ordinary-place support tests}
\releasecite{Development/AbelianVarietyTheory/Proofs/TwistedOperators/RowTestSupportFamiliesStaircaseCosets.lean}{27}{The fixed-place staircase at the actual depth}.

For the cell of \(b=b_{\beta,j}\), its mass in the quotient
\(R_j(\mathbb A_{F,f})\backslash G_j(\mathbb A_{F,f})\) is
\[
 M_{j,\beta}=
 \frac{\operatorname{vol}(K_j)}
      {\operatorname{vol}_{R_j}(R_j(\mathbb A_{F,f})\cap bK_jb^{-1})}.
\]
Indeed its parametrization by the right compact group is the compact
quotient of \(K_j\) by \(b^{-1}R_jb\cap K_j\). Haar integration
over the stabilizer gives the displayed denominator. Write \(N_j\)
for the unipotent radical of \(R_j\) and \(J_{\beta,j}\) for the
projection of the stabilizer to the middle Levi. The latter is the
term level of the definite column at this vertex. Locally, this is
the condition that a middle-block element have a lift over the
given flag whose action preserves \(g_{\beta,v,j}\mathcal O_v^{n_j+c_j}\)
in the stated frame. The tilt and sign commute with the middle
block and normalize its flag radical. At ordinary places the
orthogonal self-dual decomposition gives the full middle lattice
stabilizer. This identifies the projection with the product
of the prescribed term levels.
\releasecite{Development/AbelianVarietyTheory/Proofs/TwistedOperators/VolumePeriodProjectionSetEqFamilyLevel.lean}{296}{Projection to the definite column level}.

The Haar measure on \(R_j\) is the product of the unipotent
and middle-block measures. In the split model the two sets of
rectangular coordinates on which the Levi has nontrivial
determinant are dual, so the middle action has modulus one.
Fubini on the compact stabilizer gives a denominator proportional
to \(\operatorname{vol}_{N_j}(N_j\cap bK_jb^{-1})
\operatorname{vol}(J_{\beta,j})\), with the proportionality
independent of the cell. When a Haar measure normalized to give
the reference column level \(J_{0,j}\) volume one is replaced by
the measure giving \(J_{\beta,j}\) volume one, the conversion
factor multiplying \(M_{j,\beta}\) is
\(\operatorname{vol}(J_{\beta,j})/\operatorname{vol}(J_{0,j})\).
There is no column-level change at the endpoints and no change
on the other definite path. Thus the variable Levi volume
cancels exactly. For the remaining factor at \(v\) set
\[
 \begin{split}
 U^0_{v,j}&=N_j(F_v)\cap K_{j,v},\\
 U^\beta_{v,j}&=N_j(F_v)\cap
                 g_{\beta,v,j}K_{j,v}g_{\beta,v,j}^{-1},\\
 d_{\beta v,j}
   &=\frac{[U^0_{v,j}:U^0_{v,j}\cap U^\beta_{v,j}]}
           {[U^\beta_{v,j}:U^0_{v,j}\cap U^\beta_{v,j}]}.
 \end{split}
\]
The fixed tilt multiplies both volumes by the same determinant.
The integral diagonal sign preserves \(U^0_{v,j}\) and its
volume. Hence the remaining factor is the fixed base measure
factor times \(\prod_v d_{\beta v,j}\). This gives
\(d_{\beta v}=\prod_{j\in I}d_{\beta v,j}\) with the inverse-volume
orientation in \eqref{eq:period-principal-unfolding}.
The fixed base factors and the ordinary-place vacuum values can
be included in the common normalization of the real Fock
expansion. The measure conversion defining each term is performed
before evaluating its definite periods.
\releasecite{Development/AbelianVarietyTheory/Proofs/TwistedOperators/CellVolumeMulCoef.lean}{656}{Relative indices in the actual frames}
\releasecite{Development/AbelianVarietyTheory/Definitions/TwistedInnerProducts/IsRealFam.lean}{100}{Product of term-level measure conversions}.

\teamanchor{proof.principal-opposite-phases}
\subsubsection{The phase of the finite expansion}
Use determinant half-roots \(d_j\) of the unitary base points
\(\delta_{j,f}\) and the initial inverse staircase. A half-root
of a determinant \(u\) satisfies \(d_j/\overline d_j=u\). Its
choice drops out of \(\xi_{H,\boldsymbol\mu}(d_j)\), since the
character is trivial on the fixed-field adeles. The root of a
real translate may be taken with finite part one. In a paired
downward contraction at \(j,j+1\), changing the twist by \(\nu\)
changes the two operators by
\[
 \nu(d_j)^{-1}\nu(d_{j+1})\,
 \nu(d_{j+1})^{-1}\nu(d_{j+2})
 =\nu(d_j)^{-1}\nu(d_{j+2}).
\]
The intermediate determinant cancels with this sign under the
linear flat contraction. The product on even \(j\) telescopes,
and the final root can be taken to be one. In the reference-vacuum
comparison the factor used is
\[
 Z_{\rm base}(\boldsymbol\mu)=
 \prod_{j\ {\rm even},\,j<m}
 \frac{\xi_{H,\boldsymbol\mu}(d_{j+2})}
      {\xi_{H,\boldsymbol\mu}(d_j)}.
\]
\releasecite{Development/AbelianVarietyTheory/Proofs/TwistedOperators/UnfoldingJointDataTwistRatioAtDetRoots.lean}{78}{Paired determinant twist}
\releasecite{Development/AbelianVarietyTheory/Proofs/TwistedOperators/UnfoldedVacuaDifferByTwistValuesLargeDepth.lean}{31}{Normalization at the base points}.

On the opposite path at a selected \(v\), for a matrix
assignment \(g\) write
\[
 X_{v,\boldsymbol\mu}(g)=
   \prod_{b\ {\rm downward}}
     \frac{\xi_{H,\boldsymbol\mu,v}(\det g_{b,\mathrm{source}})}
          {\xi_{H,\boldsymbol\mu,v}(\det g_{b,\mathrm{target}})} .
\]
This is the ratio of the local determinant factors on switching
from the untwisted to the twisted model. The rightward arrows
use the other splitting character, so they contribute one.
The absolute-value factors of the integral sign assignment
\(s_v\) are one. At the retained-label places the untwisted determinant factors at
\(s_v\) are also one, because the downward factors are unramified
on those units and the rightward vertices are unchanged. At the fixed auxiliary
split places one moves from the retained reference staircase \(g^0_v\)
to the base vacuum. The determinant factors of the untwisted
model there are fixed. The remaining ratio is therefore
\[
 Z_{\rm loc}(\boldsymbol\mu)=
    \prod_{v\in S}\bigl(X_{v,\boldsymbol\mu}(s_v)
                        X_{v,\boldsymbol\mu}(g^0_v)\bigr)^{-1},
\]
using \(g^0_v=1\) at the places where the cell label is retained.
The common scalar of the reference model can be placed in the
first even factor. When expressing the moved definite column
vacua in the finite dictionary there is one further product
\(Z_{\rm move}(\boldsymbol\mu)\) of determinant twist values.
Write \(x_j\) for the determinant half-root of the prescribed
finite moving element on the definite vertex \(j\). The two
splitting characters of arrow \(j\) change by the same twist.
With the source and target conventions in the moved column
action its multiplier is
\(\xi_{H,\boldsymbol\mu}(x_j)^{-1}
   \xi_{H,\boldsymbol\mu}(x_{j+1})\).
The reference model is the same moved vacuum at the base
tuple. These character values factor out before summing the
vacuum expansion. Their product over the column arrows is
\(Z_{\rm move}\).
\releasecite{Development/AbelianVarietyTheory/Proofs/TwistedOperators/FlatDictionaryScalarLargeDepth.lean}{23}{Local determinants of signs and base moves}
\releasecite{Development/AbelianVarietyTheory/Proofs/TwistedOperators/MovedVectorsAreCombinationsOfVacuaOfLargeDepth.lean}{27}{Moved definite column vacua}.

The scalar removed from the finite expression is consequently
\[
 \gamma'_{\boldsymbol\mu}
   =Z_{\rm loc}(\boldsymbol\mu)Z_{\rm move}(\boldsymbol\mu)/
                      Z_{\rm base}(\boldsymbol\mu).
\]
In particular the ratio uses the same determinants for every term
of a cell. These are values of the given finite-order twist and
their inverses, so the simultaneous test field contains the
phase. The finite character \(\gamma_{\boldsymbol\mu}\) from the
preceding swaps is applied separately.
\releasecite{Development/AbelianVarietyTheory/Proofs/TwistedOperators/IsExpansionZero.lean}{64}{Phase in the flat expansion}
\releasecite{Development/AbelianVarietyTheory/Proofs/TwistedOperators/RowTestCellValueOfStepsLargeDepth.lean}{29}{Assembly with the moved column factor}.

\teamanchor{proof.principal-opposite-flat}
\subsubsection{The placed finite Fourier contractions}
\paragraph{The local models at the points of a cell.}
Fix the tuple, one split place \(v\in S\), and one row cell.
Write \(B_j\) for the columns of the row frame at the chosen split
component and \(P_j\) for the frame of its definite core. These are
the two frame assignments in the simultaneous choice. Write
\(\epsilon_j=s_{v,j}\), \(g_j=g_{\beta,v,j}\), \(S_0\) for the initial
staircase, and \(S_j=1\) for \(j>0\). The finite one-place point \(r\)
has coordinates \(1,\epsilon_jg_j,1,1\) at the initial, interior, and
last two vertices, respectively. The reference point \(r_0\) has
coordinates \(S_0^{-1},1,\dots,1\). In particular the matrix of
\(\rho_j=r_{0,j}^{-1}r_j\) in \(B_j\) is \(\epsilon_jg_j\) at
every vertex, with \(g_0=S_0\) and \(g_{m-1}=g_m=1\).
The two point assignments are identity off \(v\). On the other split
component use their inverse Hermitian transposes. Let \(h_j\)
be the common factor, which has the local component of the fixed
tilt at \(v\), and put \(a_j=h_jr_{0,j}\). At the last two vertices
all three points can be taken equal to the identity. Thus the
required relative product is \(a_j\rho_j=h_jr_j\).
\releasecite{Development/AbelianVarietyTheory/Proofs/TwistedOperators/EpsVEqOne.lean}{191}{Placed and reference points}
\releasecite{Development/AbelianVarietyTheory/Proofs/TwistedOperators/DictionaryPointsLocalMatricesLargeDepth.lean}{121}{Their local matrices}.

Use the matrix coordinate \(X\) of a map from vertex \(j\) to
vertex \(j+1\) on this split component. If \(\phi_b\) denotes the
given split test at the downward grid arrow \(b\), the coordinate
in \(\phi_b\) is \(X^{\mathsf T}\). On the chosen component the
global Hom lattice at these vertices restricts to all integral
matrices in \(B_j,B_{j+1}\). Indeed the inverse frames on the two
sides send the lattice-inclusion condition to entrywise
integrality. The supported split half is zero at every other
component, and the other half of the trace-dual Hom lattice is its
adjoint. Consequently the one-place model \(\mathcal M_b\) with
these frames sends the integral indicator \(O_b\) to the reference
vector \(v_b\) and \(\phi_b(X^{\mathsf T})\) to the given test
vector \(f_b\). This is the local realization of the specified
modification of the global vacuum. Equivalently, one obtains the
one-place model by applying Heisenberg translations and character
projectors at \(v\) to \(v_b\); finite Fourier inversion on the
quotients of the integral matrix lattice gives
\(\mathcal M_b(\phi_b(\,\cdot\,^{\mathsf T}))=f_b\).
\releasecite{Development/AbelianVarietyTheory/Proofs/TwistedOperators/FrameLatticePointsEqIntegralMatricesLargeDepth.lean}{34}{Split Hom lattice in the chosen frames}
\releasecite{Development/AbelianVarietyTheory/Proofs/TwistedOperators/ArrowLocalModelMovedFrames.lean}{75}{Transpose and realization of the given local vector}.

The chosen oscillator action at a supported point of local matrices
\((k_j,k_{j+1})\) is multiplication by \(D_b(k)\) and pullback by
\(X\mapsto k_{j+1}^{-1}Xk_j\). In terms of the downward grid arrow
\(b\) with \({\rm src}(b)=j+1\) and \({\rm tgt}(b)=j\), the factor is
\[
 D_b(k)=
 \frac{\kappa_{b,j+1}(\det k_{j+1})}{\kappa_{b,j}(\det k_j)}
 |\det k_j|_v^{\operatorname{rk}(B_{j+1})/2}
 |\det k_{j+1}|_v^{-\operatorname{rk}(B_j)/2}.
\]
Here \(\kappa_{b,i}\) is the block character at \(b\) divided by
the stipulated parity power of \(\eta_v\) at vertex \(i\).
The two splitting characters evaluate on single-component ideles
with precisely these values, since projection of the idempotent
idele at the chosen component gives \(\det k_i\), and the modulus
uses the inverse transpose on the grid arguments. The powers of
the absolute value are the square roots of the Jacobians on the
matrix space for this component.
\releasecite{Development/AbelianVarietyTheory/Definitions/Divisors/RankEqCardIffExistsFrameMatrixFirstColumnBlock.lean}{986}{Edge factor with the inverse transpose}
\releasecite{Development/AbelianVarietyTheory/Proofs/TwistedOperators/LocalCharacterComponentsLargeDepth.lean}{74}{Identification of splitting characters}
\releasecite{Development/AbelianVarietyTheory/Proofs/TwistedOperators/ArrowLocalModelMovedFrames.lean}{93}{Weil action on both vectors}.

In the moved frames \(B'_i=a_iB_i\epsilon_i\), with core indices
first and flag indices second, set
\[
 H(X)=\epsilon_{j+1}^{-1}X\epsilon_j,\qquad
 \mathcal M'_b(F)=\omega_b(a_j,a_{j+1})\mathcal M_b(F\circ H).
\]
Indeed the same polarized vector in the two frames satisfies
\[
 a_{j+1}(B_{j+1}X B_j^{-1})a_j^{-1}
       =B'_{j+1} H(X) (B'_j)^{-1}.
\]
The opposite-half operators are transported by the adjoint
identity for this equation. Thus \(\mathcal M'_b\) has the
translation and multiplication laws of the Schr\"odinger
model in \(B'_j,B'_{j+1}\), with no choice of scalar in this
transport. Since the two \(\epsilon_i\) are integral invertible
matrices, \(O_b\circ H=O_b\). Substituting the relative point
in the action gives
\[
 \begin{split}
 \mathcal M'_b(O_b)&=\omega_b(h_jr_{0,j},h_{j+1}r_{0,j+1})v_b,\\
 \mathcal M'_b(\Phi_b)&=\omega_b(h_jr_j,h_{j+1}r_{j+1})f_b,\\
 \Phi_b(Z)&=D_b(\epsilon g)\,
     \phi_b\bigl(g_j^{\mathsf T}Z^{\mathsf T}g_{j+1}^{-\mathsf T}\bigr).
 \end{split}
\]
In this equation \(Z\) is reindexed in the order of \(B'_j\) and
\(B'_{j+1}\). The identity
\[
 (\epsilon_{j+1}g_{j+1})^{-1}
       (\epsilon_{j+1}Z\epsilon_j^{-1})(\epsilon_j g_j)
 =g_{j+1}^{-1}Zg_j
\]
accounts for the absence of \(\epsilon\) from the argument of
\(\phi_b\).
\releasecite{Development/AbelianVarietyTheory/Proofs/TwistedOperators/ExistsLocalModelMovedFrames.lean}{77}{Transport of the polarization}
\releasecite{Development/AbelianVarietyTheory/Proofs/TwistedOperators/ArrowLocalModelMovedFrames.lean}{110}{Placed inputs in the same local model}.

At \(v\in S\) write \(k=F_v\), normalize additive Haar
measure by \(\operatorname{vol}(\mathcal O_v)=1\), and
use the given character \(\psi_v\) with annihilator
\(\mathcal O_v\). For an arrow \(j<m-1\) let \(c=m-2-j\).
A downward matrix has block sizes
\((n_{j+1},c)\) by \((n_j,c+1)\).
Start with
\(\left(\begin{smallmatrix}x&0\\0&(1_c\;\;0)\end{smallmatrix}\right)\)
in these coordinates. For \(p=\lfloor c/2\rfloor+1\) allow free entries
in the lower-left rows \(i<p\), upper-right columns \(i>p\),
and strictly upper-triangular positions of the flag block.
In the upper-right column \(p\) retain \(y\) when \(j\) is even.
In the lower-left row \(p\) retain \(u\) when \(j\) is odd.
These retained entries are outside the free positions. Denote the
resulting affine matrices by \(C_j(x,y;e)\) and \(R_j(x,u;e)\),
where \(e\) runs through the free entries, and let \(s(A)\) be
the sum of the first superdiagonal of the flag block.
For the two moved local Schr\"odinger inputs \(\Phi_1,\Phi_2\)
the two flat symbols in an interior pair \(j,j+1\) are
\[
 \begin{aligned}
 F_1(x,y)&=\int \Phi_1(C_j(x,y;e))\,
                              \psi_v(-s(C_j(x,y;e)))\,de,\\
 F_2(x',u)&=\int \Phi_2(R_{j+1}(x',u;e'))\,
                              \psi_v(-s(R_{j+1}(x',u;e')))\,de'.
 \end{aligned}
\]
Here \(\Phi_i\) includes the shifted frame action, its determinant
factor and the prescribed coordinates of the split transfer.
The sign of the row-column contraction in these frames is
\[
 \mathfrak C_v(\Phi_1,\Phi_2)(x,x')
      =\int F_2(x',u)F_1(x,y)\psi_v(-u y)\,dy\,du .
 \tag{*}
\]
In the terminal pair \(j=m-2\) there are no free entries in
the first matrix and no row integration after it. At the last
vertex write \(x'=\binom{Y}{u}\) with the extra line in the
prescribed position. The terminal formula is
\[
 \mathfrak C_v^{\mathrm{last}}(\Phi_1,\Phi_2)(x,x')
     =\Phi_2(Y)\int\Phi_1((x\;\; y))\psi_v(-u y)\,dy .
\]
The endpoint decomposition has \(z_0\) as marked line vector;
the projection of its image into the preceding line has
coordinate \(-\xi_v a\,\zeta_v\), with \(a=1\) at this
last pair. This is the minus sign in the terminal Fourier
argument. The zero Witt index on the remaining source removes
the auxiliary coordinates by evaluation at zero. The final
position of the marked row is that of \(x'\) in the normalized
top frame.
\releasecite{Development/AbelianVarietyTheory/Definitions/LatticeFrames/ContrFactor.lean}{14}{The free entries and the paired transform}
\releasecite{Development/AbelianVarietyTheory/Definitions/Foundations/SimpleSum.lean}{53}{Choice of free entries and their character}
\releasecite{Development/AbelianVarietyTheory/Proofs/TwistedOperators/LastBiExpand.lean}{16}{The placed augmentation at the last vertex}.

\paragraph{Witt columns and the two contractions.}
The frame at an even Bessel vertex of Witt index \(\rho\) is
obtained by inserting \(P_j\) on the definite core, putting
\(s\xi_v^{\rho-i}e_i\) in flag position \(i\), putting \(s\zeta\)
in the marked position \(\rho\), and putting
\(s(-h/\xi_v)(-1/\xi_v)^{\rho-1-i}e'_i\) in the reflected
position \(2\rho-i\), where \(0\le i<\rho\),
\(s=\xi_v^{-\rho}\),
\(h=-(-\xi_v)^\rho\xi_v^\rho\), and the marked line vector
\(\zeta\) is included in the core of the Witt decomposition.
The two Witt vectors in the original orthogonal flag coordinates
are \(\xi_v^i w_i\) and \((-\xi_v)^{-i}w_{2\rho-i}\). Inserting
these gives \(w_i\) and \(-w_{2\rho-i}\), and \(\zeta\) has
coordinate \(\xi_v^\rho\) in the middle line, giving \(w_\rho\).
These agree with the columns \(B'_j\). At the initial vertex the
powers of the staircase in \(B_0\) cancel against \(S_0^{-1}\)
on the flags and have exponent zero on the core.

At the intervening Fourier--Jacobi vertex with index \(\rho\)
the corresponding columns are \(P_{j+1}\) on the core, the
isotropic columns \(s\xi_v^{\rho-i}e_i\), and the reflected
columns
\[
 s(h/\xi_v)(-1/\xi_v)^{\rho-1-i}e'_i
 \qquad\text{in positions }2\rho-1-i.
\]
The two multipliers simplify to
\(\xi_v^{-i}\) and \((-\xi_v)^i\). These are the tilt
multipliers on \(e_i,e'_i\). On the reflected flag the frame
already contains \(-1\), canceled by \(\epsilon_{j+1}\).
For the following Bessel vertex the normalization is
\(s'=\xi_v s\) and its last norm satisfies \(-\xi_v^2h'=h\).
The same middle frame therefore serves both arrows of the
pair, with the same marked vectors and dual pairings. At
Witt index zero the flag formulas have no isotropic columns.
\releasecite{Development/AbelianVarietyTheory/Definitions/LatticeFrames/ExistsUnitarySupportedAtPlaceOfMulEqIdem.lean}{78}{Mixed-model frame columns}
\releasecite{Development/AbelianVarietyTheory/Proofs/TwistedOperators/BesselFrameOfGridLatticeLargeDepth.lean}{167}{Bessel columns in the chosen frame}
\releasecite{Development/AbelianVarietyTheory/Proofs/TwistedOperators/JacobiFrameOfGridLatticeLargeDepth.lean}{49}{Fourier--Jacobi columns in the same frame}.

Consider now the actual mixed-coordinate operators on
\(\mathcal M'_b(\Phi_b)\). Heisenberg modulation at an unread
entry \(Z_{ia}\) multiplies by
\(\psi_v(t(Z_{ia}-Z^0_{ia}))\), where \(Z^0\) is the matrix
with the marked flag block \((1_c\;\;0)\) and zero off the
core and this block. Its image under the actual flat map is
unchanged. This follows by pairing the respective column or row
with its dual: in the split component the Heisenberg pairing is
the trace pairing, and the above inverse scales on the dual
columns cancel the scales on the evaluated columns. The marked
columns pair with coefficient \(1\) at their diagonal positions,
and the orthogonal-to-core rows have pairing zero with the core.
For a free elementary entry \(E\), the translation
\(\Phi_b(Z+E)\) has flat image \(\psi_v(s(E))\) times the
original image. In the upper flag block this follows from
\(s(UZ)=s(Z)+\sum U_{i,i+1}\) on the affine slice when \(U\)
is upper unipotent with diagonal one. Multiplying by a simple
root on the target uses its inverse, so its factor is
\(\psi_v(-\sum U_{i,i+1})\). Left multiplication by the simple
roots accounts for the free entries of the earlier flag
columns, using the diagonal \(1\) in each such column as the
pivot. In the remaining free rectangular blocks the corresponding
matrix unit pairs orthogonally with the marked blocks and gives
the factor \(1\). A free transvection meeting the marked
column pairs only with its neighboring flag vector and gives
exactly that superdiagonal entry. These identities come from
moving translations through the defining mixed Fourier
integrals. Their coordinate changes are additive translations
or triangular shears with unit diagonal.
\releasecite{Development/AbelianVarietyTheory/Proofs/DecoratedLatticeProofs/LocalFlatMapFreeEntryTranslationLawsBesselToJaco.lean}{378}{Free translations with the Bessel marked column}
\releasecite{Development/AbelianVarietyTheory/Proofs/DecoratedLatticeProofs/FlatFourierJacobiToBesselColumnMulLaws.lean}{36}{The reverse column evaluations}
\releasecite{Development/AbelianVarietyTheory/Proofs/DecoratedLatticeProofs/FourierJacobiBesselFlatMapUnipotentLaw.lean}{39}{Flag translations for the reverse operator}.

Consequently the quotient by these relations is the indicated
partial Fourier integral \(F_1\) or \(F_2\). One can see the
whole quotient by first imposing the evaluations at the early
nonfree flag entries and then using their diagonal pivots to
perform the remaining simple-root translations. The other
nonfree entries are evaluated by modulation. Translation in the
remaining free entries integrates against \(\psi_v(-s)\).
On any finite partition by compact open rectangles the
evaluation eliminates the rectangles missing the marked values,
and the integration identifies each translated indicator with
its prescribed character times the original indicator. Taking
all such partitions proves the factorization for every
Schwartz--Bruhat input. The surviving translations and
multiplications on \(F_1,F_2\) are the Heisenberg operators in
the indicated core-plus-line frames, by applying the same
trace pairing to the retained row and column. The linking
map between the lines pairs their coordinates by \(-uy\),
using \(s'=\xi_vs\), \(-\xi_v^2h'=h\), and the fact that
the link carries the second marked vector to the negative of
the first. Thus the actual two-arrow contraction factors through
\(\mathfrak C_v\). Applying this factorization to \(O_b,O_{b'}\)
gives the reference contraction with the product of the integral
indicators. If the symbol is expanded as a sum of products,
linear descent in each core coordinate sends that expansion
to the same sum of the Heisenberg operators applied to this
reference contraction. In particular, the descended map kills
every relation between two such expansions.
\releasecite{Development/AbelianVarietyTheory/Proofs/DecoratedLatticeProofs/FlatMapFourierJacobiFactorsLocalFlat.lean}{55}{Evaluation and integration for the column descent}
\releasecite{Development/AbelianVarietyTheory/Proofs/DecoratedLatticeProofs/FlatMapLawsFactorFourierJacobiSource.lean}{59}{Evaluation and integration for the row descent}
\releasecite{Development/AbelianVarietyTheory/Proofs/TwistedOperators/FlatDictionaryDoubleStepAtPlace.lean}{13}{The placed interior contraction}.

\paragraph{Pullback at the final line.}
For \(j=m-2\), the first arrow retains the column \(y\) and
has no free entries. Its two inputs after the Bessel flat map
are therefore read in the core-plus-line frame
\((P_j,\zeta)\) and the core frame \(P_{j+1}\).
On the next arrow the source has Witt index zero. The map
called extension by zero here inserts the source core into the
full source, which at index zero is its entire space. Its model
is \(Y\mapsto\Phi_2(Y)\), with empty complementary coordinates,
and the reference model is \(1_{\rm int}(Y)\).

Write the final column space as \(T\perp L''\), where \(T\)
is the last definite row space. Its frame is
\((i_T B_m,i_{L''}z_0)\) with \(z_0=1\). The map to the
preceding marked line is the skew adjoint of \(\tau w\),
where \(\tau z_0=-\xi_v\zeta\). In particular it exchanges the
two polarized halves. Its Hermitian normalization is
\(\tau^*H_L\tau=-\xi_0\overline{\xi_0}\,H_{L''}\);
in the two line frames \(H_L=-1\) and \(H_{L''}=1\).
For a column \(y\) on the first line and a row \(u\) on
the last line, substitution in the trace pairing on the chosen
component cancels \(\xi_v\) against the character normalization
and gives \(-uy\). Therefore the pullback of a line test \(c\)
under the skew adjoint, in the Schr\"odinger model on the last
line, is represented by
\(\widehat c^{\, -}(u)=\int c(y)\psi_v(-uy)\,dy\).
Indeed translation of the new coordinate becomes multiplication
by the old dual character and conversely, and Fourier inversion
uses the same Haar measure of volume one on integral vectors,
which is self-dual for the indicated conductor \(\mathcal O_v\).
The two polarizations therefore give the terminal symbol
\(\Phi_2(Y)\widehat{\Phi_1(x,\cdot)}^{\, -}(u)\).
The fixed reindexing sends the \(T\)-rows to their original
indices and the \(L''\)-row to the final index of \(P_m\);
it preserves the integral indicator and the coordinate pairing.
This is exactly the augmented definite matrix argument at the
last vertex. The same computation for the two reference
indicators gives the product of the two integral indicators in
the definite matrix coordinates.
\releasecite{Development/AbelianVarietyTheory/Proofs/TwistedOperators/LastTwoStepsLocalModelFlatMapFourierJacobiStep.lean}{24}{The first input at Witt index zero}
\releasecite{Development/AbelianVarietyTheory/Proofs/TwistedOperators/LastTwoStepsLocalModelExtendByZeroTopStep.lean}{23}{The second input in the row summand}
\releasecite{Development/AbelianVarietyTheory/Proofs/DecoratedLatticeProofs/LineModelKappaTransformIsLocalModel.lean}{24}{The skew adjoint and the negative Fourier transform}
\releasecite{Development/AbelianVarietyTheory/Proofs/TwistedOperators/LastTwoStepsTopVertexFrameAtPlace.lean}{25}{The augmented frame and its reindexing}
\releasecite{Development/AbelianVarietyTheory/Proofs/TwistedOperators/FlatDictionaryLastTwoStepsAtPlace.lean}{74}{The placed terminal contraction}.

For integral indicators on the large matrices, the flat
operators give the integral indicators in their smaller
matrix spaces. Indeed all their fixed flag entries are zero
or one, and integration in each free entry has volume one and
trivial character there. The conductor formula
\(\widehat{1_{\mathcal O_v^n}}=1_{\mathcal O_v^n}\)
shows that the contracted interior symbol is
\(1_{M_{n_{j+1},n_j}(\mathcal O_v)}(x)
  1_{M_{n_{j+2},n_{j+1}}(\mathcal O_v)}(x')\),
and the same holds for the augmented terminal symbol. These
calculations fix the reference normalization with no residual
Haar factor at the two-arrow contraction. For a moved symbol
partition the support into rectangular compact open cosets.
Fourier transformation of one such indicator is
\[
 \widehat{1_{b+L}}(u)
    =\operatorname{vol}(L)\psi_v(-ub)1_{L^\perp}(u).
\]
Consequently the retained symbols are finite sums of products
of functions on the two matrix spaces. Translation and
modulation in either retained matrix factor act on the
standard lattice vacuum by the corresponding Heisenberg
operators in the same frame. This gives the finite dictionary
by applying these products of translations and Fourier
projectors to the reference product of indicators. A relation
between two finite expansions is killed by the descended
linear operator. Thus no independence assumption on the
chosen indicators or translates is required.
\releasecite{Development/AbelianVarietyTheory/Proofs/DecoratedLatticeProofs/FlatMapBesselToFourierJacobiLocalModel.lean}{24}{Covariances for the column descent}
\releasecite{Development/AbelianVarietyTheory/Proofs/DecoratedLatticeProofs/FlatMapFourierJacobiToBesselLocalModel.lean}{20}{Covariances for the row descent}
\releasecite{Development/AbelianVarietyTheory/Proofs/DecoratedLatticeProofs/LineContractionSplitPlaceLocalModel.lean}{26}{The paired Heisenberg action}
\releasecite{Development/AbelianVarietyTheory/Proofs/TwistedOperators/SplitPlaceDoubleStepFlatContractionExpand.lean}{26}{Product expansion relative to the integral vacuum}.

The one-staircase claim at the fixed auxiliary places uses
the conductor and periodicity on these same symbols. Write
\(e_0\) for the test depth and \(\varpi\) for a uniformizer.
Index flag vectors starting at \(1\), and put
\(d_i=\varpi^{e_0 i}\). At the initial step the source model
of the support calculation uses the staircase
\(\operatorname{diag}(1_{n_0},d_1,\ldots,d_{m-1})\).
This is the matrix in the test argument; the acting
base point used to untwist the vacuum is its inverse.
Suppose at one step the source frame uses the staircase
and write \(h\) for the target finite representative in
the unshifted local symbol. Its lattice is
\(L=h\mathcal O_v^{n_{j+1}+c}\).
A nonzero value of the joint contraction implies a point
of the affine support where the moved full matrix is integral.
Its \(i\)-th flag column \(v_i\) has flag coordinate \(1\)
at \(i\), zero below \(i\), and zero middle block if
\(i\le \lfloor(c+1)/2\rfloor\). Integrality says \(d_i v_i\in L\).

For the reverse divisibility take \(v\in L\) with
zero flag coordinates below \(i\) and zero middle block
when \(i\le \lfloor c/2\rfloor\). In the next flag column
\(i+1\) of the test argument translate by \(t v/d_i\)
for \(t\in\mathcal O_v\). It changes only free entries.
Conjugating this change to the moved coordinates multiplies
on the right by \(d_{i+1}\) and on the left by \(h^{-1}\).
The resulting change is integral at depth \(e_0\), so the
test value is invariant. The only change in the Fourier
character is the factor \(\psi_v(-t v^{(i)}/d_i)\), where
\(v^{(i)}\) denotes the \(i\)-th flag coordinate.
The pair term \(-u y\) is unchanged because both its
entries are retained. Translation of the integral gives
\(\psi_v(-t v^{(i)}/d_i)=1\) for all \(t\in\mathcal O_v\).
The conductor condition gives \(v^{(i)}\in d_i\mathcal O_v\).

Let \(U\) have columns \(v_i\) on the flag and identity on
the middle block. It belongs to the flag radical with the
middle block at position \(\lfloor(c+1)/2\rfloor\).
Then \(U^{-1}L\) contains \(d_i e_i\) on each flag line.
Choose an integral basis of \(U^{-1}L\) adapted to the block
flag with middle block of position \(\lfloor c/2\rfloor+1\).
Such a basis is obtained by successive lifts of bases of
the filtration quotients of the lattice, which are free over
the discrete valuation ring. The two divisibilities just
proved give \(d_i\mathcal O_v\) as the projected lattice on
the \(i\)-th successive flag line. Replacing its lift by
\(d_i e_i\) keeps a basis, because it changes the \(i\)-th
pivot by a unit and adds only earlier coordinates.
The middle basis can have entries only in the middle block
and the earlier flag coordinates. The resulting basis matrix
is a middle Levi element times an element of the flag
radical times \(\operatorname{diag}(1,d_1,\ldots,d_c)\).
Both radicals used in the two orders lie in \(R_{j+1}\).
Therefore \(h\) belongs to the reference staircase double
coset. Applying this argument successively, with the chosen
reference representative of that coset at each step, proves
the single allowed staircase at these places. Here the
depth is exactly the staircase depth \(e_0\) of the
transferred tensor, and the argument applies to the joint
integral even if this tensor is a finite sum.
\releasecite{Development/AbelianVarietyTheory/Proofs/DivisorProofs/OneCosetOfIsTensor.lean}{183}{Conductor detection by depth-preserving translations}
\releasecite{Development/AbelianVarietyTheory/Proofs/DivisorProofs/ExistsKOOfDet.lean}{49}{Adapted basis of the target staircase lattice}.

\teamanchor{proof.principal-opposite-ordinary}
\subsubsection{The ordinary-place support calculation}
The ordinary-place calculation uses the moved self-dual row
lattices, with the prescribed core lattice and reflection-paired
flag at each vertex. To see the ideal normalizations, write
\(\mathfrak a,\overline{\mathfrak a}\) for the two host ideals
and index the flag by \(i=0,\ldots,c_j-1\).
Before the tilt its line ideals are
\[
 I_{j,i}=\mathfrak a^{c_j-i}\overline{\mathfrak a}^{-(i+1)}.
\]
In an even row the normalized first Witt vectors are
\(\xi_0^i\) times the flag vectors. In an odd row the finite
tilt multiplies the first-half ideals by \(\xi_0^{-i}\).
Both conventions therefore give isotropic coefficient ideal
\(\xi_0^{-i}I_{j,i}\). On the reflected Witt line the
coefficient ideal is its conjugate inverse, since
\(I_{j,i}\overline I_{j,c_j-1-i}=\mathcal O_E\) and the
finite tilt on the reflected lines uses inverse conjugate
multipliers. The even middle line, with Gram \(-1\), has
ideal \((\mathfrak a/\overline{\mathfrak a})^{\rho+1}\)
if \(c_j=2\rho+1\). At either parity the normalized neighboring
isotropic ideals have quotient
\((\xi_0)\mathfrak a\overline{\mathfrak a}\),
the annihilator of the integral scalars for the trace
character on the completion. This is the primitivity of
the adjacent entries. The relation
\(\mathfrak a I_{j+1,i}=I_{j,i}\) gives the equality of
coefficient ideals under the marked column map after including
the Hom-lattice scale \(\mathfrak a\).
For Bessel to Fourier--Jacobi the column map uses unit
first flag vectors at the target. In the other direction
it uses the vectors scaled by \(\xi_0^i\), ending with the
marked vector \(\xi_0^\rho\) in the middle line.
\releasecite{Development/AbelianVarietyTheory/Definitions/Divisors/CmHostImaginaryGeneratorNeZero.lean}{133}{The flag ideals}
\releasecite{Development/AbelianVarietyTheory/Proofs/TwistedOperators/TiltSc.lean}{25}{The moved flag and core lattices}
\releasecite{Development/AbelianVarietyTheory/Proofs/TwistedOperators/BesselToJacobiBasePointCarriesRowLatticeLines.lean}{43}{Transport of the first flag ideals}
\releasecite{Development/AbelianVarietyTheory/Proofs/TwistedOperators/JacobiToBesselCarriesFlagLinesRowPeriodSucc.lean}{36}{Transport including the marked vector}.

The same ideal check gives primitivity of the last
Bessel-to-Fourier--Jacobi covector. Choose a local
generator \(g\) of the middle line ideal. The corresponding
norm in the marked coordinate is
\(-\overline{\xi_0}^{\,\rho}\xi_0^\rho\).
The two fractional factors in the pairing with the last
reflected flag are \(g\xi_0^{-\rho}\) and
\(g^{-1}\overline{\xi_0}^{-(\rho-1)}\). Their product
with this norm is \(\xi_0\), using
\(\overline{\xi_0}=-\xi_0\). Applying the Hom-lattice scale
therefore recovers the trace dual just computed. Thus the
last covector has integral scalar stabilizer. The split
calculation uses both summands of \(E\otimes_F F_v\);
at inert places the trace pairing is on the unramified
quadratic extension.
\releasecite{Development/AbelianVarietyTheory/Proofs/TwistedOperators/BesselLastVectorMemHomLatticeIffIntegral.lean}{111}{Terminal covector with its line and Hom-lattice factors}.

Here is how the support implications follow from the
local mixed-coordinate laws. Iwasawa for the self-dual
Witt lattice writes the target point as a row unipotent,
a diagonal Witt element, a core element and a compact
factor. This decomposition follows by extending primitive
isotropic frames to Witt bases of the lattice, splitting
off the paired vectors successively. Strip off the compact
factor by the unramified vacuum property and the unipotent
by the equivariant character. For the first arrow the core
factor is allowed in the row Levi, so only the diagonal
\(t=\operatorname{diag}(t_i)\) remains. Multiply the vacuum
by the character in the dual coordinate to the \(i\)-th marked
column. For a lattice translate the vacuum is fixed. After
conjugating by the torus the diagonal pairing is multiplied by
\(t_i^{-1}\), because the marked columns are \(X_0U\) for
upper unipotent \(U\) and \(U_{ii}=1\). Nonvanishing gives
\(t_i^{-1}\) in the integral scalars by trace duality and the
column ideal equality. Applying instead the elementary
transvection in position \((i,i+1)\) of the Witt flag gives
the pairing \(a\,t_i/t_{i+1}\) with \(a\) in the above
annihilator ideal. Its character is one, since the conjugated
transvection fixes the vacuum. Hence \(t_i/t_{i+1}\) is integral.
The last covector calculation in the previous paragraph
gives \(t_\rho\) integral by testing the core-to-Witt
Heisenberg translations. These conditions give all \(t_i\)
and \(t_i^{-1}\) integral. This places the diagonal in
the original stabilizer.
\releasecite{Development/AbelianVarietyTheory/Proofs/DecoratedLatticeProofs/InvDiagTorusIntegralBesselJacobiUnramified.lean}{138}{Dual diagonal tests}
\releasecite{Development/AbelianVarietyTheory/Proofs/DecoratedLatticeProofs/TorusEntryRatioIntegralAtUnramifiedPlace.lean}{202}{Simple-root tests}
\releasecite{Development/AbelianVarietyTheory/Proofs/DecoratedLatticeProofs/LastDiagonalEntryTorusIntegralOfFlatMapNeZero.lean}{277}{Last covector test}.

For the second arrow \(g^{-1}X\) has diagonal Witt
components \(t_i^{-1}\) and last component
\(\gamma^{-1}\zeta\) in the core, where \(\gamma\) is
the remaining core element. Dual coordinate tests on the
individual Witt components give \(t_i^{-1}\) integral.
The simple-root translation of column \(i+1\) by
\(t_i a\) times column \(i\) is a lattice translate
when \(a\) belongs to the same annihilator ideal.
Nonvanishing forces its character to be one, hence \(t_i\)
integral. The last dual test, now projected to the core,
gives \(\gamma^{-1}\zeta\) in the normalized marked-line
lattice. More precisely it gives
\(\gamma^{-1}(a\zeta)\) integral for every \(a\) for which
\(a\zeta\) was integral. Take a generator of this ideal.
The two integral marked vectors have the same unit norm;
their orthogonal complements in the self-dual core lattice
are again self-dual. Over an unramified quadratic extension
of odd residue characteristic these complements are
isometric, by diagonalization and surjectivity of the norm
on units. In the split case the vector and corresponding
covector form a direct summand with unit pairing, so an
integral change of bases on the two halves gives the same
orbit. This proves that \(\gamma\) is a core element fixing
the marked line times an integral stabilizer. Nonvanishing
at the second arrow thus forces the Bessel row orbit.
At an inert place the lattice Heisenberg tests use \(2\)
a unit. At a split place use each isotropic half separately;
their characters multiply to the full trace character.
These tests, and the twisting away from the places where
the test tensor has been changed, give the two conditions
used in the ordinary-place induction.
\releasecite{Development/AbelianVarietyTheory/Proofs/DecoratedLatticeProofs/JacobiToBesselFlatMapSupportProof.lean}{476}{Shift and last-line tests on the Fourier--Jacobi-to-Bessel arrow}
\releasecite{Development/AbelianVarietyTheory/Proofs/DecoratedLatticeProofs/UnitNormVectorsOneOrbitSelfDualLatticeSplit.lean}{24}{Integral orbit of the marked vector}.

\paragraph{Combining the finite places.}
Outside \(S\) the flip and the modification of the global
reference vector act through the split idempotents supported
in \(S\). They commute with Heisenberg translations at the
other finite places, so the local factors there are the
transported self-dual lattice vacua. For a lattice isometry
between the old and moved Hom lattices the identity
\(\omega(g)H(z)\omega(g)^{-1}=H(gz)\) carries the defining
translation and opposite-half multiplication invariances to
the moved lattice. In particular the Heisenberg tests in the
ordinary-place support calculation apply to the actual two
flat maps. Their dual diagonal factors are obtained by
multiplication inside the first mixed Fourier integral, their
isotropic shifts inside the second mixed integral, and their
unipotent covariance by translation of the flag integral in
these same moved lattices. At nonsplit places these formulas
use the Hermitian trace pairing of the lattice and its trace
dual. At the ordinary split places they use the two inverse
dual components. No variation of the finite cell label enters
these identities.

The local factors combine by multilinearity because their
Heisenberg operators at different places commute. More
explicitly, a finite tensor relation of functions is detected
by evaluations at tuples of points. For each finite family of
functions, finitely many evaluations span the dual of its
linear span, by successive selection of a point where a
nonzero function in the remaining kernel is nonzero. Hence
applying the descended local maps to the factors preserves the
relation. This applies at every paired position including the
last one, and then to the finite tensor expansion over the
places. At \(v\in S\) the factors \(D_b(\epsilon g)\) are
included in the placed vectors above. In comparison with the
symbols for the cells, the ratio is
\(\prod_bD_b(\epsilon)\), times the removed determinant
character at \(g=g_v^0\) at a fixed staircase place. The
two sources of this ratio are the sign assignment and whether
the cell test retains the character of its representative.
All absolute-value factors at the sign assignment are one.
This recovers \(Z_{\rm loc}^{-1}\) after taking out the fixed
untwisted product, so the normalization is the one used in
the phase calculation. The common reference bi-vacuum,
the moved definite vectors, and the Haar changes have the
factors \(Z_{\rm base}\), \(Z_{\rm move}\), and
\(\prod_v d_{\beta v}\) computed there. The ordinary exceptional-place factors give the prescribed
\(s_{\rm AA}\); the other ordinary places use the
unramified vacuum. Substitution in the convergent unfolded
integral and application of Fubini give the finite factors of
the first line of \eqref{eq:period-principal-unfolding}, with its
real tensor determined in the normalized class pairing below.
\releasecite{Development/AbelianVarietyTheory/Proofs/TwistedOperators/SingleCosetOffTwistPlacesLargeDepth.lean}{164}{The actual translated vacua off the selected places}
\releasecite{Development/AbelianVarietyTheory/Proofs/TwistedOperators/ExistsSignFlatDictionaryAtOnePlaceLargeDepth.lean}{120}{Assembly over pairs at one place}
\releasecite{Development/AbelianVarietyTheory/Proofs/TwistedOperators/FlatDictionaryRelativeBaseVacuaLargeDepth.lean}{28}{Assembly with the reference vacuum}.

\paragraph{The exceptional-place factor.}
At an exceptional prime \(w\) of \(E\), write \(q_w\) for its
residue cardinality and \(a,a'\) for the valuations of the
two host ideals \(\mathfrak a,\overline{\mathfrak a}\). Put
\(e_j=0\) on the core coordinates and
\(e_j(i)=a(c_j-i)-a'(i+1)\) on the flags indexed from zero.
These are the exponents of the diagonal representative of
the original lattice, relative to the integral matrix frame.
The undisturbed finite test on each arrow is the indicator
of matrices in \(\varpi^a M(\mathcal O_w)\). In the moved
box at the origin, a coordinate in row \(h\) and column \(k\)
of the downward matrix at \(j\) is in
\(\varpi^{b_{j,hk}}\mathcal O_w\), where
\[
 b_{j,hk}=a-e_j(k)+e_{j+1}(h).
\]
The adjacent flag exponents differ by \(a+a'\) on a tested
superdiagonal position. The sum of the exponents of each
paired retained row and column coordinate is also \(a+a'\).
The fixed terminal coordinate is zero at the origin.
Therefore the Fourier phase is trivial on the box, by the
\(\varpi^{a+a'}\mathcal O_w\) conductor. Write \(B\) for the
sum of \(b_{j,hk}\) over the integrated free and retained
coordinates. The box contributes volume \(q_w^{-B}\).

Let \(\zeta_X\) be the sum of the diagonal exponents at a
last-column vertex \(X\) and \(0\) elsewhere. Put \(r_X=\dim X\).
For an arrow \(b:s\to t\) on the opposite path the two
inverse-transpose half-densities give
\(q_w^{(\zeta_s r_t-\zeta_t r_s)/2}\). The cell weight is
\(q_w^{M_w}\), where \(M_w\) is the sum of \(e_j(h)-e_j(k)\)
over entries \((h,k)\) of the integrated unipotent radicals.
Indeed conjugation multiplies each such entry by
\(\varpi^{e_j(h)-e_j(k)}\), and the weight is the inverse volume.
Writing the vacuum in the reference trace-dual frame adds
\(\prod_{b:s\to t}q_w^{(a-a')r_s r_t/4}\).
Thus the exceptional factor in the unfolding is
\[
 s_w=q_w^{\,M_w-B+
         \sum_{b:s\to t}(\zeta_s r_t-\zeta_t r_s)/2+
         (a-a')\sum_{b:s\to t}r_sr_t/4}.
\]
Successive ranks have opposite parity, so this is positive
and is an integral power of \(\sqrt{q_w}\).
Their product is \(s_{\rm AA}\), which is nonzero and
belongs to the field with these square roots.
\releasecite{Development/AbelianVarietyTheory/Definitions/Divisors/InvVolConjLevelOne.lean}{129}{The exceptional representatives and the trace-dual correction}
\releasecite{Development/AbelianVarietyTheory/Proofs/DivisorProofs/OriginDatumValueHalfIntegerPowPlaceAA.lean}{24}{The box with trivial Fourier phase}
\releasecite{Development/AbelianVarietyTheory/Proofs/DivisorProofs/UnfoldingTermCoefEqResidueCardZpow.lean}{153}{The inverse unipotent volume}
\releasecite{Development/AbelianVarietyTheory/Proofs/DivisorProofs/FlatIntegrableAaVacOfConductor.lean}{197}{The positive normalized scalar}.


For the test, put

\[
 V_{\boldsymbol\mu,a}(x)=s_{\rm AA}\sum_{\beta}
      \left(\prod_v d_{\beta v}\right)J_{\boldsymbol\mu,a,\beta}(x).
\]
For the coefficient field \(K\) with chosen embedding \(\iota_K\),
the field condition on the test includes the values of
\(\xi_{H,\boldsymbol\mu}\). They have common finite order for each
tuple, so \(\gamma_{\boldsymbol\mu}\gamma'_{\boldsymbol\mu}\)
is \(\iota_K(e_{\boldsymbol\mu})\) for a root of unity
\(e_{\boldsymbol\mu}\in K^\times\). Substituting the expansion and
jet identity into the path formula gives
\begin{equation}
\label{eq:period-principal-test}
 \iota_K(e_{\boldsymbol\mu})c^\square_{\boldsymbol\mu}\,
 \psi_a(\Theta_{\boldsymbol\mu}^{\downarrow}j_{\boldsymbol\mu}^{\rm top}x)
     =\kappa V_{\boldsymbol\mu,a}(x).
\end{equation}
The product of the square constants appears separately from the
root of unity. The scalar \(\kappa\) is nonzero as soon as one test
has nonzero left side.

\subsubsection{Rationality of the finite sums}
In the coefficient model the column and row inputs at a class
representative are polynomials over \(K\). Transport by a rational
unitary matrix through the specified field embeddings is polynomial
substitution and multiplication by a determinant power.
Place type factors are character values. A chain of normalized
coefficient matrices over \(K\) preserves this polynomial model.
At the top, the coefficient restriction in codimension one is
the algebraic projection onto the required branching summand.
One can take its adjoint projection in the polynomial model
using the monomial pairing, with factorial diagonal entries.
Normalizing by the product of the corresponding extremal-vector
pairings divides by a nonzero element of \(K\). Thus the resulting
bilinear contraction is over \(K\).

For the finite test coefficients at a split place, taking the
Fourier sums preserves the field containing the occurring phases.
The complementary integral also preserves it when it contains
\(\sqrt{q_v}\) and the occurring determinant character values.
Here is a direct calculation. With additive matrix measure the
weight on \(\mathrm{GL}_\ell(F_v)\) is a rational measure
normalization times \(\nu(\det g)|\det g|^{n/2-\ell}\).
For the old rank \(\ell=a'\) the exponent \(s=n/2-\ell\) is
nonnegative. Composition of the test with the injective
full-rank frame map gives a compactly supported locally constant
function of the matrix \(g\). After rescaling, it is a finite
linear combination of integral matrix coset indicators.
In a coset \(D+\varpi^e M_\ell(\mathcal O_v)\), integral
elementary divisors allow \(D\) diagonal, with valuations
less than \(e\) in its nonzero positions. The block over those
positions is invertible of constant determinant valuation.
The other block in the determinant is its Schur complement,
\(\varpi^e\) times an unrestricted integral matrix translated
by an integral matrix. The character of the first block is
computed on a finite congruence quotient.
For the unrestricted \(j\)-by-\(j\) matrix integral, the ramified
character case with \(j>0\) gives zero by multiplying a column
by a unit. For an unramified character, decompose the first
column by the common valuation \(k\) of its entries and take
its primitive part to the first basis vector by an integral
change of rows. Expanding the determinant along that column
gives the recursion
\[
 I_j=\frac{1-q_v^{-j}}{1-\nu(\varpi)q_v^{-(s+j)}}I_{j-1},
 \qquad I_0=1.
\]
All factors are in the field. This argument applies at every
split swap; the final Fourier inversion again involves a
finite sum.

In the two definite paths of this test, the field calculation
uses a further normalization of the kernel entries. In the
dictionary for an arrow \(p\), the raw derivative table at a
matched pair of points has the form
\[
 J^{\rm raw}_p(\alpha)
   =c_{0,p}\prod_{\tau}\omega_{p,\tau}^{|\alpha_\tau|}
                  J^{\rm alg}_p(\alpha).
\]
The table \(J^{\rm alg}_p\) is obtained by differentiating the
finite sum of reduced torsion translates of the moved vacuum.
For the reversed-arrow convention the complex factors are
conjugated. The \(\omega_{p,\tau}\) are the chosen nonzero
CM periods, with that convention included. The symbol
\(|\alpha_\tau|\) is the derivative degree at \(\tau\).
In the coefficient operator the source determinant powers
and the substitution of frame variables give a single degree
\(d_{p,\tau}\) at each \(\tau\), as in the polynomial
calculation below. Hence the scalar
\(c_{0,p}\prod_\tau\omega_{p,\tau}^{d_{p,\tau}}\)
passes through the entire operator. Its reciprocal is the
normalization of the kernel matrix, so the normalized entries
are the entries computed from \(J^{\rm alg}_p\).
\teamanchor{proof.principal-dictionary-field}
\paragraph{The field for the actual finite dictionary.}
Use each occurring pair of finite twist character and marked
ramified set when choosing the coefficient field. For such a
pair the split tensors, residue levels and depths are fixed
before choosing the auxiliary class realization. For each
test and tuple the support of its finite unfolding consists
of finitely many row cells. At the fixed staircase places
there is just the prescribed cell. At the other split places
one can see finiteness directly from the free-coordinate
translations. Suppose the previous representative and its
inverse have bounded entries. At a point where the cell
integrand is nonzero, an affine matrix \(A\) in the two
flat symbols has \(h^{-1}A\) bounded, where \(h\) is the next
row representative. This follows from integral support
after multiplication by the previous representative.
Write \(h=bk\) by block Iwasawa, where \(k\) is integral
invertible and \(b\) is upper triangular in flag order with
a middle block. Put the core before the central flag when
these have the same row-group weight.

For each flag index \(i\), multiply an elementary free
increment \(zE_{i,i+1}\) in the flag block by \(b\).
Its nonzero column has entries only in the free positions
by triangularity, and \(h^{-1}\) carries it to
\(k^{-1}zE_{i,i+1}\). The fixed test period and the bound on
the previous representative give one nonzero ideal of
values of \(z\) for which this preserves the test. The
phase change is \(\psi_v(zb_{ii})\), since the paired
retained entries stay fixed. Nonvanishing forces this
phase to be one for every \(z\) in that ideal. The
conductor hence bounds the flag diagonals \(b_{ii}\).
In the row Levi component, the inverse flag columns are
bounded by those of \(b^{-1}A\), because the affine matrix
has flag pivots \(1\) and zero entries below those pivots.
This holds also for the central block when a flag shares
the core weight. In that case the Levi block has the form
\(\left(\begin{smallmatrix}C&d\\0&a\end{smallmatrix}\right)\).
Its inverse flag column bounds \(a^{-1}\) and
\(C^{-1}d\,a^{-1}\). After removing \(C\) by the row
Levi action, the bound on \(a\) bounds \(C^{-1}d\).
Removing the flag radical now leaves a matrix and its
inverse in bounded sets. The set of matrices with both
bounds is compact. An open right level has finitely
many cosets intersecting this compact set. Induction
from the fixed initial staircase proves the required
finite support, at every selected place.
\releasecite{Development/AbelianVarietyTheory/Proofs/DivisorProofs/SumOKEx.lean}{199}{Support and depth after the split changes of path}
\releasecite{Development/AbelianVarietyTheory/Proofs/DivisorProofs/FiniteUnfSupportOfNeOne.lean}{67}{Bounding the flag and the middle block}
\releasecite{Development/AbelianVarietyTheory/Proofs/DivisorProofs/FiniteUnfSupportOfNeOne.lean}{288}{Depth-preserving increments in the actual affine slice}
\releasecite{Development/AbelianVarietyTheory/Proofs/DivisorProofs/FiniteUnfSupportOfNeOne.lean}{375}{Finite induction on the representatives}.

For a supported cell take the prescribed expansion by
cosets at its least positive depth of expansion. At fixed
depth the products of the coset indicators on the split
halves are disjoint. The compact support of the symbol
therefore bounds the possible nonzero products by a
finite quotient on each of those halves. The constant
coefficient of one product is the value of the local
symbol at its representative. These coefficients are
over \(K\) by the Fourier and matrix integrals above,
with the specified characters and phases.
The class spaces at the term levels are finite because
the groups there are definite. The terminal column
representative is the image of the last row representative,
so it needs no new choice. For these finitely many classes,
pure summands and coefficient indices, take the normalized
matrix entries on both paths. The derivatives after pulling
back the self-dual Hom lattices are those of the finite
sums of torsion translates used below in
\S\ref{sec:period-normalized-entry}. The Gaussian
normalization removes the same fixed CM period powers
on every monomial of the indicated homogeneous weight.
The fourfold theta transport there gives algebraicity
of each normalized entry, with the transpose and
conjugation on the row path. Adjoining these finite
lists, the character values and the required square
roots of residue cardinalities yields the specified
coefficient field on the representative family.
\releasecite{Development/AbelianVarietyTheory/Proofs/TwistedOperators/FinitePureDataOfCosetIndicatorExpansionsDepth.lean}{95}{Finiteness at the prescribed depth}
\releasecite{Development/AbelianVarietyTheory/Proofs/TwistedOperators/NormalisedDefiniteKernelEntriesIsAlgebraic.lean}{68}{The normalized class entries}
\releasecite{Development/AbelianVarietyTheory/Proofs/TwistedOperators/NormalFormJets.lean}{43}{Adjoining the finite lists on the representative family}.

In passing to the selected grid the corresponding tuples
have the same twist and marked ramification. Carry the
definite spaces, lattices, local frames, rational coefficient
frames and finite moving elements by their specified
identifications. The initial row corner is \(X'\), with
the distinguished lattice in its identity frame on both
sides. The terminal embedding is the same orthogonal-summand
embedding. This carry also uses the same normalized
vacua, constants \(c_{0,p}\) and period parameter. With
these choices the split transfer gives the same tensor
in the same local matrix arguments. Hence taking its
flat Fourier symbols gives identical functions, with the
same supported cells, levels and least depths. In particular
the finite coefficient expansions can be compared in the
same quotients. At the class points the two derivative
tables then evaluate the same Heisenberg translates of
the same normalized vacuum, with the same two finite
moves and real monomials. The arithmetic entries in the
selected grid are thus the previously adjoined numbers.
This gives the values of \(V_{\boldsymbol\mu,a}(x)\) in \(\iota_K(K)\) in
\eqref{eq:period-principal-test}.
\releasecite{Development/AbelianVarietyTheory/Proofs/TwistedOperators/CarrySecond.lean}{260}{Transport of frames and finite moves}
\releasecite{Development/AbelianVarietyTheory/Proofs/TwistedOperators/CarrySecond.lean}{321}{Transport of the normalized vacuum}
\releasecite{Development/AbelianVarietyTheory/Proofs/TwistedOperators/TrJetOfTrFlip.lean}{449}{Transport of the support and both series of entries}
\releasecite{Development/AbelianVarietyTheory/Proofs/TwistedOperators/RowTestRationalSumsOfFiniteSumsMeaningful.lean}{51}{Field membership on the selected tests}.


\teamanchor{proof.principal-opposite-pairing}
\subsubsection{The normalized definite class pairing}
\paragraph{The coefficients at the initial points.}
Write \(C_j\) and \(R_i\) for the definite vertices of the
ascending column and row, respectively. In this paragraph a
finite point \(X_j\) is on \(C_j\), and \(Y_i\) is on \(R_i\).
At the terminal vertex put \(X_m=\iota(Y_m)\), using the
stabilizer embedding adjoining the identity on the line.
The two initial coefficient functions are
\[
 \begin{split}
 u_0(X_0)&=\rho^{\rm fr}_{C_0}{}^{-1}
                      \Psi_a^{\,\mathbb C}(X_0),\\
 v_0(Y_0)&=\rho^{\rm fr}_{R_0}{}^{-1}
                    \operatorname{read}\bigl(A_{\boldsymbol\mu}
                           (j_{\boldsymbol\mu}^{\rm top}x,Y_0)\bigr).
 \end{split}
\]
Here \(\rho^{\rm fr}\) is the prescribed coefficient change to
the orthonormal real frame, and \(\operatorname{read}\) takes the
listed row weight to its dual coefficient convention.
The column coefficient is the image of the given polynomial
\(\Psi_a\) at the finite point. The row coefficient
\(A_{\boldsymbol\mu}(y,Y_0)\) is the evaluation of the
coefficient section \(y\) of \(X'\) at the image of \(Y_0\)
under the fixed identification \(R_0\simeq X'\), followed by the
coefficient comparison. This comparison extends the coefficient
embedding by polynomial substitution and renames the indices
under the rank identification. Thus \(A_{\boldsymbol\mu}\) is
the coefficient used in the arithmetic sum.

In the automorphic realization of these two functions the
archimedean components enter by the inverse real group actions.
In fact after writing \(g=g_fg_\infty\) the two automorphic values
are
\[
 \Lambda_{C,0}\bigl(\rho_C(g_\infty^{-1})u_0(g_f)\bigr),
 \qquad
 \Lambda_{R,0}\bigl(\rho_R(g_\infty^{-1})v_0(g_f)\bigr).
\]
For the column this follows by evaluating the fixed real
coefficient vector against \(\Psi_a^{\,\mathbb C}\). For the
row, in the polynomial realization the fixed real vector is
paired with the inverse of the same coefficient comparison,
then with the inverse real action at \(g_\infty\). The
comparison is injective because it is scalar extension of the
coefficient space followed by an injective renaming of the
matrix variables. Its inverse on the image therefore yields
exactly the initial row value, including for the original
input \(j_{\boldsymbol\mu}^{\rm top}x\). Finite and real parts
commute, so no finite determinant character is introduced by
this frame change. The two functionals \(\Lambda_{C,0}\) and
\(\Lambda_{R,0}\) are fixed for the family.
\releasecite{Development/AbelianVarietyTheory/Definitions/Divisors/InvVolConjLevelOne.lean}{366}{Initial column and row coefficient evaluations}
\releasecite{Development/AbelianVarietyTheory/Proofs/TwistedOperators/ExistsRealModelSq.lean}{170}{The column coefficient realization}
\releasecite{Development/AbelianVarietyTheory/Proofs/TwistedOperators/ExistsRealModelSq.lean}{264}{The row coefficient realization}
\releasecite{Development/AutomorphicArithmetic/LocalEmbeddings/JetsTest.lean}{34}{Initial real functionals in the two frames}.

\paragraph{Normalized recursion at a pair of finite points.}
For a finite product summand \(d\) let \(J_{p,d}^{\rm raw}(X,Y)\)
be the Fock derivative table for definite arrow \(p\) in its
coefficient frames. If \(a_X,a_Y\) are the two fixed finite
moves from the standard self-dual lattices, the table is read at
\((Xa_X,Ya_Y)\) on the original vacuum with the finite
Heisenberg operators of \(d\) applied to it. Indeed those
moves are on the right of the group arguments of the kernel,
and they commute with the real creation operators. At the
opposite-sign row arrow the matrix indices in the Fock
monomial are interchanged and the CM table is conjugated.
These are the derivative table and the sign convention in
the prescribed coefficient operator.

For clarity, let \(a,b\) be the two ranks, and at each complex
embedding let \(n_\tau\) be the \(a\)-th entry of the polynomial
part of the target highest weight, with indices starting at
one. Given a source coefficient polynomial \(u(H)\), the
coefficient operator first multiplies it by
\(\prod_\tau\det(H_\tau)^{n_\tau}\). It then substitutes
\(H_\tau=Z_\tau^{\mathsf T}T_{\tau,[a]}\), where \(Z_\tau\)
is \(b\) by \(a\) and \(T_{\tau,[a]}\) consists of the first
\(a\) columns of the target polynomial variables. Replace
each \(Z^\alpha\) by its entry in the actual derivative table.
This is the jet operator in the fixed polynomial coefficient
realization. The first \(a\) entries of the target polynomial
weight are the source polynomial weight plus \(n_\tau\),
and the remaining entries are zero. Thus the result is
already in the required target coefficient space, and the
total \(Z_\tau\) degree is the sum \(d_{p,\tau}\) of the target
polynomial weight. This proves directly that division by
\(b_p=c_{0,p}\prod_\tau\omega_{p,\tau}^{d_{p,\tau}}\) yields
the prescribed normalized operator
\(\mathcal J_{p,d}(X,Y)\). Both the real period and its finite
column or row reading use the same polynomial and the same
points. The nonzero \(b_p\) is independent of the finite
points, and on every column arrow it is independent of the
tuple because the normalized reference vacuum and its
periods are fixed before twisting.
\releasecite{Development/AbelianVarietyTheory/Definitions/Divisors/KernelNormConstantSpec.lean}{62}{The row convention}
\releasecite{Development/AbelianVarietyTheory/Definitions/Divisors/KernelNormConstantSpec.lean}{295}{The prespecified normalized matrix}
\releasecite{Formalism/LocalToGlobal/KernelLattices/ThDistC.lean}{23}{The polynomial substitution}
\releasecite{Development/AbelianVarietyTheory/Proofs/TwistedOperators/FockComponentTransportedPointsEqJetOperator.lean}{43}{The actual transported column points}
\releasecite{Development/AbelianVarietyTheory/Proofs/TwistedOperators/LastColumnArrowRealAverageFunctional.lean}{78}{The actual transported row points}.

One can see the real integration giving this same operator
before the finite class sum. At each embedding write
\(k\in U(a)\) and \(h_\infty\in U(b)\) in the two coefficient
frames. The oscillator on the coefficient variables substitutes
\(Z\mapsto h_\infty^{\mathsf T}Zk^{-\mathsf T}\).
If \(\ell_s,\ell_t\) are the last entries of the source and
target weights in that convention, its scalar is
\(\det(k)^{n_\tau-\ell_s}\det(h_\infty)^{\ell_t}\).
In the column frame let \(c\) be the center of the block
inserted below the common entries. The infinitesimal
character is read as highest weight by subtracting the
half-sum of the positive roots. The last target entry is
\(c+a/2\), and each common entry is decreased by
\((b-a)/2\). Hence \(n_\tau-\ell_s=-c-b/2\).
These are the two splitting exponents with the positive
oscillator half-ranks, in their source and target order.
In the reversed row frame let \(c\) instead be the center
of the inserted block before reversing its sign and let
\(\varepsilon(n)\) be the parity of \(n\).
The two row exponents are
\[
 c-\varepsilon(b)/2-b/2,\qquad
 -c+\varepsilon(a)/2+a/2.
\]
The latter is the last target entry. Subtracting the parity
shifts and the difference of the half-ranks in the common
entries gives the former as \(n_\tau-\ell_s\).
The source coefficient action in the integral supplies
\(\det(k)^{\ell_s}\) and substitutes \(H\mapsto kH\).
The two determinants on \(k\) therefore give exactly
the power of \(\det H\) in the jet operator.

For fixed initial real functional \(\Lambda\), express it
as a linear combination of finitely many polynomial evaluations.
These evaluations span its dual since a polynomial vanishing at
all invertible matrices is zero. For one such evaluation at
\(H=H_0\), absorb \(\det(k)^{n_\tau}\) by using
\(f(H)=u(H)\prod_\tau\det(H_\tau)^{n_\tau}\), and divide at
\(H_0\) by the nonzero \(\prod_\tau\det(H_{0,\tau})^{n_\tau}\).
The average of
\(P(Zk^{-\mathsf T}) f(kH_0)\), taken coefficientwise in \(Z\),
is an intertwiner from the source polynomial representation to
the rectangular polynomial space. This follows by changing
variables in Haar, and extends from the compact group to its
complexification. By the Cauchy decomposition, the space of
these intertwiners is an irreducible representation of the target
general linear group at each embedding. It is spanned by the
substitutions \(f(H)\mapsto f(Z^{\mathsf T}T_{[a]})\) at
invertible \(T\). Indeed one such substitution is nonzero by
specializing \(Z\) on its first \(a\) rows, and its target
translates give exactly these substitutions as \(T\) varies.
The same argument for the product of the compact groups gives
the assertion at all embeddings. Applying the finite derivative
functional to this linear combination recovers the jet
operator evaluated by a new real functional. The target
determinant and substitution are its inverse real action.
Multiplying the new functional by \(b_p\) passes from the
raw to the normalized finite operator.
\releasecite{Development/AbelianVarietyTheory/Proofs/DecoratedLatticeProofs/MvZOneEqMapUnivSubst.lean}{90}{The compact average of the rectangular substitution}
\releasecite{Development/AbelianVarietyTheory/Proofs/TwistedOperators/ColumnArrowKernelRealEqDetCharMulFockKernel.lean}{126}{The column exponents}
\releasecite{Development/AbelianVarietyTheory/Proofs/TwistedOperators/DenseOfPolynomialStdGaussianMemSchwartz.lean}{26}{The column averaging}
\releasecite{Development/AbelianVarietyTheory/Proofs/TwistedOperators/CompactPairRealAverageFockEqJetOperator.lean}{317}{The row exponents}.

Thus, with probability Haar on the real source group, at
each arrow the exact recursion for a pure summand and a
fixed real polynomial pair \(P\) is
\[
 \begin{split}
 &\int_{G_{p,s}(F_\infty)}
    K_{p,d,P}(Xk,h)\,
             \Lambda_{p,s}\big(\rho_{p,s}(k^{-1})u\big)\,dk\\
 &\hspace{2em}=
 \Lambda_{p,t}\big(\rho_{p,t}(h_\infty^{-1})
                  \mathcal J_{p,d}(X,h_f)u\big).
 \end{split}
\]
The new functional depends on the real polynomial, the preceding
functional and the fixed normalized vacuum, uniformly in the
finite input. Iterate along both paths, starting with
\(u_0(X_0)\) and \(v_0(Y_0)\). The terminal average over
\(G_{R_m}(F_\infty)\), where \(X_m=\iota(Y_m)\), is the
bilinear invariant pairing of the last two representations.
The dual type of \(R_m\) occurs with multiplicity one in the
restriction of the \(C_m\)-representation. Averaging the two
terminal functionals with probability Haar therefore gives
\(\kappa(P)\mathcal B\), where \(\mathcal B\) is the specified
algebraic branching contraction. Here \(\kappa(P)\) depends
only on \(P\) and the fixed real data.
\releasecite{Development/AutomorphicArithmetic/LocalEmbeddings/JetsGrid.lean}{152}{The column and row recursions at matched points}
\releasecite{Development/AutomorphicArithmetic/ThetaCorrespondence/ThetaPairLifts/LocalTypeDictionary/PairingOfPointsExpansion.lean}{18}{The last compact average}.

\paragraph{The prescribed term levels and their class weights.}
First refine all the compact open levels used in a cell to
levels fixing each factor in its finite expansion. Such levels
exist by local constancy of the input coefficients and
Schwartz vectors, since the expansion is finite. On this
refinement the real-average recursion gives a finite sum
with weight \(\operatorname{vol}(K')/
|G(F)\cap X K' X^{-1}|\) for each class \(X\) at each
integrated vertex. Probability Haar in the real variable is
the one just used in the recursion. The fine-level calculation
also proves invariance under left rational changes of
representative, by moving the real component of such an
element through the compact average.

To descend to the prescribed levels use the entire finite
expansion of the cell. At an interior column vertex a
prescribed level element is the middle projection of a
stabilizer of the chosen row double coset. Lift it to that
stabilizer and change the row variable by its unipotent
component. The flag characters from the two
adjoining flat kernels cancel by the paired translations of the
retained row and column. To check the determinant in this
substitution, write \(d=\det k_v\) on the split component.
The lift and the projected core matrix have the same
determinant \(d\), because the flag lift is unipotent and
conjugating by the cell representative preserves the
determinant. It is a unit. The two half-densities are
one and the parity factors from the same vertex occur
in opposite powers. Thus the path factor at an interior
column vertex \(j\) is
\(\mu_{j,v}(d)/\mu_{j+1,v}(d)\). The same factor is the
product of the scalars on the two translated core vacua.
At the initial column vertex only the factor
\(\mu_{1,v}(d)^{-1}\) occurs, and \(K^+\) fixes \(\Psi_a\).
At an interior row vertex \(i\) the two edges give
\(\chi_{m-i,v}(d)/\chi_{m-i+1,v}(d)\), again the product
of the two vacuum factors. The typed path gives invariance
with precisely these path factors, which persist through
the square changes of variables. At the last row vertex
make the same change also in the column vertex by \(\iota\).
The orthogonal complement is the fixed line, so this
preserves the terminal Fourier argument and branching
contraction. The last rightward and downward factors there
give \(\mu_{m,v}(d)/\chi_{1,v}(d)\). In each case the
product of vacuum factors replaces the determinant normally
outside the substituted test arguments. Hence the unchanged
full tensor equality for the placed path becomes the
unchanged sum of composed maps. This uses linearity of the
Schr\"odinger realization and of the normalized operator in
its derivative table. Away from the selected split places
the level stabilizes the reference lattice and the vacuum
factors on units are trivial.
\releasecite{Development/AbelianVarietyTheory/Proofs/TwistedOperators/UnfoldingExpansionProductFunctionLocalInvariant.lean}{42}{Core lift and the interior column determinant factors}
\releasecite{Development/AbelianVarietyTheory/Proofs/TwistedOperators/ExpansionProductInvariantInnerVertexLevel.lean}{17}{Interior projected column level}
\releasecite{Development/AbelianVarietyTheory/Proofs/TwistedOperators/UnfoldingExpansionProductFunctionFixedByKplus.lean}{219}{Initial column level}
\releasecite{Development/AbelianVarietyTheory/Proofs/TwistedOperators/ExpansionTensorSumInvariantGridRowLevel.lean}{14}{Interior row level}
\releasecite{Development/AbelianVarietyTheory/Proofs/TwistedOperators/ExpansionSumProdLastVertexMovedEq.lean}{113}{The two terminal determinant factors}
\releasecite{Development/AbelianVarietyTheory/Proofs/TwistedOperators/ProductFunctionTotalTensorInvariantLastArrows.lean}{143}{Joint terminal level}.

At the first row vertex the comparison has a character.
For \(k\) in the prescribed first row level let \(k_v\)
be its split coordinate. This level lies in the actual
\(X'\) level, in the frame used for the initial row
coefficient. Its character on that coefficient is
\(\chi_{\boldsymbol\mu}(k)=\prod_v A_v(k_v)\), where
\[
 A_v(k_v)=
 \begin{cases}
  \displaystyle\prod_{c=1}^m \chi_{c,v}(\det\bar k_{v,c}),
               &v\text{ a type component},\\
  (\xi_{\boldsymbol\mu,v}^{R})^{-1}
                      (\det\bar k_{v,\mathrm{tail}}),
               &v\text{ a tuple component},\\
  1,&\text{otherwise}.
 \end{cases}
\]
Here the type product omits the initial \(r\)-block, and
the tail block at a tuple place has size \(mt\). These
values are the diagonal-block determinants in the same
residue matrices, since the first row frame is the
prescribed \(X'\)-frame. The original last downward
step test at this vertex transforms under the local
path action with eigenvalue \(A_v(k_v)^{-1}\).
This follows by multiplying the block diagonal
characters and canceling the determinant weights
in the split Weil action on units. The vertex is
exterior to every square being changed. The same
local eigenvalue therefore persists in the opposite
tensor by performing every square change with the
exterior substitution held fixed.

More explicitly, let \(D^{\rm first}_v(k_v)\) be the
determinant factor of the first rightward edge
in its split frames. Its half-densities are one
on this level. Replacing \(Y_0\) by \(Y_0k\) in
the raw function arguments removes
\(\prod_vD^{\rm first}_v(k_v)\) from the path
action. In the standard oscillator model, this
same change is by \(a_{Y_0}^{-1}ka_{Y_0}\)
on the vacuum and its Heisenberg translates, where
\(a_{Y_0}\) is the fixed first row moving element.
The scalar on the vacuum is
\(\prod_vD^{\rm first}_v(k_v)\): its determinant in
the moved frame is still \(\det k_v\). Away from the selected
split components the unit determinant factors are trivial.
The trace-dual lattice
condition guarantees preservation of the vacuum
line. Applying the same argument substitution to
all factors in the expansion now gives the multiplier
\(\prod_v A_v(k_v)^{-1}\) for the full normalized
kernel at \(Y_0k\). The frames and degree normalizations
are unaffected by \(k\). Multiplying by the initial
row coefficient, which has multiplier
\(\chi_{\boldsymbol\mu}(k)\), gives invariance of the
paired summand on the prescribed \(Y_0\)-class.
\releasecite{Development/AbelianVarietyTheory/Proofs/TwistedOperators/FlippedTensorEigenvalueLaws.lean}{122}{Original and transferred first row level character}
\releasecite{Development/AbelianVarietyTheory/Proofs/TwistedOperators/FirstVertexWeilOrbitValue.lean}{381}{Block determinant and vacuum scalar comparison}
\releasecite{Development/AbelianVarietyTheory/Proofs/TwistedOperators/RowShift0Model.lean}{87}{The moved first row vacuum}
\releasecite{Development/AutomorphicArithmetic/ThetaCorrespondence/ThetaPairLifts/LocalTypeDictionary/FirstRowInvariantUnderTwist.lean}{18}{Cancellation on the paired class value}.

A coarse class above \(X\), on passage to \(K'\subseteq K\),
has fine classes indexed by the orbits of
\(G(F)\cap X K X^{-1}\) acting on \(K/K'\).
The sum of the reciprocal stabilizer orders is
\([K:K']/|G(F)\cap X K X^{-1}|\) by orbit--stabilizer.
Multiplying by \(\operatorname{vol}(K')\) and using
\(\operatorname{vol}(K)=1\) gives the desired weight
at the actual term level. Thus the full class sum is
\[
 \kappa(P)\sum_{X,Y}
 \frac{\displaystyle\sum_d c_d\,\mathcal B
       \bigl(\mathcal J_{C,d}u_0(X_0),\,
             \mathcal J_{R,d}v_0(Y_0)\bigr)}
 {\displaystyle\prod_{j=0}^{m-1}|\Gamma_{X_j}|\,
                \prod_{i=0}^{m}|\Gamma_{Y_i}|}.
\]
Here \(\mathcal J_{C,d}\) and \(\mathcal J_{R,d}\) are the
composed normalized arrow operators at the indicated points.
The terminal column point is \(\iota(Y_m)\), so it contributes
no independent class mass. The two chains, initial coefficients
and stabilizer products in this expression are exactly those in
\(J_{\boldsymbol\mu,a,\beta}(x)\).
The finite expansion enters only through its total tensor:
the same evaluation argument as for the local contractions
passes linear relations through the composed operators. It
therefore gives the value for the prescribed arithmetic
expansion as well as for any refinement used in the integrals.
\releasecite{Development/AbelianVarietyTheory/Definitions/Divisors/InvVolConjLevelOne.lean}{276}{Term class representatives and weights}
\releasecite{Development/AutomorphicArithmetic/LocalEmbeddings/JetsGlue.lean}{240}{Regrouping the fine classes}
\releasecite{Development/AutomorphicArithmetic/ThetaCorrespondence/ThetaPairLifts/LocalTypeDictionary/PairingOfPointsExpansion.lean}{55}{The paired class expansion}.

The preceding class calculation shows that every term
as a functional of \(P\) is proportional to \(\kappa(P)\).
Use this as the common functional \(\ell\) when it is
nonzero. If it vanishes identically, choose any nonzero
\(\ell\), since in that case the multiple is zero in every
term. Choose a polynomial tensor \(\sigma_0\) with
\(\ell(\sigma_0)=1\). At a real
quotient point and for a real input tensor \(P\), the
placed comparison gives a finite definite Fock combination
\(\sigma_{P}(p)\). Because the reference bi-vacuum and both
initial real types are fixed, its functional value
\(\ell(\sigma_P(p))\) is independent of the finite inputs.
The compact Fock reduction projects to the central weights
of the two fixed real types. Their scalar circle weights fix
the degree at each embedding, so only finitely many real
Fock coefficients enter that functional. Its evaluation
commutes with the finite Heisenberg operators. This explains
the same linear functional for all terms, with the
finite determinants already removed. If some reference
definite period is nonzero, divide the placed identity
for that term by its fixed value. This writes the scalar
as a measurable, integrable function of the real quotient
point by absolute convergence of the unfolded integral.
If all the reference definite periods vanish, every
placed term vanishes and one can take the integrated
coefficient to be zero. Denote the resulting real
quotient integral by \(c(P)\).
If \(\sum_P b_P P\) is the original fixed real tensor,
the representative of the real output is therefore
\[
 \sigma=\frac{\sum_P b_P c(P)}{s_{\rm AA}}\,\sigma_0.
\]
The denominator is the nonzero exceptional-place product
computed above. This is the real combination when \(s_{\rm AA}\) is written explicitly in
the first line of \eqref{eq:period-principal-unfolding}.
Using the same \(\sigma\) in the class calculation replaces
\(\kappa(P)\) by \(\kappa=\sum_P \sigma(P)\kappa(P)\).
Thus the scalar from the real data is common across the tuples,
tests, cells and initial coefficients. This proves
\(A_{\boldsymbol\mu,a,\beta}(x)=\kappa
J_{\boldsymbol\mu,a,\beta}(x)\) in the original term
measures.
\releasecite{Development/AbelianVarietyTheory/Proofs/TwistedOperators/ColumnRealDatumFockCombinationLargeDepth.lean}{58}{Real central characters in the Fock reduction}
\releasecite{Development/AbelianVarietyTheory/Proofs/TwistedOperators/LastRowRealIteratedIntegralEqConstMulFinite.lean}{214}{The common real combination with the exceptional vacuum factor}
\releasecite{Development/AbelianVarietyTheory/Proofs/TwistedOperators/JetsOfOrthogonalSum.lean}{19}{Linearity in the real tensor}.

\teamanchor{proof.principal-square-ratios}%
\subsubsection{Ratios of square constants}
To compare principal-level tests across tuples one needs
\(c_{ij,\boldsymbol\mu}/c_{ij,\boldsymbol\mu_0}\in\iota_K(K^\times)\)
for each square. Write \(\lambda_{\boldsymbol\mu}\) for the
comparison factor on the standard definite triangle whose two
normalized class sums are defined over \(K\).
Section~\ref{sec:period-definite-triangle} proves
\(\lambda_{\boldsymbol\mu}/\lambda_{\boldsymbol\mu_0}
\in\iota_K(K^\times)\) by evaluating a fixed complementary pair of
minor products on that triangle. Equation
\eqref{eq:period-bottom-constant-field} gives the ratio on any of
the prescribed bottom triangles in the cell.

In the comparison of the square with this triangle, put
\(\mathcal R_{\boldsymbol\mu}(s)=
\Lambda_{\boldsymbol\mu}(s)/\Lambda_{\boldsymbol\mu_0}(s)\) for the
specified Hecke \(L\)-values of that square. Put
\(p_j=\lfloor(m-j-1)/2\rfloor\), \(q_j=\lfloor(m-j)/2\rfloor\),
and \(n_0=\lfloor(t+r+6)/2\rfloor\) from the difference of
its two intermediate ranks. The
required identity in the actual normalizations is
\begin{equation}
\label{eq:period-square-ratios}
\begin{split}
 \frac{c_{ij,\boldsymbol\mu}}{c_{ij,\boldsymbol\mu_0}}
 \prod_{k=0}^{q_j-1}\mathcal R_{\boldsymbol\mu}(n_0+k+1)
  ={}& a_{ij,\boldsymbol\mu}
       \frac{\lambda_{\boldsymbol\mu}}{\lambda_{\boldsymbol\mu_0}}
       \prod_{k=0}^{p_j-1}\mathcal R_{\boldsymbol\mu}(n_0+q_j-k).
\end{split}
\end{equation}
where \(a_{ij,\boldsymbol\mu}\in\iota_K(K^\times)\) is the product of
the measure and local step factors. For nonzero \(L\)-values the
common factors cancel. When \(m-j\) is odd, \(p_j=q_j\) and the two
products have the same indices. When \(m-j\) is even, \(p_j=q_j-1\)
and the only uncanceled factor is
\(\mathcal R_{\boldsymbol\mu}(n_0+1)\) on the left. Consequently
only this latter ratio in even cells needs a field calculation.
\releasecite{Development/AbelianVarietyTheory/Proofs/DivisorProofs/SquareRatioInEmbeddingRange.lean}{23}{Triangle and critical-value factors}
\releasecite{Development/AutomorphicArithmetic/ThetaCorrespondence/ThetaPairLifts/ThirdPlaceSeries/SquareRatioAtPointOfLeadingTerm.lean}{19}{Application at the fitted points}.

\subsubsection{The square and the descending triangles}
Keep the square, the tuple family and the coefficient embedding fixed.
For a triangle at stage \((p,q)\), the two source ranks are
\(A_0+2p\) and \(B_0+2q\), where
\(A_0=R_{i+1,j+1}-2p_j\) and \(B_0=R_{ij}-2q_j\).
The negative indices of the two sources and the target are
\(p,q,p+q\), respectively. In particular the half-difference of the
source ranks is \(e(p,q)=n_0+q_j-p_j+p-q\).
The two splitting characters stay fixed along each Witt tower,
since a hyperbolic plane changes each rank by two. Their quotient at
every stage is \(\nu_{ij,\boldsymbol\mu}\).
\releasecite{Development/AbelianVarietyTheory/Proofs/DivisorProofs/NegIdxOfIsStd.lean}{24}{Stage and splitting characters}.

Let \(D_{p,q}(\boldsymbol\mu)\) be the quotient of the triangle's
nonzero comparison constants at \(\boldsymbol\mu\) and \(\boldsymbol\mu_0\).
This quotient is independent of the choice of common polynomial pair
representing the comparison. In fact, if the two matched operators
are \(\lambda_{\boldsymbol\mu} J_{\boldsymbol\mu}p_*\) and
\(J_{\boldsymbol\mu}p'_*\) for fixed polynomial maps and injective
\(J_{\boldsymbol\mu}\), the proportionality of the two values on a
fixed polynomial pair implies
\begin{equation}
   p_*(P)=\frac{c_{\boldsymbol\mu}}{\lambda_{\boldsymbol\mu}}p'_*(P').
   \label{eq:period-triangle-pair-normalization}
\end{equation}
For \(p'_*(P')\ne0\) this forces \(c_{\boldsymbol\mu}/\lambda_{\boldsymbol\mu}\)
to be independent of the tuple. This lets the successor obtained at
one step be used as the input of the next.
\releasecite{Development/AbelianVarietyTheory/Proofs/TwistedOperators/RatInOfUnit.lean}{104}{Comparison of stage constants}.

On the actual row grid, the top equality from the matched,
regularized theta integrals takes the form
\[
  \frac{c_{ij,\boldsymbol\mu}}{c_{ij,\boldsymbol\mu_0}}
     = a^{\rm top}_{\boldsymbol\mu}D_{p_j,q_j}(\boldsymbol\mu),
  \qquad a^{\rm top}_{\boldsymbol\mu}\in\iota_K(K^\times).
\]
The two hyperbolic-plane comparisons, on the same stages and with
the stage measures, give
\begin{align*}
 D_{p,q}(\boldsymbol\mu)\mathcal R_{\boldsymbol\mu}(e(p,q)+1)
      &= b_{p,q,\boldsymbol\mu}D_{p,q-1}(\boldsymbol\mu) &&(q>0),\\
 D_{p,0}(\boldsymbol\mu)
      &= b_{p,0,\boldsymbol\mu}\mathcal R_{\boldsymbol\mu}(e(p,0))
                    D_{p-1,0}(\boldsymbol\mu) &&(p>0).
\end{align*}
The powers in the finite factors of the two comparisons are
\(-e(p,q)\) and \(1-e(p,0)\), respectively. The \(L\)-functions
initially omit places above the fixed base level. At its auxiliary
prime the denominator character is ramified and the numerator is
unramified. At the type primes the denominator has conductor exponent
two and the numerator has exponent at most one. Thus the quotient
is ramified above every omitted place, so both step identities use
the finite primitive \(L\)-functions on their indicated arguments.
\releasecite{Development/AutomorphicArithmetic/ThetaCorrespondence/ThetaPairLifts/LocalTypeDictionary/ThetaSquareNormComparison.lean}{17}{Top comparison}
\releasecite{Development/AutomorphicArithmetic/ThetaCorrespondence/ThetaPairLifts/LocalTypeDictionary/WittTowerNormStepOfSmallRank.lean}{19}{Step on the smaller source}
\releasecite{Development/AutomorphicArithmetic/ThetaCorrespondence/ThetaPairLifts/LocalTypeDictionary/WittTowerNormStepOfLargeRank.lean}{173}{Step on the larger source}
\releasecite{Development/AbelianVarietyTheory/Proofs/DivisorProofs/PartialHeckeLFunctionSquareEqPrimitive.lean}{26}{Euler factors at the base level}.

The field assertion about \(b_{p,q,\boldsymbol\mu}\) has an elementary
local calculation once the two matched step identities hold. In the
split chart at \(v\), write \(\nu_v\) for the quotient character on
that factor and \(h_v\) for its conductor exponent. The local factor
in the step, in the convention using \(\nu_v^{-1}\) on units, is
\[
  \nu_v(s_*)q_v^{-h_v s_\varepsilon}\nu_v(\varpi_v)^{h_v}
    \sum_{u\in(\mathcal O_v/\varpi_v^{h_v})^\times}
           \nu_v(u)^{-1}\psi_v(u/\varpi_v^{h_v}),
\]
where \(s_*\) is the fixed sign and \(s_\varepsilon\) is
\(-e(p,q)\) or \(1-e(p,0)\) in the two respective steps. An
unramified twist at a base place leaves the Gauss sum on units fixed,
and the change of the uniformizer factor is a character value in the
field. At the remaining places the relevant conductor exponent is at
most the specified depth. The additive character is trivial on
\(\mathcal O_v\), so the phases have order dividing \(q_v^{h_v}\).
These roots and the local character values belong to the coefficient
field of the grid. The two fractional-ideal factors in the step are
finite products of local uniformizer values and their inverses. The
change of a translating vector contributes a finite-order twist value
in the field. Taking ratios on the matched nonzero tests proves
\(b_{p,q,\boldsymbol\mu}\in\iota_K(K^\times)\). The measure factor in
each step is common to both tuples and cancels. This deduction uses
only the field part of the local calculation.
\releasecite{Development/AbelianVarietyTheory/BlockTheory/DivisorBase/WeilTensorModels.lean}{121}{Gauss factors}
\releasecite{Development/AbelianVarietyTheory/Proofs/DivisorProofs/BasePlaceCharacterFactor.lean}{206}{Base places}
\releasecite{Development/AbelianVarietyTheory/Proofs/DivisorProofs/GaussSumStepCharacterTwistPlaceDvd.lean}{30}{Auxiliary places}.

At stage \((0,0)\), transport the reached triangle to the standard
definite triangle used for the rational class sums. First carry the
Hermitian spaces by their isometries, then move the self-dual lattices
and their frames. The common Haar and real changes multiply the
comparison constants by a tuple-independent scalar. In the finite
transport the two vacuum factors scale by \(r_{\boldsymbol\mu}\) and
\(r'_{\boldsymbol\mu}\). The change in the comparison constant is
\(r'_{\boldsymbol\mu}/r_{\boldsymbol\mu}\), with ratio to its base value
in \(\iota_K(K^\times)\). The frame changes have the corresponding
property from the finite character quotients. Writing
\(\lambda_{\boldsymbol\mu}\) for the factor on this standard triangle
gives
\[
   D_{0,0}(\boldsymbol\mu)
       =a^{\rm bot}_{\boldsymbol\mu}
                 \frac{\lambda_{\boldsymbol\mu}}{\lambda_{\boldsymbol\mu_0}},
   \qquad a^{\rm bot}_{\boldsymbol\mu}\in\iota_K(K^\times).
\]
\releasecite{Development/AbelianVarietyTheory/Proofs/DivisorProofs/DefiniteBottomLatticeChangeConstantsAwayUnit.lean}{28}{Lattice and vacuum transport}
\releasecite{Development/AutomorphicArithmetic/ThetaCorrespondence/ThetaPairLifts/LocalTypeDictionary/WittTowerConstantRatioEqBottom.lean}{26}{Transport to the standard triangle}.

Run the smaller-source steps first. The successive arguments on the
left are \(n_0+1,\ldots,n_0+q_j\). After these steps run the larger-source
steps, whose arguments on the right are \(n_0+q_j,\ldots,n_0+q_j-p_j+1\).
Multiplying the identities cancels each intermediate comparison
quotient and yields \eqref{eq:period-square-ratios}, with
\(a_{ij,\boldsymbol\mu}=a^{\rm top}_{\boldsymbol\mu}a^{\rm bot}_{\boldsymbol\mu}
    \prod b_{p,q,\boldsymbol\mu}\), the product over these steps.
Since all \(L\)-values are nonzero there, the ranges cancel completely
when \(p_j=q_j\), and leave exactly \(\mathcal R_{\boldsymbol\mu}(n_0+1)\)
on the left when \(p_j=q_j-1\). In the product of square constants
used for the downward coordinates, each remaining Hecke factor has
exponent \(-1\) and is indexed by its cell \((i,j)\).
\releasecite{Development/AbelianVarietyTheory/Proofs/DivisorProofs/TowerOfThetaMatching.lean}{99}{Descent identity}
\releasecite{Development/AutomorphicArithmetic/ThetaCorrespondence/ThetaPairLifts/LocalTypeDictionary/TowerOfThetaMatching.lean}{57}{Comparison on the same grid}.

\teamanchor{proof.period-square-real-quotient}
\subsubsection{The real quotient for the square}
At one real place use the oriented Hermitian--skew-Hermitian
Schr\"odinger models of the two-leg kernels. The two Hermitian
source signatures are \((a-p,p)\) and \((q,a'-q)\).
In the convention that the skew-Hermitian target form is multiplied by
\(i\), its signature is \((R,S)=(a+a'-p-q,p+q)\).
Write \(m_\alpha,m_\beta\) for the exponents in the embedding of
the triangle model, so the characters of \(\mathbb C^\times\)
have factors \((z/|z|)^{m_\alpha}\) and
\((z/|z|)^{m_\beta}\). Put \(\kappa_\infty=m_\alpha-m_\beta\).
The grid gives \(\kappa_\infty>a+a'\) and \(\kappa_\infty\) even. The latter follows because
\(a,a'\) have the same parity and the restrictions of the
characters to \(\mathbb R^\times\) are the corresponding
powers of the sign character.

The polynomial oscillator action in a signed orthonormal frame
uses the operators
\[
 (B_\ell,A_\ell)=
 \begin{cases}
 (x_\ell,\partial_\ell),&\text{the two frame signs agree},\\
 (-\partial_\ell,x_\ell),&\text{the two frame signs differ}.
 \end{cases}
\]
The matrix unit indexed by \(\ell,\ell'\) acts by
\(B_\ell A_{\ell'}+\tfrac12\delta_{\ell,\ell'}\)
before adding the splitting-character trace. This is the Fock
action obtained by intertwining the prescribed Heisenberg
operators with the Gaussian line. On compact diagonals it gives
\((\tfrac{R-S}{2},\ldots,\tfrac{R-S}{2};
  \tfrac{S-R}{2},\ldots,\tfrac{S-R}{2})\) on the source vacuum,
and the analogous target vacuum weight. These are the signs of
the oscillator used in \cite[(3.3)--(3.4)]{paul-trapa-characters}.
For the large source the splitting pair is
\((\beta,\beta/\alpha)\), with the first entry used on the
target. For the small source it is
\((\alpha,\alpha/\beta)\). The source splitting twists therefore
make invariance equivalent to the \(\det^{\kappa_\infty/2}\) and
\(\det^{-\kappa_\infty/2}\) coinvariants of the respective untwisted
oscillator actions. The target trace shifts are
\(m_\beta/2\) and \(m_\alpha/2\).

The ordered pair of signatures
\[
 ((a-p,p),(a'-q,q))
\]
defines one \(\theta\)-stable parabolic \(\mathfrak q\) of
\(\mathfrak u(R,S)_{\mathbb C}\).
We have \(\lvert R-S\rvert<\kappa_\infty\), \(R\ge a-p,a'-q\), and
\(S\ge p,q\). Paul--Trapa's Table~1 and Theorem~4.1(1),(2)
give the two untwisted small theta lifts on this same ordered
Levi. With the genuine character convention of their
Section~2.8, they are
\[
 A_{\mathfrak q}(\det^{\kappa_\infty/2}\boxtimes 1),\qquad
 A_{\mathfrak q}(1\boxtimes\det^{-\kappa_\infty/2}).
\]
Tensoring with a character of the full group commutes with
cohomological induction, by tensoring the defining complexes
and restricting the character to the Levi. Thus the two
representations agree after applying their respective target
splitting twists, since the two resulting ordered pairs of
exponents are \((m_\alpha/2,m_\beta/2)\).
The parity of the ranks is the same on both legs, so the
metaplectic genuineness conventions in this comparison agree.
The theorem supplies nonzero irreducible quotients, including
the compact cases and the cases \(p=0\) or \(q=0\).
\cite[\S2.8, Table~1 and Theorem~4.1(1),(2)]{paul-trapa-characters}

Choose surjective quotient homomorphisms from the two Fock
polynomial modules to this representation, and tensor them
over the real places to obtain \(q_A,q_B\).
Howe's polynomial version of duality gives a unique maximal
proper submodule in each source coinvariant module. These
maximal submodules are the kernels of \(q_A,q_B\).
For unitary groups the compact groups
\(\mathrm U(R)\times\mathrm U(S)\) are connected, and each
vector is compact finite. A Lie-stable subspace is therefore
stable under the compact group, by exponentiation on the
finite-dimensional compact orbit of each vector.
The resulting irreducible real module and both quotient
maps can be chosen once for the prescribed frames and
infinity types. The finite twist leaves those types fixed.
\cite[Theorem~2.1]{p2flip-howe-real}

The polynomial quotients are the compact-finite parts
of the continuous irreducible quotients of the ordinary
real Schwartz models. In fact the signed frames give
compatible compact subgroups for the oscillator, so
the polynomial modules are the modules in the
algebraic correspondence. Coincidence with the smooth
correspondence on both unitary pairs gives nonzero
continuous maps to the source character tensored
with the smooth globalization of the computed target
module. Their restrictions to the polynomial space
are nonzero by density of the Hermite functions in
the Schwartz topology. Polynomial multiplicity one
then allows their normalization to the chosen
\(q_A,q_B\). The one-dimensional source twist and
the target twist enter these continuous actions by
the same characters as in the polynomial computation.
The comparison uses polynomials with fixed finite
tensors, and continuity on real Schwartz space is
used to select an initial polynomial on which the
large integral has nonzero value.
\cite[Proposition~5.1]{bao-sun-smooth-theta}
\cite[Theorems~1 and~2.1]{p2flip-howe-real}

To apply these maps to the grid, identify the polynomial
of the first arrow with the first target block in the
orthogonal sum of the two exterior vertices. On the second
arrow, before concatenating columns, transpose the matrix
of Fock variables. In the large route the transpose has
factor \(\epsilon_i\epsilon'_j\) on the variable with
indices \(i,j\) in the original arrow, where
\(\epsilon_i,\epsilon'_j\) are its target and source
frame signs. In the small route the original second-arrow
source frame is negated, so the factor is the product of
its signed frame signs. Indeed the Hermitian adjoint
conjugates the coefficients and multiplies the transposed
matrix unit by this sign product. In the large route the
second conjugation in the signed model restores the original
coefficients. In the small route pullback to the negated
second-arrow space gives the signed transpose directly.
These identities hold first for the Gaussian and then for
each creation operator by differentiation of the Heisenberg
translation. Orthogonal direct sums give the product of
the two Gaussian polynomials in the two sets of columns.
Let \(\iota_A,\iota_B\) be the resulting maps from the
two route polynomial tensor spaces to the two Fock
polynomial spaces. They are onto by the signed variable
permutations and the polynomial tensor product in disjoint
variables. The identities hold with the finite tensors
fixed throughout, including finite sums of pure tensors.
\releasecite{Development/AbelianVarietyTheory/Proofs/TwistedOperators/FockVectorProductMulPolynomials.lean}{221}{Large-route Fock substitution}
\releasecite{Development/AbelianVarietyTheory/Proofs/TwistedOperators/FockVectorProductMulPolynomials.lean}{272}{Small-route Fock substitution}

The comparison relation
\[
 q_A(\iota_A P)=q_B(\iota_B P')
\]
has onto projections and is proper, since the common
real module is nonzero. In any product realization of an
irreducible adelic quotient, the polynomial maps at fixed
finite input vanish on the corresponding local kernels.
Consequently the two maps on such an input have the form
\(q_A(\iota_A P)\otimes r_A(\varphi)\) and
\(q_B(\iota_B P')\otimes r_B(\varphi')\), in common
product coordinates. This factorization follows by
expressing the real-isotypic module as its irreducible real
factor tensored with the span of the intertwining maps into
the quotient. Indeed evaluation of this tensor product is
onto and injective, since the real endomorphism ring is
\(\mathbb C\). Schur's lemma in this form applies to
the countable-dimensional real polynomial module.
\releasecite{Development/AutomorphicArithmetic/ThetaCorrespondence/ThetaPairLifts/LocalTypeDictionary/TwistedQuotientReal.lean}{18}{Factoring the fixed infinity-type maps}

\subsubsection{Transport of the quotient across the two polarizations}
For each leg, the rational Darboux transformation of the
square intertwines the source and target actions on its
algebraic Schwartz space. The transpose and sign on the
small-source Darboux model give the opposite polarization.
Its source relations are the augmentation differences
for that transported action and the Lie derivatives for the
actual small-source splitting character. The polynomial
quotient maps on the two models have the same kernel after
transport. In detail, take a finite tensor stage containing
the exceptional places and all finite components that
depart from the vacua. In the Hermitian model the sum of the
local quotient kernels is the kernel of the tensor of the
local surjections. Each local kernel is the unique maximal
proper submodule containing the source relations. At the
real place this follows by polynomial Howe duality, and at
a finite place by the corresponding irreducible quotient
and uniqueness for the smooth Weil representation. A tensor
product of these finite-length coinvariant modules has only
one irreducible quotient. To see uniqueness, use a composition
series on the factors. An irreducible quotient of the product
is the exterior tensor product of simple factors. Restricting
to each leg gives a map from that local coinvariant module to
a sum of copies of its simple factor, so it must kill the
corresponding local kernel. The sum of those local kernels
already has simple quotient. For split factors one may use
the maximal-kernel assertion on the product directly.
The transported kernels are compatible with passing to a
larger finite tensor stage, since the extra vacuum maps are
equal and nonzero. Comparing the two surjective algebraic
quotient maps on each leg now identifies the kernels.
The real Darboux operator takes the Gaussian and its
creation operators to those in the corresponding
transported Cartan frame. The finite intertwiners act
on the corresponding finite Schwartz spaces; they
introduce no new places away from the compatible
lattice vacua. The smaller-leg comparison also
applies to arguments with its prescribed full-rank
finite support.

If both pairs of quotient maps use a common target, the
two legwise identifications differ by a single target
intertwiner. Its graph is stable for the real Lie algebra
and the finite target group. It is therefore multiplication
by a nonzero scalar in the common irreducible target
coordinates. This supplies the scalar when transporting
the Hermitian-triangle equality to the square, using the
unchanged Schwartz tests in each Darboux model.
\releasecite{Development/AutomorphicArithmetic/ThetaCorrespondence/ThetaPairLifts/LocalTypeDictionary/ThetaKernelMatchOnOppositeParabolic.lean}{30}{Kernel on the first polarization}
\releasecite{Development/AutomorphicArithmetic/ThetaCorrespondence/ThetaPairLifts/LocalTypeDictionary/ThetaKernelMatchOnOppositeParabolic.lean}{55}{Kernel on the opposite polarization}

\subsubsection{Constant terms on the chosen hyperbolic planes}
At an upper stage write the source ranks as \(a>a'\), the target
rank as \(d=a+a'\), and \(e=(a-a')/2\). Write \(T_A,T_B\) for the
larger and smaller source integrals. In the mixed-model convention
the smaller source has its form sign reversed, so the signatures
are \((a-p,p)\), \((q,a'-q)\), and \((d-p-q,p+q)\).
The rank hypotheses \(e>2(p+q)+1\) and \(a'>2(p+q)+1\)
hold throughout the rectangle of stages. Indeed the two differences
are bounded below by \(n_0-2(p_j+q_j)-1\) and
\(r-2(p_j+q_j)-1\), which are positive by \(k\ge m+3\) and
\(t\ge r+5\). The infinity-type weight difference is
\(2a_\tau>d\), with \(a_\tau\) as in the critical-value
calculation below. Use the actual splitting characters
\(\alpha,\beta\) and set \(\nu=\alpha/\beta\). The planes have
isotropic bases \((x,y)\) in the source being reduced and
\((x',y')\) in the target, with the prescribed pairings equal to \(1\).
The orthogonal complements carry the lower lattices and Cartan
involutions. The removed basis vectors are the two outer vectors of
the corresponding block at the exceptional places. These conditions
are the lattice and frame conditions used in the evaluation maps below.
\releasecite{Development/AbelianVarietyTheory/BlockTheory/DivisorBase/WeilTensorModels.lean}{532}{Adapted hyperbolic planes}.

For the target parabolic of \(Ex'\), take the constant term by
averaging over its rational unipotent quotient with total mass \(1\).
In the mixed Schr\"odinger model the trivial Fourier character cuts
out the isotropic vectors in the source. The zero vector gives the
target-only evaluation \(E_0\) at the origin of the removed mixed
coordinates. It yields the lower theta integral if the source is
unchanged. When the source also loses a plane, both nonzero shifts
and their reduced evaluation maps can be described on the smaller
tensor space. Denote these by \(u_A(l)=l\,x\otimes y'\),
\(u_B(l)=\bar l\,y\otimes x'\), and \(E_A,E_B\), respectively.
They are evaluation at zero after Fourier transform on the
isotropic coordinates complementary to the corresponding shift.
The mixed-model Fourier expansion follows by Poisson summation on
the removed coordinates, with the self-dual measures for the fixed
additive character. Its zero orbit and its nonzero orbit have the two
Levi characters
\[
 \alpha(t)|t|_E^{a'/2},\qquad \beta(t)|t|_E^{a/2}.
\]
Since the real exponents differ, these two components are unique.
Indeed, comparing at \(t=1\) and at an idele of norm other than \(1\)
determines the values of both components on the lower Levi.
\releasecite{Development/AbelianVarietyTheory/Definitions/Divisors/CTSmallCone.lean}{42}{Mixed evaluations}
\releasecite{Development/AbelianVarietyTheory/Definitions/Divisors/CTSmallCone.lean}{86}{Two Levi characters}.

For an isotropic line in a source of rank \(n\), the determinant of
conjugation by \(t\in\mathbb A_E^\times\) on the unipotent radical
has absolute value \(|t|_E^{n-1}\). The \(E\)-linear coordinates
contribute \(|t|_E^{n-2}\), and the center, which is one-dimensional
over \(F\), contributes \(|N_{E/F}t|_F=|t|_E\). The rational nonzero
isotropic vectors form one orbit by Witt's theorem. Unfold its sum by
the stabilizer of \(x\), with quotient integration in the order
unipotent, lower source, ideles, and compact group. This gives the
factor \(|t|_E^{1-n}\,d^\times t\). For the larger source one uses
\(t=l^{-1}\), and \(E_A\) has weight \(|t|_E^{(d-2)/2}\)
and splitting character \(\nu^{-1}(t)\). The Mellin weight is thus
\(\nu(l)|l|_E^{a-1-(d-2)/2}=\nu(l)|l|_E^{e}\).
For the smaller source the second polarization uses \(t=l\),
and \(E_B\) has weight \(|t|_E^{d/2}\) and character \(\nu(t)\).
The Mellin weight is \(\nu(l)|l|_E^{d/2+1-a'}=\nu(l)|l|_E^{e+1}\).
The compact average is taken with the fixed product of normalized
compact measures. For \(J=A,B\), the resulting operator is
\[
 C_J\Phi=\int_K\int_{\mathbb A_E^\times}
 E_J\rho(u_J(l))\omega(k)\Phi\,\nu(l)|l|_E^z\,d^\times l\,dk,
 \qquad z=e\ (J=A),\quad z=e+1\ (J=B).
\]
All source and target Haar changes in this unfolding depend on the
geometry and the fixed additive character. The fixed rank bounds give
\(z\geq2\), so the finite product here is in its absolute convergence
range. The exchange of the global kernel sum and integral also uses the
stage's absolute theta estimates, including the smaller-source support
condition at the extra split place \(v_1\).
\releasecite{Development/WeilBound/ThetaLifts/SmallCones/Unfold.lean}{65}{Line-parabolic quotient measure}
\releasecite{Development/WeilBound/ThetaLifts/SmallCones/ConeLawAB.lean}{61}{Larger-source mixed weight}
\releasecite{Development/WeilBound/ThetaLifts/SmallCones/ConeLawAB.lean}{115}{Smaller-source mixed weight}
\releasecite{Development/WeilBound/ThetaLifts/SmallCones/ConeCore.lean}{22}{Absolute unfolding and compact average}.

For the smaller source the switch to \(u_B\) must be made before
unfolding. To see the correction, expand the theta sum in the two
polarizations and average over their common isotropic coordinates. The
intersection is spanned by \(x\otimes x'\), \(x\otimes W_1\), and
\(X_1\otimes x'\), where \(X_1,W_1\) are the two complements.
Orthogonality for the mixed-model Fourier characters leaves the
indices \(\xi x\otimes y'\) on one side and \(\xi y\otimes x'\)
on the other, with \(\xi\in E\). The nonzero rational vectors of a
fixed isotropic line are reindexed on the second side by
\(\xi\mapsto\bar\xi^{-1}\). Subtracting the zero indices leaves a
multiple of the unshifted \(E_A\) term, since \(E_B\Phi=0\) for the
smaller-source tests. The unshifted correction lies in the
\(\alpha(t)|t|_E^{a'/2}\) component. It does not change the
\(\beta(t)|t|_E^{a/2}\) component used in this comparison.
At \(v_1\), the upper small input is in the split-coordinate vanishing
subspace used in the smaller-source estimate. Its condition is
preserved by the target action and by the projection associated to
\(E_B\), so the lower input in the cone integral still satisfies it.
\releasecite{Development/WeilBound/ThetaLifts/SmallCones/TwoExp.lean}{121}{Fourier restriction to each rational line}
\releasecite{Development/WeilBound/ThetaLifts/SmallCones/GenSmallOf.lean}{48}{Poisson correction and split components}
\releasecite{Development/WeilBound/ThetaLifts/LocalTheory/UnramifiedCohomology/FLocBOf.lean}{84}{Smaller-source support and quotient at the extra split place}.

\subsubsection{Local scalars on the matched quotients}
Away from the ramified places, let \(\mathfrak b,\mathfrak c\)
be the line ideals for \(x,x'\) on the two adapted planes. Write
\(\mathfrak d_E^{-1}\) for the absolute codifferent and use a bar for
conjugation of ideals. The two line ideals for \(y,y'\) are
\(\overline{\mathfrak b}^{-1}\) and
\(\mathfrak d_E^{-1}\overline{\mathfrak c}^{-1}\) in the Hermitian
and skew-Hermitian lattice conventions. Thus the two shift conditions
on \(l\) cut out the ideals
\[
 \mathfrak a_A=\mathfrak b\mathfrak d_E^{-1}\overline{\mathfrak c}^{-1},
 \qquad
 \mathfrak a_B=\overline{\overline{\mathfrak b}^{-1}\mathfrak c}.
\]
For the component at a place \(w\) of \(E\) these conditions on the
standard vacuum are \(\operatorname{ord}_w(l)\geq m_w\), where
\(m_w=\operatorname{ord}_w\mathfrak a_J\). With
\(\operatorname{vol}(\mathcal O_{E_w}^{\times})=1\) integration gives
\[
 \sum_{h\geq m_w}\nu_w(\varpi_w)^h q_w^{-zh}
   =\nu_w(\varpi_w)^{m_w}q_w^{-zm_w}L_w(z,\nu_w).
\]
The ideal norm powers are common geometry factors across twists. The
character evaluates to \(\nu(\mathfrak b\mathfrak c)\) for the first
cone and \(\nu(\mathfrak b\mathfrak c)^{-1}\) for the second. This
uses \(\nu\circ\mathrm{conj}=\nu^{-1}\) and triviality on the
codifferent factor. In this setting the relative discriminant is the
unit ideal, so the absolute codifferent is extended from \(F\), on
whose ideles \(\nu\) is trivial. At places omitted from the vacuum
calculation the character is ramified over every place of \(E\), so
the Euler product in these factors is the finite primitive \(L(z,\nu)\).
\releasecite{Development/WeilBound/ThetaLifts/LocalTheory/ConeShifts/IdealId0.lean}{17}{Line ideals and their evaluations}
\releasecite{Development/WeilBound/ThetaLifts/LocalTheory/ConeShifts/FubMain.lean}{145}{Reassembly of the Euler product}.

At an exceptional split place put \(k=F_v\) and identify
\(E_v=k\times k\) using the selected place \(w\), with character
\((\vartheta,\vartheta^{-1})\) and conductor exponent \(h>0\).
Use the measures \(\operatorname{vol}(\mathcal O_k)=1\) and
\(\operatorname{vol}(\mathcal O_k^\times)=1\) for the additive and
multiplicative groups. The local additive character is unramified.
The two standard sections in the matrix model have their respective
square blocks in \(\mathrm{GL}(\mathcal O_k)\), weighted by
\(\vartheta^{-1}\) of the determinant. All entries in the other
columns lie in \(\varpi^h\mathcal O_k\). In the common quotient they
are both the parabolic section supported by
\(Z\in \varpi^h\operatorname{Mat}_{a'\times a}(\mathcal O_k)\) on
the opposite cell \(\left(\begin{smallmatrix}1&0\\Z&1\end{smallmatrix}\right)\),
where they have value \(1\). Indeed the diagonal block minors there
are identities and the remaining columns are \(Z\) or its transpose
with sign. Off the open cell the two complementary projected block
minors are singular, so both sections vanish. Parabolic covariance
then gives equality at every group element. This verifies the common
nonzero test in both the upper and lower quotients.
\releasecite{Development/WeilBound/ThetaLifts/LocalTheory/StandardValuations/SmallFSing.lean}{134}{Two sections on the opposite cell and its complement}
\releasecite{Development/WeilBound/ThetaLifts/LocalTheory/StandardValuations/ExistsVol.lean}{101}{Transport to the fixed frames}.

The partial Fourier evaluations on these sections remove two pivot
positions from the same square block. In the cone evaluation the
determinant on the chart is minus the product of the two removed
pivots times the determinant of the lower block. The translation
sets one pivot to the first coordinate \(s\) of \(l\). The
Fourier variable pairs with the other coordinate \(t\) of \(l\).
The minus sign in the determinant and the Fourier sign on that
coordinate cancel under the simultaneous change of the integration
variable. The other integrated entries are one row of the untouched
block or one column after transposition. Their total volume is
\(q_v^{-hn_{\rm other}}\), where \(n_{\rm other}=a'\) for the larger
source step and \(n_{\rm other}=a\) for the smaller source step.
Exactly the same volume multiplies the plain evaluation on the
other section. Set \(D_\vartheta(x)=1_{\mathcal O_k^\times}(x)
\vartheta(x)^{-1}\) and
\(\widehat D_\vartheta(t)=\int_k D_\vartheta(y)\psi_v(ty)\,dy\).
After the common volume is canceled, the cone multiplier is
\[
 \int_{k^\times\times k^\times}
  D_\vartheta(s)\widehat D_\vartheta(t)\,
  \vartheta(s)\vartheta(t)^{-1}|st|^z\,d^\times s\,d^\times t.
\]
The integral defining \(\widehat D_\vartheta(t)\) is zero for
\(\operatorname{ord}_v t<-h\), by additive orthogonality on cosets
modulo \(\varpi^h\). It is zero for \(\operatorname{ord}_v t>-h\),
by multiplication by a unit in the last conductor subgroup on which
\(\vartheta\) is nontrivial, using all units when \(h=1\).
For \(u\in\mathcal O_k^\times\), partition into the residue classes of
the units to get
\[
 \widehat D_\vartheta(\varpi^{-h}u)=q_v^{-h}\vartheta(u)
   \sum_{x\in(\mathcal O_k/\varpi^h)^\times}
   \vartheta(x)^{-1}\psi_v(\varpi^{-h}x).
\]
Substitution in the multiplier gives
\(q_v^{h(z-1)}\vartheta(\varpi)^h\) times this Gauss sum before its
factor \(q_v^{-h}\) and its \(u\)-character. Equivalently, the
multiplier is exactly \(\varepsilon_v(1-z,\vartheta,\psi_v)\) in the
finite convention of the comparison. It is nonzero by Fourier
inversion and the character's transformation law on the shell.
\releasecite{Development/WeilBound/ThetaLifts/LocalTheory/StandardValuations/ConeVal.lean}{158}{Removed pivots and surviving row or column}
\releasecite{Development/WeilBound/ThetaLifts/LocalTheory/StandardValuations/ExistsVol.lean}{85}{Epsilon convention}
\releasecite{Development/WeilBound/ThetaLifts/LocalTheory/StandardValuations/ExistsVol.lean}{162}{Smaller-source volume}
\releasecite{Development/WeilBound/ThetaLifts/LocalTheory/StandardValuations/ExistsVol.lean}{219}{Larger-source volume}.

To use these scalar calculations after taking a constant term, the
plain and cone maps must descend to the common upper and lower
target quotients. They intertwine the target line parabolic with the
Levi character selected above. At an unramified place the scalar
comparison on the spherical line then holds on its parabolic span.
Indeed the parabolic times the maximal compact gives the full target
group, the compact translates of the vacuum are scalar, and its
nonzero image is cyclic in the irreducible quotient. At a ramified
split place the equivariant maps are proportional. The split matrix
argument discards the two removed columns by the unipotent relations
and a unit on which \(\vartheta\) is nontrivial. The surviving
full-rank stratum has the twisted coinvariants of the transitive
space of full-rank rectangular matrices. Integrating over the
invertible source block identifies this representation, up to a
target determinant twist and inverse transpose, with normalized
parabolic induction of the unitary determinant character and the
trivial representation. Bernstein's irreducibility theorem and
the countable-dimensional form of Schur's lemma give proportionality.
The nonzero matched matrix sections computed above determine the
scalar. At \(v_1\) in the smaller-source step one also needs descent
from the full Schwartz space to the quotient. The change of variables
in the cone integral kills the augmentation relations from the
source action, so the descent agrees with the one obtained using the
small support subspace, which still surjects to that quotient.
\releasecite{Development/WeilBound/ThetaLifts/LocalTheory/UnramifiedCohomology/UnramifiedDeterminedBySphericalHolds.lean}{100}{Cyclicity of the spherical comparison}
\releasecite{Development/WeilBound/ThetaLifts/LocalTheory/DixmierCommutants/LocalMultiplicityOneIrreducibleGL.lean}{215}{Ramified Levi restriction}
\releasecite{Development/WeilBound/ThetaLifts/LocalTheory/LocBOf.lean}{101}{Descent of the small cone at the extra split place}.

The real quotient maps are proportional on the polynomial vectors
with the prescribed Fock data. In the highest-weight model the
Levi weight and the compact diagonal weight on the deleted pair of
coordinates determine the possible image of the upper highest-weight
vector up to one scalar. The unipotent covariance equations and the
mixed-root equations then propagate that scalar from the highest-weight
vector to the whole polynomial module. For the maps at hand,
nonvanishing of the plain map follows by surjectivity of the mixed
evaluation and the lower quotient maps. Nonvanishing of the cone
map uses a Schwartz bump at infinity in the removed variables,
followed by polynomial density on the prescribed quotient. The real
cone and evaluation maps, their quotient normalizations and their
characters at infinity are common to all finite twists under
comparison. Their ratio is therefore a common nonzero scalar.
\releasecite{Development/WeilBound/ThetaLifts/LocalTheory/ToralCohomology/ParZero.lean}{175}{Highest-weight restriction}
\releasecite{Development/WeilBound/ThetaLifts/LocalTheory/LocAOf.lean}{100}{Nonzero quotient maps on the larger-source comparison}
\releasecite{Development/WeilBound/ThetaLifts/LocalTheory/LocBOf.lean}{132}{Nonzero quotient maps on the smaller-source comparison}.

Upon tensoring, the vacuum scalar at each extra unramified place
used to present a tensor cancels its inverse from factoring out the
Euler product away from that finite set. Thus for a matched
polynomial pair \((\Phi_A,\Phi_B)\) the reduced pairs are
\[
 \big(C_A\Phi_A,\,k_A E_0\Phi_B\big),\qquad
 \big(k_B E_0\Phi_A,\,C_B\Phi_B\big)
\]
in the larger and smaller source steps, respectively. Here \(k_J\)
is a common nonzero geometric and real factor times the corresponding
line-ideal character factor, the product of finite epsilon factors,
and \(L(z,\nu)\). This assertion uses the quotient maps of the
prescribed test models on every tensor factor. Their product on an
upper tensor equality gives the indicated lower tensor equality,
and linearity gives the statement for the finite sums making up the
matched pair.
\releasecite{Development/WeilBound/ThetaLifts/LocalTheory/LocAOf.lean}{121}{Assembly of the normalized larger-source maps}
\releasecite{Development/WeilBound/ThetaLifts/LocalTheory/LocBOf.lean}{162}{Assembly of the normalized smaller-source maps}.

\paragraph{Using the matching constant.}
Let \(c,c_1\) be the upper and lower matching constants on the
specified quotients. For the smaller removal the
\(\beta(t)|t|_E^{a/2}\) coefficients of the larger and smaller upper
integrals restrict to \(C_0 T_{A,1}(E_0\Phi_A)\) and
\(C_1 T_{B,1}(C_B\Phi_B)\), respectively. The lower matching
identity gives \(k_B T_{A,1}(E_0\Phi_A)=c_1 T_{B,1}(C_B\Phi_B)\).
For the larger removal the matching coefficients are those of
\(\alpha(t)|t|_E^{a'/2}\). They restrict to
\(C_1 T_{A,1}(C_A\Phi_A)\) and \(C_0 T_{B,1}(E_0\Phi_B)\),
with \(T_{A,1}(C_A\Phi_A)=c_1 k_A T_{B,1}(E_0\Phi_B)\).
Here \(C_0,C_1\) refer separately to the geometric constants for the
two cases and are independent of the finite twist. Surjectivity of
the target-only mixed evaluation on polynomial data and nonvanishing
of the lower lifts give a test with nonzero chosen coefficient in
each case. Canceling that coefficient gives
\[
 c\, k_B (C_1/C_0)=c_1\quad\text{for the smaller removal},\qquad
 c=(C_1/C_0) k_A c_1\quad\text{for the larger removal}.
\]
These equations compare actual matching constants from the
chosen coefficient. The argument of
Section~\ref{sec:period-triangle-uniqueness} promotes the
coefficient equality to equality on the matched polynomial
tests. The nonzero starting lift is the definite one in
Section~\ref{sec:period-definite-triangle}.
\releasecite{Development/WeilBound/ThetaLifts/ConeHypotheses/KeyAOf.lean}{37}{Smaller-source nonzero matching coefficient}
\releasecite{Development/WeilBound/ThetaLifts/ConeHypotheses/KeyAOf.lean}{88}{Larger-source nonzero matching coefficient}
\releasecite{Development/WeilBound/ThetaLifts/ConeHypotheses/CtSmallCone.lean}{88}{Induction using the end stages}.

\teamanchor{proof.period-square-inert-quotient}
\subsubsection{The normalized quotient at the inert places}
Write \(L/k_v\) for the unramified quadratic extension
at an inert place \(v\). The triangle assumptions give odd
residue characteristic, the conductor-zero additive
character, and self-dual lattices. The exceptional set
is split, so both splitting pairs are unramified here.
The unramified characters \(\alpha_v,\beta_v\) of
\(L^\times\) agree and have product one. Indeed a
uniformizer of \(k_v\) is a uniformizer of \(L\), and on
it the two values are \((-1)^{a'}\) and \((-1)^a\).
In particular the character on the source determinant
in each Weil action is trivial.

The target is the quasi-split unitary group \(H\) of
dimension \(2m_v=a+a'\). Use the integral Siegel parabolic
\(P=MN\) and the stabilizer \(K\) of the target self-dual
lattice. The decomposition \(H=PK\) can be taken in an
integral Witt basis for this lattice. Its Levi
\(M=\mathrm{GL}_{m_v}(L)\) acts on \(N\) by
\(B\mapsto tBt^*\) in Hermitian matrix coordinates
(after multiplying by a fixed trace-zero unit if
necessary). For diagonal \(t\), each diagonal coordinate
contributes \(|N_{L/k_v} t_i|_{k_v}\).
Each off-diagonal pair is one \(L\)-coordinate with
multiplier \(t_i\overline{t_j}\) and hence contributes
\(|N_{L/k_v} t_i|_{k_v}|N_{L/k_v} t_j|_{k_v}\).
The elementary unipotents have determinant one in this
action. Thus
\[
 \delta_P(t)=|\det t|_L^{m_v}.
\]

Pass to the mixed Schr\"odinger model associated with
the target Witt basis. Evaluation \(e_A\) at zero
on the large source and \(e_B\) on the small source
are invariant under their source unitary groups.
The Heisenberg self-dual lattice lines on the two legs
correspond to scalar multiples of the characteristic
functions of the lattices in this model. In particular
the evaluations on the given vacua \(u_A,u_B\) are
nonzero. Normalize each evaluation by its value on
the corresponding vacuum. The standard Schr\"odinger
formulas give the two left \(P\)-characters of
\(f_D(h)=e_D(\omega_D(1,h)u_D)/e_D(u_D)\), \(D=A,B\).
They are
\[
 \chi_A(t)=\beta_v(\det t)|\det t|_L^{a/2},\qquad
 \chi_B(t)=\alpha_v(\det t)|\det t|_L^{a'/2}.
\]
In fact the Siegel unipotent acts at zero by the
value one of its quadratic phase. The Levi acts by
the indicated determinant character and absolute
Jacobian factor. Both \(f_D\) are right \(K\)-invariant
and \(f_D(1)=1\). Their product of left characters
is precisely \(\delta_P\).

For Haar probability on \(K\) define
\[
 \ell_D(\phi)
  =\int_K\frac{e_D(\omega_D(1,k)\phi)}{e_D(u_D)}\,dk,\qquad
 \Omega_D(h)=\ell_D(\omega_D(1,h)u_D).
\]
These are integrals over finite quotients on smooth
vectors. The compact integral pairing
\(\int_K f_A(k)f_B(k)\,dk\) intertwines the right
actions on the two induced representations, since the
product of left characters is \(\delta_P\).
Indeed this is the integral on \(P\backslash H\) of
the density of left weight \(\delta_P\).
Its transformation law follows by changing variables
in the Iwasawa formula. It follows by moving the
right translation to the other argument that
\(\Omega_A(h)=\Omega_B(h^{-1})\).
The inverse of a Cartan double coset \(KhK\) equals
\(KhK\), since the Weyl group acts by independent
inversions of the uniformizer coordinates in the
even-dimensional quasi-split unitary group.
Thus \(\Omega_A=\Omega_B=\Omega\), with \(\Omega(1)=1\).
\releasecite{Development/WeilBound/ThetaLifts/SphericalLattices/FunAsm.lean}{113}{The two normalized spherical coefficients}

For an arbitrary vector \(\phi\), put
\[
 F_{D,\phi}(h)=\frac{e_D(\omega_D(1,h)\phi)}{e_D(u_D)}.
\]
Let \(\chi=\alpha_v=\beta_v\). The images of the two
maps \(F_{D,\bullet}\) are precisely the evaluation
representations \(R(V_A,\chi)\) and \(R(V_B,\chi)\)
in the degenerate principal series, in the normalization
of \cite[\S1]{tr3yamana}. To apply Proposition~1.4
there, the size of the Siegel Levi is \(m_v\) in our notation
and the small source dimension is \(a'<m_v\).
The large source is isometric to the small source plus
\(m_v-a'\) hyperbolic planes. Indeed every self-dual
Hermitian lattice for the unramified extension has an
orthogonal integral basis with unit norms: find a vector
of unit norm in the nondegenerate reduction, split off
its line, and repeat. Each norm unit can be made one
using surjectivity of the norm on units. A hyperbolic
plane also has a self-dual lattice, and the dimensions
of \(V_A\) and of \(V_B\) plus those planes agree.
The parity condition \(\chi|_{k_v^\times}=\epsilon_{L/k_v}^{a'}\)
holds on both sources. Yamana's Proposition~1.4(1),(2)
identifies \(R(V_B,\chi)\) with the irreducible small-source
coinvariant module. Parts (4),(5) give a surjection
\(R(V_A,\chi)\twoheadrightarrow R(V_B,\chi)\).
The full large-source coinvariant module has only one
irreducible quotient by local Howe duality. Its maximal
kernel is thus the kernel of the composed surjection.
\cite[Proposition~1.4(1),(2),(4),(5)]{tr3yamana}
\cite[Theorem~1.2]{p2flip-gan-takeda}

For each leg send a smooth vector \(\phi\) to
\(h\mapsto\ell_D(\omega_D(1,h)\phi)\).
On the small leg this is a nonzero map from the
irreducible \(R(V_B,\chi)\), since \(\Omega(1)=1\).
On the large leg, pair \(F_{A,\phi}\) against
\(R(V_B,\chi)\) using the compact integral above. The
pairing is nonzero on \(f_A,f_B\), giving \(1\).
Thus the resulting map to the smooth contragredient of
\(R(V_B,\chi)\) is onto and has the unique maximal kernel.
Evaluation of this dual pairing on right translates
of \(f_B\) separates the quotient, since the translates
of the nonzero vector \(f_B\) span \(R(V_B,\chi)\).
In particular the large coefficient map has precisely
that same kernel. Both coefficient maps now have
irreducible image and send their prescribed vacua to
the function \(\Omega\). Hence the images agree as the
span of the right translates of \(\Omega\). This normalizes
the two maps in the exact spherical model on all smooth
inputs. The restricted tensor product uses this common
vector at every inert place.
\releasecite{Development/WeilBound/ThetaLifts/SphericalLattices/FinInert.lean}{120}{The common coefficient representation}

For the global square-integrable comparison, the resulting local
kernel is killed by the theta map into \(L^2\). With the other
factors fixed, this local map factors through the admissible
finite-length source coinvariants. Its image carries the positive
invariant quotient form induced by \(L^2\). Such a smooth
finite-length module is semisimple. Indeed for any stable
submodule \(M\) and any vector fixed by a compact open \(J\),
orthogonally project the vector to the finite-dimensional
space \(M^J\). By \(J\)-averaging the remainder is orthogonal
to all of \(M\). This gives an invariant direct complement.
The unique irreducible quotient and local multiplicity one
now identify the square-integrable image as the same quotient
or zero. This also proves the required vanishing of the maximal
local kernel at the restricted spherical place.

\subsubsection{Growth bounds for the triangle tests}
The source integrals and the constant terms below use an absolute
kernel estimate uniform in the real vector for fixed finite data.
Write \(p,q\) for the two source real Witt indices. The imposed
rank gap is \(e=(a-a')/2>2(p+q)+1\). For two Schwartz tests on a Schr\"odinger
space use the Heisenberg coefficient
\[
 B_{\Phi_1,\Phi_2}(z)
    =\int (\rho(z)\Phi_1)(u)\overline{\Phi_2(u)}\,du.
\]
Identify the rational symplectic space with the \(E\)-matrix space
of size \(n\) by \(d\), where \(n\) is the source rank. Expanding
the product of the two rational theta sums and applying Poisson
summation in one variable expresses it as the sum of the coefficients
over rational \(z\), with the fixed additive covolume factor. Indeed,
organize the double sum by the difference of its two indices and
expand the sum in the second index by Fourier inversion. The central
Heisenberg phase for rational indices is trivial. Weil covariance
replaces the coefficient of the acted tests at \(z\) by that of
the original tests at \(g^{-1}z(h^{-1})^{\mathrm t}\). For fixed
finite data the finite coefficient is supported in a compact set,
since it is an integral of compactly supported locally constant
functions against an additive character. Its real coefficient has
Schwartz decay of any order, bounded by a product of continuous
Schwartz seminorms of the real tests. This follows by exchanging
derivatives and polynomial weights in the displayed integral, or by
taking its partial Fourier transform in the integration variable.
\releasecite{Development/WeilBound/ThetaLifts/FrameEstimates/SqAllC.lean}{62}{The rational coefficient majorant}
\releasecite{Development/WeilBound/ThetaLifts/FrameEstimates/SquareIntegrableWeilRange.lean}{96}{The Heisenberg coefficients}

For unrestricted tests use the original test in both inputs. For
the restricted source at \(v_1\), first write its local test as a
finite sum supported in compact open patches \(D_i\) of the split
full-rank matrix locus with \(D_i+D_i\) still full rank. Such a
patch about \(Y\) exists because \(2Y\) has full rank and full rank
is open. For one patch reflect the adelic test \(u\mapsto\Phi_1(-u)\) at all places
in the second input. It has the same theta sum, since reflection
commutes with the Weil action and permutes the rational summation
indices. In the split model the Heisenberg translation in one
component of \(z\) shifts the matrix variable, and the other component
gives a phase. Nonzero support of the product in the coefficient
implies that the translation belongs to \(D_i+D_i\), since the
second test is reflected. The translated component therefore has
full row rank. This implies full row rank of the rational matrix
in the sum, since the matrices multiplying it in Heisenberg
covariance are invertible. Bound the original restricted test by
summing the bounds of its finitely many patches. The restriction
here is at the actual split place used in the cone calculation.
\releasecite{Development/WeilBound/ThetaLifts/FrameEstimates/ThetaSumNeg.lean}{80}{Parity in the restricted square bound}
\releasecite{Development/WeilBound/ThetaLifts/FrameEstimates/HermitianFormDetNeZero.lean}{76}{The full-rank support patches}

Use adelic Siegel sets on the source and the target through
Borel's reduction for reductive groups over a number field
(\href{https://www.numdam.org/item/PMIHES_1963__16__5_0.pdf}{\S\S3.4--3.5,
 5.1--5.2, and 5.6}).
In rational
Witt bases write their nonnegative decreasing diagonal parameters
as \(s_1,\dots,s_l\) and \(t_1,\dots,t_{l'}\), repeated with
the opposite signs and zeros in the full bases. Bounded shifts in
the chambers affect the estimates by constants. Moving compact
positive unipotent factors to the right of the diagonal keeps them
in a bounded set by contraction on the positive chamber.
The parameters
act by the same positive scales at all infinite embeddings of
the number field, with the remaining factors in fixed compact
sets. Let \(S=\sum_i s_i\), \(T=\sum_j t_j\), and \(D=[E:\mathbb Q]\).
Absorb the compact factors and the rational basis changes into
the constants and the fixed fractional ideal containing the
matrix entries in the coefficient sum. Each entry has dilation
\(e^{w_{ij}}\), with \(w_{ij}\) the sum of its signed source and
target parameters. For any fixed fractional ideal \(\mathfrak a\)
in \(E\), lattice counting in its Minkowski embedding gives
\[
 \sum_{x\in\mathfrak a}(1+e^w\|x\|_\infty)^{-N}
        \ll_N \max(1,e^{-Dw}) \qquad (N>D).
\]
This follows by counting points in balls using a bounded fundamental
parallelepiped and summing over dyadic annuli. Apply a sufficiently
large coefficient seminorm and split its decay among the entries.
The negative weights satisfy
\(\sum_{i,j}\max(0,-w_{ij})\le dS+nT\).
Thus the kernel is bounded by \(C\exp(D(dS+nT)/2)\).

In the supported sum, the row whose source dilation is \(e^{s_1}\)
has some nonzero rational entry. The norm of a nonzero element
of the fixed fractional ideal has a positive lower bound in the
Minkowski embedding. After dilation this costs at least a constant
times \(e^{s_1-t_1}\) in the real coefficient, interpreting \(t_1=0\)
if there is no target parameter. Reserve an additional Schwartz
power to obtain an extra factor \(e^{-M\max(0,s_1-t_1)}\) in
the bound for the kernel, for any prescribed \(M\). This gives
arbitrary source decay at the expense of a larger target exponent
because \(S\le l s_1\) and \(t_1\le T\). The anisotropic source
case has an empty parameter list and a compact quotient.
\releasecite{Development/WeilBound/ThetaLifts/FrameEstimates/FrameBasis.lean}{42}{Rational diagonal weights}
\releasecite{Development/WeilBound/ThetaLifts/FrameEstimates/SoSemi.lean}{24}{The source and target Siegel estimates}

For the larger source the Siegel density in the source parameters
is, up to bounded factors,
\[
 \exp\!\left(-D\sum_{i=1}^l(n-2i+1)s_i\right).
\]
The decreasing parameters give
\(\sum_i(n-2i+1)s_i\ge(n-l)S\). The hypothesis
\(d+2l<2n\) therefore makes the unrestricted bound integrable
in the source parameters. It holds for the larger source of the
triangle and for its target-only lowered lift, since \(l\le p\)
and \(e>2(p+q)+1\). For the smaller source the supported bound
gives an integrable source bound by its arbitrary decay. This is
also valid after a target-only reduction onto supported data.

The resulting source integral of the absolute kernel grows at
most as \(e^{C_1 T}\) for a fixed \(C_1\). This is a polynomial
height bound. To see that the reduction controls the target
height even before reduction, take a nonzero rational leading
isotropic vector \(x\) in the target Witt frame for a reduced
representative \(\gamma h=a(t)k\). The adelic product norm of
\(h^{-1}\gamma^{-1}x=k^{-1}a(t)^{-1}x\) is
\(O(e^{-D t_1})\) by compactness of \(k\). The rational vector
\(\gamma^{-1}x\) has product norm at least one. Multiplying by
the matrix norm of \(h\) proves that \(e^{D t_1}\) is bounded by
a constant times the adelic height of \(h\), in a compatible
matrix height convention. Since \(T\le l't_1\), this gives the
required polynomial bound. The fixed finite tests contribute
only constants. On compact target sets the same bounds give an
integrable source majorant times a continuous real seminorm. They
yield continuity in the real Schwartz test and justify differentiation
in the target real action. Every real derivative is another real
Schwartz test with the same finite data. In particular the height
exponent works uniformly for all derivatives. A compact unipotent
average preserves these estimates.
\releasecite{Development/WeilBound/ThetaLifts/FrameEstimates/DomWeil.lean}{13}{Domination in the unrestricted source range}
\releasecite{Development/WeilBound/ThetaLifts/FrameEstimates/ThickFrames.lean}{147}{Uniform seminorm domination and integration}

\teamanchor{proof.period-triangle-uniqueness}
\subsubsection{The triangle identity from a Levi coefficient}\label{sec:period-triangle-uniqueness}
Write \(V_A,V_B,W\) for the three spaces at one triangle, put
\(d=a+a'\) and \(e=(a-a')/2\), and write \(G=U(W)\). At each real
place the signatures, in the chosen form conventions, are
\((a-p,p)\), \((q,a'-q)\), and \((d-p-q,p+q)\). The two Weil
splittings use \((\beta,\beta/\alpha)\) and
\((\alpha,\alpha/\beta)\), in the order source character, target
character. In particular the integrated source representations are
trivial in these splittings. The rank and weight bounds used below are
\(e>2(p+q)+1\), \(a'>2(p+q)+1\), and
\(m_\alpha-m_\beta>d\) at all real places, where \(m_\chi\) is the
exponent of \(z/|z|\) in the chosen embedding. Let \(\mathscr D\)
be the space of differences \(T_A(\Phi)-T_B(\Phi')\) of the two
theta integrals on \([G]=G(F)\backslash G(\mathbb A_F)\), with
polynomial real and Schwartz finite data. Impose the prescribed support
at \(v_1\) on the \(B\) tests. The two source tests may be chosen
independently. In particular \(T_A(\Phi)-cT_B(\Phi')\) belongs to
\(\mathscr D\) for every complex scalar \(c\).
\releasecite{Development/AbelianVarietyTheory/BlockTheory/DivisorBase/WeilTensorModels.lean}{360}{Triangle hypotheses and splitting characters}
\releasecite{Development/WeilBound/ThetaLifts/EigenvalueEndgame/PlaneRemovalConfiguration.lean}{59}{Space of differences}

For a rational isotropic line \(L\) in \(W\), let \(P_L\) be its
stabilizer and \(N_L\) its unipotent radical. The line constant term
\(f_L\) is the average over \([N_L]\) with mass one. With a chosen
dual isotropic line the Levi is
\(\operatorname{Res}_{E/F}\mathrm{GL}_1\times U(W_1)\), where
\(W_1=L^\perp/L\). The two line coefficients have the unnormalized
characters \(\alpha(t)|t|_E^{a'/2}\) and
\(\beta(t)|t|_E^{a/2}\). The modulus of \(P_L\) on this factor is
\(|t|_E^{d-1}\), so in the normalized constant term the two
exponents are \((1-a)/2\) and \((1-a')/2\), respectively.
The coefficients are uniquely determined by \(f_L\),
by evaluating at two different norms, and are left \(N_L\)-invariant.
For the latter assertion, compare the two evaluations after replacing
the argument by \(ng\). The torus normalizes \(N_L\), so the
constant terms are unchanged in both evaluations.

Suppose that for every matched polynomial pair the selected coefficient
of \(T_A(\Phi)-cT_B(\Phi')\) vanishes on the embedded \(U(W_1)\).
Here the selected coefficient is the \(\beta\) coefficient when
\(q>0\), and the \(\alpha\) coefficient when \(q=0<p\). A right
translate by the target real maximal compact and by an arbitrary
finite target element is still a matched polynomial pair. The real
and finite quotient maps intertwine these actions, the real compact
preserves the polynomial space, and the smaller test support at
\(v_1\) is preserved by the target action. Right translation commutes
with the line constant term and with projection onto each of its two
characters. Thus the coefficient vanishes on \(U(W_1)K\), where
\(K\) is the real maximal compact times the finite adelic group.
Adelic Iwasawa decomposition \(G(\mathbb A_F)=P_L(\mathbb A_F)K\)
and the covariance of the coefficient now give its vanishing at every
argument. This uses the coefficient identity for all matched pairs,
including the translates of the given pair.
\releasecite{Development/WeilBound/ThetaLifts/EigenvalueEndgame/PlaneRemovalConfiguration.lean}{70}{Iwasawa decomposition and stability of matching}
\releasecite{Development/WeilBound/ThetaLifts/EigenvalueEndgame/TowerCodimensionTwo.lean}{188}{Propagation from the Levi}

We use three consequences of the mixed-model line expansion for the
other coefficient. When \(q=0<p\), the \(\beta\) coefficient
vanishes on all of \(\mathscr D\). The definite smaller source has
only the zero orbit, which gives its \(\alpha\) coefficient. The
larger source's \(\beta\) coefficient comes from the zero orbit and
is a lift from \(V_A\) to \(W_1\). Its target has real negative index
\(p-1\), and \((d-2)+2p<2a\). The real vanishing calculation below
applies to this convergent lift. When \(p=0<q\), the \(\alpha\)
coefficient vanishes on all of \(\mathscr D\). The definite larger
source has only the \(\beta\) coefficient, and the smaller source's
zero-orbit lift to \(W_1\) has target negative index \(q-1<q\).
Use the restricted support at \(v_1\) for this lift. If its target-only
evaluation does not already have that support, one uses a representative
of the upper test with the same integral whose evaluation has the
lower support. This is the local quotient support calculation at
\(v_1\). At this place the source character is nontrivial, and source
augmentation relations within the prescribed support integrate to zero.

In the split matrix chart at \(v_1\) a supported test is a Schwartz
function on the open full row-rank locus. In coordinates for the
removed plane the target-only evaluation is, up to its fixed nonzero
measure factor, the map
\[
 E_0\phi(Y_1)=\int_{F_{v_1}^n}\phi(0,Y_1,z)\,dz.
\]
The first coordinate is set to zero by the trivial Fourier character,
and the dual coordinate is integrated by the partial Fourier transform.
The image of full-rank tests is the Schwartz space on matrices \(Y_1\)
of rank at least \(n-1\). Indeed for every such \(Y_1\) one can choose
a column \(z\) completing it to rank \(n\). Small compact open boxes
about these matrices give product charts lifting test functions, and
a compact open partition handles every Schwartz function on this
locus. Source actions intertwine the evaluation, since the Jacobian
in \(z\) changes the source determinant power from \(d/2\) to
\((d-2)/2\). In the rank \(n-1\) quotient the source stabilizer has
arbitrary determinant. The nontrivial unitary source character on
a uniformizer rules out its twisted coinvariant by the rank argument
below. Hence a finite sum of source augmentation terms can be
subtracted from \(E_0\phi\) to obtain full-rank support. Lift every
term of that subtraction to the full-rank upper test space using
the surjectivity just proved. Subtract the corresponding source
augmentations above, which preserve the upper integral since the
tests and their translates lie in the supported subspace. The
adjusted upper test has the required lower support after evaluation.
\releasecite{Development/WeilBound/ThetaLifts/EigenvalueEndgame/EigenKilling/EigenKillSet.lean}{613}{Lowering the split support}

When \(p,q>0\), choose a plane removal on the larger source and
the target, keeping \(V_B\) fixed. For every \(k\in K\) the
restriction of the \(\alpha\) coefficient at \(h_1k\) to
\(h_1\in U(W_1)(\mathbb A_F)\) belongs to the lower difference
space, with ranks \(a-2,a'\) and indices \(p-1,q\). The nonzero
orbit on the larger source gives its cone evaluation. The cone
integral converges at \(e\), preserves polynomial real data through
its compact average, and sends a restricted tensor to a lower
restricted tensor times a scalar. The zero orbit on the smaller
source gives its target-only evaluation with the same support
representative just described. All these assertions concern separately
chosen tests on the two sides. The rank and weight inequalities
persist after this larger-source removal.
\releasecite{Development/WeilBound/ThetaLifts/EigenvalueEndgame/EigenKilling/AlphaZeroOf.lean}{50}{The two absent coefficients and the lower difference}
\releasecite{Development/WeilBound/ThetaLifts/EigenvalueEndgame/EigenKilling/SubMemKer.lean}{170}{Changing the representative within the smaller support}

For the assertion about the real cone, use the orthogonal symplectic
decomposition into the lower tensor space and its complement. The
adapted Cartan involutions preserve this decomposition. In compatible
Fock coordinates, the Gaussian is a product and each polynomial
vector is a finite sum of products of polynomials in the two sets of
variables. The removed-coordinate shift acts on the complement.
The mixed evaluation commutes with the lower Heisenberg action, so
on each product it preserves the lower vector up to the scalar of
the complementary test. Indeed a continuous operator commuting with
the real Schr\"odinger phase multipliers has distribution kernel equal
to a multiplication operator on the diagonal. Commuting also with
translations forces that multiplier to be constant.
The compact source orbit of the upper polynomial lies in the space
of polynomials of bounded degree times the Gaussian, which is finite
dimensional. All the reduced shifted vectors therefore lie in a
fixed finite-dimensional space of lower polynomials. Absolute
integration in the cone preserves this space. The finite local
integrals are Schwartz by compact support in the evaluated variables,
with the unchanged lower vacuum outside a finite set after its scalar
has been factored out. The absolute convergence of the idele product
at \(e\ge2\) permits summing the reduced tensors. This gives the lower
polynomial tests used in the induction.
\releasecite{Development/WeilBound/ThetaLifts/LocalTheory/ArchimedeanBumps/SlicePol.lean}{143}{Finite-dimensional lower polynomial space}

\paragraph{The real weight obstruction in the zero orbit.}
Consider a convergent theta integral from signature \((P,Q)\) to
\((R,S)\) with \(\min(R,S)<\min(P,Q)\). The determinant splitting
on the source has exponent \(m\), with \(|m|>R+S\). Fixing the
finite test and the target point gives a continuous functional on the
real Schwartz space, invariant under the real source. Restrict to the
Fock polynomials at one real place and fix the vectors at the other
places. For a rank-one trace-one matrix \(v w^{\mathrm t}\) in one
compact source block, the complexified infinitesimal action, after
the constant term is removed, is the derivation
\[
 D=\sum_{i,j,c}v_i w_j
     \begin{cases}
       x_{ic}\partial_{jc}, & c\text{ has the same sign as the block},\\
       -x_{jc}\partial_{ic}, & c\text{ has the opposite sign}.
     \end{cases}
\]
The shift removed from this operator is half the difference between
the numbers of same-sign and opposite-sign target columns. Source
invariance says that the functional intertwines \(D\) with the
scalar \(-m/2\) minus that shift. This scalar has sign opposite to
\(m\). If it is positive, choose the block with fewer same-sign
target columns. For an arbitrary linear polynomial \(\ell\), choose
\(w\ne0\) annihilating the coefficient columns of \(\ell\) of
that sign, and choose \(v\) with \(w^{\mathrm t}v=1\). Such \(w\)
exists because the source block has larger dimension. The action of
\(D\) on the space generated by the orbit of \(\ell\) has weights
\(0,-1\). The powers of \(\ell\) have only nonpositive weights, so
the functional kills them. For a negative required scalar,
choose the block with fewer opposite-sign columns and exchange the
roles of \(v,w\). The weights of the corresponding powers are then
nonnegative. Powers of linear forms span the homogeneous polynomials
by polarization. Density of the polynomial vectors in the real
Schwartz space gives the asserted vanishing. This calculation uses
the actual source splitting on each lower zero-orbit lift, whose
exponent is \(m_\alpha-m_\beta\) or its negative.
\releasecite{Development/WeilBound/ThetaLifts/EigenvalueEndgame/EigenKilling/FockActionFirstEq.lean}{78}{Rank-one action on powers of linear forms}
\releasecite{Development/WeilBound/ThetaLifts/EigenvalueEndgame/EigenKilling/RealVanishOf.lean}{60}{Passage to the invariant Schwartz functional}

\paragraph{Iteration of an \(\alpha\) constant term.}
Suppose \(q>0\) and the full line constant term of a difference
has only its \(\alpha\) coefficient. Witt's extension theorem and
left rational invariance move this property to any rational isotropic
line. When choosing a plane adapted to a new lattice frame, conjugate
the real Cartan involutions and right-translate the target argument
accordingly. On the source, a right change of variables in the
integral carries the Fock data to the conjugate Cartan model. The
finite support condition stays at the same place. Thus the resulting
function is a difference of the same kind and retains its one-character
line constant term. With \(p=0\), the \(\alpha\) coefficient is zero
by the calculation above. If \(p>0\), take the larger-source removal.
For each \(k\in K\) write \(f_{1,k}\) for the resulting lower
difference on \(U(W_1)\).

Choose a second rational isotropic line in \(W_1\). The two successive
averages equal the average over the radical of the parabolic of the flag
formed by the first line and their two-plane \(U\). That radical is
also \(N_U\) followed by the upper triangular unipotent of the
\(\mathrm{GL}_2(\mathbb A_E)\) factor in the Levi of the two-plane.
Both decompositions are semidirect products over \(F\), so Fubini
with probability quotient measures identifies the averages. Write
\(C(g)\) for this flag average. Fix all other data, and let \(t\)
be the idele equal to \(2\) at the complex places and \(1\) at finite
places. Let \(a_2\) act on the two-plane as \(\operatorname{diag}(t,t^{-1})\).
The upper \(\alpha\) law and the two lower laws imply
\[
 C(a_2^i g)=A(g)+\lambda^i B(g),\qquad
 \lambda=\alpha(t)\beta(t^{-1})|t|_E^{(a'-(a-2))/2},\qquad
 |\lambda|<1\quad(i\in\mathbb Z).
\]
This identity holds for every \(g\). Indeed, write
\(g=n m(s)j(h_1)k\) in the first Iwasawa decomposition. The outer
unipotent drops out of the first average, and the lower function is
\(\alpha(s)|s|_E^{a'/2}R(h_1)f_{1,k}\), with \(R\) denoting right
translation. Apply the two lower laws to its line average. The
\(\alpha\) factors at \(t,t^{-1}\) cancel and the two different
norm powers give the displayed \(\lambda\). In the rank-one argument
below the same result holds on every right translate, so
\(C(a_2g)=C(g)\). The relation \((\lambda-1)B(g)=0\) gives vanishing
of the lower \(\beta\) coefficient, at every argument. Thus each
\(f_{1,k}\) has only its \(\alpha\) coefficient. Induction on \(p\),
applying line cuspidality and orthogonality below at each stage,
makes \(f_{1,k}=0\). Iwasawa decomposition then gives a zero full
line constant term of the upper difference. The induction applies to
the entire difference space and needs no equality on lower matched
pairs for these auxiliary removals.
\releasecite{Development/WeilBound/ThetaLifts/EigenvalueEndgame/SignatureSplitting/TwoStepOf.lean}{22}{Two successive constant terms}
\releasecite{Development/WeilBound/ThetaLifts/EigenvalueEndgame/SignatureSplitting/SpanPairIff.lean}{207}{Rank-one restriction of the flag average}
\releasecite{Development/WeilBound/ThetaLifts/EigenvalueEndgame/TowerCodimensionTwo.lean}{115}{Induction on the larger-source index}

Here is the rank-one argument in a form that also applies to the
rank-two Fourier calculation below. Let \(F_2\) be a smooth,
finite-level function on \(\mathrm{SL}_2(E)\backslash\mathrm{SL}_2(\mathbb A_E)\)
of uniform moderate growth, meaning the same height exponent works
for all its real right derivatives, with constants depending on the
derivative. Suppose its upper unipotent average \(C_2\) obeys
\(C_2(a_2^i x)=A_2(x)+\lambda^i B_2(x)\) with \(0<|\lambda|<1\).
All right derivatives of the average have bounded growth on the sets
\(a_2^i K_2\) for \(i\ge0\) and compact \(K_2\). To see this,
differentiate the recurrence, or express its two coefficients in
terms of the values at \(1,a_2\); compactness bounds those values
uniformly on \(K_2\). Rank-one adelic reduction covers the quotient
by \(a_2^i K_2\), \(i\ge0\), for one compact \(K_2\). This is the
rank-one Siegel set, with the norm-one diagonal directions reduced
modulo units and the remaining positive radial direction discretized
by powers of \(2\).

For this compact \(K_2\) and the fixed finite level choose a compact
set of representatives of \(E\backslash\mathbb A_E\) with finite
parts in a sufficiently small compact open additive subgroup.
Additive approximation gives these representatives, with bounded
archimedean parts obtained modulo a fractional ideal. The finite
parts act trivially at the chosen level after conjugation by \(K_2\).
The remaining, archimedean unipotents conjugate across \(a_2^i\)
by multiplication with \(2^{-2i}\). Thus comparison with the
unipotent average by the mean value estimate bounds a derivative of
\(F_2\) at \(a_2^i k\) by a constant plus \(2^{-2i}\) times the
bound for its first further derivatives on \(a_2^i K_2'\), with
\(K_2'\) compact. The estimate holds for every compact \(K_2\) and
all derivatives, since real differentiation preserves the finite
level. If every derivative is \(O(2^{Mi})\), this improves their
common exponent to \(\max(0,M-2)\). After finitely many applications
\(F_2\) is bounded. Its average is then bounded for negative powers
also, forcing \(B_2=0\). With \(C_2=0\) the same argument gives
boundedness directly. In the application, \(F_2\) is the restriction
of the \(N_U\)-average of a fixed right translate of \(f\) to the
embedded \(\mathrm{SL}_2(\mathbb A_E)\). The average is invariant
under \(\mathrm{SL}_2(E)\), which normalizes the rational unipotent.
The theta estimates give uniform moderate growth because for fixed
finite data the exponent is independent of the real Schwartz vector,
and real derivatives act on that vector. The compact unipotent
integral and the algebraic group embedding preserve this condition.
\releasecite{Development/WeilBound/ThetaLifts/EigenvalueEndgame/SignatureSplitting/ConvCont.lean}{140}{Growth reduction for the rank-one unipotent average}
\releasecite{Development/WeilBound/ThetaLifts/EigenvalueEndgame/SignatureSplitting/AeFacts.lean}{18}{The adelic rank-one reduction}
\releasecite{Development/WeilBound/ThetaLifts/FrameEstimates/CuspRapid.lean}{30}{Uniform moderate growth of the tested theta integrals}

\subsubsection{Line and parabolic constant terms}
Let \(f\in\mathscr D\) have zero constant term along one rational
isotropic line of \(W\). Witt's theorem and rational invariance
give the same vanishing for every line. For a rational isotropic
\(j\)-plane \(U\), write \(f_U\) for the constant term along the
parabolic stabilizing \(U\). Choose a dual frame for \(U\), and
write \(m_U(s)\) for the resulting embedding of
\(\mathrm{GL}_j(\mathbb A_E)\) in the Levi. Proceed by induction on
\(j\). For each target argument \(h\) the function
\(s\mapsto f_U(m_U(s)h)\) is left \(\mathrm{GL}_j(E)\)-invariant.
Its constant terms along the proper maximal parabolics stabilizing
prefixes of the chosen frame are zero. Indeed the average over
\(N_U\) and then the indicated rectangular upper unipotent is the
average over the radical of the flag parabolic. One can take the
average for the smaller isotropic plane first in the same radical.
All quotient measures here have mass one. Consequently the lower
dimensional vanishing supplies all those constant terms.

An invariant function on \(\mathrm{GL}_j(\mathbb A_E)\) with these
zero averages is determined by its generic Whittaker integrals on
all right translates. For this implication, expand successively along
the upper unipotent column groups of dimensions \(j-1,j-2,\ldots,1\).
At each stage rational \(\mathrm{GL}_r(E)\) acts transitively on
nonzero characters of the vector quotient \((E\backslash\mathbf
A_E)^r\). Fix the representative nonzero on its last entry. The
zero Fourier term at stage \(r\) is an integral of the constant term
for the rectangle \((r,j-r)\), with the already selected characters
on the remaining upper triangular group. It vanishes by the previous
averages. Thus if the last integral, with a fixed nontrivial character
on successive simple roots, is zero on every translate, Fourier
uniqueness on each compact vector quotient gives vanishing backwards
through the expansions. This uses only uniqueness of Fourier
coefficients of continuous functions on compact quotients.
\releasecite{Development/WeilBound/ThetaLifts/EigenvalueEndgame/CuspidalAverages/CuspOfLineOfKeys.lean}{76}{Whittaker testing of constant terms}

We show that the required Whittaker integrals are zero on \(f_U\).
Choose a finite place \(v\) split in \(E\) and write \(k=F_v\),
\(\varpi\) for a uniformizer and \(Q_v\) for the residue cardinality.
At a chosen factor of \(E_v=k\times k\), choose the split
Schr\"odinger chart with source rank \(n\) and matrix variable
\(Y\in\operatorname{Mat}_{n,d}(k)\). A local functional \(\ell\)
is obtained by varying the \(v\)-component of a fixed tensor and
taking its source theta integral, the \(N_U\)-average and the
Whittaker integral. The global integrals are linear on this test
space, and the unipotent integrations are compact. For the smaller
source take \(v\ne v_1\), so this variation keeps the prescribed
support. In the split chart the source and target actions are
\[
 \omega(g,1)\phi(Y)=\chi_W(\det g)|\det g|_v^{d/2}\phi(g^{\mathrm t}Y),
 \qquad
 \omega(1,h)\phi(Y)=\chi_X(\det h)|\det h|_v^{n/2}\phi(Yh),
\]
where \(\chi_X,\chi_W\) are the two characters of this source in the
chosen splitting. Invariance under the source and covariance of the
Whittaker integral give the laws
\begin{align*}
 \ell(\phi(AY))&=\xi(\det A)\ell(\phi),
 &&\xi=(\chi_W|\cdot|_v^{d/2})^{-1},\\
 \ell(\phi(Y(1+e_i u^{\mathrm t})))&=\psi_v(u_{i+1})\ell(\phi),
 && u_l=0\ \text{for }l\le i,\quad 1\le i<j.
\end{align*}
Here the frame occupies the first \(j\) columns. A column shear has
determinant one, preserves \(U\), is trivial on \(U^\perp/U\), and
induces the indicated upper triangular unipotent on the frame. On the
other split component it acts by inverse transpose, which fixes the
dual-component part of \(U\). This proves the displayed laws also
when the shear changes columns outside the frame. Conjugating to an
arbitrary rational frame gives the same laws in its coordinates.
\releasecite{Development/WeilBound/ThetaLifts/EigenvalueEndgame/CuspidalAverages/ActsOfReads.lean}{34}{Split actions and determinant powers}
\releasecite{Development/WeilBound/ThetaLifts/EigenvalueEndgame/CuspidalAverages/ActsOfReads.lean}{481}{The global functional in a split chart}

Write \(y_i\) for the columns of \(Y\). The first shear law forces
\(\ell\) to be supported on dependence of \(y_1,y_2\). At a matrix
with these columns independent, take \(w\) with \(w(y_1)=0\) and
\(w(y_2)=1\), and choose \(t\) with \(\psi_v(t)\ne1\).
The operation of left multiplication by \(1-t y_1w\) combined
with right multiplication by \(1+t e_1(wY)\) fixes this matrix and
has character \(\psi_v(t)\) on the functional. The matrices in the
operation are fixed at the chosen point. Both multiplications are
unipotent and commute. Their combined linear action preserves a
lattice, as is seen by taking a basis adapted to the resulting
nilpotent operator. Arbitrarily small compact open balls about the
point are therefore preserved. The value of the functional on each
such ball is zero. A finite disjoint compact open refinement on the
support gives the claimed vanishing away from the dependence locus.
Equivalently, the functional factors through restriction to this
closed locus. For locally constant compactly supported functions,
restriction to a closed subspace is onto by compact open extension,
and its kernel consists of functions supported in the complement.

When \(j\ge3\), on the dependence locus with \(y_1\ne0\) write
\(y_2=s y_1\). Adding \(t y_2-s t y_1\) to the third column fixes
the matrix and has character \(\psi_v(t)\) by the second shear law.
The compact open argument kills this stratum. On the remaining
closed locus \(y_1=0\), adding \(t y_1\) to the second column fixes
every point and again has nontrivial character. The functional is
zero. Hence the rank at least three Whittaker integrals of each
theta image vanish separately.
\releasecite{Development/WeilBound/ThetaLifts/EigenvalueEndgame/CuspidalAverages/LocalFieldT2Space.lean}{196}{The adjacent simple-root calculation}

For \(j=2\) choose \(v\) outside all the finite supports of the
tests, the frame denominators and the nonintegral places of the given
argument \(h\). Take it split with unramified splitting characters.
There are infinitely many such places. The local vacuum is a scalar
multiple of the integral matrix characteristic function. In the
chosen chart the frame change and the \(v\)-component of \(h\) are
integral with integral inverses, so right multiplication by them
preserves that function. The element \(b_i=\operatorname{diag}
(\varpi^i,\varpi^{-i})\) at the chosen factor of \(E_v\), with all
other components equal to one, has determinant one in
\(\mathrm{GL}_2(\mathbb A_E)\). Its action on the vacuum moves the
column bounds to \((-i,i)\), with no Weil determinant factor.

Here is the calculation for these bounds. Let \(B(p_0,u_0)\) be
the value of \(\ell\) on the box with columns one and two in
\(\varpi^{p_0}\mathcal O_v^n\) and \(\varpi^{u_0}\mathcal O_v^n\),
respectively, and all other columns integral. Put
\(D(p_0,u_0)=B(p_0,u_0)-B(p_0+1,u_0)\) and
\(R=\xi(\varpi)Q_v^{n-d}\). Let \(c_v\) be the least integer
such that \(\psi_v\) is trivial on \(\varpi^{c_v}\mathcal O_v\).
Then
\begin{align*}
 B(p_0,u_0)&=0 && (u_0-p_0<c_v),\\
 D(p_0,u_0)&=R D(p_0+1,u_0) && (u_0-p_0-1\ge c_v),\\
 D(p_0,u_0)&=Q_v D(p_0,u_0+1) && (u_0-p_0\ge c_v).
\end{align*}
For the first equality take a shear of column two by column one
with parameter in \(\varpi^{c_v-1}\mathcal O_v\) and nontrivial
character. It preserves the box. For the other equalities partition
the first-column shell by the first row whose valuation is \(p_0\).
The corresponding test with all other entries of the first column
at depth \(p_0+1\) occurs with weight \(Q_v^{n-i}\) for pivot row
\(i\), by row shears distributing the later entries among the residue
classes. These shears preserve all the other column bounds, and
have determinant one. Scale the pivot row by \(\varpi\). In the
new box the pivot is at \(p_0+1\), the other first-column entries
still have depth \(p_0+1\), and the other entries of the pivot row
each gain a depth. Restricting the other first-column entries to
depth \(p_0+2\) gives a factor \(Q_v^{n-1}\) by the same row shear
count. Loosening the extra depth in the pivot row in the other
columns contributes \(Q_v^{-(d-1)}\) in this expression, by the
column shears. The shear affecting column two has a parameter of
depth at least \(u_0-p_0-1\), so its character is one. Together
with the determinant factor this proves the second equality. To
prove the third, loosen the column-two entry of the pivot row from
depth \(u_0+1\) to \(u_0\) using the column shear with parameter
of depth \(u_0-p_0\). This contributes \(Q_v\). On the dependence
locus the other column-two entries automatically have depth at
least \(u_0+1\) because the other first-column entries have depth
at least \(p_0+1\). Thus the resulting pivot test agrees with the
one in \(D(p_0,u_0)\) on the support of \(\ell\).

At \(u_0-p_0=c_v\) one has \(D(p_0,u_0)=B(p_0,u_0)\). The last
two recurrences give
\(D(u_0-c_v,u_0)=Q_vR D(u_0+1-c_v,u_0+1)\).
Sum the shells, using the second recurrence at each depth down
from \(u_0-c_v\). The resulting value at \((-i,i)\) is a scalar
multiple of
\[
 y_i(R)=(Q_vR)^{-i}\sum_{j=0}^{2i-c_v}R^j,
\]
where the sum is zero for \(2i<c_v\). For the two sources the
parameters are \(r_v=\nu_v(\varpi)Q_v^e\) and \(r_v^{-1}\),
where \(\nu=\alpha/\beta\). The relation
\(y_i(r_v^{-1})=r_v^{c_v}y_i(r_v)\) holds by reversing the sum,
and \(|r_v|=Q_v^e>Q_v\). Hence \(|y_i(r_v)|\) tends to infinity.
The same profile holds for finite sums of pure tensors and for
their difference, since the same split place and frame can be
used. In particular the Whittaker coefficient along \(b_i h\) is
\(C_h y_i(r_v)\), for a scalar \(C_h\) independent of \(i\).
\releasecite{Development/WeilBound/ThetaLifts/EigenvalueEndgame/CuspidalAverages/ColOpSingle.lean}{160}{The two shell recurrences and their boundary}
\releasecite{Development/WeilBound/ThetaLifts/EigenvalueEndgame/CuspidalAverages/ActsOfReads.lean}{629}{The reciprocal parameters at the same split place}

On the other hand \(s\mapsto f_U(m_U(s)h)\), restricted to
\(\mathrm{SL}_2(\mathbb A_E)\), has zero upper unipotent average
by the line case and the same flag calculation. It is bounded on
the \(\mathrm{SL}_2(E)\) quotient by the rank-one argument above,
which applies since it is smooth of uniform moderate growth and
finite level. Its Whittaker integrals have absolute value bounded
by the same bound, with mass-one unipotent measure. It follows that
\(C_h=0\), and taking \(b_0=1\) gives the required vanishing at
\(h\). This completes the induction on \(j\). Every proper
rational parabolic of \(U(W)\) stabilizes a nontrivial isotropic
flag. Its constant term is an iterated average including one of
these plane averages. Thus \(f\) is cuspidal.
\releasecite{Development/WeilBound/ThetaLifts/EigenvalueEndgame/CuspidalAverages/CuspOfLineOfKeys.lean}{169}{Induction on the Witt index of the averaging plane}

\subsubsection{Orthogonality of cuspidal differences}
Let \(f\in\mathscr D\) be cuspidal. It is smooth, of finite level,
and of uniform moderate growth. The last assertion follows by
applying the theta kernel estimates to derivatives of the real
Schwartz vector with its finite tensor unchanged. Cuspidality and
uniform moderate growth give rapid decrease on Siegel sets. In the
usual proof, a zero unipotent mean bounds the value by derivatives
along the radical by integration on its compact quotient. Conjugating
these derivatives across the positive radial part of a Siegel set
gains an inverse power of a large simple-root coordinate. Choose the
radical of the maximal parabolic for that coordinate. Each root of
this radical contains that simple root. Applying the same estimate
to the derivatives and using their uniform moderate growth gives any
inverse power. The finite level lets one take the finite unipotent
coordinates in a compact open set acting trivially. The bounded
parts of the Siegel sets and the finitely many cusp representatives
only change the constants. Thus the rapid decrease is a bound in
the quotient height, the infimum of the adelic heights of rational
translates. This also gives integrability after multiplying by an
automorphic function with a fixed polynomial height bound.
\releasecite{Development/WeilBound/ThetaLifts/FrameEstimates/CuspRapid.lean}{24}{Growth and rapid decrease on the triangle}

For \(X=V_A\) or \(V_B\), put \(H_X=U(X)\) and define
\[
 A_X(f;\Psi)(g)=\int_{[G]} f(h)\overline{\theta_X(\Psi;g,h)}\,dh.
\]
For fixed \(g\) the kernel has polynomial height growth in \(h\),
locally uniformly for \(g\) in compact sets. Rapid decrease gives
absolute convergence, continuity and rational left invariance of
\(A_X\). The theta bounds for the two tested source integrals give
a polynomial height bound on the source integral of the absolute
kernel. Thus Fubini yields, with the given source Haar measure,
\[
  \langle T_X(\Psi),f\rangle_{[G]}
   =\int_{[H_X]} \overline{A_X(f;\Psi)(g)}\,dg
\]
for the tested \(\Psi\). We will show that for every Schwartz test
the adjoint function is constant on \([H_X]\).
\releasecite{Development/WeilBound/ThetaLifts/EigenvalueEndgame/OrthogonalDecomposition/RapidBound.lean}{118}{Fubini for the adjoint pairings}

Work at a finite place \(v\ne v_1\), outside the ramification set,
whose underlying rational prime splits completely in \(E\). The
value \(\nu_v(\varpi)\) at either split factor is different from
\(1\). In fact, for the corresponding degree-one prime \(\mathfrak w\)
of \(E\) choose \(b\) with \(\mathfrak w^h=(b)\). The (2h)-th power
of the unramified value, by the Hecke character relation on \(b\),
is, up to inversion,
\(\prod_\tau(\tau(b)/\overline{\tau(b)})^{m_\alpha(\tau)-m_\beta(\tau)}\).
All embeddings are in the Galois field \(E\). The conjugate prime
ideals of \(\mathfrak w\) are distinct, and all the displayed
exponents are nonzero. Taking a valuation at a prime in this product
rules out a root of unity. This argument applies at every such
place, even though the coefficients of \(f\) need not be vacuum
vectors there.
\releasecite{Development/WeilBound/ThetaLifts/LatticeBasics/GoodNe.lean}{42}{Nontrivial uniformizer values at degree-one places}

We use two facts about the split matrix model. First, on
\(\operatorname{Mat}_{N,d}(F_v)\) a functional with right covariance
\(\ell(\phi(\cdot h))=\eta(\det h)\ell(\phi)\), where
\(\eta(\varpi)\notin\mathbb R_{>0}\), vanishes on each stratum of
rank less than \(d\). For a matrix of rank \(r\), the orbit of
right translation is the space of surjections onto its image, with
stabilizer containing a \(\mathrm{GL}_{d-r}\) on its kernel. The
stabilizer therefore contains an element of determinant \(\varpi\).
On a transitive orbit, a relatively invariant distribution is
possible only if the character on the stabilizer equals a ratio
of positive Haar modular characters. One proves this by pulling the
distribution back to the group by averaging test functions along
stabilizer fibers, where the only relatively invariant functional is
weighted Haar integration. On the rank stratum the image varies in
a Grassmannian, and local choices of a frame give product charts
with the orbit as fiber. Compactly supported locally constant test
functions in each chart are finite sums of products. The stabilizer
criterion therefore applies on the whole stratum. The filtration by
closed rank bounds proves the vanishing, using restriction of
Schwartz functions to the closed stratum. On square matrices the
only surviving contribution is the open orbit \(\mathrm{GL}_d\).
Its functional is a multiple of Haar integration weighted by
\(\eta(\det h)^{-1}\). In particular left multiplication on the
test variable has the same covariance as right multiplication.
\releasecite{Development/WeilBound/ThetaLifts/LatticeBasics/LocDistE.lean}{77}{The rank filtration for the local distribution}

Second, the analogous source coinvariant on \(n\)-by-\(d\) matrices
is represented by full row-rank tests. On the other strata the left
stabilizer contains a determinant-\(\varpi\) element, and the same
argument applies to \(\xi=(\chi_W|\cdot|_v^{d/2})^{-1}\).
On the full-rank stratum, integrating the test over the fibers of
the row-space map with weight \(\xi^{-1}\) gives the coinvariant.
Indeed each fiber is freely transitive under \(\mathrm{GL}_n(F_v)\),
and on compact open charts weighted Haar integration is the
one-dimensional twisted coinvariant on that fiber. Local sections
of the row-space map and a compact open partition lift every smooth
section over the Grassmannian. With \(Y_0=(I_n\;0)\), the resulting
map in the target group model is
\[
  \phi\longmapsto
  \left[h\longmapsto \chi_X(\det h)|\det h|_v^{n/2}
        \int_{\mathrm{GL}_n(F_v)}
         \phi(A Y_0 h)\xi(\det A)^{-1}\,dA\right].
\]
For a block lower triangular element with diagonal blocks \((A,D)\)
the multiplier is
\[
 (\chi_X/\chi_W)(\det A)\chi_X(\det D)
 |\det A|_v^{-(d-n)/2}|\det D|_v^{n/2}.
\]
The two powers form the square root of the parabolic modulus. The
two sources thus give the normalized inductions with unitary
determinant characters \(\alpha\) on the block of size \(a\) and
\(\beta\) on the block of size \(a'\), in the two block orders.
Each inducing representation is irreducible and unitary, so
Bernstein's irreducibility theorem applies. The two irreducible
inductions have the same class by commutativity of normalized
parabolic induction in the Grothendieck ring of the general linear
groups. Denote their common local module by \(P_v\).
\cite[\S8.2, Corollary (a)]{p2bernstein-unitary-induction}
\cite[Theorem~1.9]{p2zelevinsky-induction-ring}
\releasecite{Development/WeilBound/ThetaLifts/EigenvalueEndgame/OrthogonalDecomposition/LocDecompOf.lean}{151}{Using the common split quotient in a cuspidal difference}

For this place fix the other tensor factors in a finite expression
for \(f\) as a difference. The local orbit maps of the two source
integrals factor through \(P_v\) by source invariance. Each image
is simple or zero, and their finite sum is a semisimple isotypic
module. The submodule generated by the cuspidal vector \(f\) is
therefore a finite sum of copies of \(P_v\) inside the space of
rapidly decreasing functions. Consequently \(f\) is a finite sum
of values of local equivariant maps from the Schwartz space for
either source, each with rapidly decreasing outputs. This step
allows one to form the adjoint pairing with the opposite source
locally for every varied test, with convergence for each of those
pairings.

For \(A_X\) use the equivariant maps from the other source \(X'\),
of rank \(n'=d-n\). Fixing all places of the adjoint test except
\(v\), the local pairing is bilinear after conjugating the test
for \(X\). Stack the \(n\)-by-\(d\) and \(n'\)-by-\(d\) matrices
to view it as a functional on square matrices. Simultaneous target
right translation preserves the pairing. The Weil target scalars
on the two tests are \(\overline{\chi_X}|\cdot|_v^{n/2}\) and
\(\chi_{X'}|\cdot|_v^{n'/2}\). After removing these scalars, the
functional has right covariance
\(\eta=(\chi_X/\chi_{X'})|\cdot|_v^{-d/2}\).
Since \(\chi_X/\chi_{X'}=\chi_W\) is unitary and nontrivial on the
uniformizer, \(\eta(\varpi)\notin\mathbb R_{>0}\). By the square
matrix calculation the same left covariance holds on the first
block, the other block fixed. This cancels exactly the conjugated
source factor \(\overline{\chi_W}|\cdot|_v^{d/2}\). Right
translation of \(g\) by any element of \(H_X(F_v)\) consequently
leaves \(A_X\) unchanged. The square matrix calculation uses
exactly \(d\) rows from the two blocks together.
\releasecite{Development/WeilBound/ThetaLifts/EigenvalueEndgame/OrthogonalDecomposition/PairLocOf.lean}{26}{Cancellation of the adjoint Weil factors}

The rational group and the groups at these degree-one places
generate a dense subgroup of \(H_X(\mathbb A_F)\). On
\(\mathrm{SU}(X)\), strong approximation off one such place applies
because the group is simply connected, absolutely almost simple,
and split and noncompact there. For the determinant, use a section
of the norm-one torus \(S\) obtained by scaling a rational
anisotropic line in \(X\). The compact group \([S]\) is generated
densely by the same local determinant groups. Otherwise there
would be a nontrivial continuous character trivial on them. Its
pullback by adelic Hilbert~90, \(x\mapsto x/\bar x\), would be a
nontrivial unitary Hecke character \(\omega\) of \(E\) trivial at
the degree-one places outside a finite set. The quotient
\(\zeta_E(s)/L(s,\omega)\) then has an absolutely convergent,
nonzero Euler product near \(s=1\), since the remaining infinite
set consists of primes of norm at least the square of their
rational prime. This contradicts the simple pole of \(\zeta_E\)
and holomorphy of \(L(s,\omega)\) there. This proves
density on the determinant as well. Local right invariance implies
local left invariance because conjugation by a fixed adelic element
preserves each single-place subgroup. Continuity and rational left
invariance now make \(A_X(f;\Psi)\) constant.
\cite[Theorem~7.12]{p2platonov-strong-approximation}
\cite[\S4]{tate-thesis}
\releasecite{Development/WeilBound/ThetaLifts/LatticeBasics/ConstOfGood.lean}{114}{Strong approximation and the norm-one determinant}
\releasecite{Development/WeilBound/ThetaLifts/LatticeBasics/ConstIdele.lean}{134}{The degree-one determinant argument}

\paragraph{Eliminating the adjoint pairings.}
Suppose first that \(X\) is isotropic. The rank-one constant-term
formula for a theta lift of a cusp form shows that the constant term
of \(A_X(f;\Psi)\) along an isotropic line of \(X\) transforms under
its idele torus by \(\overline{\chi_W(t)}|t|_E^{d/2}\), in the
conjugate Weil splitting. In
the mixed model the central unipotent integral retains isotropic
vectors of \(W\). The nonzero orbit unfolds against the constant
term of \(f\) along the stabilizer of an isotropic line, and its
integral vanishes by cuspidality. The zero vector gives evaluation
at zero in the removed coordinates, on which the dilation has the
indicated weight. This is the computation of the line case in
Wu's Proposition~3.1, with the two spaces interchanged and the
conjugate Weil representation used for the test. Its proof uses
cuspidality and rapid decrease of the input to unfold the nonzero
orbit, so it applies directly to the cuspidal difference. A constant
function equals its mass-one line average and obeys this law only
if zero. Thus \(A_X=0\), and the adjoint identity gives
\(\langle T_X(\Psi),f\rangle=0\) on the tested data.
\cite[Proposition~3.1 and (3.4)]{p2wu-unitary-tower}

Suppose instead that \(X\) is definite, in a nondefinite triangle.
Its adelic quotient is compact. Suppose a tested vector \(\Psi_2\)
has \(\langle T_X(\Psi_2),f\rangle\ne0\). Then the constant
\(A_X(f;\Psi_2)(1)\) is nonzero by Fubini. On the doubled group
\(U(W\oplus(-W))\) use the Siegel--Weil section obtained from
\(\Psi_1\otimes\overline{\Psi_2}\), with arbitrary polynomial
\(\Psi_1\). Its inducing character is \(\chi_X\) and its normalized
parameter is \(s_0=(n-d)/2<0\). The tensor product of the original
Weil model and its conjugate cancels the source splitting character
and identifies with the restriction of the doubled model after
multiplication by \(\chi_X(\det h')^{-1}\) on the second target
factor. Products of the theta sums therefore give the doubled
kernel, with this determinant correction, at \(\iota(h,h')\).
The anisotropic Siegel--Weil formula in the unitary case \(n<d\)
gives a holomorphic value at \(s_0\), equal to twice the doubled
theta integral with source quotient measure of mass one. If
\(v_X=\operatorname{vol}([H_X])\) for the source measure in
\(T_X\), this gives
\[
 \begin{split}
 J_{s_0}(h)&=\frac{2}{v_X}A_X(f;\Psi_2)(1)T_X(\Psi_1)(h),\\
 J_s(h)&=\int_{[G]} E(s,\iota(h,h'))\,
                  \chi_X(\det h')^{-1} f(h')\,dh'.
 \end{split}
\]
The identification uses the compatible doubled splitting.
Compactness of \([H_X]\), polynomial kernel bounds in the other
variable and rapid decrease justify both integrations in the
identity. The Eisenstein series is the usual flat extension of the
section at \(s_0\). Its polynomial real test is finite under the
doubled maximal compact in the corresponding Fock model, and its
finite data have finite level.
\cite[\S3.1, Theorem~7.1 and \S11.3]{gqt}
\cite[Theorems~2.1 and 2.2]{p2yamana-singular-siegel-weil}
\releasecite{Development/WeilBound/ThetaLifts/LatticeBasics/OrthAn.lean}{73}{The definite adjoint identity and its measures}

The functions \(J_s\) have zero constant terms. To see this first
in the half-plane of absolute convergence, unfold the Siegel sum
against \(f\). An isotropic maximal subspace of \(W\oplus(-W)\)
transverse to the two copies of \(W\) is the graph of an isometry
and gives the open orbit. For any other subspace \(M\), the
intersection \(U=M\cap(0\oplus W)\) is a nonzero isotropic
subspace. Indeed the two intersections with the coordinate copies
have equal dimension, since isotropy makes each projection into
the perpendicular of the intersection on its side. The second
projection is \(U^\perp\), by the same dimension count. The
unipotent radical stabilizing \(U\) acts trivially on
\(U^\perp/U\) and therefore preserves \(M\). Its determinant on
\(M\) is one. Its average in the unfolding is thus the zero
parabolic constant term of \(f\). The open orbit gives
\[
 J_s(h)=\int_{G(\mathbb A_F)} F_s(\iota(1,z))\,
            \chi_X(\det z)^{-1}f(hz)\,dz,
\]
with \(F_s\) the inducing section. Averaging in \(h\) gives zero
along each proper parabolic, again by cuspidality of \(f\).
The section on the open orbit is bounded in a right half-plane by
a large negative power of the matrix height of \(z\). In the
graph coordinates, the Siegel height is the inverse of the
Pl\"ucker norm of the graph. Its local factors are comparable to
the products of the larger of \(1\) and each Cartan singular value,
using both components at split places. These products bound the
matrix norms of \(z\) and \(z^{-1}\). Cartan integration on each
local group gives polynomial volume growth for these heights, and
the finite factors converge once the power is sufficiently large.
The closed-orbit interchanges follow by majorizing the defining
series with the positive series of the real part of the parameter
and pairing its polynomial growth with rapid decrease.

Meromorphic continuation of Eisenstein series with values in the
space of smooth functions of uniform moderate growth continues
each constant term of \(J_s\). Pairing with the rapidly decreasing
\(f\) after a restriction along \(\iota(h,h')\) is a continuous
linear functional on each fixed growth space. This remains true
after the additional compact unipotent average in \(h\). The
constant terms are consequently zero also at the holomorphic
point \(s_0\). The argument works with the actual (possibly
non-quasi-split) unitary group in the two copies of the doubling
embedding. It uses the Siegel parabolic of the split doubled space.
\cite[Theorem~2.3(1)]{p2bernstein-lapid-eisenstein}
\releasecite{Development/WeilBound/ThetaLifts/LatticeBasics/UnfoldC.lean}{109}{The open and nonopen doubling orbits}
\releasecite{Development/WeilBound/ThetaLifts/LatticeBasics/CuspLift.lean}{52}{Continuation of the zero parabolic means}

The nonzero adjoint constant would thus force every
\(T_X(\Psi_1)\) in the polynomial image to be cuspidal. Remove
a plane from the other (isotropic) source and the target, leaving
\(X\) unchanged. Assume the strong induction for the triangle
identity has supplied a nonzero lower polynomial lift for this
definite source. This is a smaller value of \(p+q\). The target-only
evaluation is onto on polynomial data and gives the sole line
coefficient of a definite-source integral, with a nonzero geometric
factor. A preimage of the lower nonzero lift therefore has nonzero
line constant term above, a contradiction. This use of nonvanishing
works at the \(p=0<q\) endpoint of the auxiliary induction on \(p\)
as well. Its further plane removal still has smaller total index
than the triangle under consideration. It follows that the
pairings to the definite-source image vanish on the tested data.
\releasecite{Development/WeilBound/ThetaLifts/LatticeBasics/OrthAn.lean}{333}{Contradiction using the lower definite-source lift}
\releasecite{Development/WeilBound/ThetaLifts/ConeHypotheses/PolDense.lean}{22}{The lower-stage nonvanishing premise available in the induction}

\paragraph{Conclusion of the comparison.}
The two pairings of a cuspidal difference with its defining lifts
vanish, so its squared Petersson norm is zero. Smoothness and the
full support of the quotient measure give pointwise vanishing.
Thus equality of the selected line coefficient on the embedded
Levi for all matched pairs implies equality of the two theta
integrals with that scalar on every matched polynomial pair. The
converse is obtained by taking the line average and projecting to
the selected character. The proof handles the remaining
constant-term contributions on the full space of differences
before pairing in the cuspidal subspace. In particular the final
difference has zero residual and continuous projections in
\(L^2([G])\). The nonzero lower coefficient specifies the common scalar with the
normalization of the cone calculation. To pass to the matched
Schwartz pairs, approximate their real tests by simultaneously
matched polynomial pairs with the finite tensors fixed, retaining
the smaller-source support. The real Schwartz continuity of both
integrals for those finite tensors gives the same identity in the
limit, pointwise on \([G]\).
\releasecite{Development/WeilBound/ThetaLifts/EigenvalueEndgame/TowerCodimensionTwo.lean}{73}{The positive norm conclusion}
\releasecite{Development/WeilBound/ThetaLifts/ConeHypotheses/CtSmallCone.lean}{176}{Passage to the matched Schwartz class}

\teamanchor{proof.period-square-regularization}
\subsubsection{Regularization and the scalar of the square}
\label{sec:period-square-regularization}
The comparison on supported smaller-source tests extends to
arbitrary finite tensors through a finite sum of source
translates with nonzero total coefficient. For any finite
collection of tensors and source translations, choose an
additional completely split place \(v\) outside their
exceptional places and the ramification set. In the chosen
split Schr\"odinger chart the smaller source rank is \(n=a'\).
The finite vacuum is a scalar times
\(\mathbf1_{\operatorname{Mat}_{n,d}(\mathcal O_v)}\), and the
translations at \(v\) preserve this vacuum. Let \(q_v\) be the
residue cardinality and let \(U\) range over the subspaces of
\(k_v^n\), with \(L_U\) their inverse images in
\(\mathcal O_v^n\). For codimension \(c\), choose \(g_U\) with
\(g_U^{\mathrm t}L_U=\mathcal O_v^n\) and
\(\det(g_U)\) of valuation \(-c\). Under the source Weil
action, the translated indicator is
\(\mathbf1_{\{Y:\operatorname{im}Y\subset L_U\}}\) times
\(\chi_W(\varpi_v)^{-c}q_v^{cd/2}\), where \(\chi_W\)
is the unramified source-action character.

Multiply that translate by
\((-1)^c q_v^{c(c-1)/2}\chi_W(\varpi_v)^c q_v^{-cd/2}\)
and sum over \(U\). The lattice of subspaces has M\"obius
function \((-1)^c q_v^{c(c-1)/2}\) on an interval of
codimension \(c\). This follows by summing these values over
the interval using the Gaussian binomial identity, whose
recurrence counts subspaces according to their intersection
with a fixed hyperplane. Summing over \(U\) containing the
image of \(Y\) modulo \(\varpi_v\) therefore gives the
indicator that this image has dimension \(n\). The sum
\(Z_v\) of source action operators has augmentation
\[
 d(Z_v)=\prod_{j=0}^{n-1}
          \bigl(1-\chi_W(\varpi_v)q_v^{j-d/2}\bigr)\ne0.
\]
The inequality \(n<d/2\) and unitarity of \(\chi_W\) give
nonvanishing. This produces the required smaller-source
support and hence the Schwartz continuity and integrability
of the regularized integral. In a translate of any one of the
finite tensors, the \(v\)-component of the original
translation can be dropped because it fixes the vacuum.
The other components commute with \(Z_v\).
\releasecite{Development/AutomorphicArithmetic/ThetaCorrespondence/ThetaPairLifts/LocalTypeDictionary/ExistsCombinationMemRegularSubspace.lean}{58}{The finite operator on the actual tensor model}
\releasecite{Development/AutomorphicArithmetic/ThetaCorrespondence/ThetaPairLifts/LocalTypeDictionary/TwistRegulariserIsGood.lean}{17}{Simultaneous source translations}.

In the common target quotient, applying \(Z_v\) to the
smaller-source vector multiplies its image by \(d(Z_v)\)
because each summand is a source translate. For a matched
pair in the same quotient, the normalized regularized
integral is \(d(Z_v)^{-1} I_B(Z_v\Phi_B)\) in the chosen
measure. The factor from the supported comparison is
independent of the auxiliary split place used for this
calculation. Apply both finite operators at once to a
test with nonzero common quotient and nonzero image under
the larger-source integral. Their augmentations are nonzero.
The resulting test is supported at both places, so evaluating
the two equalities on this pair gives the same scalar. This
uses the common quotient maps for the two supported realizations.
If the original smaller-source integral already converges
absolutely, changing variables by each source translate
gives \(I_B(Z_v\Phi_B)=d(Z_v)I_B(\Phi_B)\) directly.
The integrals of finite source augmentation differences are
zero after the same operation, by the preceding commutation
and change of variables. This proves the regularized identity
from the supported comparison whenever its two quotient maps
and its common nonzero image are available.
\releasecite{Development/AutomorphicArithmetic/ThetaCorrespondence/ThetaPairLifts/LocalTypeDictionary/RegularisedIntegralComparison.lean}{22}{Augmentation and comparison}
\releasecite{Development/AutomorphicArithmetic/ThetaCorrespondence/ThetaPairLifts/LocalTypeDictionary/RegConstEqAwayUnitMulFactor.lean}{22}{Agreement with the convergent integral}.

The relative scalar from changing the reference sections of
the triangle to the actual square can be checked before
factoring the triangle \(L\)-values. For a fixed square let
\(q_A,q_B\) be its maps to the same irreducible real polynomial
module, with the two infinity types fixed across the tuples.
For the two legs, let \(r_{A,\boldsymbol\mu},r_{B,\boldsymbol\mu}\)
be the finite maps in the product quotient. The images of
the two real and finite pairs are
\(q_A(P)\otimes r_{A,\boldsymbol\mu}(\Phi_A)\) and
\(q_B(P')\otimes r_{B,\boldsymbol\mu}(\Phi_B)\).
Such a product description follows from the real quotient
once each finite datum gives a real intertwiner out of the
same irreducible module. Take the span of those intertwiners
as the finite factor. Schur's lemma and simplicity identify
its evaluation tensor with the span of the images. In the
countable-dimensional real module, the scalar form of Schur's
lemma follows since an endomorphism with no eigenvalue would
give an uncountable linearly independent family of resolvents
of a fixed nonzero vector. Indeed, clearing the denominators
in a nonzero finite relation between resolvents gives a
nonzero polynomial in the endomorphism annihilating that
vector. The polynomial is a product of invertible shifts
times a nonzero scalar, a contradiction. The real module and the two
polynomial maps come from the same infinity types, so the
factorization may be used with fixed real coordinates.
\releasecite{Development/AutomorphicArithmetic/ThetaCorrespondence/ThetaPairLifts/LocalTypeDictionary/TwistedQuotientReal.lean}{18}{Common real quotient}
\releasecite{Development/AutomorphicArithmetic/ThetaCorrespondence/ThetaPairLifts/LocalTypeDictionary/TwistedSquareLocalMatching.lean}{28}{Product quotient}.

For the reference triangle sections \(\Phi_A^0,\Phi_B^0\)
the finite maps are proportional with ratio
\(\gamma_{\boldsymbol\mu}\). Choose \(P,P'\) with the same
nonzero real quotient, and with the reference smaller-source
integral nonzero after identifying the triangle coordinates.
Write its triangle polynomial coordinate as \(p_F'(P')\),
and write \(p_F(P)\) for the larger-source coordinate.
The actual quotient Haar measures differ from the triangle
measures by factors \(s_A,s_B\). If \(b_{\boldsymbol\mu}\)
is the regularized integral scalar on the product quotient,
injectivity of the triangle realization gives
\[
 s_A\lambda_{\boldsymbol\mu}p_F(P)
       =b_{\boldsymbol\mu}\gamma_{\boldsymbol\mu}s_B\,p_F'(P').
\]
Since the triangle polynomial maps and the chosen vectors
are independent of \(\boldsymbol\mu\), the ratio
\(m_{\rm ref}=b_{\boldsymbol\mu}\gamma_{\boldsymbol\mu}s_B/
             (s_A\lambda_{\boldsymbol\mu})\)
is independent of \(\boldsymbol\mu\).

On the actual finite flip tests the required normalized
finite comparison reads
\[
 s_A r_{A,\boldsymbol\mu}(\Phi_A)
   =\kappa_{\boldsymbol\mu}\gamma_{\boldsymbol\mu}
                 s_B r_{B,\boldsymbol\mu}(\Phi_B).
\]
The two prescribed model identifications use one common Haar
and real factor, the finite twist on the moving frame and the
local characters of the complementary-frame diagonal units.
In those identifications
\[
 \kappa_{\boldsymbol\mu}=C\,
   \xi_{\boldsymbol\mu}(y)
   \left(\prod_{v\in S_{\boldsymbol\mu}}
                        \nu_{\boldsymbol\mu,v}(x_v)\right)^{-1},
\]
where \(C\ne0\) is common, \(y\) is the prescribed moving
idele, \(x_v\) is the fixed local frame unit and
\(S_{\boldsymbol\mu}\) is the tested split support.
The relative scalar \(\kappa_{\boldsymbol\mu}/
\kappa_{\boldsymbol\mu_0}\) is in \(\iota_K(K^\times)\).
Indeed \(S_{\boldsymbol\mu_0}\subseteq S_{\boldsymbol\mu}\).
On the old places the ratio of the two \(\nu_v(x_v)\)
is the finite twist value. At the added places the
additional factors are values of the reference finite local
characters on units and of the finite twist. The coefficient
field includes these values on the chosen support, and
\(\xi_{\boldsymbol\mu}(y)\) is a finite twist value in the
same field. This argument checks the field consequence
of the frame identity for the actual square.
\releasecite{Development/AutomorphicArithmetic/ThetaCorrespondence/ThetaPairLifts/LocalTypeDictionary/ThetaFlipInvolutiveCorrespondence.lean}{164}{The finite frame and Haar factors}
\releasecite{Development/AbelianVarietyTheory/Proofs/DivisorProofs/ExistsPowClose.lean}{82}{Local unit character ratios}
\releasecite{Development/AbelianVarietyTheory/Proofs/DivisorProofs/XiAwayIm.lean}{16}{The moving idele}.

In the tied real relation the quotient maps are normalized
by equality. The finite frame identity and the scalar
of the reference triangle give
\[
 c^0_{ij,\boldsymbol\mu}
   = b_{\boldsymbol\mu}\kappa_{\boldsymbol\mu}
       \gamma_{\boldsymbol\mu}s_B/s_A
   = m_{\rm ref}\kappa_{\boldsymbol\mu}\lambda_{\boldsymbol\mu}.
\]
The witness and the real comparison below give this
same ratio for the original square constants. In
particular \(a^{\rm top}_{\boldsymbol\mu}\) is
\(\kappa_{\boldsymbol\mu}/\kappa_{\boldsymbol\mu_0}\).
\releasecite{Development/AutomorphicArithmetic/ThetaCorrespondence/ThetaPairLifts/LocalTypeDictionary/RegConstEqAwayUnitMulFactor.lean}{78}{Scalar extraction on one nonzero triangle vector}.

\teamanchor{proof.period-square-witness}
\subsubsection{A nonzero evaluation on the prescribed square tests}
Fix a tuple and a square. Use its two middle-vertex
quotient measures, and place the Hermitian triangle on
the same source lattices and the sum of the two exterior
lattices. Write \(I_A\) for the large-source integral and
\(T\) for the exceptional set of the placed triangle.
The bijection from the selected split places above
\(T\) to \(T\) uses the components prescribed by the grid.
The designated auxiliary split place \(v_0\) of the square
belongs to \(T\). The place used to regularize the small
integral lies outside \(T\).

First choose a polynomial Fock vector \(P\) and a finite
pure tensor \(\varphi\), with the prescribed vacua
outside \(T\), for which \(I_A(P\otimes\varphi)\) is nonzero
as a function. The nonzero Hermitian-triangle comparison
gives a nonzero value of the large integral on the
Schwartz space. Hermite polynomial density and its continuity
with finite data fixed then give a polynomial and a finite
pure tensor. To restrict the exceptional support to \(T\),
argue by contraposition. At any \(v\notin T\) the common
vacuum image is nonzero and cyclic in the irreducible
local quotient. For fixed other factors the large integral
depends only on the local quotient image. This follows by
using the same matching small-source vector for two tests
with the same quotient image. That small-source vector can
be chosen with the required full-rank support at the good
split regularizing place. Replacing a vacuum by a local
target translate translates the resulting function.
Consequently if the function vanishes with the vacuum it
vanishes with every local vector. Induction on the finite
set of exceptional factors outside \(T\) proves the
claimed support reduction.
\releasecite{Development/AbelianVarietyTheory/Proofs/DivisorProofs/ExistsINeZeroKind.lean}{196}{Polynomial nonvanishing with exactly the permitted exceptional set}

The evaluation may be taken in the block subgroup
\(\mathcal H\) of the target, the product of the two exterior
unitary groups. In fact, \(\det(\mathcal H(\mathbb A_F))\)
is the full norm-one idele group. Use the action by the
given norm-one idele on a nondegenerate line in a positive-rank
exterior group and the identity on its complement. Every
target element is therefore \(s h\) with \(s\in\mathrm{SU}(\mathbb A_F)\)
and \(h\in\mathcal H(\mathbb A_F)\).
At the specified split place \(v_0\), strong approximation
gives density of \(\mathrm{SU}(F)\mathrm{SU}(F_{v_0})\) in
\(\mathrm{SU}(\mathbb A_F)\). If all evaluations on \(\mathcal H\)
were zero for all allowed \(\varphi\) with \(P\) fixed,
the value at \(\gamma k h\) would also be zero. The rational
factor \(\gamma\) drops out by invariance of the theta
distribution. Moving \(k\) to the right of \(h\) changes only
the factor of \(\varphi\) at \(v_0\). Continuity in the target
then gives zero at every \(s h\), contrary to the first step.
There are thus \(\varphi\) and
\(h=\operatorname{diag}(g,g_1)\in\mathcal H(\mathbb A_F)\)
with nonzero value.
\cite{p2platonov-strong-approximation}
\releasecite{Development/AbelianVarietyTheory/Proofs/DivisorProofs/ExistsINeZeroKind.lean}{223}{Restriction to the exterior block subgroup}

At each \(v\in T\) replace its local test by one coming from
the actual flip, and correct the evaluation by a block translate.
For this, take a small-source preimage \(\Psi\) of the given
large-source local quotient class in the complementary
Schr\"odinger chart. The quotient in these charts is the
space of the complementary zeta sections. We use the
following consequence of the split frame calculation.
Let \(w\) be the chosen split component above \(v\).
In the large-source chart the placed mixed-coordinate
change has the form \(Z\mapsto g_0Zc_0\), followed by
the prescribed column permutation and completion
isomorphism. In the chosen completion, write \(C_1\)
for the frame of the first exterior vertex and \(C_2\)
for the negative of the frame of the second exterior
vertex. Write \(b_W,b_A\) for the triangle target and
large-source bases, \(\gamma_A\) for the large middle
grid frame, and \(J_A\) for its Hermitian Gram matrix
in the ambient coordinates. Then the matrices, with
all entries transported by the chosen completion
isomorphism, are
\[
 c_F=b_W^{-1}\operatorname{diag}(C_1,C_2),\qquad
 c_0=c_F^{-\mathsf T},\qquad
 g_0=-\gamma_A^{\mathsf T}J_A^{\mathsf T}b_A.
\]
The row change is integral invertible by self-duality:
the pairing matrix of two integral bases of the
same self-dual lattice is unimodular.
The partial Fourier transforms on the
complementary blocks cancel, so the image of the large
flip test in this chart is the primitive-frame section
of the first partial Fourier transform of the small
test in the placed coordinates.
This uses the inverse Fourier sign on the complementary
block and the old middle level's inverse-transpose Haar
measure. Normalize the mixed model by this Fourier
inversion so that the chart identity has prefactor one.
\releasecite{Development/AbelianVarietyTheory/Proofs/TwistedOperators/ExistsMixedLagrangianChartIntegralRowChange.lean}{426}{The integral row frame and the two mixed charts}
\releasecite{Development/AbelianVarietyTheory/Proofs/DivisorProofs/ThetaTripleRealizationSquareClass.lean}{458}{The chosen column frame and its sign}

Let \(J\) be the prescribed first column block of the exterior
sum and choose a large integer \(M\).
Put
\[
 D_M=\operatorname{diag}(\varpi^M I_J,\varpi^{-M}I_{J^c}).
\]
For each exterior block, use the unitary element with
component \(C_j(D_M|_j)^{-\mathsf T}C_j^{-1}\) on the
chosen half and inverse adjoint on the other half.
The two components give an element of \(\mathcal H(F_v)\).
In the large-source Schr\"odinger chart it acts by right
substitution by \(c_0 D_M c_0^{-1}\) with a nonzero scalar
\(r\). Indeed the operator of this substitution and the
given Weil operator induce the same conjugation on the
two Heisenberg generators. Their ratio is a nonzero
scalar by Stone--von Neumann. Put \(N=c_0D_M\).
Use for the first partial Fourier transform of the small
test, in the transported coordinates, the function
\[
 \Xi(Y)=
 \frac{\nu(\det g_0)\nu(\det N)
             |\det g_0\det N|_v^{(a+a')/2}}
      {r\,u\,s_\nu}\,\Psi(YN^{\mathsf T}).
\]
Here \(u\) is the frame-permutation factor
\((\nu(\det Q)|\det Q|_v^{(a+a')/2})^{-1}\) for the
prescribed permutation \(Q\), and \(s_\nu>0\) is the ratio
of the transported inverse-transpose old-level Haar
measure to the volume-one maximal-compact Haar measure
on the small-source linear group. All these factors
are nonzero. Indeed the frame-permutation factor follows
from replacing the completion \(h\) by \(Qh\) and using its
determinant weight. For the two moved frames the large
zeta section acquires the inverse of the numerator above.
The integrand on the small side is evaluated at
\(Y N^{-\mathsf T}\); this cancels \(N^{\mathsf T}\) in
\(\Xi\). The factors \(r,u,s_\nu\) then restore the
zeta section of \(\Psi\) in its original measure and
frame. The equality of sections is an equality of
local quotient classes.

For the depth check, choose \(M\) so that the support of
\(\Psi(Yc_0^{\mathsf T})\) lies in the box
\(\varpi^{-M}\mathcal O_v\) and it is invariant under
translation by the box \(\varpi^M\mathcal O_v\).
Take \(n=\max(2M,\mathrm{depth}_v)\).
The pullback by \(D_M\) is supported in the lattice
with \(J\)-columns in \(\varpi^{-n}\mathcal O_v\) and
\(J^c\)-columns integral. It is invariant by
\(\varpi^n\) times this lattice: the changes in the
arguments of \(\Psi(Yc_0^{\mathsf T})\) have valuations at
least \(M\) and \(n-M\) respectively. The inverse finite
Fourier transform on \(J\) gives an integral test
\(\Phi_w\) of depth \(n\). The additive character and
the local transported lattice use the same dual pairing
as the split transfer. Applying that transfer to
\(\Phi_w\), then taking the block translate, recovers
the original local class in the required quotient.
\releasecite{Development/AbelianVarietyTheory/Proofs/DivisorProofs/SquareFlipReachesLocalQuotientClasses.lean}{49}{Local quotient reach by admissible flipped tests}
\releasecite{Development/AbelianVarietyTheory/Proofs/DivisorProofs/ExistsDepthSplitThetaBigEqSmallBlockTranslate.lean}{38}{Depth, frame movement and normalizing scalar}

Replace a local component at a time. If \(\phi_v^{\rm new}\)
is the unshifted flip datum, its image after action by
\(k_v\in\mathcal H(F_v)\) equals the old quotient image.
Right covariance gives the same nonzero value with the
new tensor at \(h k_v\). The new evaluation remains block
diagonal, and subsequent replacements leave earlier
finite components in their new unshifted forms. This
produces the local symbols \(\Phi_w\) of each square and a
nonzero evaluation on the pure tensor determined by them.
The depths in this existence argument may exceed the
minimum depth of the row tests.

To identify this tensor with the actual two-arrow tests,
use the split idempotent which is one at the chosen
components over \(T\) and zero at the other components.
The local tensor decomposition gives a finite sum
\(\tau\) on the large route and a finite sum \(\tau'\) on the
small route, having symbols \(\operatorname{Flip}_w\Phi_w\)
and \(\Phi_w\) on the prescribed halves. Off \(T\) use the
specified lattice vacua. In the large Hermitian
Schr\"odinger model, put the local mixed-polarization
intertwiners in the exact frames of these halves. Each
takes the local vacuum to its lattice indicator times its
value at zero. This value is nonzero. The common Gaussian
at infinity is multiplied by the product of these values
so that the two models agree on the initial vacuum.
Applying translations and multiplication operators of
the two complementary local Lagrangians then preserves
the equality. They act by the same Heisenberg phase in
both models, and their translates of the local vacuum
span the local Schwartz space. Thus the equality for
the local symbols on the chosen halves extends to the
actual two-arrow tensor datum on the full adelic space.
This proves the equality for the Gaussian with exactly
these finite factors. The route Fock substitutions are
surjective and commute with the creation operators,
so we can realize the specified \(P\) on \(\tau\) with
the finite factors unchanged. The real relation has
onto projections, which supplies a matching small-route
polynomial on \(\tau'\).
\releasecite{Development/AbelianVarietyTheory/Proofs/DivisorProofs/ExistsGaussianDatumFlippedFiniteTensorBigRoute.lean}{246}{Gaussian with the finite flip tensors}
\releasecite{Development/AbelianVarietyTheory/Proofs/DivisorProofs/RealPolynomialVectorReachedByRealDatum.lean}{108}{All Fock polynomials on those same finite tensors}

The product of the two large-route kernels for this
datum is the large Hermitian theta kernel times the
nonzero splitting prefactor at \(\operatorname{diag}(g,g_1)\).
In the compatible frames the source conjugation
pushes the selected grid quotient Haar measure to the
triangle measure. The integral over a different
fundamental domain agrees by the rational invariance
of the summed kernel. We obtain
\[
 \int_{[U_A]} \ker_A(\Theta_A;g,g_1,u)\,du\ne0
\]
on the matched route tests \(\Theta_A,\Theta_B\)
just constructed. Take a finite-source translate
regularizer \(z\) for the corresponding small-route
test, with nonzero augmentation \(d(z)\).
Its normalized integral is the one in the square
comparison. Hence
\[
 0\ne d(z)\int_{[U_A]}\ker_A(\Theta_A;g,g_1,u)\,du
   =c_{ij,\boldsymbol\mu}
      \int_{[U_B]} (z\ker_B)(\Theta_B;g,g_1,u)\,du .
\]
This proves nonvanishing both of the square constant
and of the translated small-route evaluation, in the
measures used in the matching identity.
\releasecite{Development/AbelianVarietyTheory/Proofs/DivisorProofs/ExistsMatchedFlipThetaIntegral.lean}{39}{The selected square evaluation}
\releasecite{Development/AutomorphicArithmetic/ThetaCorrespondence/ThetaPairLifts/LocalTypeDictionary/ConstK.lean}{15}{Nonzero square constant}

\teamanchor{proof.period-square-scalar}
\subsubsection{The constant for the real relation of the row tests}
Keep the real relation \(\mathscr R\) used in the grid
row tests. The normalization argument also uses the
relation \(\mathscr R_0\) given by equality in the two
fixed real quotient maps. Both relations have onto
projections. They use the same real vacua in every
arrow, so the condition of having paired finite flip
symbols is the same for both. Write
\(c_t=c_{ij,\boldsymbol\mu}\) for the square constant
on \(\mathscr R\), with \(t=\boldsymbol\mu\), and
\(c_t^0\) for the regularized comparison constant on
\(\mathscr R_0\).

To see why the local finite matching supplies one
scalar for all paired tests, fix \(t\) and a product
coordinate system with the real maps \(q_A,q_B\)
of the relation being used. For one nonzero real
quotient vector \(x\), choose a large polynomial
with image \(x\). The pure-tensor matching gives
a small polynomial with image \(x'\) such that
\[
 x\otimes r_A(\tau)=x'\otimes r_B(\tau')
\]
for every finite pair with the prescribed flip
symbols. The small polynomial can be chosen once
here, since it is chosen in the onto real map before
the matching finite tensors are inserted. Apply
a complex-linear functional on the real quotient module taking
\(x\) to \(1\). This gives
\(r_A(\tau)=\lambda_t r_B(\tau')\), where \(\lambda_t\)
is its value at \(x'\). It is independent of the
paired finite test and its depth. Consequently
the product quotient satisfies
\[
 p_A(\Theta_A)=\lambda_t p_B(\Theta_B)
\]
on every test in the real relation. For the tied
relation \(\mathscr R_0\), the same argument uses
the placed Hermitian triangle and its reading in
the prescribed frames. Multiply the comparison
scalar of the common product quotient by
\(\lambda_t\). The kernel realization has a common
splitting prefactor, and it transports the two
source measures to the given grid measures. This
gives \(c_t\), respectively \(c_t^0\), in the
regularized comparison.
\releasecite{Development/AutomorphicArithmetic/ThetaCorrespondence/ThetaPairLifts/LocalTypeDictionary/SquareIdentityOfMatching.lean}{187}{Uniform finite scalar for the original real relation}
\releasecite{Development/AbelianVarietyTheory/Proofs/DivisorProofs/RealMatchingTiedRegularizationConstant.lean}{180}{The tied comparison scalar in the actual measures}

For fixed exterior arguments let \(T_B(z)\) be the
integrated translated small test when integrable.
The normalized value \(J=T_B(z_0)/d(z_0)\) is
independent of the chosen good-place regularizer
\(z_0\), and \(T_B(z)=d(z)J\) for every \(z\) with an
integrable translated kernel. Indeed place an auxiliary
regularizer \(w\) of nonzero augmentation at a good
split place outside the support of \(z_0\) and outside
the finite exceptional set of the finite translation
sum \(z\). Strip the compact factors of \(z\) outside
its finite exceptional set, preserving its augmentation
and translated function. The resulting sum commutes
with \(w\). Change variables in the middle quotient
and use absolute convergence for the two orders of
each product of translates to get
\[
 d(w)\,T_B(z)=d(z)\,T_B(w),\qquad
 d(w)\,T_B(z_0)=d(z_0)\,T_B(w).
\]
The latter identity uses disjoint support for \(w,z_0\).
Cancelling the two nonzero augmentations proves the
assertion. Thus the original scalar identity may be
applied to the regularizer used for the tied relation.
\releasecite{Development/AutomorphicArithmetic/ThetaCorrespondence/ThetaPairLifts/LocalTypeDictionary/RegularizedIntegralSmallConstant.lean}{19}{One normalized small-source value for every integrable translate}
\releasecite{Development/AutomorphicArithmetic/ThetaCorrespondence/ThetaPairLifts/LocalTypeDictionary/TwistedConstantExistence.lean}{50}{The original relation for all convergent comparisons}

For each paired finite test, the large integrated
kernel at fixed exterior arguments is a linear
functional on the large-route real polynomial tensor
space. The common nullspace, over all such tests and
arguments at any one \(t\), is
\[
 \mathcal K=\ker(q_A\iota_A).
\]
Indeed the large kernel is the evaluation of the
large-source lift on \(q_A(\iota_A P)\otimes y\),
times its nonzero splitting prefactor, where \(y\)
is the fixed finite quotient image for the test.
The quotient map of the theta integral is either zero
or injective by irreducibility. The nonzero grid
evaluation just constructed gives one finite test at
each \(t\) with \(y\ne0\) and a nonzero integral map.
If all paired grid evaluations vanish for a given
\(P\), the block-translation and local replacement
arguments applied to this fixed \(P\) give vanishing
of its global theta integral on that test. Injectivity
and cancellation of \(y\) give \(q_A(\iota_A P)=0\).
Conversely that quotient equality gives vanishing
for every paired large test.
Thus the nullspace is the same for all \(t\), and it
is proper.
\releasecite{Development/AbelianVarietyTheory/Proofs/DivisorProofs/SquareConstantProportionalRegular.lean}{343}{Common large-route nullspace}

Choose \((P,Q)\in\mathscr R\) and
\((P_0,Q)\in\mathscr R_0\). Use the very same
finite tensors and quotient measures for the two
large tests and their common small test. Take the
regularizing finite translate supplied by the
comparison on \(\mathscr R_0\).
For a large-route evaluation functional \(L\), both
identities now use the same normalized translated
small-route value \(J\), giving
\[
 L(P)=c_tJ,\qquad L(P_0)=c_t^0J.
\]
It follows that \(c_t^0P-c_tP_0\in\mathcal K\).
If \(c_t^0=0\), every \(P_0\) would belong to
\(\mathcal K\) by the onto projection of
\(\mathscr R_0\). The nonzero evaluation thus gives
\(c_t^0\ne0\). Choose one pair as above with
\(P_0\notin\mathcal K\), and keep it fixed for all \(t\).
The class of \(P\) in the quotient by \(\mathcal K\)
is then \(c_t/c_t^0\) times the fixed nonzero class
of \(P_0\). Consequently
\[
 c_t=\ell_{ij}c_t^0,\qquad
 \ell_{ij}\in\mathbb C^\times
\]
with \(\ell_{ij}\) independent of the tuple.
\releasecite{Development/AbelianVarietyTheory/Proofs/DivisorProofs/SquareConstantProportionalRegular.lean}{412}{Transport between the two real relations}

In the ratio \(c_t/c_{t'}\) the factor \(\ell_{ij}\)
cancels. The reference triangle's coefficient formula
is applied to the \(c^0\) constants with the measure
changes already included.

For a listed smooth cusp form, apply the regularized
integral comparison inside the exterior integral
in the prescribed quotient measures. The joined
integrals on the chosen path are absolutely
convergent after the finite-source translates,
and the translated small lift is multiplied by
the augmentation of the regularizer by change
of variables. Cancel its nonzero augmentation.
This gives
\[
 \mathcal L_{\rm right,down}(\Theta_A)
   =c_t\,\mathcal L_{\rm down,right}(\Theta_B).
\]
For a square inside a longer path, decompose its
finite test with the two exterior paths fixed.
Each summand has two finite symbols related
by the prescribed flip, and the corresponding
Fock polynomials meet the original real relation.
Linearity and the same convergence bound at
that prefix give the identity for their sum.
Replacing the small route by the large route
therefore multiplies the expression by \(c_t^{-1}\).
\releasecite{Development/AbelianVarietyTheory/Proofs/DivisorProofs/CommutesWithOfIsSquareConstOfTriple.lean}{27}{Cancellation after integrating the source translates}
\releasecite{Development/AbelianVarietyTheory/Proofs/TwistedOperators/GridSquareFlipSmallRouteEqInvSmulBigRoute.lean}{314}{Use of the inverse constant along the path}

\subsubsection{The definite triangle}\label{sec:period-definite-triangle}
At the bottom use the three standard Hermitian spaces of ranks
\(A_0,B_0,A_0+B_0\), their self-dual lattices and their prescribed
finite sections. The two sources and the common target are definite
at every infinite place. The larger-source jet uses the
negative-sign Fock convention. Put \(h_0=(A_0+B_0)/2\) and
\(e_0=(A_0-B_0)/2\). Then \(A_0>B_0>0\) and
\(A_0<A_0+B_0\le 2A_0\). The target has Witt index zero, so there
is no smaller global member of this Witt tower. The doubling
parameter for the larger source is \(s_0=B_0/2\). At the unramified
finite places the constant source vector gives an exponent string
\((A_0-1)/2-k\), \(0\le k<A_0\). The local standard factor
in the doubling convention uses \(\nu_{ij,\boldsymbol\mu}^{-1}\).
Consequently at \(s=s_0\) the arguments in that factor are
\[
  s_0+\tfrac12+\tfrac{A_0-1}{2}-k=h_0-k,\qquad 0\le k<A_0 .
\]
The exceptional set here consists of the base-level places and the
places under the finite twist. At the new places
\(\nu_{ij,\boldsymbol\mu_0}\) is unramified and the local component
of \(\xi_{\boldsymbol\mu}\) is ramified. Ramification occurs on both
split factors, since the character is trivial on the diagonal.
The base-level factors are ramified as well, so the omissions
remove only trivial primitive factors. Since
\(\nu_{ij,\boldsymbol\mu}^{-1}=\nu_{ij,\boldsymbol\mu}\circ c\)
for CM conjugation \(c\), reindexing the remaining Euler factors by
\(w\mapsto\bar w\) gives the primitive functions for \(\nu_{ij,\boldsymbol\mu}\).
Analytic continuation gives
\(\prod_{s=1-e_0}^{h_0}\Lambda_{\boldsymbol\mu}(s)\)
at the doubling point. This is the finite factor in the anisotropic
Rallis formula \cite[S\S7.1, 11.3--11.7]{gqt}.
\releasecite{Development/AbelianVarietyTheory/Proofs/TwistedOperators/BottomExistsIntegrable.lean}{104}{Standard triangle}
\releasecite{Development/AbelianVarietyTheory/Proofs/DivisorProofs/PartialLFunctionEqNormalizingFactorDefiniteBotto.lean}{86}{Exceptional places and exponent string}
\releasecite{Development/AbelianVarietyTheory/Proofs/DivisorProofs/SquareNuUnramifiedOffPairPartialLPrimitive.lean}{82}{The inverse square character in the standard factor}.

For the larger-source finite section one can check the nonzero
exceptional doubling contribution using its support. At infinity
take the product of the powers of the top \(A_0\)-minors of the
Fock matrices. The powers are \(a_\tau-h_0\), which are nonnegative
integers by the grid gap. Right substitution acts on each such minor
by the determinant of the source matrix. The vacuum exponent and
the source splitting character contribute the inverse determinant
power, so this section is fixed at infinity. At each exceptional
split place write \(Z\) for the coordinate matrix in the polarization
given by the selected place of \(E\) over \(v\). The first section
has integral \(Z\) with its top \(A_0\)-minor invertible over
\(\mathcal O_v\). If \(Z\) and \(Zg_v^{-1}\) both meet this
condition, their top minors differ by \(g_v^{-1}\), which then
belongs to \(\mathrm{GL}_{A_0}(\mathcal O_v)\). The other split
component is its adjoint inverse, by unitarity. Thus outside the
lattice stabilizer the section and its translate have disjoint
support in this polarization. Their local inner product is zero.
Equivalently one can apply a finite Fourier projector on a common
finite quotient of the two supports. On the lattice stabilizer the
determinant character of the section cancels the splitting factor,
so the section is fixed. At these exceptional places, the local coefficient in the doubling
integral is therefore supported on the compact stabilizer and is the
nonzero squared norm of the vector there. The Siegel section height
is one on this stabilizer. Its integral is the Haar volume times the
squared norm. The other places contribute the spherical factors above.
In particular, with the prescribed doubling and section matching, the
larger-source lift is nonzero by the norm formula and the nonvanishing
of the Hecke product.
\releasecite{Development/AbelianVarietyTheory/Proofs/DivisorProofs/DefiniteBottomLocalDoublingIntegralConverges.lean}{377}{Doubling calculation on the prescribed section}
\releasecite{Development/AutomorphicArithmetic/ThetaCorrespondence/ThetaPairLifts/LocalTypeDictionary/BigThetaNormEqRallisInnerProduct.lean}{19}{Norm on the definite triangle}.

In the common oscillator quotient, take a fixed pair of maps
\(p_*,p'_*\) on the two archimedean polynomial spaces, both
surjective onto the quotient \(\mathcal P\). They depend on the
infinity types, which are fixed across the tuples. For each tuple
the matched finite data lie on the same line in the common finite
quotient. Thus after taking the definite theta integrals the two
maps factor as \(\lambda_{\boldsymbol\mu}J_{\boldsymbol\mu}p_*\)
and \(J_{\boldsymbol\mu}p'_*\). The kernel of
\(J_{\boldsymbol\mu}\) is a subrepresentation of the irreducible
quotient. Nonvanishing of the larger-source lift gives
\(\lambda_{\boldsymbol\mu}\ne0\) and injectivity of
\(J_{\boldsymbol\mu}\), and \(p'_*\ne0\) then gives nonvanishing of
the other lift. This step uses the local matching on the prescribed
finite data, the polynomial quotient and the global integral
comparison on those quotients.
\releasecite{Development/AutomorphicArithmetic/ThetaCorrespondence/ThetaPairLifts/LocalTypeDictionary/ThetaBottomFinitePlaceMatching.lean}{15}{Prescribed finite sections in the common quotient}
\releasecite{Development/AutomorphicArithmetic/ThetaCorrespondence/ThetaPairLifts/LocalTypeDictionary/BottomTripleIntegralFactorizationOfReal.lean}{21}{Polynomial and finite factors}
\releasecite{Development/AutomorphicArithmetic/ThetaCorrespondence/ThetaPairLifts/LocalTypeDictionary/FirstOccurrenceThetaLiftNeZero.lean}{15}{Nonvanishing of both integrals}.

For each source leg, choose a compact open subgroup of its finite
adelic group fixing the sections for every tuple. At level given by
this subgroup the class set is finite by definiteness. Integration
over the adelic quotient is a single Haar volume times a sum of
compact real averages at class representatives, each divided by
the size of the representative's finite stabilizer. The compact
average of the Fock action is the projection to the invariant
subspace after twisting by the determinant character. To see its
algebraicity on a homogeneous polynomial space, restrict to the
dense subgroup of compact unitary matrices with algebraic entries.
The invariant space is the intersection of the kernels of
\(\rho(g)-1\), and the complementary space is the span of their
images. Finitely many of these matrices suffice to give the same
intersection and span. Both spaces are defined over algebraic
numbers, so the projection onto the first along the second has
algebraic coefficients. Averaging over the central circles leaves
a single multidegree at the infinite places, fixed across the
tuples. In the smaller-source leg the surviving determinant
semi-invariants are spanned by products of maximal minors.
\releasecite{Development/AutomorphicArithmetic/ThetaCorrespondence/ThetaPairLifts/LocalTypeDictionary/BottomClassSumIntegralNeg.lean}{24}{Definite class formula}
\releasecite{Development/AutomorphicArithmetic/ThetaCorrespondence/ThetaPairLifts/LocalTypeDictionary/TwistedRealAverageKernelIsAlgebraic.lean}{22}{Real averaging in a fixed multidegree}
\releasecite{Development/AutomorphicArithmetic/ThetaCorrespondence/ThetaPairLifts/LocalTypeDictionary/TwistRealProjectionAverageMonomial.lean}{15}{Projection on the smaller source}.

On the larger-source leg the finite Fourier projection expresses
the kernel on the surviving multidegree as a finite algebraic
linear combination of jets of finite Heisenberg translates of a
standard vacuum. The CM jet comparison divides them by a common
nonzero scalar, the vacuum comparison constant times the product
of the local CM periods to the surviving degrees, with conjugation
for the negative sign. On the smaller-source leg one uses the
CM jet comparison for the surviving products of minors, again with
a common period. The two Haar volumes are included in the two
periods. At the target use a compact open subgroup fixing the
finite vectors for all tuples. Its finite class set gives a finite
list of evaluations on the surviving monomial spaces. Changing
representative acts by polynomial substitution in the rational
Hermitian matrices through the chosen embeddings, together with
their determinant characters. These coefficients are in a finite
number field. Thus adjoining the finitely many normalized class
coefficients at the target representatives gives the field of the
two class sums. This field is included at the same stage as the
Hecke-ratio field in the construction of the tested coefficient
field. The CM jet comparisons here must be taken with the stated
vacuum and polynomial normalizations.
\releasecite{Development/AutomorphicArithmetic/ThetaCorrespondence/ThetaPairLifts/LocalTypeDictionary/ClassSumKernelIsAlgebraic.lean}{19}{Larger-source kernels}
\releasecite{Development/AutomorphicArithmetic/ThetaCorrespondence/ThetaPairLifts/LocalTypeDictionary/JetsDivPeriodIsAlgebraic.lean}{57}{Fixed-degree CM period}
\releasecite{Development/AutomorphicArithmetic/ThetaCorrespondence/ThetaPairLifts/LocalTypeDictionary/TwistKernelDivPeriodIsAlgebraic.lean}{15}{Minor jets}
\releasecite{Development/AutomorphicArithmetic/ThetaCorrespondence/ThetaPairLifts/LocalTypeDictionary/BottomFieldIsNumberField.lean}{18}{Coefficient field of both sums}.

Evaluation on the finite adeles distinguishes the resulting
functions. In fact, \(U(F)\) is dense in \(U(F_\infty)\).
The Cayley transform identifies the open subset where \(-1\) is
not an eigenvalue with an open subset of the space of
skew-Hermitian matrices. The \(F\)-rational matrices are dense
there by simultaneous approximation at the infinite places,
and this open subset is dense in the real unitary group.
For \(g=(g_\infty,g_f)\), take \(x_\ell\in U(F)\) with
\(x_{\ell,\infty}\to g_\infty\) and evaluate at
\(x_\ell(1,x_{\ell,f}^{-1}g_f)\). Left rational invariance and
continuity show that a function vanishing on finite adeles vanishes
at \(g\). The same argument applies to linear combinations of the
two operators.

\paragraph{The quotient of the bottom constants in the prescribed field.}
\teamanchor{proof.period-bottom-constant-field}
Fix the bottom triangle \(Y\) whose normalized class sums were
adjoined to \(K\) for this cell. Write \(I^Y_{\boldsymbol\mu}\) and
\((I^Y_{\boldsymbol\mu})'\) for its two integrals, with the common
quotient maps \(p_*,p'_*\) above. Use the product \(P\) of the powers
of the top minors used to prove nonvanishing and take \(P'\) to be
the product of the complementary minors, with the same powers.
Both are polynomials with integer coefficients in the prescribed
Fock variables. In the standard compact frames, the complementary
minor identity on a unitary matrix relates the two column blocks
by conjugation, the determinant, and the sign of the ordered
complement. On the surviving determinant types this gives the
intertwiner of the rectangular Schur spaces in the two quotient
maps, with the determinant twist on the second block. Schur's
lemma gives
\[
   p_*P=\rho\,p'_*P'\ne0
\]
for a scalar \(\rho\in\mathbb C^\times\) that includes the sign.
The spaces, frames and degrees at infinity are fixed throughout
the tuples, so \(\rho\) is fixed. Since \(J_{\boldsymbol\mu}\) is
injective, \((I^Y_{\boldsymbol\mu})'(P')\) is nonzero. Choose a
finite target point \(b_{\boldsymbol\mu}\) at which it is nonzero.
With the two fixed class periods \(\Omega,\Omega'\), the two values
\[
 \frac{I^Y_{\boldsymbol\mu}(P)(b_{\boldsymbol\mu})}{\Omega},
 \qquad
 \frac{(I^Y_{\boldsymbol\mu})'(P')(b_{\boldsymbol\mu})}{\Omega'}
\]
belong to \(\iota_K(K^\times)\). Their quotient is
\(\rho\lambda_{\boldsymbol\mu}\Omega'/\Omega\). Dividing by the
corresponding quotient at the base tuple shows that
\(\lambda_{\boldsymbol\mu}/\lambda_{\boldsymbol\mu_0}
\in\iota_K(K^\times)\).
\releasecite{Development/AbelianVarietyTheory/Proofs/DivisorProofs/RealQuotientMinorEqConstSmulComplementMinor.lean}{180}{Products of complementary minors}
\releasecite{Development/AbelianVarietyTheory/Proofs/DivisorProofs/ExistsConstantTripleCoefficient.lean}{95}{The common quotient for any prescribed bottom triangle}.

Now let \(X\) be the bottom triangle used in the square and let
\(c_0(\boldsymbol\mu)\) be its comparison constant. Both triangles
have the same three ranks, with identity Gram matrices in the
rational coordinates and the canonical definite frames at
infinity. Identify their three spaces by these coordinates. At
each selected split factor, move the finite frame
\(\beta_Y\) to \(\beta_X\) by \(\beta_X\beta_Y^{-1}\).
This matrix and its inverse are integral. On the conjugate factor
take the adjoint inverse, giving a unitary move that preserves
the self-dual lattice. At all other places the move is the identity.
Write \(g,g',h\) for the resulting two source moves and the common
target move at the given tuple. Let \(x,x',y\) be their determinant
ideles on the selected split halves, extended by \(1\) elsewhere.
In a split polarization the normalized Weil action on a moved
section multiplies the pullback by
\(\chi_T(x)^{-1}\chi_S(y)\), for source and target characters
\(\chi_S,\chi_T\). If \(\alpha,\beta\) are the two parity-twisted
characters of this cell, the larger and smaller leg multipliers
are therefore
\[
    \kappa=(\beta/\alpha)(x)^{-1}\beta(y),
    \qquad
    \kappa'=(\alpha/\beta)(x')^{-1}\alpha(y).
\]
The finite tests have the same defining support in the new frames.
Changing variables by \(g^{-1},(g')^{-1}\) in the two source
integrals preserves quotient Haar measure, and both target
arguments undergo the same right translation by \(h^{-1}\).
Consequently the transported comparison constant is
\(\lambda_{\boldsymbol\mu}\kappa/\kappa'\). In this quotient the
common parity character cancels. For
\(\nu=\mu_{0,j}\xi_{\boldsymbol\mu}/\chi_{m-i}\) it gives
\[
       u^{\rm fr}_{\boldsymbol\mu}
       =\kappa/\kappa'=\frac{\nu(x)\nu(x')}{\nu(y)}
       \ \in\ \iota_K(K^\times).
\]
Indeed the local values of these three characters are in the
prescribed coefficient field. For a determinant in \(E_w^\times\)
its character value equals the value on a sufficiently close
element of \(E^\times\), since the local characters are smooth.
Only the finitely many selected split halves contribute.
\releasecite{Development/AbelianVarietyTheory/Proofs/DivisorProofs/DefiniteBottomTripleTransportSameRanksConstants.lean}{278}{Identifying the standard spaces}
\releasecite{Development/AbelianVarietyTheory/Proofs/DivisorProofs/ExistsSupportedUnitariesMovingBottomTriples.lean}{110}{Integral moves on selected halves}
\releasecite{Development/AbelianVarietyTheory/Proofs/DivisorProofs/ExistsDefiniteBottomTripleTransferredFrames.lean}{39}{The two split Weil multipliers}
\releasecite{Development/AbelianVarietyTheory/Proofs/DivisorProofs/DefiniteBottomTripleWeilMoveRightTranslate.lean}{143}{Common right translation}
\releasecite{Development/AbelianVarietyTheory/Proofs/DivisorProofs/DefiniteBottomTripleBasisChangeFactorMemField.lean}{92}{Membership of the frame factor in the prescribed field}.

With the frames aligned, the prescribed vacua span the same
vacuum line on each leg. Replacing either normalized vacuum
multiplies its section by one fixed scalar, since the section
operator for each tuple uses the same support, conductor and Weil
action. The source measures on the two aligned presentations also
differ by fixed factors. These changes only multiply the comparison
by a nonzero constant independent of the tuple. Applying
\eqref{eq:period-triangle-pair-normalization} on \(X\) to the
transported pair and to the fixed pair defining \(c_0\) yields
\begin{equation}
 \frac{c_0(\boldsymbol\mu)}{c_0(\boldsymbol\mu_0)}
   =\frac{u^{\rm fr}_{\boldsymbol\mu}}{u^{\rm fr}_{\boldsymbol\mu_0}}
      \frac{\lambda_{\boldsymbol\mu}}{\lambda_{\boldsymbol\mu_0}}
   \ \in\ \iota_K(K^\times).
 \label{eq:period-bottom-constant-field}
\end{equation}
\releasecite{Development/AbelianVarietyTheory/Proofs/DivisorProofs/DefiniteBottomTripleConstantsEqOfSameLattices.lean}{94}{Fixed vacuum factors}
\releasecite{Development/AbelianVarietyTheory/Proofs/DivisorProofs/DefiniteBottomTripleConstantsEqOfSameLattices.lean}{165}{Fixed Haar factors}
\releasecite{Development/AbelianVarietyTheory/Proofs/DivisorProofs/BottomConstantOfBottomPair.lean}{26}{Transport between prescribed triangles}
\releasecite{Development/AbelianVarietyTheory/Proofs/DivisorProofs/BottomTriangleConstantAt.lean}{41}{Bottom constant ratio}.

\subsubsection{Hecke values in the square comparison}
Write \(E/F\) for the CM extension in the grid. At cell \((i,j)\)
the two splitting characters carry the same parity twist by \(\eta\).
Their quotient, with \(\xi_{\boldsymbol\mu_0}=1\), is
\[
  \nu_{ij,\boldsymbol\mu}
    =\frac{\mu_{0,j}\xi_{\boldsymbol\mu}}{\chi_{m-i}}
    =\nu_{ij,\boldsymbol\mu_0}\xi_{\boldsymbol\mu}.
\]
It is unitary and trivial on \(\mathbf A_F^\times\). In the CM type
\(\phi_\tau=\overline{\operatorname{emb}_\tau}\) the infinity character
at \(\tau\) is \((z/\bar z)^{a_\tau}\), where
\(a_\tau=e_\tau(j+1)+\operatorname{cen}_\tau(m-i)\) is the sum specified
by the signed weight grid. The grid gap gives
\(a_\tau>N_{\mathrm{amb}}+1\), where
\(N_{\mathrm{amb}}=1+r+m(t+r+6)\) is the ambient rank bound.
The finite-order character \(\xi_{\boldsymbol\mu}\) has trivial
infinity component, so these exponents also apply at the base tuple.
\releasecite{Development/AbelianVarietyTheory/Definitions/Divisors/FinitePrimesUnder.lean}{233}{Square character}
\releasecite{Development/AbelianVarietyTheory/Proofs/DivisorProofs/SquareNuTwistedTupleIsCMValueHeckeCharacter.lean}{62}{Infinity type and finite twists}
\releasecite{Development/AbelianVarietyTheory/Proofs/TwistedOperators/NormalFormToricRatio.lean}{37}{Weight bound}.

Here \(\Lambda_{\boldsymbol\mu}(s)\) is the finite primitive Hecke
\(L\)-function. Its factor at \(w\) is
\((1-\nu_{ij,\boldsymbol\mu,w}(\varpi_w)(\mathrm Nw)^{-s})^{-1}\)
if the character is trivial on the local units, and is \(1\) otherwise.
This convention uses each tuple's own conductor. The archimedean
factor is \(\prod_\tau\Gamma_{\mathbb C}(s+a_\tau)\), with
\(\Gamma_{\mathbb C}(u)=2(2\pi)^{-u}\Gamma(u)\). The dual factor at
\(1-s\) has arguments \(1-s+a_\tau\), so the critical integers are
\[
   1-\min_\tau a_\tau\ \le s\le\ \min_\tau a_\tau .
\]
In particular, \(1\le s\le N_{\mathrm{amb}}+1\) is in this interval.
For the two diagonal ranks write
\(R_{ij}=r+jt+i(r+6)\). Both \(t\) and \(r+6\) are odd. Hence
\(n_0=(t+r+6)/2\ge1\), and
\(n_0+q_j\le R_{ij}+n_0=(R_{ij}+R_{i+1,j+1})/2\le N_{\mathrm{amb}}\).
In an even cell \(q_j\ge1\), so this bounds the remaining argument
\(n_0+1\) as well. The finite factors converge absolutely to a
nonzero product at every argument \(s\ge n_0+1\) used in
\eqref{eq:period-square-ratios}.
\releasecite{Development/AbelianVarietyTheory/Definitions/Divisors/EntireEqOfHalfPlane.lean}{53}{Primitive Euler factors}
\releasecite{Development/AbelianVarietyTheory/Statements/PlaceShiftStatements/AuxiliaryPrimeNeTypePrime.lean}{63}{Rank identities}.

The triangle norm calculation also uses nonvanishing at integers
from \(1-e_0\) to the half-sum of the diagonal ranks, where
\(e_0=n_0+q_j-p_j\) is half the difference of the two bottom ranks.
For a unitary Hecke character with nontrivial restriction to the
ideles of norm one,
Tate's theorem gives an entire completed \(L\)-function and relates
its value at \(s\) to the inverse character at \(1-s\) by a
nonvanishing epsilon factor. The completed \(L\)-function of the
trivial character has a simple pole at \(s=1\).
\cite[Chapter~XV, \S4]{tate-thesis}
These assertions apply to \(\nu=\nu_{ij,\boldsymbol\mu}\), \(\nu^2\)
and \(\nu^{-1}\), none of which can be a power of the norm because
each has unequal paired infinity exponents.
At the boundary, let \(S\) contain all ramified finite places. For real
\(\sigma>1\), the Euler-product expansion and the inequality
\(3+4\cos\theta+\cos(2\theta)=2(1+\cos\theta)^2\ge0\) give
\[
  \big|\zeta_E^S(\sigma)^3 L^S(\sigma,\nu)^4
                L^S(\sigma,\nu^2)\big|\ge1.
\]
The zeta factor has a simple pole at \(1\) and the last factor is
holomorphic there. Thus \(L^S(1,\nu)\ne0\). Reinstating the unramified
factors in \(S\) preserves nonvanishing. The same argument applies to
\(\nu^{-1}\). If \(1-e_0\le s\le0\), then \(1-s\ge1\) and both
infinity factors in the functional equation are finite and nonzero,
since \(e_0\le N_{\mathrm{amb}}<\min a_\tau\). The conductor power and
root number are nonzero. This gives nonvanishing throughout the
integer range for the triangle.
\releasecite{Development/AbelianVarietyTheory/PublishedTheorems/Divisors₂/HeckeTateFunctionalEquationCM.lean}{19}{Hecke functional equation}.

For algebraicity, use Blasius's critical-value formula for an
algebraic Hecke character \(\chi\) of \(E\) with unequal infinity
exponents at every complex place. At a critical integer \(s\) it gives
\[
  L(s,\chi)\ \sim_{\overline{\mathbb Q}}\
    d_F^{1/2}(2\pi i)^{[F:\mathbb Q]s}
    p(\chi^{-1,c},\Phi_\chi).
\]
Here \(L\) is the finite primitive Euler product, \(c\) is complex
conjugation of \(E\), \(d_F\) is the absolute discriminant, and
\(p\) is the CM period. The CM type \(\Phi_\chi\) takes the embedding
with the smaller infinity exponent from each pair. The relation
\(\sim_{\overline{\mathbb Q}}\) means equality up to a nonzero algebraic
factor in the chosen complex component. For a fixed CM type \(\Phi\),
the CM periods multiply under product of characters with critical
infinity type. This follows from the diagonal map of the Shimura
torus to its square.
\cite{blasius-hecke,blasius-period-relations}
\cite[\S1.2, Proposition~1.1 and Theorem~1.1]{lin-cm-periods}

Apply this with \(\chi=\nu_{ij,\boldsymbol\mu}\) and
\(\chi_0=\nu_{ij,\boldsymbol\mu_0}\). The conjugate inverses equal
the original characters, since \(\chi(x)\chi(x^c)=\chi(\mathrm N_{E/F}x)=1\).
They have the same CM type \(\Phi\), the conjugate of the type
\((\phi_\tau)_\tau\) used to write the positive exponents.
Take \(d\ge1\) such that \(\xi_{\boldsymbol\mu}^d=1\).
For both characters all positive powers have distinct paired
exponents \(ha_\tau,-ha_\tau\) and the same \(\Phi\). Multiplicativity
and \(\chi^d=\chi_0^d\) give
\[
  \left(\frac{p(\chi,\Phi)}{p(\chi_0,\Phi)}\right)^d
       \in\overline{\mathbb Q}^{\times}.
\]
Taking the actual quotient of the two complex periods gives an
element of \(\overline{\mathbb Q}^{\times}\), since it satisfies the
corresponding polynomial of degree \(d\). The common discriminant and
powers of \(2\pi i\) cancel in the two critical-value formulas. Thus
\(\mathcal R_{\boldsymbol\mu}(s)\in\overline{\mathbb Q}^{\times}\)
for every \(1\le s\le N_{\mathrm{amb}}+1\).

The coefficient field is chosen using the field generated by these
ratios for the finitely many tuples, cells and integers. A twist key
specifies the finite-order character and its prescribed primes.
A grid key specifies the auxiliary prime, its selected place and
the grid characters, all over the fixed host. For every realizable
pair of keys take a basis of the compositum of the ratio field with
the fields for the finite local coefficients and other grid data.
The selection imposes the union of these finite bases before
realizing the tested grid, using the common family thresholds.
One uses a finite Galois field containing these algebraic numbers
and its involution is the restriction of complex conjugation under
the chosen inclusion \(\iota_K\). Transport between records with
the same keys identifies \(\mu_{0,j}\), \(\chi_{m-i}\) and the
matched finite-order twists, and hence their primitive
\(L\)-functions. The realized ratios therefore belong to
\(\iota_K(K^\times)\), in particular the factor
\(\mathcal R_{\boldsymbol\mu}(n_0+1)\) in each even cell.
\releasecite{Development/AbelianVarietyTheory/Proofs/TwistedOperators/NormalFormToricRatio.lean}{149}{The finite field of ratios}
\releasecite{Development/AbelianVarietyTheory/Definitions/TwistedInnerProducts/TwistPlacesOffPrimeQZero.lean}{223}{Twist keys}
\releasecite{Development/AbelianVarietyTheory/Definitions/AutomorphicMembership/GKeyFin.lean}{24}{Grid keys}
\releasecite{Development/AbelianVarietyTheory/Proofs/IwasawaDecomposition/GridExistenceOfPieces.lean}{52}{Field constraint on the actual grid}
\releasecite{Development/AbelianVarietyTheory/Proofs/IwasawaDecomposition/EventuallyLargeGrid.lean}{14}{The common thresholds}
\releasecite{Development/AbelianVarietyTheory/Proofs/IwasawaDecomposition/ConjugateEmbApply.lean}{163}{The complex embedding and involution}.

\subsection{Bessel periods and Poincar\'e series}
\label{sec:bessel}

\teamanchor{proof.bessel-real-test}
The bottom inclusion \(X_0\subset Y_0\) has codimension \(2r_S+1\),
where \(r_S=m/2-1\). The Bessel subgroup \(\mathcal S_w\) for \(w\) is
the semidirect product of the unipotent group \(N_w^{\mathrm{un}}\) of a
length-\(r_S\) isotropic flag with the copy of \(U(X_0)\) stabilizing the
flag and a marked vector of negative norm. Write \(\psi_w\) for the
generic character of \(N_w^{\mathrm{un}}\) determined by the flag and
vector. The pairing with a bottom row form \(f_0\) is
\[
 \int_{[U(X_0)]}\phi(h)
    \int_{[N_w^{\mathrm{un}}]} f_0(nh)\,\overline{\psi_w(n)}\,dn\,dh .
\]
The product contracts the chosen coefficient vectors in the bilinear
convention, using the prescribed quotient Haar measures. This is the
Gan--Gross--Prasad Bessel period for a split odd-dimensional
orthogonal complement. In the invariant-functional convention,
\(\overline{\psi_w}\) is the one-dimensional factor on the Bessel
subgroup, trivial on the diagonal \(U(X_0)\). Integrating first on
the unipotent quotient gives the displayed Fourier coefficient.
\cite[\S\S12, 23]{ggp}
The ambient \(Y_0\) is common to all words. The
construction scales a reference marked vector of norm \(-T_{\rm ref}\)
by \(\lambda_w^{-1}\), with \(\lambda_w\overline{\lambda_w}T_w=T_{\rm ref}\),
to obtain one of norm \(-T_w\). The character, stabilizer and column
level for \(w\) enter its period separately.

At a real place the compact column type \(\rho_w\) has parameter
\(\alpha_1>\cdots>\alpha_k\). The row discrete series \(\sigma_\infty\)
has decreasing entries \((\beta_1,\ldots,\beta_{k-1},-b_0,\ldots,-b_{m-1})\).
The \(\beta_i\) carry positive signature labels and the \(-b_i\)
alternate, starting with a positive label. The column and row entries
have half-integral differences and satisfy
\(\alpha_i>\beta_i>\alpha_{i+1}\) and \(\alpha_k>-b_0\), as in the
\releasecite{Formalism/PeriodIntegrals/BesselPeriodGGP.lean}{27}{local interlacing conditions}.
Write \(S=\mathcal S_w(\mathbb R)\), embedded in the product of the row
and column groups by \((hn,h)\), and extend \(\psi_w\) trivially across
the column group. Jiang--Sun--Zhu's multiplicity-one theorem for real
unitary Bessel groups says that, for irreducible Casselman--Wallach
representations \(\tau,\lambda\) of the two groups and a generic Bessel
character \(\xi\), the continuous intertwining space
\(\operatorname{Hom}_S(\tau\widehat\otimes\lambda,\xi)\) has dimension
at most one. Here \(\lambda\) is inflated across the unipotent radical.
\cite[Theorem~A]{p2jiang-sun-zhu}

Smooth matrix coefficients of real discrete series with compact
center belong to the Harish--Chandra Schwartz space. Both
\(\sigma_\infty\) and the compact group representation \(\rho_w^\vee\)
are square-integrable with compact centers, so the same assertion
holds for their product.
\cite[\S7.2]{p2ichino-real-theta}
For an odd-dimensional orthogonal complement that is the sum of
hyperbolic planes and a nondegenerate line, Beuzart-Plessis shows
that Harish--Chandra Schwartz functions on the product group restrict
to \(L^1(S)\). For an irreducible tempered product \(\tau\boxtimes\lambda\),
his continuous sesquilinear Bessel form on smooth vectors is nonzero
exactly when \(\operatorname{Hom}_S(\tau^\infty
\widehat\otimes\lambda^\infty,\xi)\ne0\). For Schwartz matrix
coefficients, inversion in the unimodular subgroup gives the
ordinary integral of \(\langle v,(\tau\boxtimes\lambda)(s)u\rangle
\overline{\xi(s)}\) in our pairing convention.
\cite[\S6.2, Lemma~6.5.1(ii), \S7.1 and Theorem~7.2.1]{p2beuzart}
In the prescribed Haar measures its application to the actual compact
column type gives, entry by entry, the absolutely convergent integral
\[
 \mathcal B_\infty(u,v)
 =\int_{U(X_0)(\mathbb R)\times N_w^{\mathrm{un}}(\mathbb R)}
  \big\langle v,\sigma_\infty(hn)u\big\rangle
  \overline{\psi_w(n)}\,\rho_w(h)^{-1}\,dh\,dn .
\]
Xue's real Bessel criterion applies to generic packets for a unitary
pair of ranks \(d\) and \(d+2r_S+1\) whose complement has \(r_S\)
hyperbolic planes and a marked line of sign \((-1)^d\), with the
associated generic Bessel character. For discrete-series parameters
\(\varphi_{\rm row}=\bigoplus_a\chi_a\) and
\(\varphi_{\rm col}=\bigoplus_i\mu_i\), write \(e_a,e'_i\) for the
component-group generators of these characters. The continuous Bessel
Hom space has dimension one exactly when both packet characters satisfy
\[
 \eta_{\rm row}(e_a)=\epsilon(\chi_a\otimes\varphi_{\rm col},\psi_{\mathbb C})
 \quad\text{and}\quad
 \eta_{\rm col}(e'_i)=\epsilon(\varphi_{\rm row}\otimes\mu_i,\psi_{\mathbb C})
 \qquad\text{for all }a,i .
\]
The epsilon factors use \(s=1/2\) and the trace-zero character
\(\psi_{\mathbb C}(z)=\exp(2\pi i\operatorname{Tr}_{\mathbb C/\mathbb R}(iz))
=\exp(2\pi(\bar z-z))\), which also fixes the even-rank Whittaker
normalization. In odd rank \(d'\) use the unique datum on
\(U((d'+1)/2,(d'-1)/2)\).
\cite[\S\S1.1--1.2, Theorem~1.1]{p2xue-bessel}
The row and dual compact column lie in generic discrete-series
packets. For \(d=k\) even, negating both Hermitian forms gives the
required positive marked line. For \(r_S>0\), scalars \(t_a\) on the flag
basis and \(\overline{t_a}^{-1}\) on its dual scale the consecutive
character coordinates by \(t_a/t_{a+1}\) and the last by \(t_{r_S}\).
They fix the column group and marked line and realize the prescribed
generic character by conjugation.

For an irreducible discrete series of \(U(p,q)\), with distinct
Harish--Chandra weights in \(\mathbb Z+(p+q-1)/2\), Atobe attaches
the character \(\chi_x(z)=(z/|z|)^{2x}\) to each weight \(x\).
In his normalization, a weight at decreasing position \(j\) has
component sign \((-1)^{j-1}\) for a positive signature label and
\((-1)^j\) for a negative one. Negating the form flips each label,
and passing to the dual negates the weights, reverses their order
and preserves the signature label attached to each negated weight.
In even rank, changing to the other Whittaker datum reverses the
component sign at every generator.
\cite[\S2.2, Theorem~2.1(4),(6),(7) and Theorem~A.4]{p2atobe-hc}
For \(a\in\mathbb Z+\tfrac12\), the archimedean root number of
\((z/|z|)^{-2a}\) at \(s=1/2\) for \(\psi_{\mathbb C}\) above is
\(\operatorname{sgn}(a)\). Taking \(a=-x\) gives
\(\epsilon(\chi_x,\psi_{\mathbb C})=(-1)^{[x>0]}\).
\cite[Proposition~2.1]{p2ggp-examples}

To match the even-rank datum, one can test the positive-definite
pair \(U(d+1)\supset U(d)\). Write any compact dual column weights
in decreasing order as \(-\gamma_d>\cdots>-\gamma_1\). Choose the
compact row weights \(\delta_1>\gamma_1>\delta_2>\cdots>
\gamma_d>\delta_{d+1}\) in the other parity. This is the
Harish--Chandra form of the compact branching rule, so the Hom space
has dimension one. Xue's epsilon product on the generator of
\(-\gamma_i\) then has \(i\) negative factors, and Atobe's positive
signature label at position \(d+1-i\) gives the same sign
\((-1)^{d-i}=(-1)^i\). This distinguishes the two even-rank
normalizations. Both sources use the odd quasi-split form with one
more positive than negative direction.
\cite[Introduction and Theorem~1.1]{p2xue-bessel}
\cite[Appendix~A.3--A.4]{p2atobe-hc}
Thus after the form reversal the packet signs at \(\beta_i\),
\(-b_j\) and the dual column entry \(-\alpha_i\) are
\((-1)^i\), \(1\) and \((-1)^{i-1}\), respectively. Multiplying
the root numbers for the summands counts positive row-minus-column
differences, respectively \(k-i\), \(0\) and \(i-1\). Since \(k\)
is even, both equations of the criterion hold.
By the nonvanishing equivalence above and smooth cyclicity, there
are \(u,v\) among the translates of any unit vector \(v_0\) in the
lowest \(K\)-type with \(\mathcal B_\infty(u,v)\ne0\).

For the real orbital norm, fix the section on the row group used in
the Gram estimates. It transforms on the left by \(\rho_w\) and \(\psi_w\).
It has positive finite norm on the Bessel quotient with the
specified quotient measure and is reproduced by right convolution
with \(d(\sigma_\infty)\overline{\langle v_0,\sigma_\infty(g)v_0\rangle}\),
where the formal degree uses the same ambient Haar measure.
These properties identify the real test used in those estimates.
Their proof in the release uses a direct nonvanishing calculation on
translates of \(v_0\).
\releasecite{CompanionTheory/HarishChandraModules/BesselParameterIntegrals/BesselOrbitConvolutions/BesselNonvanishingL2.lean}{34}{Direct nonvanishing argument using the chosen section}.
Its
\releasecite{CompanionTheory/AutomorphicTheory/LocalBesselUniqueness/BesselModels/LocDDQ.lean}{19}{comparison with the chosen local factor}
then factors \(\mathcal B_\infty(u,v)\) as
\(\xi\mapsto\langle\ell_w(v),\xi\rangle\ell_w(u)\) and identifies the
local factor as a multiple of evaluation by the nonzero Bessel
functional \(\ell_w\). This calibration enters the proof of the full
Gram estimates in Section~\ref{sec:trace}.

\teamanchor{proof.full-bessel-realization}
At the fixed bottom row level and type, let \(\mathcal H\) be the
finite-dimensional Hilbert space containing the Bessel Poincar\'e series
for all words. From \(\phi\in H_w^{\rm full}\), the construction forms a
seed \(\Phi_{w,\phi}\) on \(G_0=U(Y_0)\) with the Bessel character and the
specified local factors. Let \(x_w\) be its chosen finite placement,
\(K_2\) the ambient finite level with character \(\chi_2\),
and \(W_w\) the real section just described. Write
\(T_\phi(\xi;1,h_f)\) for the column coefficient evaluated at \(h_f\)
on the vector \(\xi\) in the real type. The seed has finite support
modulo the Bessel group in \(N_w^{\rm un}(\mathbb A_{F,f})\,
U(X_0)(\mathbb A_{F,f})\,x_wK_2\), with the column group in its
specified embedding. For \(y_f=n h_f x_w k_2\) it has value
\[
 \Phi_{w,\phi}(y_\infty,y_f)=
   \psi_{w,f}(n)\chi_2(k_2)\,
   T_\phi(W_w(y_\infty);1,h_f).
\]
It is zero off that support. The finite level relation for the
construction is as follows. If
\(x_w^{-1}n h_f x_w\in K_2\), then
\(x_w^{-1}h_f x_w\in K_2\) and the level character of the
first product is \(\psi_{w,f}(n)\chi_2(x_w^{-1}h_f x_w)\).

\teamanchor{proof.period-bessel-seed-overlap}
Here is the finite congruence calculation for the chosen flag.
Let \(g=g_0g_1\) be the modulus, \(\Lambda_f\) the finite
adelic lattice of \(Y_0\), \(K_\Lambda\) its integral stabilizer,
and \(e_a,f_a\) the rational isotropic vectors with
\(\langle f_a,e_b\rangle=\delta_{ab}\). The finite column group
fixes these vectors and the marked \(z_w\).
The rational lattice splits along the hyperbolic spaces and their
orthogonal complement. Its coefficients on the \(e_a\)'s are
\(I_a=\mathfrak d^{-(r_S-a)}\) for indices \(0\le a<r_S\),
where \(\mathfrak d=\mathfrak d_{F/\mathbb Q}\mathcal O_E\).
The coefficients on the \(f_a\)'s are \(I_a^{-1}\).
The orthogonal projection onto the complementary space preserves
\(\Lambda_f\). If \(nh_f\in K_\Lambda\), apply this projection
to \(nh_f y\) for \(y\) in the complementary lattice.
As \(h_f\) preserves the complement and \(n-1\) maps it into
the span of the \(e_a\)'s, the projection gives \(h_f y\).
The hyperbolic coordinates are fixed by \(h_f\).
The inverse of \(nh_f\) has the form \(n'h_f^{-1}\)
because the column group normalizes the unipotent group.
Applying the argument to the inverse gives \(h_f\in K_\Lambda\)
and then \(n\in K_\Lambda\). Since \(x_w\in K_\Lambda\)
and \(K_2\subset K_\Lambda\), this applies to the overlap above.
\releasecite{Formalism/LocalToGlobal/SiegelLatticeTheory/HaarTransversals/HidaSetting.lean}{207}{The rational flag and its lattice}
\releasecite{Formalism/LocalToGlobal/SiegelLatticeTheory/HaarTransversals/UnitaryHaarProducts/HermitianETransportInvariant.lean}{240}{Projection to the complementary lattice}.

Write \(e^0_a,f^0_a,z^0\) for the integral frames of \(K_2\),
and \(\beta_i,\beta_i^\vee\) for the ordered residual vectors and
their duals there. The equations of \(K_2\) modulo \(g\Lambda_f\)
are triangular on the \(e^0_a\)'s, have increment in their span
on \(z^0\), and permit on \(\beta_i\) increments in that span
and the preceding \(\beta\)'s. Use the given simultaneous frame
approximations. At each prime dividing \(g\), \(e^0_a,f^0_a\)
reduce to \(e_a,f_a\) and \(x_w z^0\) reduces to \(z_w\).
The element \(x_w\) fixes \(e_a,f_a\), so the first two
congruences hold for \(u=x_w^{-1}nx_w\) by the flag
conditions defining \(N_w^{\rm un}\).
Moreover \(\langle e_a,x_w\beta_i\rangle\) is divisible by
\(g\) by the marked-frame orthogonality. Subtracting the
corresponding \(f_a\)-components gives a vector in the exact
orthogonal to every \(e_a\) with the same reduction. Its
increment under \(n\) lies in their span by definition of the
unipotent group. This increment has integral coefficients, as
the \(f_a\)'s are integral dual vectors at these primes.
Transporting back gives the third congruence for \(u\), in
fact with increment in the span of just the \(e^0_a\)'s.
At primes away from \(g\) only the integral stabilizer condition
is needed. Hence \(u\in K_2\).
\releasecite{Formalism/LocalToGlobal/SiegelLatticeTheory/HaarTransversals/HidaSetting.lean}{118}{The triangular level}
\releasecite{Formalism/LocalToGlobal/SiegelLatticeTheory/HaarTransversals/UnitaryHaarProducts/HermitianETransportInvariant.lean}{548}{The transported flag, marked vector and residual congruences}.

The numerator of the finite level character on \(u\)
is the sum of the pairings
\(\langle f^0_a,u\,\operatorname{next}(e^0,z^0)_a\rangle\)
and \(\langle\beta_i^\vee,u\beta_{i+1}\rangle\) for existing
successor indices. Here \(\operatorname{next}\) uses the next
flag vector, and the marked vector at the last position.
Call the sum \(A\). Write \(B\) for the corresponding rational
flag sum \(\sum_a\langle f_a,n\,\operatorname{next}(e,z_w)_a\rangle\),
the numerator of the finite unipotent character. The residual
pairings of \(A\) are zero modulo \(g\) by the congruences and
the dual Gram matrix. Take a single integer \(C\), prime to \(g\),
clearing the denominators in the given frame congruences.
In particular \(C(x_wf^0_a-f_a)\) and
\(C(x_w\operatorname{next}(e^0,z^0)_a-
                  \operatorname{next}(e,z_w)_a)\) lie in \(g\Lambda_f\).
Expanding each difference of pairings gives
\(C^2(A-B)\in g\widehat{\mathcal O}_E\).
Also \(A\) is integral, and \(B\) belongs to
\(\mathfrak d^{-1}\widehat{\mathcal O}_E\).
For the latter statement, \(n-1\) preserves the finite
lattice and its entry from \(I_{a+1}e_{a+1}\) to \(I_ae_a\)
lies in \(I_aI_{a+1}^{-1}=\mathfrak d^{-1}\). At the last
index the marked vector is integral and the same conclusion
uses \(I_{r_S-1}=\mathfrak d^{-1}\).
The additive character is trivial on this trace-dual ideal.
It is therefore trivial on \(A-B\) and on \(C^2(A-B)/g\).
A B\'ezout identity for \(C^2\) and \(g\) gives triviality on
\((A-B)/g\). Thus \(\chi_2(u)=\psi_{w,f}(n)\).
Now \(x_w^{-1}h_fx_w=u^{-1}(x_w^{-1}nh_fx_w)\) is in \(K_2\),
and multiplicativity of \(\chi_2\) gives the stated level relation.
\releasecite{Formalism/LocalToGlobal/SiegelLatticeTheory/HaarTransversals/UnitaryHaarProducts/HermitianETransportInvariant.lean}{786}{The trace-dual ideal for the rational flag}
\releasecite{Formalism/LocalToGlobal/SiegelLatticeTheory/HaarTransversals/UnitaryHaarProducts/HermitianETransportInvariant.lean}{935}{and the character after transport}.

For two decompositions \(n h_f x_w k_2=n' h_f' x_w k_2'\), apply
this relation with \(h=(h_f')^{-1}h_f\) and
\(n''=(h_f')^{-1}(n')^{-1}n h_f'\), for which
\(x_w^{-1}n''h x_w=k_2'k_2^{-1}\). The column transforms
under \(h\) with multiplier \(\chi_2(x_w^{-1}h x_w)\),
and the Bessel character on \(n''\) is
\(\psi_{w,f}(n)/\psi_{w,f}(n')\).
Thus the two expressions for the seed agree.
\releasecite{Formalism/LocalToGlobal/SiegelLatticeTheory/HaarTransversals/PoincareInHH.lean}{156}{The prescribed Poincar\'e seed}
\releasecite{Formalism/LocalToGlobal/SiegelLatticeTheory/HaarTransversals/UnitaryHaarProducts/HermitianETransportInvariant.lean}{935}{The finite level relation}.

The archimedean Bessel covariance and the finite character
formula are inverse on the rational unipotent by the product
formula, and the rational action of the stabilizing column
group agrees on the two coefficient realizations. The seed is
therefore left invariant by \(\mathcal S_w(F)\).
Summing \(\Phi_{w,\phi}(\gamma g)\) over
\(\gamma\in\mathcal S_w(F)\backslash G_0(F)\) gives
the operator \(\mathcal P_w\phi\) in the same-word and mixed Gram
statements.

\teamanchor{proof.period-bessel-seed-fold}
The orbital proof uses the weighted integrability of the
chosen real section on the family with \(r\ge1801\) and
successive real \(b\)-gaps at least \(16r^{19}\).
\releasecite{Formalism/LocalToGlobal/SiegelLatticeTheory/HaarTransversals/ConcreteSettingExists.lean}{13}{The real gap hypothesis}.
Set
\[
 H_w(x)=\inf_{s\in \mathcal S_w(F_\infty)}
                        \prod_{\tau}\|s_\tau x_\tau\|_{\mathrm{op}},
 \qquad N=4r^2\sqrt{2q},
\]
where \(q\) is the maximum real negative index of \(Y_0\).
The local argument in that proof gives
\[
 \int\|W_w(x)\|H_w(x)^N\,d\bar x<\infty
\]
on the real Bessel quotient, and a bounded compactly supported
right convolution reproducing \(W_w\).
The integrable matrix-coefficient projector above also reproduces
\(W_w\) with the given Haar measure.
\releasecite{Development/AbelianVarietyTheory/Proofs/FlagMeasures/WuDecayOfGap.lean}{7}{The weighted estimate and reproduction}
use the actual frame and quotient measure.

Here is how the weighted bound gives absolute periodization.
Write \(G=G_0(\mathbb A_F)\), \(\Gamma=G_0(F)\) and
\(\Delta=\mathcal S_w(F)\). Use the compact convolution at infinity
and the normalized \(K_2\) average at finite places. They give a
compact set \(C\subset G\) and a constant \(a\) such that
\[
 |\Phi(x)|\le a\int_C|\Phi(xy)|\,dy,
 \qquad \Phi=\Phi_{w,\phi}.
\]
For a compact set \(B\) define
\(m_B(x)=\#\{\gamma\in\Gamma:x^{-1}\gamma x\in B\}\).
If the transporter from \(z\) to \(xC\) contains one element
\(\gamma_0\in\Gamma\), every other such element satisfies
\(x^{-1}\gamma\gamma_0^{-1}x\in CC^{-1}\).
Unfolding the nonnegative integral of \(|\Phi|\) against this
transporter gives
\begin{equation}
 \label{eq:period-bessel-absolute-fold}
 \sum_{\delta\in\Delta\backslash\Gamma}|\Phi(\delta x)|
 \le a\,m_{CC^{-1}}(x)\int_{\Delta\backslash G}|\Phi(z)|\,dz .
\end{equation}

On the support of the seed, first take the quotient by the
adelic stabilizing group. The finite coordinate is represented
by \(x_wK_2\). The rational quotient of the stabilizer is compact:
the unipotent rational quotient is compact, and the column
unitary group is definite. Thus representatives on the support
have finite coordinate in one compact set and real coordinate
\(c_\infty x\), where \(c_\infty\) is in a fixed compact set and
\(x\) can be chosen with \(\prod_\tau\|x_\tau\|_{\rm op}\le2H_w(x)\).
The quotient measure is obtained by integrating first on these
compact stabilizer fibers. The norm of the column coefficient
is bounded on its finite class set. For every compact \(B\),
an element counted by \(m_B\) has finite entries in a fixed
fractional ideal of \(E\), and its entries at \(\tau\) have
absolute value at most \(C_B\|x_\tau\|_{\rm op}^2\).
Indeed it is a conjugate by \(c_\infty x\) of an element of \(B\),
and the operator norms of \(x_\tau\) and its inverse agree
in the unitary frames. Counting integral lattice points in boxes
in the \(r^2\) complex matrices gives a constant times
\(\prod_\tau\|x_\tau\|_{\rm op}^{4r^2}\).
Since \(N\ge4r^2\), this and the weighted bound make both
\(\int_{\Delta\backslash G}|\Phi|\) and
\(\int_{\Delta\backslash G}m_B|\Phi|\) finite.

Unfold one factor of the square of the absolute series on
\(\Gamma\backslash G\). Its integral is
\(\int_{\Delta\backslash G}|\Phi(x)|
              \sum_\delta|\Phi(\delta x)|\,dx\).
Equation~\eqref{eq:period-bessel-absolute-fold} makes this finite.
Thus the periodization represents an \(L^2\) function. It has
the right character of \(K_2\). The integrable archimedean
projector can be interchanged with the periodization in \(L^2\),
using the \(L^2\) absolute series and invariance of Haar. It
fixes the result by the convolution of the seed, so the result
lies in the selected Petersson space. In its pairing with a
row form, absolute convergence and left rational invariance
join the summation domains into the quotient by \(\Delta\).
The seed's left equivariance there gives the Bessel integral
with the indicated column coefficient.
\releasecite{Formalism/LocalToGlobal/SiegelLatticeTheory/HaarTransversals/PoincareSeriesSquareIntegrable.lean}{52}{The absolute folding estimate}
\releasecite{Formalism/LocalToGlobal/SiegelLatticeTheory/HaarTransversals/UnitaryHaarProducts/InHHClosed.lean}{61}{the realization in the selected Petersson space} and
\releasecite{Formalism/LocalToGlobal/SiegelLatticeTheory/FlagStabilizerTensors/LinkSFinal.lean}{97}{the orbit unfolding}.

In the direct application of the Gram bounds we take their
operators on the prescribed columns. The change from the
arithmetic column to the adelic column reads the same values
on the transported finite cosets with the character of
\(x_w^{-1}h_fx_w\). These values determine the column by the
class covariance. The invariant Hermitian coefficient
pairings and class masses give the isometric ramified
inclusion \(\jmath_w:\mathcal M_w\hookrightarrow H_w^{\rm full}\).
The simultaneous comparison supplies the common row
isometry. In the notation of that comparison,
\(s:E_{\rm row}\hookrightarrow\mathcal H\) is
conjugate linear and isometric.
\releasecite{Formalism/LocalToGlobal/SiegelLatticeTheory/FlagStabilizerTensors/STiesR.lean}{19}{The actual column levels and maps}
\releasecite{Development/AutomorphicArithmetic/ThetaCorrespondence/ThetaPairLifts/ThirdPlaceSeries/AnchorPairs.lean}{24}{and the corresponding bottom pairing}. The conjugate-linear first-variable convention gives the
bilinear form
\[
 \mathcal E_w(\phi,F)
   =\big\langle s(F),\mathcal P_w\jmath_w\phi\big\rangle_{\mathcal H}.
\]

Write \(\pi_w\) for the tensor product of the compact \(U(k)\)-types
at the real places. Let \(v_N(w)\) be the volume of the unipotent
fundamental domain times the volume of the common ambient level, and
\(v_H(w)\) the volume of the real compact factor of the Bessel
stabilizer. Compact Schur orthogonality with probability measure
contributes \(1/\dim\pi_w\). Indeed averaging a rank-one operator
over that factor commutes with \(\pi_w\) and has the trace of the
original operator. Restoring the measure multiplies by \(v_H(w)\).
The quotient and level factors give \(v_N(w)\). The full operators
and reference measures are fixed by
\begin{equation}
\label{eq:period-gram}
 c_0(w)=\frac{v_N(w)v_H(w)}{\dim\pi_w}>0,\qquad
 L_w=c_0(w)^{-1/2}\mathcal P_w,\qquad
 K_{wv}=L_w^*L_v .
\end{equation}
Put \(d_w=\dim H_w^{\rm full}\) and \(r_w=\dim\mathcal M_w\). The
single-word estimate for \(K_{ww}\) and the value of \(r_w\) in
Section~\ref{sec:trace} will give a nonzero adelic reference pairing.

\teamanchor{proof.principal-test-separation}%
\subsubsection{The auxiliary grid and the detection problem}
Write \(Y_i=V_{i,0}\) for the first row of the auxiliary grid and
\(C_i\) for its core spaces. The ranks are \(r+i(r+6)\) and
\(r+1+i(r+5)\), respectively. The initial core is
\(C_0=Y_0\perp\langle1\rangle\). The inclusion
\(\iota:U(Y_0)\hookrightarrow U(C_0)\) fixes a norm-one vector \(z_0\)
in its orthogonal complement. The last core is identified with the
definite column vertex. The last row and core form the Bessel pair
used in \eqref{eq:period-principal-test}. At a real place the initial
row has negative index \(m/2\). These groups and the test family
are those of the auxiliary grid.

Fix the level, the Weil data and the quotient measures of the grid.
Let \(t\) be a nonzero element of the bottom cuspidal multiplicity
space at principal level and trivial finite level character.
Fix a nonzero vector \(w\) of its prescribed smooth real
representation. The weak separation assertion asks for a prescribed
finite coset tuple \(\Phi\), a permitted column \(c\) at the last
core, real Fock polynomials \(\rho\) and a real column covector
\(\xi\) for which the last Bessel period of
\(c_\xi\) and the first-row lift of \(t(w)\) is nonzero. Here \(c_\xi\)
is the scalar form obtained by evaluating the coefficient of \(c\)
against \(\xi\). The finite tuple may depend on \(t,w\).
\releasecite{Development/AbelianVarietyTheory/Definitions/TwistedInnerProducts/OffQOfSizeRows.lean}{119}{Statements of core detection, core span and ascent}.

\teamanchor{candidate.fixed-level-period}%
\subsubsection{The bottom corank-one period}
\textbf{AI-proposed candidate (detection at the prescribed principal level).}
For each size exponent \(\varepsilon\) and depth constant
\(C_F\ge6510\), there is a common threshold \(x_0(\varepsilon,C_F)\).
For every permitted size datum of exponent \(\varepsilon\) and
scale \(x\ge x_0(\varepsilon,C_F)\), take any compatible row level
and character family with an admissible base tuple, any auxiliary
grid at constants \((1,C_F,b)\) for that family, and its admissible
Weil data and row quotient Haar measures. Let \(M_Y\) and \(M_C\)
be models of the prescribed bottom row and core discrete-series
types. Let \(t\ne0\) be any cuspidal multiplicity map for \(M_Y\)
at the full row principal subgroup on \(\Lambda_Y\) of modulus
\((g_0g_1)\), with trivial finite character. For every nonzero
vector \(w\) in the tensor product of the smooth real modules of
\(M_Y\), there exist a cuspidal multiplicity map \(y\) for \(M_C\)
of trivial finite character at
\(K_C=K(\Lambda_Y\perp\mathcal O_E z_0,\mathfrak m)\), and a
\(K_\infty\)-finite vector \(v\) in its smooth tensor model, such
that \(a=y(v)\) satisfies
\begin{equation}
 \label{eq:principal-core-detection}
 \int_{[U(Y_0)]} a(\iota h)t(w)(h)\,dh\ne0 .
\end{equation}
Here \(S_0\) consists of the chosen place \(v_0\) over the least
split auxiliary rational prime \(q_0\) of the grid and the places
over \(g_0g_1p_{\rm de}p'_{\rm de}\). With one specified split
prime \(\mathfrak w_v\) over each \(v\in S_0\), the core modulus
is \(\mathfrak m=\prod_{v\in S_0}\mathfrak w_v^{d_v}\),
where \(d_v=2\) away from \(v_0\) and
\[
 d_{v_0}=1+
 \left\lfloor
 \frac{C_F[F:\mathbb Q]^2(r+1)^6
    \log\bigl(b(r+1)g_0g_1p_{\rm de}p'_{\rm de}|\Delta_E|\bigr)}
      {\log q_0}
 \right\rfloor .
\]
The modulus, row measure and real types are fixed before \(t,w\).
\releasecite{Development/AbelianVarietyTheory/Definitions/TwistedInnerProducts/OffQOfSizeRows.lean}{96}{Quantifiers for detection at the grid level}
\releasecite{Development/AbelianVarietyTheory/Definitions/Divisors/BasePlaceNotMemChiGridBadPlaces.lean}{78}{Core modulus}
\releasecite{Development/AbelianVarietyTheory/Definitions/Divisors/LeastSplitAuxiliaryPrimeThetaChainSpec.lean}{250}{Depth at the auxiliary place}.

The real branching pair uses the dual of the row discrete series
because the period is bilinear. Xue gives its local dimension
criterion. \cite[Theorem~1.1]{p2xue-bessel}
The global corank-one period theorem concerns a fixed Hermitian
Arthur parameter and uses central value nonvanishing. Its existence
result with one rank fixed takes the larger-rank parameter as data.
The displayed detection fixes each lower form at the stated level
and seeks an upper witness in a prescribed real type and level.
\cite[Theorems~1.1.5.1 and~1.1.5.3]{depth-endoscopic-ggp}
The argument below proves this first step for the admitted grid.
The next step is to span the prescribed core test space by
descending chains, then ascend with the permitted row tests.

For \(U(p,q)\) and the subgroup \(U(p-1,q)\) fixing a
positive line, He identifies the discrete part of the restriction
of each discrete series. It is the direct sum of the discrete series
of the subgroup whose Harish--Chandra parameters satisfy the signed
interlacing rule, each with multiplicity one.
\cite[Theorem~0.1]{depth-he}
To verify the signed interlacing, merge the core and
row-dual parameters in decreasing order. The prescribed
word is
\[
 a_0,\ldots,a_m,\quad
 v_2,v_2-\tfrac12,\ldots,v_k,v_k-\tfrac12,\quad
 -b_1,\ldots,-b_m.
\]
The \(a_i\) are core entries with signs alternating from
positive. All entries in the middle block are positive,
with \(v_i\) from the row dual and \(v_i-\tfrac12\) from
the core. The final block consists of row-dual entries
with signs alternating from positive. The strict ordering
uses \(a_m>v_2\), \(v_i-v_{i+1}\ge1\) and \(v_k+b_1\ge1\).
Since \(m\) is even, \(a_m\) has positive sign. Thus two
successive entries have the same sign precisely when
they belong to different groups, including at the two
block boundaries. This is the signed interlacing condition
in He's theorem for deletion of a positive line.
\releasecite{Development/AbelianVarietyTheory/Proofs/TwistedOperators/InterlacedSignedParameters.lean}{8}{Merged real parameter word}
\releasecite{Development/AbelianVarietyTheory/Proofs/TwistedOperators/BranchingHypothesesGridParametersLargeScale.lean}{80}{Matching the word to the actual models}.

Put \(d=[F:\mathbb Q]\) and \(n=r+1\). In the adapted real frames,
write \(|z|_v\) for the positive Euclidean norm of a vector of
\(C_0\) at a real place \(v\), normalized by \(|z_0|_v=1\).
The needed coefficient \(k_\infty\) is the tensor of matrix
coefficients of translated lowest-compact-type vectors in the
specified ambient discrete series. It must reproduce the value of
\(f=t(w_*)\) for a fixed nonzero real row vector \(w_*\). In the
real Haar factors of \(dh\), the required identities and bounds are
\begin{align}
 \int_{U(Y_0)_\infty} k_\infty(\iota h)f(xh)\,dh_\infty
   &= c_\infty f(x), \qquad c_\infty\ne0,
   \label{eq:principal-core-reproducing}\\
 \int_{U(Y_0)_\infty}|k_\infty(g\iota h)|\,dh_\infty
   &\le e^Q|c_\infty|
      \prod_{v\mid\infty}|g_v z_0|_v^{-(2n+5)} .
   \label{eq:principal-core-transversal}
\end{align}
The first variation along the row Lie algebra has an
\(L^1\)-bound with constant \(e^Q|c_\infty|\). Moving that variation
to the coefficient in \eqref{eq:principal-core-reproducing} gives
the row Lipschitz constant \(e^Q\|f\|_\infty\). The required
uniform range is \(Q=C d n^4\log b\) with an absolute \(C\).
The vector before translation must also have sufficient
matrix-coefficient decay to form a cuspidal Poincar\'e series.
The lower bound on the Hilbert projection and the
transversal and first-variation bounds use the same translated
vector. The reproducing identity follows from branching and
Schur orthogonality.

At one real place let \(H=U(Y_0)(\mathbb R)\), let \(\pi\)
be the ambient representation restricted to \(H\), and let
\(\sigma\) be the row dual. He gives a unique Hilbert copy
\(j\sigma\) of \(\sigma\) inside \(\pi\), with \(j\) an
isometry. Take the specified translated ambient vector \(v\),
with \(j^*v\ne0\) and
\(k_v(h)=\langle v,\pi(h)v\rangle\) in \(L^1(H)\). Inner
products are linear in the second argument. Let \(u_0\)
be a unit lowest-compact-type vector in \(\sigma\).
The real weight gaps of the row imply that its diagonal
coefficient is integrable. Write \(\mathcal D\) for the
linear span of its \(H\)-translates, which is dense in
\(\sigma\). All coefficients with both vectors in
\(\mathcal D\) are integrable by left and right translation.

For \(a,c\in\mathcal D\) let
\(e_{a,c}(h)=d_\sigma\overline{\langle a,\sigma(h)c\rangle}\),
with \(d_\sigma\) the formal degree for the chosen Haar
measure. Schur orthogonality gives
\(e_{a,c}*e_{a',c'}=\langle c,a'\rangle e_{a,c'}\) and
\(e_{a,c}^*=e_{c,a}\). On any unitary representation
\(\tau\) of \(H\), integrate \(\tau\) against these \(L^1\)
functions. If \(\tau(e_{u_0,u_0})\) has nonzero image,
choose a unit vector \(z\) in its image. The matrix-unit
identities make \(a\mapsto \tau(e_{a,u_0})z\) an isometry
on \(\mathcal D\). Left translation of \(e_{a,u_0}\) by
\(g\in H\) gives \(e_{\sigma(g)a,u_0}\), so the isometry
extends to a copy of \(\sigma\) in \(\tau\). Applied to
the orthogonal complement of \(j\sigma\), uniqueness shows
that \(\pi(e_{u_0,u_0})\) vanishes there. Factoring each
\(e_{a,c}\) through this projection gives the same
vanishing for all matrix units. On \(j\sigma\) they act
by the usual rank-one operators.

Taking the coefficient against \(v\) gives
\[
 \int_H \overline{k_v(h)}\,\langle a,\sigma(h)c\rangle\,dh
     =d_\sigma^{-1}\langle a,j^*v\rangle\langle j^*v,c\rangle .
\]
The integral extends continuously to all \(a,c\) by
\(k_v\in L^1(H)\). With \(c=j^*v/\|j^*v\|\), this is
the reproducing eigenvector with eigenvalue
\(\|j^*v\|^2/d_\sigma\). The vector is smooth since \(j^*\)
is a bounded \(H\)-intertwiner and \(v\) is smooth. Transporting
back to the row model conjugates the two coefficients in
the bilinear convention. Multiplying the identities over
the real places gives \eqref{eq:principal-core-reproducing}
in the prescribed measures.

The smooth tensor in the row multiplicity space is algebraic.
Its map \(t\) extends to the Hilbert tensor, as follows. Freeze
all but one real factor. Point evaluation is continuous on the
remaining smooth Fr\'echet space by the defining Sobolev
condition on multiplicity maps. Every such cuspidal value is a
bounded function by rapid decay for finite-level smooth cusp
forms of moderate growth, finite under the center.
See Miller--Schmid~\cite[Theorems~A and~2.23--2.24]{p2miller-schmid-decay}.
The uniform boundedness
principle applied to all point evaluations therefore gives
continuity of this one-variable map into bounded functions and
hence into \(L^2([U(Y_0)])\), since the quotient has finite volume.
The pullback of the \(L^2\) inner product is invariant on the
compact-finite module. Such a form is a nonnegative scalar
multiple of the model form: on each finite-dimensional compact
isotypic space compare the forms by the Riesz operator, then
invariance makes the combined operator an endomorphism of the
irreducible Harish--Chandra module. Density of compact-finite
vectors in the smooth topology extends the identity to the
smooth factor. On pure tensors, applying the identity in each
factor in turn gives
\[
 \|t(u_1\otimes\cdots\otimes u_d)\|_2^2
     =c_t\prod_j\|u_j\|^2 .
\]
Separate polarization gives the same scalar times the Hilbert
inner product on all pairs of algebraic tensors. The scalar is
positive for \(t\ne0\), since continuous functions zero in
\(L^2\) vanish. Thus \(t\) extends as a scalar isometry and
intertwines the real group on the Hilbert tensor.

First apply the integrated Schur
operator on the Hilbert space of the row representation and
use the automorphic intertwiner. The pointwise integral against
the smooth bounded row form is continuous by the \(L^1\)-bound
on the kernel. By Fubini against compactly supported continuous
functions on the quotient it agrees with the Hilbert-space
integral. The resulting equality of continuous functions holds
at each \(x\).
\releasecite{Development/AutomorphicArithmetic/ThetaCorrespondence/RealBands/BranchDefs.lean}{11}{Quantitative real input}
\releasecite{Development/AutomorphicArithmetic/ThetaCorrespondence/RealBands/BranchInput.lean}{60}{Integrable coefficients at the actual parameter gaps}
\releasecite{Development/AutomorphicArithmetic/ThetaCorrespondence/FockModels/RealBranch.lean}{140}{Quantitative real projection}
\releasecite{Development/AutomorphicArithmetic/ThetaCorrespondence/RealBands/BranchPack.lean}{107}{The prescribed real types and frames}
\releasecite{Development/AbelianVarietyTheory/Definitions/TwistedInnerProducts/DepthPays.lean}{94}{Real reproducing normalization}.

\subsubsection{Quantitative real branching and the period kernel}\label{sec:principal-real-bounds}
Work at one real place. Write \(G=U(p,q)\), \(H=U(p-1,q)\), and
\(N=p+q=n\) for the ambient dimension. Use the adapted frame with
signature matrix \(J=\operatorname{diag}(\epsilon_1,\ldots,\epsilon_N)\).
The inclusion fixes the positive coordinate vector \(z=e_{i_0}\).
Let \(\lambda\) and \(\mu\) be the ambient and row-dual Harish--Chandra
parameters, respectively, in decreasing order with their signs. We use
the following size properties of the grid:
\[
 \lambda_i-\lambda_{i+1}\ge 2,\qquad
 |\lambda_i|,|\mu_j|\le B=b^2,\qquad
 \log(8NB)\le 6\log b .
\]
Within either parameter the gap between entries of opposite signs is
at least \(1000N^3\). Indeed the \(a_i\) have adjacent gap at least
\(1000N^3\), as do the \(b_i\), and \(a_m\) is separated from the middle
block by at least \(1000N^3\). In the row dual the positive middle block
is separated from the first negative entry in the final block by the
gap \(b_2-b_1\), in addition to \(v_k+b_1\ge 2\). The middle block in
each parameter has successive spacing at least \(2\). The parameters
are integral after subtracting the respective half-sum of positive
roots. The signature includes at least one positive and one negative
coordinate in each group. These are the real grid estimates at the
chosen scale.
\releasecite{Development/AbelianVarietyTheory/Proofs/TwistedOperators/BranchingHypothesesGridParametersLargeScale.lean}{34}{Gaps and size of the actual real grid}

Let \(P=j^*\), and take unit vectors \(v_0\) and \(u_0\) in the lowest
compact types of \(\pi\) and \(\sigma\). Use Haar measures \(dg\) and
\(dh\) in the two groups and denote the corresponding formal degrees
by \(d_\pi\) and \(d_\sigma\). Write \(\|\cdot\|_F\) and
\(\|\cdot\|_{\mathrm{op}}\) for the Frobenius and Euclidean operator
norms of matrices. All implicit constants below are absolute. Set
\(T=8NB\) and \(L=T^2\).

\emph{Derived vectors in the lowest compact type.}
Put \(E_{ab}=d\pi(e_{ab})\) on the smooth space. Infinitesimal
unitarity gives \(E_{ab}^*=\epsilon_a\epsilon_b E_{ba}\) there. Let
\(C_G=\sum_{a,b} E_{ab}E_{ba}\) and let \(C_K\) be the same sum over
equal-sign pairs. For any smooth vector \(y\),
\[
 \sum_{a,b}\|E_{ab}y\|^2=\langle y,(2C_K-C_G)y\rangle .
\]
On the lowest type the highest weight is
\(\lambda+\rho_{\mathrm{nc}}-\rho_K\), with each coordinate bounded in
absolute value by \(B+N/2\). The space generated from this type by at
most \(a\) derived factors is a quotient of a direct sum of its tensor
products with \(\mathfrak g_{\mathbb C}^{\otimes j}\), \(j\le a\), as
a compact module. The weights of the initial type lie in the convex
hull of its compact Weyl orbit. Each adjoint weight has coordinates of
absolute value at most \(1\). Hence each highest weight \(\nu\) of the
generated space satisfies \(|\nu_i|\le B+N/2+a\). The Casimirs have
eigenvalues \(\|\nu+\rho_K\|^2-\|\rho_K\|^2\) and
\(\|\lambda\|^2-\|\rho_G\|^2\), respectively. For \(a\le N+2\), the
square root of the upper bound for \(2C_K-C_G\) is at most \(L\).
Cauchy--Schwarz now gives
\begin{equation}
 \|d\pi(X_1)\cdots d\pi(X_j)v_0\|
 \le L^j\prod_{i=1}^j\|X_i\|_F\qquad (j\le N+2).
 \label{eq:real-branch-derived}
\end{equation}
\releasecite{Development/WeilBound/CoreStatements/OrbitDecayEstimates/CurveDerivativeDecay/KFiniteWordBoundIndefUnitary.lean}{20}{Derived vector estimate}

\emph{A moment in the transverse coordinate.}
Let \(F(g)=\langle j u_0,\pi(g)v_0\rangle\). The required moment is
\begin{equation}
 \int_G \|g^{-1}z\|^2 |F(g)|^2\,dg
 \le T^{8N}\int_G |F(g)|^2\,dg=T^{8N}/d_\pi .
 \label{eq:real-branch-transverse-moment}
\end{equation}
Use \(\Pi=\pi\otimes (\mathbb C^N)^*\) and put
\(x_a(y)=y\otimes e_a^*\). On coordinate tuples \(x\) its action is
\((\Pi(g)x)_a=\sum_l(g^{-1})_{la}\pi(g)x_l\). Its invariant Hermitian
form is \([x,y]=\sum_a\epsilon_a\langle x_a,y_a\rangle\).

The Harish--Chandra module of \(\Pi\) is the direct sum of the
discrete-series modules \(\pi_i\) with parameters \(\lambda-e_i\),
in the same signed chamber. To see this, use the upper triangular
\(\theta\)-stable Borel in the ordered frame. The inducing character
of \(\pi\) on its compact Cartan has weight \(\lambda-\rho_G\).
Apply the finite-dimensional tensor identity for cohomological
induction to the Borel filtration of \((\mathbb C^N)^*\). Its graded
characters have weights \(-e_i\). Every shifted parameter is regular
in the original chamber, since the ambient gaps are at least \(2\).
Good-range concentration thus makes the degree-\(\dim(\mathfrak n
\cap\mathfrak k)\) long exact sequences short exact. Each graded factor
is the corresponding discrete series. Their infinitesimal characters
are distinct, since each \(\lambda-e_i\) remains decreasing and the
tuples are different. The primary decomposition for the center splits
the filtration. This uses the tensor identity and filtered complex
in~\cite[Proposition~3.6 and (14)]{real-branch-prudhon} and the
good-range Borel construction in~\cite[Theorems~5.35, 5.99 and~8.2,
and Chapter~XI]{real-branch-knapp-vogan}. The smooth spaces give the
same direct sum by Casselman--Wallach globalization
\cite[Theorems~1.1 and~10.6]{real-branch-bernstein-kroetz}.

Write \(p_i\) for the projections and \(R_i\) for the maps onto the
smooth unitary models of \(\pi_i\). The summands are orthogonal for
\([\cdot,\cdot]\). Indeed the total quadratic Casimir is symmetric
for this form and its eigenvalues on the summands are real and
distinct. On each summand the form is a nonzero real multiple of
the unitary inner product, by irreducibility and Schur's lemma on
compact-finite vectors. Nondegeneracy follows from the direct sum
and orthogonality. Rescale \(R_i\) so that
\begin{equation}
 [x,p_i y]=s_i\langle R_i x,R_i y\rangle,\qquad s_i\in\{1,-1\}.
 \label{eq:real-branch-tensor-form}
\end{equation}

The matrix operator \(\mathcal E=(E_{ab})\) on tuples satisfies
\(C_{\Pi}=C_G I+NI-2\mathcal E\), by expanding the action on the
tensor with the dual standard representation. Thus its eigenvalue
on \(p_i\) is \(\lambda_i+(N-1)/2\). Let \(Q_G(u)\) be the Capelli
column determinant with column \(l\) shifted by \(u-l+1\). Its
infinitesimal-character eigenvalue is
\(\prod_j(u+\lambda_j-(N-1)/2)\). Let \(A(u)\) be the matrix whose
\((a,l)\) entry is the signed Capelli minor omitting row \(l\) and
column \(a\), with consecutive column shifts starting from \(u\).
The Capelli comatrix identity gives
\[
 A(u)(\mathcal E+(u-N+1)I)=Q_G(u)I .
\]
This is the specialization of the antisymmetrizer identity in
\cite[Proposition~2.7(4), Remarks~2.8(1) and~2.11]{real-branch-mno} under
\(T(u)\mapsto I+\mathcal E/u\). Group the product by the row of
its last factor, after fixing the column sequence consisting of
all columns except \(a\) followed by the final column. This gives
the stated sign and Capelli shifts. The eigenvalue of \(Q_G(u)\)
also follows by evaluation on a highest weight vector: the diagonal
contribution in column \(l\) is \(u-l+1+\lambda_l-\rho_{G,l}\).
The largest displaced column in each other summand gives a raising
operator acting on the highest weight vector after scalar factors.
Evaluating the comatrix identity on each eigenspace and taking
\(u_i=-\lambda_i+(N-1)/2\) gives
\[
 p_i=\frac{A(u_i)}{\prod_{k\ne i}(\lambda_k-\lambda_i)}.
\]
The denominators have absolute value at least \(1\). By
\eqref{eq:real-branch-derived}, each signed minor applied to \(v_0\)
has norm at most \((N-1)!(B+N+L)^{N-1}\le T^{4N}\). In particular,
\(|s_i|\|R_i x_l(v_0)\|^2\le T^{4N}\) by
\eqref{eq:real-branch-tensor-form} and the \((l,l)\) minor.

The \((i_0,i_0)\) minor is the Capelli determinant of the complement
of \(z\). On smooth vectors the bounded intertwiner \(P\) commutes
with its derived action. Comparing the two Harish--Chandra shifts
of rank \(N\) and \(N-1\) therefore gives
\begin{equation}
 P\bigl((p_i x_{i_0}(y))_{i_0}\bigr)
   =\gamma_i P y,qquad
 \gamma_i=\frac{\prod_j(\mu_j-\lambda_i+1/2)}
                  {\prod_{k\ne i}(\lambda_k-\lambda_i)},\qquad
 |\gamma_i|\le T^N .
 \label{eq:real-branch-subgroup-minor}
\end{equation}
Put \(\xi=j u_0\). The \(L^1(H)\) matrix unit
\(e=d_\sigma\overline{\langle u_0,\sigma(\cdot)u_0\rangle}\) from
the preceding argument acts on \(\pi|_H\) as the orthogonal
projection onto \(\mathbb C\xi\). It acts on \(\pi_i|_H\) as an
orthogonal projection \(Q_i\), by the self-adjoint idempotent
identities. Choose a smooth vector \(\eta\) with
\(\langle\eta,\xi\rangle=1\) and put \(a=x_{i_0}(\eta)\). Since
\(H\) fixes this coordinate, integrating the equality in
\eqref{eq:real-branch-tensor-form} on \(\Pi(h)y\) against \(e(h)\)
gives, for smooth tuples \(y\),
\[
 \langle\xi,(p_i y)_{i_0}\rangle
    =s_i\langle b_i,R_i y\rangle,\qquad b_i=Q_i R_i a .
\]
With \(y=a\), \eqref{eq:real-branch-subgroup-minor} yields
\(s_i\|b_i\|^2=\gamma_i\). This averaging takes place in the
unitary Hilbert spaces of the summands. For each \(l\) we obtain
\[
 (g^{-1})_{l i_0} F(g)
   =\sum_i s_i\langle b_i,\pi_i(g) R_i x_l(v_0)\rangle .
\]
Schur orthogonality bounds the square integral of each summand by
\(T^{5N}/d_{\pi_i}\). Harish--Chandra's formal-degree formula on
the same group and for the same Haar measure is proportional to
\(\prod_{a<b}|\lambda_a-\lambda_b|\), with a common factor for all
these parameters. For the compact central factor, Haar integration
of the squared coefficient contributes the same factor in every
representation, since the central character is unitary. Decreasing
one coordinate by \(1\) changes only
\(N-1\) factors and each ratio is at most \(2\), because the original
differences have absolute value at least \(2\). Thus
\(d_\pi/d_{\pi_i}\le 2^{N-1}\le T^N\).
\cite[(3.10), (3.12) and Theorem~(8.1)]{ar4-atiyah-schmid}
Cauchy--Schwarz over \(i\), followed by summation over \(l\),
introduces at most \(N^3\) in the bound \(T^{6N}/d_\pi\).
This proves \eqref{eq:real-branch-transverse-moment}.
\releasecite{Development/WeilBound/RepresentationTheoryGL/HolomorphicDiscreteSeries/SchurEigenProjections/HolomorphicMoments/MomentHigherEstimateHolds.lean}{28}{Weighted Schur estimate}

\emph{A translate normalized by its reproducing scalar.}
For a matrix unitary group \(D\) of matrix size \(r_D\), let
\(B_D(\delta)=\{g\in D:\|g-I\|_F\le\delta\}\) and
\(D(R)=\{g\in D:\|g\|_F\le R\}\). Any Haar measure on it satisfies
\begin{equation}
 \operatorname{vol}D(R)
 \le(1+2R^2/\delta)^{2r_D^2}\operatorname{vol} B_D(\delta).
 \label{eq:real-branch-volume}
\end{equation}
Indeed \(g^{-1}=Jg^*J\) gives \(\|g^{-1}\|_F=\|g\|_F\). A maximal
\(\delta/R\)-separated subset of \(D(R)\) gives a cover by its left
translates of \(B_D(\delta)\) and the Euclidean packing bound in
real dimension \(2r_D^2\) gives the number of points. With
\(\delta=1/4\) and dyadic shells, \eqref{eq:real-branch-volume} gives
\begin{equation}
 \int_D\|g\|_F^{-u}\,dg
 \le 2\,33^{2r_D^2}\operatorname{vol}B_D(1/4),
 \qquad u\ge 4r_D^2+1.
 \label{eq:real-branch-power-integral}
\end{equation}

Let \(Y=2T^{8N}\). By \eqref{eq:real-branch-transverse-moment}, the
region \(S=\{g:\|g^{-1}z\|^2\le Y\}\) carries square integral at
least \(1/(2d_\pi)\). This region is left \(H\)-invariant. A vector
\(u\) of Hermitian norm one is the image of \(z\) under a matrix
\(b_u\) with \(\|b_u\|_F^2\le N+2\|u\|^2\). To construct it, take
\(u=u_++u_-\) in the positive and negative subspaces, with lengths
\(a_+,a_-\) satisfying \(a_+^2-a_-^2=1\). First map \(z\) to the
direction of \(u_+\) by an element of the compact group, then act
in the span of the two directions by the block with diagonal
\(a_+\) and off-diagonal \(a_-\). The compact map alone suffices
when \(u_-=0\).

Apply this with \(u=g^{-1}z\) for \(g\in S\). Since \(g b_u\in H\),
every left coset in \(S\) has a representative of Frobenius norm
at most \(R_0=(N+2Y)^{1/2}\). Put \(R_G=2R_0\). The left translates
of each representative by \(B_H(1/4)\) belong to \(G(R_G)\).
Quotient integration with \(dg\) and \(dh\) bounds the quotient
volume of \(H\backslash S\) by
\(\operatorname{vol}G(R_G)/\operatorname{vol}B_H(1/4)\). It also
identifies the square integral on each fiber with
\(\|P\pi(g)v_0\|^2/d_\sigma\), by Schur on \(H\). Consequently some
representative \(g_*\) has \(\|g_*\|_F\le R_0\) and
\begin{equation}
 c=\|P\pi(g_*)v_0\|^2/d_\sigma\ge
  \frac{\operatorname{vol}B_H(1/4)}
       {4d_\pi\operatorname{vol}G(R_G)} .
 \label{eq:real-branch-normalized-lower}
\end{equation}
Here the factor \(4\) allows choosing a value at least half the
average, and the fiber integral is independent of the representative.

For \(\delta=(4L)^{-1}\), the logarithm of \(g\in B_G(\delta)\)
belongs to the real Lie algebra and has Frobenius norm at most
\(2\delta\). This follows from the convergent logarithm series and
\(g^{-1}=Jg^*J\). Integrating the derived orbit along the exponential
and using \eqref{eq:real-branch-derived} shows that
\(|\langle v_0,\pi(g)v_0\rangle|\ge 1/2\) on this ball. Schur on
\(G\) gives \(d_\pi\operatorname{vol}B_G(\delta)\le 4\).
As \(R_G\le T^{O(N)}\), \eqref{eq:real-branch-volume} and
\eqref{eq:real-branch-normalized-lower} yield
\begin{equation}
 \|g_*\|_F\le T^{O(N)},\qquad
 c\ge \operatorname{vol}B_H(1/4)\exp(-C_0N^3\log T)
 \label{eq:real-branch-local-scale}
\end{equation}
for an absolute \(C_0\).

\emph{Decay of the lowest-type coefficient.}
For a group of matrix size \(M\) in this argument, let \(\beta\) be its
Harish--Chandra parameter with opposite-sign gaps at least
\(1000M^3\). If \(\tau\) is the representation and \(V\) is its
lowest compact type, we claim
\begin{equation}
 \|P_V\tau(g)|_V\|_{\mathrm{op}}
  \le \exp(100M^3)\|g\|_F^{-40M^2} .
 \label{eq:real-branch-lowest-decay}
\end{equation}
Use the polar decomposition \(g=k\exp X\), with \(X\) Hermitian
and supported in the off-sign blocks. The compact factor is
isometric and preserves \(V\). The unitary group is connected
reductive with finite center on its derived group. Sun's bottom-layer
formula applies to the compression of \(\exp X\) for the upper
triangular \(\theta\)-stable Borel and the Cartan character of weight
\(\beta-\rho_G\). The bottom type is the lowest compact type, with
highest weight \(\beta+\rho_{\mathrm{nc}}-\rho_K\), and occurs with
multiplicity one. The weight is dominant for the induced compact
ordering because the entries of \(\beta\) are decreasing with
successive gaps at least \(1\).
\cite[Theorem~6.3 and (25)--(26)]{real-branch-sun}

In the formula take \(x=\exp(X/2)\). For \(l\) in the compact
group, the map \(H_{\mathfrak q}\) extracts the diagonal from the
upper triangular factor of \(l x\) modulo the compact real form
of the complex group. Thus \(\widetilde H_{\mathfrak q}(l x)\) is
the positive diagonal \(D\) in the factorization
\(\exp(l X l^{-1})=u D u^*\) with \(u\) upper unitriangular.
The scalar \((\delta_{\mathfrak n}\chi)(D)\) in Sun's formula has
weight \(\beta+\rho_G\), since \(\delta_{\mathfrak n}\) contributes
\(2\rho_G\). If \(D_i(\exp Z)\) denotes the principal minor on the
coordinates \(i+1,\ldots,M\), this scalar is
\[
 \prod_{i=1}^{M-1}D_i(\exp Z)^{-(\beta_i-\beta_{i+1}+1)},
 \qquad Z=l X l^{-1}.
\]
Indeed the trailing determinant is the product of the corresponding
diagonal entries of \(D\), and \(\det D=1\). The other factor in
the integral is the conjugate of the projection onto a highest
weight line. With compact probability measure, its average is
\(I_V/\dim V\). It follows by positivity that the norm of the
compression is bounded by the supremum of the displayed scalar,
with the factor \(\dim V\) in Sun's formula canceled by this
average.

To bound a minor, view it as the expectation of \(e^U\) in the
spectral probability measure of the exterior-power derived operator
of \(Z\) on the unit wedge for the trailing set. If
\(t=\|X\|_{\mathrm{op}}\), then \(|U|\le Mt\). Its mean is zero
because the diagonal of \(Z\) is zero. Its second moment is
\(\sum_{a\le i<b}|Z_{ab}|^2\), as seen by replacing one vector in
the wedge. For \(\theta=(\max(1,Mt))^{-1}\), use
\(e^v\ge 1+v+v^2/4\) for \(|v|\le 1\), followed by
\(\log(1+h)\ge h/2\) for \(0\le h\le 1\). Convexity of the
logarithm of the moment-generating function then gives
\[
 \log D_i(\exp Z)\ge \tfrac18\theta\sum_{a\le i<b}|Z_{ab}|^2.
\]
The sum of these variances with weights \(\beta_i-\beta_{i+1}+1\)
has coefficient \(\beta_a-\beta_b+b-a\) at \(|Z_{ab}|^2\).
Every nonzero entry has opposite signs, so the sum is at least
\(1000M^3\|Z\|_F^2/2\). For \(t\ge 1\) the compression norm is
therefore at most \(\exp(-125M^2t/2)\). For \(t<1\) it is at
most \(1\) by unitarity. The inequality
\(\|g\|_F\le\sqrt M e^t\) proves \eqref{eq:real-branch-lowest-decay}.
This argument also supplies the row integrability used above by
\eqref{eq:real-branch-power-integral}, with \(M=N-1\).
\releasecite{Development/AbelianVarietyTheory/Proofs/TwistedOperators/LowestKtypeCoefficientExplicitPowerDecay.lean}{124}{Lowest-type power decay}

\emph{Transversal integrals and first variation.}
Let \(v=\pi(g_*)v_0\), \(k(g)=\langle v,\pi(g)v\rangle\),
\(G_0=\|g_*\|_F\), and \(D_0=40N^2\). Conjugating
\eqref{eq:real-branch-lowest-decay} by \(g_*\) gives
\[
 |k(g)|\le M_0\|g\|_F^{-D_0},\qquad
 M_0=\exp(100N^3)G_0^{2D_0} .
\]
For \(a\in G\), put \(\ell=\|a z\|\ge 1\). By the same block
construction write \(a=b_{az}h_1\) with \(h_1\in H\). Setting
\(x=b_{az} h\), we have \(\|x\|_F\ge\ell\) and
\(\|h\|_F\le \sqrt{N+2}\,\ell\|x\|_F\). For
\(u=4(N-1)^2+1\), this gives
\[
 \int_H |k(a h)|\,dh
 \le M_0(N+2)^{u/2}\ell^{-D_0+2u}
       \int_H\|h\|_F^{-u}\,dh.
\]
In particular \(D_0-2u\ge 2N+5\).

For \(X\) in the Lie algebra of \(H\) with \(\|X\|_F\le 1\),
let \(A_h(t)=k(\exp(-tX)h)\). The bound for the second derived
word in \eqref{eq:real-branch-derived} and
\(\|\operatorname{Ad}(g_*^{-1})X\|_F\le G_0^2\) give
\(\sup_t |A_h''(t)|\le L^2G_0^4\). For \(|t|\le 1\) we also
have \(|A_h(t)|\le M_0 e^{D_0}\|h\|_F^{-D_0}\), since the
operator norm of \(\exp(tX)\) is at most \(e\). Taylor's formula
at \(t=\|h\|_F^{-D_0/2}\) bounds
\[
 |A_h'(0)|\le
 (2M_0 e^{D_0}+L^2G_0^4/2)\|h\|_F^{-D_0/2} .
\]
Both integrals are bounded using \eqref{eq:real-branch-power-integral}.
Dividing by the same \(c\) of \eqref{eq:real-branch-local-scale},
we get the stated transversal and first-variation bounds with a
single \(\exp(C_1N^3\log T)\). For the product over real places
they give one \(Q\le C d n^3\log b\), with
\(d=[F:\mathbb Q]\). The weaker exponent \(C d n^4\log b\) in
the depth argument also suffices. Using the reproducing identity at
\(x\exp(tX)\), translation puts the variation on the kernel as
\(k(\exp(-tX)h)-k(h)\). Its \(L^1\)-norm is at most \(|t|\) times
the first-variation bound, by integration along that curve and
left invariance of Haar measure. This gives the stated Lipschitz
bound for \(f\) and supplies the transversal estimate for the
Poincar\'e seed.
\releasecite{Development/AbelianVarietyTheory/Proofs/TwistedOperators/TranslateCoefficientIntegralBoundsOfUnitaryLaws.lean}{166}{Transversal integration}
\releasecite{Development/AbelianVarietyTheory/Proofs/TwistedOperators/TranslateCoefficientDerivativeIntegralBoundsOfLa.lean}{225}{Integrated first variation}

\subsubsection{The Poincar\'e test at the prescribed depth}

For the level and the row quotient, put
\(A=r^2 d\, g_0g_1\,|\Delta_E|^{1/[E:\mathbb Q]}\).
The prescribed \(\mathfrak m\) must allow a real number \(0<s\le1/2\)
with
\begin{equation}
 \label{eq:principal-core-depth}
 e^Q A(2s)^{1/m}\le\tfrac12,\qquad
 2e^{Q+(6n+9)d}\le R,\qquad
 R=s^d\sqrt{\mathrm N\mathfrak m}>1 .
\end{equation}
Take \(s=\tfrac12(2e^Q A)^{-m}\), so the first bound is
an equality. With \(L=\log b\), the size hypotheses give
\(d,m\le n\), \(L\ge1\), \(n\ge C+1\) at sufficiently large
scale, and \(\log A\le n^3+21L\). For the last inequality,
\(r^2d\le n^3\), \(\log(g_0g_1)\le20L\) and
\(\log|\Delta_E|^{1/[E:\mathbb Q]}\le L\). The prescribed
depth at the chosen split auxiliary place, with
\(\mathrm N\mathfrak w_{v_0}\ge q_0\), gives
\(\log\mathrm N\mathfrak m>C_F d^2 n^6 L\). Hence
\[
 \log 2+Q+(6n+9)d+d\log(1/s)
 \le(2C+42)d^2n^5L
 <\tfrac12\log\mathrm N\mathfrak m .
\]
This gives the remaining bounds in \eqref{eq:principal-core-depth}.
The short-lattice estimate of Section~\ref{sec:principal-short-lattice} says that if
some nonzero \(y\) of the transported rational row lattice
\(\mathfrak a\Lambda_Y(x)\), for a nonzero integral ideal
\(\mathfrak a\), satisfies
\[
 \prod_{v\mid\infty}|x_v^{-1}y|_v
       <s^d\sqrt{\mathrm N\mathfrak a},
\]
then \(|f(x)|\le e^Q\|f\|_\infty A(2s)^{1/m}\). The
estimate uses reduction to a rational unipotent radical.
Cuspidality makes the average along that radical zero and
the small real representatives at the row level bound the
variation away from \(f(x)\).
\releasecite{Development/AbelianVarietyTheory/Proofs/TwistedOperators/NormCuspFormLeOfShortLine.lean}{67}{Small-lattice estimate}
\releasecite{Development/AbelianVarietyTheory/Proofs/TwistedOperators/GridDepthPaysEveryLoss.lean}{148}{Depth of the actual grid}.

Here is the nonvanishing argument from these inputs. Choose \(x\)
with \(|f(x)|>\|f\|_\infty/2\). Such an \(x\) exists since
the nonzero intertwiner \(t\) and its irreducible real source give
\(f\ne0\). The small-lattice estimate implies the reverse
weak inequality for every \(\mathfrak a,y\) at this \(x\).
By orthogonal addition of the norm-one line it follows that
\[
 \prod_{v\mid\infty}|\iota(x_v)^{-1}z|_v\ge R
 \quad
 (0\ne z\in \mathfrak m\Lambda_C(\iota x)).
\]
Indeed, if the \(Y_0\) component is nonzero, its product gives
the bound. A nonzero vector on the added line has coefficient in
\(\mathfrak m\), so the product is at least
\(\sqrt{\mathrm N\mathfrak m}\).

Use the seed \(k=k_\infty 1_{K_C}\) in the Poincar\'e series
\[
 P_x(g)=\sum_{\gamma\in U(C_0)(F)}
          k\bigl(\iota(x)^{-1}\gamma g\bigr).
\]
Here is the justification on the ambient smooth representation.
At a real place the translate in the first coefficient vector
changes the constant in \eqref{eq:real-branch-lowest-decay}.
The bound with exponent \(40n^2\) then holds when the second
vector is either in the lowest type or its chosen translate.
For a fixed second \(K_\infty\)-finite vector an exponent
\(20n^2\) suffices. Indeed the lowest type generates the
irreducible Harish--Chandra module under the derived action.
For a derivative of order \(j>0\) of a translated lowest-type
coefficient, use the values on a finite stencil in real
exponential coordinates, at scale \(h\le1\), exact on
polynomials of degree at most \(2j-1\). Such a stencil is
obtained by interpolation on a \(2j\)-point grid in each
variable. The absolute value of the derivative is bounded by
\(C_j(h^{-j}\|g\|_F^{-40n^2}+h^j)\).
The first term uses the power decay, uniform for each right
shift in the stencil. The remainder uses derivatives of total
order \(2j\), bounded independently of \(g\) by unitarity and
smoothness of the lowest-type vector. Taking
\(h=\|g\|_F^{-20n^2/j}\) proves the assertion for every fixed
derived word and hence for each fixed compact-finite vector.

For any \(g\), two distinct rational terms surviving the
finite indicator have their quotient, in the order
\(a b^{-1}\) for the real arguments \(a,b\) of their
coefficients, in the discrete real image of
\[
 U(C_0)(F)\cap\iota(x_f)K_C\iota(x_f)^{-1},
\]
conjugated by \(\iota(x_\infty)^{-1}\).
Choose \(\delta>0\) so that every nonidentity element
of this subgroup at infinity has at least one component
of Frobenius distance \(\delta\) or more from the identity.
In the box \(\|a_v\|_F,\|b_v\|_F\le R_v\), the quantity
\(\|a_v b_v^{-1}-1\|_F\) is at most
\(R_v\|a_v-b_v\|_F\). Thus the products of real balls of
radii \(\delta/(2R_v)\) centered at the distinct arguments
have disjoint interiors. Volume comparison bounds the number
of terms by
\(\prod_{v\mid\infty}(1+2R_v^2/\delta)^{2n^2}\).
Dyadic summation with exponent \(20n^2\) proves uniform
boundedness and convergence of the series for every fixed
compact-finite second vector and for the chosen diagonal
seed. Convergence is locally normal: on compact right
translates the matrix norms are comparable with constants
independent of the rational term. The same bound for every
fixed derived word allows differentiation term by term on
the compact-finite module.

For a compact-finite second vector, unfold the constant
term of this series along a proper rational parabolic.
The quotient of its unipotent radical by the rational points
is compact, and the absolute bound above permits unfolding.
Each resulting integral over the adelic unipotent radical
vanishes by the real discrete-series constant-term theorem
at a real place. The theorem applies to the lowest-type
vector and the compact-finite second vector. Translation of
the first vector moves the translate into the argument of
the matrix coefficient and leaves that theorem applicable.
The level is \(K_C\) and the right compact and derived
actions on the Poincar\'e series agree with those on the
second vector.

These bounded, cuspidal functions lie in the cuspidal Hilbert
space, since the adelic quotient has finite volume. The estimates
for all fixed derived words identify differentiation with right
differentiation in \(L^2\) as well, by dominated convergence.
The resulting Harish--Chandra-module map pulls the invariant
Hermitian form back to a scalar multiple of the given form, by
irreducibility. It therefore extends to a bounded map from the
ambient Hilbert representation. This map intertwines the real
group. Indeed the compact-finite vectors are analytic for each
real one-parameter subgroup. Their images have the same bounds
on all powers of each infinitesimal generator, up to the single
operator norm of the map. The intertwining follows first by
power series and then along the entire one-parameter subgroup
by analytic continuation. These subgroups generate the connected
real group.

At the prescribed finite level, local Sobolev embedding on
each real quotient chart gives the smooth representative of the
image of every smooth vector. It also bounds evaluation at any
fixed adelic point by a finite sum of Hilbert norms of derived
vectors, with the order depending only on the real dimension.
This proves the evaluation condition in the actual multiplicity
space. The same bound at \(xg_\infty\) follows by evaluating
\(\pi(g_\infty)u\) at \(x\). Commuting the fixed Sobolev
operators with \(\pi(g_\infty)\) via the adjoint action gives
one moderate-growth exponent for all right derivatives, by
applying the bound to each derived \(u\). The center acts by
the given infinitesimal character. Compact averaging approximates
each smooth vector in every Sobolev seminorm, hence its image
locally uniformly by the compact-finite images. The unipotent
quotients in the constant terms are compact, so cuspidality
persists pointwise. This gives a smooth cuspidal multiplicity
map \(y\) on the algebraic tensor of the prescribed local
models, of trivial character at \(K_C\).
\href{https://www.sciencedirect.com/science/article/pii/0022123675900348}{Harish-Chandra, Theorem~18.1.}
\releasecite{Development/AutomorphicArithmetic/ThetaCorrespondence/RealBands/ExistsMultTEqPoin.lean}{75}{Cuspidal Poincar\'e series at the prescribed level}.

In the pairing with \(f\), unfold using right cosets by
\(\iota U(Y_0)(F)\). For each coset represented by \(\gamma\)
the finite part of the integral over \(U(Y_0)(\mathbb A_F)\)
has support either empty or a left translate of
\(K'=\iota^{-1}(K_C)\) in \(U(Y_0)(\mathbb A_{F,f})\).
The orthogonal lattice sum gives the row congruence modulo
\(\mathfrak m\). For a unitary automorphism of a self-dual
lattice, the inverse is the adjoint. Taking the adjoint of
the inverse congruence therefore gives the congruence modulo
\(\overline{\mathfrak m}\) as well. Each prime of \(E\) over
\(g_0g_1\) divides \(\mathfrak m\) or its conjugate, and
\((g_0g_1)\) is squarefree in \(\mathcal O_E\). Thus
\(\mathfrak m\cap\overline{\mathfrak m}\subset(g_0g_1)\)
and \(K'\) lies in the bottom principal subgroup. For the
stabilizer coset, left translation
by \(x\), \(K'\)-invariance and
\eqref{eq:principal-core-reproducing} give the term
\(\operatorname{vol}(K')c_\infty f(x)\).
\releasecite{Development/AbelianVarietyTheory/Proofs/TwistedOperators/UnipotentPeriodAdelicClausesOfReal.lean}{173}{Real and finite factors of the stabilizer term}
\releasecite{Development/AbelianVarietyTheory/Proofs/TwistedOperators/ThetaRowGridBottomSelfDualLattices.lean}{16}{The bottom lattices and levels}
\releasecite{Development/AutomorphicArithmetic/ThetaCorrespondence/RealBands/ExistsMultTEqPoin.lean}{75}{Cuspidal Poincar\'e series and its multiplicity map}.

For every other supported coset, put \(z=\gamma z_0\).
The support condition implies \(z-z_0\in\mathfrak m
\Lambda_C(\iota x)\) by applying the congruence to \(z_0\).
The stabilizer property makes \(z-z_0\) nonzero and distinct
for distinct cosets. Since \(z\) has Hermitian norm one,
\(|\iota(x_v)^{-1}z|_v\ge1\). The triangle inequality gives
\[
 |\iota(x_v)^{-1}z|_v
 \ \ge\ (1+|\iota(x_v)^{-1}(z-z_0)|_v)/3 .
\]
By \eqref{eq:principal-core-transversal}, the absolute tail is
bounded by \(\|f\|_\infty e^Q|c_\infty|\operatorname{vol}(K')\)
times the sum of
\(\prod_{v\mid\infty}((1+|\iota(x_v)^{-1}y|_v)/3)^{-(2n+5)}\)
over nonzero \(y\) of \(\mathfrak m\Lambda_C(\iota x)\).

To bound this sum, group the vectors into boxes indexed by
\(j_v\ge0\) with \(2^{j_v}\le1+|\iota(x_v)^{-1}y|_v
<2^{j_v+1}\). Set \(D_{\mathbf j}=\prod_v2^{j_v+1}\).
In the norm on \(\prod_v\mathbb C^n\) obtained by taking the
maximum of the component norms scaled by \(2^{-j_v-1}\), the
box lies in the unit ball and two distinct points are at
distance at least \(D_{\mathbf j}^{-1/d}\). This follows from
the lower product bound for their difference. Packing disjoint
balls of half that radius bounds the count by
\((1+2D_{\mathbf j}^{1/d})^{2nd}\le
3^{2nd}D_{\mathbf j}^{2n}\).
The product of the quantities \(1+|\iota(x_v)^{-1}y|_v\)
is at least \(R\), so one power of that product supplies \(R^{-1}\).
Summing the remaining weights over the boxes gives
\[
 \frac{3^{(4n+5)d}\,2^{2nd}}{R}
     \sum_{\mathbf j\in\mathbb Z_{\ge0}^{d}}2^{-4\sum_v j_v}
 =\frac{(3^{4n+5}2^{2n}\,16/15)^d}{R}
 \le\frac{e^{(6n+9)d}}{R}.
\]
This bound also justifies absolute unfolding. The last two
inequalities of \eqref{eq:principal-core-depth} make the tail
at most \(\|f\|_\infty|c_\infty|\operatorname{vol}(K')/2\),
strictly smaller than the absolute value of the stabilizer term.
The resulting period is nonzero.

Let \(v_*\) be the tensor of the translated lowest-type vectors
in the diagonal seed. Use compact approximate identities of
finite Peter--Weyl support for the product of the adapted
maximal compact groups. Their averages of \(v_*\) converge in
the smooth topology. Averaging the corresponding right
translates of \(P_x\) commutes with the locally normal sum
and gives the series of these compact-finite second vectors.
By continuity of evaluation and bounded continuity of \(P_x\),
the resulting pointwise identity is \(P_x=y(v_*)\).

For an arbitrary nonzero \(w\) in the smooth row tensor, the
period against \(P_x\) defines a continuous functional on the
row Hilbert tensor. Indeed \(t\) is bounded into \(L^2\), and
\(P_x\circ\iota\) is bounded on the finite-volume row quotient.
The tensor of the local irreducible unitary representations
is irreducible for the product real group, so the span of the
real translates of \(w\) is Hilbert-dense. The functional is
nonzero on \(w_*\), hence nonzero on some real translate of
\(w\). Right invariance of the quotient measure moves the
inverse translate to \(P_x\). The translated ambient form
still belongs to \(y\) at level \(K_C\) and is bounded and
continuous. Average it on the right by finite Peter--Weyl
approximate identities in the maximal compact subgroup of
the prescribed ambient model. One may take products of
nonnegative compact kernels with integral one. The resulting
ambient vectors are compact-finite tensors, and the pairings
converge by bounded convergence against the integrable
function \(t(w)\). Some compact-finite vector therefore gives
\eqref{eq:principal-core-detection}.
\releasecite{Development/AbelianVarietyTheory/Proofs/TwistedOperators/PeriodNonvanishingOfArchimedeanInputs.lean}{65}{Passage to a prescribed initial real vector}.

\teamanchor{proof.principal-short-lattice-reduction}%
\subsubsection{Reduction at a point with a short lattice}\label{sec:principal-short-lattice}
Write \(d=[F:\mathbb Q]\), \(g=g_0g_1\) and
\(\operatorname{rd}(E)=|\Delta_E|^{1/(2d)}\). At an adelic
point \(x\), put \(L=Y_0(E)\cap x_f\widehat{\Lambda_Y}\)
in the finite adeles of the row space.
At the real places give \(L\) the Euclidean norms from
\(x_v^{-1}\) and the adapted frames. In these norms the
original Hermitian form is represented by a self-adjoint
involution \(J_v\). The finite lattice \(L\) is self-dual
for the original \(E/F\)-Hermitian form.

For an \(\mathcal O_E\)-lattice \(D\) of rank \(k>0\) with
induced Hermitian norms, use the size
\[
 \operatorname{sz}(D)
 =\left(2^{dk}|\Delta_E|^{-k/2}\operatorname{covol}(D)\right)^{1/(2dk)} .
\]
The covolume uses one complex embedding over each real place
of \(F\), and quotients use the orthogonal quotient norms.
Write \(\ell(D)\) for the minimum size of a nonzero sublattice
and \(h(D)\) for the maximum size of a nonzero torsion-free
quotient. Put \(\ell=\ell(L)\).
For any nonzero integral ideal \(\mathfrak b\) and
\(0\ne y\in\mathfrak b L\), write \(L\cap Ey=\mathfrak a y\).
Saturation and projectivity over \(\mathcal O_E\) give
\(\mathfrak b L\cap Ey=\mathfrak b\mathfrak a y\).
Thus \(1\in\mathfrak b\mathfrak a\) and
\(\mathrm N\mathfrak a\le1/\mathrm N\mathfrak b\). The size of
this line gives
\[
 \prod_{v\mid\infty}\|y\|_v
 =\operatorname{sz}(\mathfrak a y)^d(\mathrm N\mathfrak a)^{-1/2}
 \ \ge\ \ell^d\sqrt{\mathrm N\mathfrak b}.
\]
In particular, failure of the product lower bound with
\(s^d\sqrt{\mathrm N\mathfrak b}\) forces \(\ell<s\).
\releasecite{Formalism/HarderNarasimhan/StatementPl.lean}{29}{Size and quotient norms}
\releasecite{Development/AbelianVarietyTheory/Proofs/TwistedOperators/LeastSizePowMulSqrtAbsNormLeProdNorms.lean}{51}{The line calculation}.

\paragraph{Covering integral homomorphisms.}
Let \(D_1,D_2\) be Hermitian \(\mathcal O_E\)-lattices
of ranks \(a,b>0\), with \(\ell(D_1)\ge l>0\) and
\(h(D_2)\le u\). Give the real space of complex-linear
placewise homomorphisms from \(D_1\) to \(D_2\) its sum of
Hilbert--Schmidt squares. The lattice
\(\mathcal L=\operatorname{Hom}_{\mathcal O_E}(D_1,D_2)\)
has covering radius at most
\[
 \frac{ab\sqrt d}{2}\,\operatorname{rd}(E)\frac{u}{l}.
\]
Indeed, the trace pairing of \(\mathcal L\) with the reverse
homomorphisms is \(\operatorname{Tr}_{E/\mathbb Q}\operatorname{tr}(BT)\).
Its trace dual is
\(\mathfrak D_E^{-1}\operatorname{Hom}_{\mathcal O_E}(D_2,D_1)\),
as one sees in integral bases at the finite places. It follows
that the Euclidean dual elements are \((2B_v^*)_v\) for \(B\)
in this trace dual, because the field trace is twice the sum
of the real parts at the chosen complex embeddings.

For such a nonzero \(B\), of rank \(j\), its image in
\(\mathfrak D_E^{-1}D_1\) has size at least
\(l/\operatorname{rd}(E)\). Indeed the different times the
image lies in \(D_1\), and multiplication by the different
changes its size by \(\operatorname{rd}(E)\).
The size of \(D_2/\ker B\) is at most \(u\). Taking
determinants of \(B_v:D_{2,v}/\ker B_v\to\operatorname{im}B_v\)
and comparing the two covolumes gives
\[
 \Big(\prod_v|\det B_v|\Big)^{1/(dj)}
 \ \ge\ \frac{l}{u\operatorname{rd}(E)}.
\]
The \(dj\) nonzero singular values then give
\(\big(\sum_v\|B_v\|_{\mathrm{HS}}^2\big)^{1/2}
\ge\sqrt d\,l/(u\operatorname{rd}(E))\) by the
arithmetic--geometric mean inequality. Banaszczyk's
transference bound \(\mu(\mathcal L)\lambda_1(\mathcal L^\vee)
\le\dim_{\mathbb R}(\mathcal L_{\mathbb R})/2=dab\)
proves the stated bound. For either rank zero the
homomorphism space has radius zero.
See Banaszczyk~\cite[Theorem~(2.2)]{depth-banaszczyk}.
\releasecite{Development/AbelianVarietyTheory/Proofs/FlagMeasures/LatticePairVolumeBounds.lean}{27}{Integral homomorphism lattice}
\releasecite{Development/AbelianVarietyTheory/Proofs/FlagMeasures/CoveringRadiusHomSubquotientSublattice.lean}{46}{Application to the induced quotient norms}.

\paragraph{The short isotropic step.}
Suppose \(\ell<1/2\). The Arakelov Harder--Narasimhan
filtration of \(L\) has saturated steps and strictly
increasing successive sizes. The normalized degree
divided by rank equals \(-\log\operatorname{sz}\).
At a step \(R_i\) the largest quotient size is its last
successive size, and the least size of \(L/R_i\) is the
next successive size. These assertions follow by comparing
the degree polygon with the chords ending or starting at
the vertex for \(R_i\). They also give that the first
successive size is \(\ell\).
See Bost--Chen~\cite[\S\S1.2--1.3]{lt4bostchen}.

Every step of last size smaller than \(1/2\) spans an
isotropic space for the original Hermitian form. Indeed this
form induces an integral map \(R_i\to\overline{R_i}^{\,*}\)
of operator norm at most \(1\) at each real place, since
\(J_v\) is isometric. The integral metric dual has least
size \(1/h(R_i)\) by degree additivity and duality. If the
map were nonzero, its determinant on the quotient by its
kernel would give a geometric mean singular value at least
\(1/h(R_i)^2>1\). This contradicts the operator norm bound.
Consequently the number \(p\) of successive sizes below
\(1/2\) is at most the signature bound \(q\), and \(p\ge1\).
The bottom row has negative index \(m/2\). Use \(q=m\) as
an upper bound for that index.

Write \(s_1=\ell<s_2<\cdots<s_p<s_{p+1}=1/2\) for these
sizes with the cutoff appended, and put \(\rho=(2\ell)^{1/q}\).
The product of the \(p\) ratios \(s_i/s_{i+1}\) is
\(2\ell\). Some ratio is therefore at most \(\rho\).
For this step \(R=R_i\) and \(W=ER\), use
\(h=s_i\) and \(l=s_{i+1}\). Then
\[
 h(R)\le h,\qquad \ell(L/R)\ge l,\qquad
 h/l\le\rho,\qquad h^2\le\rho/4 .
\]
The last bound follows from \(h\le\rho/2\) and \(0<\rho<1\).
\releasecite{Development/AbelianVarietyTheory/Proofs/TwistedOperators/AutomorphicPointExistsIsotropicReductionOfSmallS.lean}{60}{The saturated gap used in the release}.

\paragraph{A fundamental set for the unipotent radical.}
Put \(P=L\cap W^\perp\) and \(k=\dim_E W\).
The integral form identifies \(L/P\) with
\(\overline R^{\,*}\). Indeed, self-duality gives all integral
covectors on \(L\), and \(R\subset L\) is a saturated
summand. At a real place the quotient norms agree because
\(V_v=W_v\oplus V_{0,v}\oplus J_vW_v\) is Euclidean
orthogonal, with \(V_{0,v}\) stable under \(J_v\).
In these coordinates the radical of the parabolic
stabilizing \(W\) has Lie algebra of matrices of the form
\[
 \begin{pmatrix}
  0&A&Z\\
  0&0&-J_{0,v}A^*\\
  0&0&0
 \end{pmatrix},
 \qquad Z^*=-Z ,
\]
where \(J_v\) identifies \(J_vW_v\) with \(W_v\).
The first layer \(A\) maps \(V_{0,v}\) into \(W_v\).
Its displayed lift has Hilbert--Schmidt norm
\(\sqrt2\|A\|_{\mathrm{HS}}\). The center consists of
the \(Z\) blocks. All triple matrix products vanish.

The covering bound on
\(\operatorname{Hom}_{\mathcal O_E}(P/R,R)\) gives an
integral approximation to the first layer, since
\(P/R\subset L/R\) is isometric onto its image and
\(\ell(P/R)\ge l\). Any such integral first-layer map
extends to a rational element preserving \(L\) in
the Lie algebra. To see this, extend it to an
endomorphism \(D\) of \(L\), with image in \(R\)
and killing a complement of \(P\). The adjoint
\(D^\dagger\) for the original Hermitian form preserves
\(L\) by self-duality, kills \(P\), and has image in
\(W^\perp\) because \(D\) kills \(R\).
Then \(D-D^\dagger\) is the required lift.
The radius for the first-layer error modulo the center,
in the total Hilbert--Schmidt norm, is at most
\[
 \sqrt2\frac{(r-2k)k\sqrt d}{2}
       \operatorname{rd}(E)\frac h l
 \ \le\ \frac{r^2\sqrt d}{8\sqrt2}\operatorname{rd}(E)\rho.
\]
For the center, use the homomorphisms from \(L/P\) to
\(R\), for which the source has least size
\(1/h(R)\ge1/h\). Approximate half the desired center
and take \(D-D^\dagger\). The adjoint is isometric in
the chosen norms. The resulting radius is at most
\[
 k^2\sqrt d\,\operatorname{rd}(E)h^2
 \ \le\ \frac{r^2\sqrt d}{16}\operatorname{rd}(E)\rho .
\]
\releasecite{Development/AbelianVarietyTheory/Proofs/FlagMeasures/NilradicalFirstLayerCoveringRadiusLe.lean}{107}{First-layer approximation}
\releasecite{Development/AbelianVarietyTheory/Proofs/FlagMeasures/NilradicalCenterIntegralCoveringRadius.lean}{201}{Center approximation}.

Let \(Y\) be a real logarithm in the radical. Apply the
first radius bound to \(Y/(2g)\). Write
\(Y=2gY_0+A+2gZ_1\) with \(Y_0\) rational and integral
on \(L\), \(Z_1\) central and \(\|A\|\) at most \(2g\)
times the first radius. Apply the second radius bound
to \((2gZ_1+[A,2gY_0]/2)/g\). This gives a rational
center \(Z_0\) integral on \(L\) and a real central error
\(Z'\) with \(\|Z'\|\) at most \(g\) times the second
radius. The two-step Baker--Campbell--Hausdorff identity
gives
\[
 \exp Y=\exp(2gY_0+gZ_0)\exp A\exp Z',
 \qquad
 \|A\|+\|Z'\|\le \frac{r^2\sqrt d}{4}\,
                         g\operatorname{rd}(E)\rho .
\]
The first exponential is congruent to the identity
modulo \(g\) on \(L\). Indeed its finite sum is
\(1+2gY_0+gZ_0+2g^2Y_0^2\), and the same argument
applies to its inverse. This gives the principal level
on \(L\) even at primes over \(2\).
\releasecite{Development/AbelianVarietyTheory/Proofs/FlagMeasures/SmallFundamentalSetRealUnipotent.lean}{74}{Integral rational exponentials}
\releasecite{Development/AbelianVarietyTheory/Proofs/FlagMeasures/SmallFundamentalSetRealUnipotent.lean}{135}{Two-step fundamental set}.

For \(n\in N_W(\mathbb A_F)\), first move its finite
component into \(x_f K_Y((g))x_f^{-1}\) by left
multiplication by a rational element of \(N_W\). The
polynomial logarithm identifies its finite-adelic points
with those of its Lie algebra, where the approximation
follows from the density of \(F\) in its finite adeles. Use the
preceding real factorization on the remaining real
component. Its rational factor has finite components in
the same conjugate principal subgroup. Expressed in the
frames transported by \(x\), the outcome is
\[
 n x=\gamma x\Big(\big(\exp A_v\exp Z'_v\big)_v,k_f\Big),
 \qquad k_f\in K_Y((g)),
\]
and the sum of the two real norms over the places is at most
\(r^2 d\,g\operatorname{rd}(E)\rho/4\) by Cauchy--Schwarz.
For any cuspidal row \(f\) invariant under this finite
level and with real Lipschitz constant \(L_f\) in the
stated norms, average \(f(x)-f(nx)\) over
\(N_W(F)\backslash N_W(\mathbb A_F)\) with probability
Haar. Its second term averages to zero by cuspidality.
If \(\ell<s\le1/2\), the variation estimate gives
\[
  |f(x)|\le L_f r^2 d\,g\operatorname{rd}(E)(2s)^{1/q}.
\]
Combined with the line calculation, this is the
short-lattice implication at the prescribed bottom
level used in the Poincar\'e argument.
\releasecite{Development/AbelianVarietyTheory/Proofs/TwistedOperators/AutomorphicPointExistsIsotropicReductionOfSmallS.lean}{141}{The adelic point factorization}
\releasecite{Development/AbelianVarietyTheory/Proofs/TwistedOperators/NormCuspFormLeOfShortLine.lean}{67}{Integration of the variation}.

\teamanchor{proof.principal-core-span}
\subsubsection{The span of descending core chains}
An admissible descending core chain is a sequence of rapidly
decreasing forms \(a_i\) on \(U(C_i)\) with
\(a_i=\Theta^{\rm core}_{i+1\to i}(\Xi_i)a_{i+1}\).
Each \(\Xi_i\) lies in the lattice-model depth space of its
arrow and the last \(a_m\) is \(c_\xi\) for a permitted
column. We show that the \(K_\infty\)-finite core forms
of \eqref{eq:principal-core-detection} lie in the algebraic
span of the \(a_0\) of these chains.

\paragraph{A common finite datum on each core arrow.}
Write \(G_i=U(C_i)\), \(n=r+1\), and \(h=r+5\), so that
\(\dim C_i=n+ih\). Let \(J_i\) be the given finite core level and
let \(S\) be the set of split places at its depth. Denote by
\(\chi_i\) the inverse of the grid splitting character with index
\(m-i\). Both splitting characters on the arrow oriented for the
upward pairing are \(\chi_i\). On a chosen split half the integral
variable for the descending arrow lies in the completed module
\[
 X_v=\operatorname{Hom}_{\mathcal O_v}
       (\Lambda_{i+1,v},\Lambda_{i,v}).
\]
The other half is its dual for the \(\psi\)-pairing. The two lattices
are dual because the vertex lattices are self-dual and the trace
coefficient at the chosen split place is a unit in these coordinates.
Multiplying both lattice frames by their specified polarization ideals
does not change the vertex stabilizers. When exchanging the legs of
the pair use \(x\mapsto x^*\), the actual Hermitian adjoint, and change
the sign on the opposite leg. Adjoint reverses the two projections
of the split algebra. It interchanges source and target norms, and
the two sign changes in the symplectic form cancel. In particular it
transports \(\psi\), the self-dual measure, and evaluation at zero.
These observations identify the space of downward tests with the
Schwartz space in the upward pairing using these frames.

Put \(I_v=\varpi_v^{d_v}\mathcal O_v\) at \(v\in S\). Choose
\(x_0:\Lambda_{i+1}\to\Lambda_i\) surjective modulo
\(\mathfrak m\) on the chosen halves. To do so, choose any projection
in bases of the two free completed modules at those places and lift
their residues through
\(\operatorname{Hom}(\Lambda_{i+1},\Lambda_i)/
  \mathfrak m\operatorname{Hom}(\Lambda_{i+1},\Lambda_i)\).
The homomorphism module is projective and tensoring with the finite
quotient commutes with \(\operatorname{Hom}\). The Chinese remainder
theorem thus gives a single integral \(x_0\). In the split
Schr\"odinger coordinate use
\(f_v={\bf1}_{x_{0,v}+I_vX_v}\). Use the unit lattice vacuum
at all other finite places and in the adapted real frames use the
Gaussian. The finite split choice is the projection onto the selected
coset by the characters of \(I_v^{-1}X_v^\vee/X_v^\vee\):
with the average normalized by the order of this finite group,
\[
 {\bf1}_{X_v}(x)
  \operatorname{avg}_{\lambda\in I_v^{-1}X_v^\vee/X_v^\vee}
          \psi_v\bigl(\langle x-x_{0,v},\lambda\rangle\bigr)
       ={\bf1}_{x_{0,v}+I_vX_v}(x).
\]
Thus it has trivial translation character on \(I_vX_v\) and
trivial multiplication character on \(X_v^\vee\). On the larger
dual lattice it has the eigencharacter
\(\lambda\mapsto\psi_v(\langle x_{0,v},\lambda\rangle)\).
Exchanging the legs gives the downward depth condition of
\(\mathrm{Dep}\) with its selected half. Real polynomial creation
operators commute with this finite operation. This constructs the
same finite datum for every real polynomial.
\releasecite{Development/AbelianVarietyTheory/Proofs/TwistedOperators/CoreChainAdjointLatticeExchange.lean}{337}{The completed lattices in the two orientations}
\releasecite{Development/AbelianVarietyTheory/Proofs/TwistedOperators/ExistsCosetFockVacuumDatum.lean}{83}{Finite Fourier projection and the Fock vacuum}.

The normalized split action on these matrices is pullback by the
two-sided action \(x\mapsto k_i x k_{i+1}^{-1}\).
The scalar factors in the Schr\"odinger representation are a positive
square root of the change in volume and, in the downward orientation
\((S,T)=(C_{i+1}^{-},C_i)\), the character
\(\chi_T(\det k_{i+1})^{-1}\chi_S(\det k_i)\).
The absolute determinant is one on the level. Both determinants are
\(1\bmod I_v\) in integral frames. The grid splitting characters
are trivial on \(1+I_v\). In particular, all depth-fixed vectors
are fixed under the two levels, since every
\((I_v X_v)\)-periodic function supported on \(X_v\) is fixed by
these pullbacks. At an unselected split place the same computation
uses \(I_v=\mathcal O_v\) and unramified characters, and the vacuum is
the line of \({\bf1}_{X_v}\). At a nonsplit place the lattice
stabilizer acts in the normalized lattice model by pullback, fixing
the basic lattice vacuum. Here the extension is unramified and the
residue characteristic is odd, since the host is split at the places
above \(2\). The standard splitting with the unramified parity
characters agrees with the lattice splitting on the two self-dual
lattice stabilizers. Changing those parity characters by the
unramified splitting quotients multiplies the action by their
values on unit lifts of the determinants, which are one. Such a
lift exists by Hilbert~90, since a lift can be multiplied by an
element of \(F_v^\times\) to make it a unit in an unramified
extension. At a dyadic split place use translations and modulations
on the individual halves, as in the displayed split formula.
This explains why the full \(2\Lambda\) translation condition in the
adelic lattice model gives the usual vacuum condition off the split
halves and does not enlarge the split invariant space. The same
argument applies to the first-row levels with
\(\chi_S,\chi_T\) the corresponding row splitting characters, whose
possible quotient by the unramified \(\eta\) preserves all these
trivialities.
\releasecite{Development/AbelianVarietyTheory/Proofs/TwistedOperators/DependenceFixed.lean}{402}{Actual character and principal action}
\releasecite{Development/AbelianVarietyTheory/Proofs/TwistedOperators/DependenceFixed.lean}{500}{Application to both families}
\releasecite{Development/AbelianVarietyTheory/Proofs/DivisorProofs/RelIndexNeZeroOfFinite.lean}{153}{Dyadic split condition}.

Let \(b(g)\) be the oscillator coefficient of this tensor against
itself, with any one of its real polynomial factors. At a selected
place \(v\), allowing arbitrary \(g\in G_i(F_v)\), its coefficient
vanishes unless the two cosets
\(x_{0,v}+I_vX_v\) and \(g(x_{0,v}+I_vX_v)\) meet. The
determinant multiplier never changes this support. Suppose they meet.
Then \(a=g b_0\) with \(a,b_0\) both in the chosen coset. Choose a
right inverse \(y\) to \(x_{0,v}\) modulo \(I_v\). The two integral
matrices \(a y, b_0 y\) are \(1\bmod I_v\), so they have integral
inverses. Hence \(g=(a y)(b_0 y)^{-1}\) belongs to \(J_{i,v}\).
Conversely on \(J_{i,v}\) the coefficient is
\(\|f_v\|^2>0\). The other split half of \(g\) is forced by
unitarity and the dual lattice, so this is the exact specified
one-sided congruence subgroup, viewed inside the unitary group.
Off \(S\) the finite datum is spherical. This proves support
on \(G_i(F_\infty)J_i\) inside the semilocal domain
where the components off \(S\) lie in the spherical compact.
\releasecite{Development/AbelianVarietyTheory/Proofs/TwistedOperators/CosetSupportedCoefficientCongruenceSubgroup.lean}{46}{Support read via the dual half eigencharacter}
\releasecite{Development/AbelianVarietyTheory/Proofs/TwistedOperators/CoreChainPreStepsExistOfDepthCongruence.lean}{35}{Use on the oriented core arrow}.

\paragraph{The real polynomial space and the placed zeta integral.}
In the determinant-cover coordinates at a real place, subtract the
common splitting center \(a_i=-c_{m-i}\). In the orientation of the
upward adjoint, write the centered parameter of the smaller
representation as \((A,B;C,D)\), with \(A,C>0\) and \(B,D<0\).
The partner has centered parameter
\[
 (A,\ ((h-1)/2,\,(h-3)/2,\ldots,-(h-1)/2),D;\ C,B).
\]
An overall reversal of positive and negative forms uses the
contragredient convention for both members. Every old entry has
absolute value at least \((h+3)/2\). In the grid this follows by
subtracting the floor of the midpoint of the new block. The distinct
centers are spaced by \({\rm gap}\), while the old entries lie in the
base lists and the other blocks of length \(h\). The formula
for the signed partner exchanges the signs of the old entries on
one side of this midpoint and inserts the new block in the positive
list. To check the choice of list, the opposite insertion would put
an entry of the block on both signs when compared with the grid
partner. Distinct positive and negative entries of that partner
are separated by more than \(2\,{\rm bigN}^3\).
The determinant exponents are integers on the original even-rank
groups. After the shift all real parameters for this pair are
half integral, as required on the equal-parity determinant covers.
\releasecite{Development/AbelianVarietyTheory/Proofs/TwistedOperators/CoreChainRealParametersCorrespond.lean}{76}{The actual centered signed parameters}
\releasecite{Development/AbelianVarietyTheory/Proofs/TwistedOperators/TestGridCommonDataPropertiesHold.lean}{33}{Splitting centers and parity}.

The displayed upward pair is the one in the real theta correspondence \cite[Theorem~2.7(b)]{p2paul-real-theta}. To check
first occurrence in its tower use the real first-occurrence criterion \cite[Theorem~6.1]{p2ichino-real-theta}. The centered half-integral smaller parameter has
no central string, so in that notation \(k_\pi=0\) and
\((r_\pi,s_\pi)=(|A|+|D|,|B|+|C|)\). The tested signatures on
that tower are \((r_\pi+h+l,s_\pi+l)\). At \(l=0\)
the condition on \(C^\pm_\pi(h/2)\) is satisfied since these sets
contain only entries in absolute value less than \(h/2\).
The condition for \(t=h/2>0\) excludes \(l<0\).
The reverse direction uses the complementary removal criterion \cite[Proposition~3.4(ii)]{p2paul-real-theta}.
The inserted string in the larger parameter is maximal, since
the neighboring half-integral entries of the string are absent.
Deleting it gives the smaller signature and the prescribed
parameter at first occurrence. A global smaller member in either
direction has a smaller real member at each real place. The local
vanishing there followed by the constant-term formula for theta
lifts gives cuspidality of the two first-occurrence lifts. This
uses the tower property in its constant-term form, including the
vanishing of the nonzero rational orbits of smaller dimension by
cuspidality of the input.

The local pairing for the actual adapted Fock model has a nonzero real integral \cite[Proposition~7.1 and Remark~7.2]{p2ichino-real-theta}
\[
 \int_{G_i(F_v)}
    \langle\omega_v(b,1) P_1 f_{{\rm vac},v},P_2 f_{{\rm vac},v}\rangle
    \overline{\langle\pi_{i,v}(b)u_1,u_2\rangle}\,db
\]
on polynomial and compact-finite data with this splitting and this
signature. Here the inner products are linear in their first
variable and \(\pi_{i,v}\) is the member used in the adjoint
convention. The oriented model agrees by the adjoint exchange
and, if using the opposite sign of the form, the twist of
Ichino's Remark~7.3. The polynomial space is unchanged by the
unitary change from an adapted orthonormal frame to the oscillator
coordinates. Polarizing successively in the oscillator slot and
in the \(\pi_{i,v}\) slot shows that the value is nonzero for
some \(P_1=P_2=P_v\), \(u_1=u_2=u_v\). The conclusion uses
absolute convergence at the critical point. For the section
at \(s\) near \(h/2\), the real doubling height weight is
\(H_v(b)^{s-h/2}\), with \(0<H_v\le1\). The real doubled
compact acts on a finite-dimensional span of the polynomial
tensor. Iwasawa decomposition applied to its section at \(h/2\)
therefore bounds its coefficient by \(C H_v^{(n+(i+1)h)/2}\).
On \({\rm Re}\,s\ge0\) the integrand is bounded by a constant
times the absolute discrete-series coefficient. It is integrable
by the integrability criterion for discrete-series coefficients \cite[Theorem~III.1]{p2milicic-integrability}, since the noncompact
root gaps exceed \(2\,{\rm bigN}^3\) and
\(\tfrac14\sum_\beta|(\alpha,\beta)|=n+ih-1\)
for a root \(\alpha\) of type \(A_{n+ih-1}\), with roots of
length squared two. This bound also yields a holomorphic integral
on \({\rm Re}\,s>0\). The product over real places uses the tensor
of the \(P_v\) and the \(u_v\), and its Haar normalization differs
from their product by a positive scalar.
\releasecite{Development/AbelianVarietyTheory/Proofs/TwistedOperators/EvenThetaStepRealZetaNonvanishing.lean}{40}{Real product and the holomorphic strip}
\releasecite{Development/AbelianVarietyTheory/Proofs/DecoratedLatticeProofs/RealMatrixCoeffFockVectorsLeDoublingHeightPow.lean}{59}{Height exponent of the actual doubled section}.

\paragraph{Semilocal norm on a whole Hecke label space.}
Let \(\mathcal A_i\) be the finite-dimensional space of cuspidal
occurrences of the specified smooth real type with right \(J_i\)-fixed
values and trivial finite character. Choose the multiplicity inner
product by Schur's lemma from the Petersson inner product on values
and the invariant product on the common real type. Thus
for \(a\in\mathcal A_i\) and \(u=\bigotimes_{v\mid\infty}u_v\)
from the real integral, the lower coefficient of \(a(u)\) at \(bk\)
with \(b\in G_i(F_\infty)\), \(k\in J_i\), is
\(\|a\|^2\prod_v\langle\pi_{i,v}(b_v)u_v,u_v\rangle\).
The doubled height at \(k\) is one, as
\(\operatorname{diag}(k,1)\) preserves the doubled lattice.
The semilocal integral is consequently
\[
 \operatorname{vol}(J_{i,S})
      \Bigl(\prod_{v\in S}\|f_v\|^2\Bigr)
      \|a\|^2\,Z_\infty(h/2).
\]
Here the spherical Haar factors off \(S\) have volume one,
and any conversion from product to chosen Haar measure gives
a positive constant. All factors on this line are nonzero for
every nonzero \(a\). The same finite datum and real witness can
be used throughout the multiplicity space.

For the unramified coefficient use spherical unitary base-change
labels over the places \(w\) of \(E\) outside \(S\). The label \(B\)
of a bottom constituent is a multiset of size \(n\) at each \(w\).
The elementary unitary Satake bound places its entries in the
annulus \(q_w^{-(n-1)/2}\le|\alpha|\le q_w^{(n-1)/2}\).
With the standard normalization this is the convex-hull bound by
the exponents of the trivial spherical representation. It follows
by bounding every double-coset operator on a spherical unit vector
by its Haar volume and applying the Satake transform to rays of
dominant coweights. Spherical theta transfer adds at arrow \(j\)
the string \(q_w^{(h-1)/2},\ldots,q_w^{-(h-1)/2}\) times the
common entry \(\chi_j(w)\). The original label is unchanged here
because the two characters are the same on each oriented core
pair. For split \(v\) this is normalized parabolic induction
with a character on the excess general linear factor. For inert
\(v\) the same formula is read on base-change parameters from
spherical Hecke matching \cite[Proposition~A.5]{tr3liu}. Every inert
place of the grid has odd residue characteristic.

At arrow \(i\) put \(s_*=h/2\) and omit the places over \(S\)
from all partial products. Its unramified numerator at \(s_*\)
is
\[
 L^S((h+1)/2,B\otimes\chi_i^{-1})
 \prod_{0\le j<i}\prod_{\ell=0}^{h-1}
         L_E^S(h-\ell,\chi_j\chi_i^{-1}).
\]
The first Euler product is absolutely convergent and nonzero since
\((h+1)/2-(n-1)/2>1\). Every quotient character is unitary and
nontrivial, since its real exponents differ. Tate's continuation
and the nonvanishing at one for a nontrivial unitary Hecke character
give the factors at \(h-\ell=1\). The remaining Hecke products
are nonzero by absolute convergence. Removing the finitely many
Euler factors multiplies by nonzero numbers, since the characters
are unitary. The denominator is
\(\prod_{j=1}^{n+ih}L_F^S(h+j,\eta_{E/F}^{\,h+j})\)
and is likewise finite and nonzero.

The Witt index of \(C_{i+1}\) is at most its real negative index,
which is at most \(m/2<h\). The doubling Siegel--Weil integral
is in the convergent range since
\(\dim C_{i+1}-{\rm Witt}(C_{i+1})>\dim C_i\).
The unfolding at \(s_*\) gives the squared norm of the upward
lift as the semilocal integral times this quotient of partial
products, times a positive measure factor. One can use the
convergent-range formula \cite[\S8.2]{p2ichino-real-theta}, with the
semilocal coefficient of the actual occurrence. To see why this
works for a sum of occurrences with one label, average the right
translations off \(S\) between spherical vectors. On the whole
label space each averaged operator acts as its scalar spherical
eigenvalue. For a finite set of these places the coefficient
therefore factors as the product of normalized spherical
coefficients times the semilocal coefficient just computed.
Unfold first in the half-plane of absolute convergence and then
continue the common Euler quotient using the displayed products.
The compact support at \(S\) and the real weighted bounds give
the continuation of the semilocal integral through \(s_*\).
This argument treats cross coefficients between occurrences
without needing to choose a global simple tensor for \(a(u)\).
Thus every nonzero \(a\) in a label descended from the bottom has
nonzero upward lift with the placed tensor.
\releasecite{Development/AbelianVarietyTheory/Proofs/TwistedOperators/EvenThetaStepSemilocalZetaNeZeroOfReal.lean}{63}{Semilocal reduction in the core measure}
\releasecite{Development/AbelianVarietyTheory/Proofs/TwistedOperators/EvenThetaChainLValuesNeZeroAndParamsInjective.lean}{95}{The common Euler quotient}
\releasecite{Development/AbelianVarietyTheory/Proofs/TwistedOperators/UpwardThetaLiftNeZeroOfLValueNeZero.lean}{38}{Application to the upward adjoint}.

\paragraph{Algebraic span of the downward chains.}
Fix the finite datum on each arrow. The one-dimensional real
Howe intertwining line in the two prescribed real types
factorizes the downward kernel operator as \(D_i\otimes T_i(P)\).
Here \(D_i:\mathcal A_{i+1}\to\mathcal A_i\) is one linear
map on multiplicities and \(T_i(P)\) is contraction of the
nonzero real oscillator pairing against \(P\). This factorization
can first be taken on Harish--Chandra modules. The global lift is
in the target real isotypic component by the real correspondence
and the first-occurrence constant terms. The fixed finite tensor
preserves the two levels by the preceding calculation.
The real contraction lands in the compact-finite module for
compact-finite \(P\) and upper vectors by compact equivariance.
Its image span as both arguments vary is nonzero. That span is
stable under the lower Lie algebra, since differentiating its
action differentiates \(P\) by a quadratic creation and annihilation
operator and leaves a polynomial. It is stable under the lower
maximal compact for the same equivariance. Irreducibility
of the lower Harish--Chandra module makes this algebraic span
the entire lower compact-finite space. Equivalently, apply this
argument at each real place and take their algebraic tensor product.

Fubini for the double theta integral against two cusp forms
makes the upwards operator \(D_i^*\otimes T_i(P)^*\) in the
chosen Petersson measures. The absolute convergence for this
swap follows by multiplying the polynomial height bound for the
theta kernel by the rapidly decreasing cusp forms on both
quotients. Positivity of the squared norm just proved implies that
\(D_i^*a\ne0\) for each nonzero \(a\) of a label transferred
from the bottom, and the spherical relations put it in its
transferred label space. Distinct labels stay distinct, since one
can cancel the added common multiset at each unramified place.
Decompose \(\mathcal A_0\) into its mutually orthogonal simultaneous
Hecke eigenspaces. Iteration shows that the adjoint of
\(D_0\cdots D_{m-1}\) is injective on their direct sum.
Thus \(D_0\cdots D_{m-1}\) is onto \(\mathcal A_0\).

For \(a\in\mathcal A_0\) choose one preimage \(a_m\) under
this composite. Expand a given lower compact-finite vector as a
finite sum of products \(T_0(P_0)\cdots T_{m-1}(P_{m-1})u_m\),
using the real span at each stage in turn. Each term evaluated
at \(a\) is an actual sequence of descending lifts of
\(a_m(u_m)\), with the Schwartz test obtained by the adjoint
exchange at each arrow. The intermediate terms are cusp forms
with the prescribed levels. At the definite core \(C_m\), the
real irreducible type is the dual of the fixed
finite-dimensional column coefficient, and the prescribed column
space is the corresponding full space of equivariant coefficient
functions at \(J_m\). For clarity, if the coefficient is \(E_m\)
and \(a_m:E_m^\vee\to\mathcal A_{\rm cusp}(G_m)^{J_m}\), then
\(c(g)=\sum_\alpha a_m(e_\alpha^\vee)(g)e_\alpha\) is the
column with \(c_\xi(g)=a_m(\xi)(g)\).
The expansion therefore uses only permitted terminal columns and
the actual smaller Haar quotients on the way down. This proves the
core span for the bottom space tested by the corank-one period.
\releasecite{Development/AbelianVarietyTheory/Proofs/DecoratedLatticeProofs/KFiniteSpanLeSpanRealIntertwinerValues.lean}{57}{Stability of the polynomial span}
\releasecite{Development/AbelianVarietyTheory/Proofs/TwistedOperators/ThetaLiftUpwardCuspSpanAndAdjoint.lean}{81}{Petersson adjoint with the actual kernels}
\releasecite{Development/AbelianVarietyTheory/Proofs/TwistedOperators/CoreChainCompositesExistLiftCombinations.lean}{75}{Simultaneous multiplicity preimage and real span}
\releasecite{Development/AbelianVarietyTheory/Proofs/TwistedOperators/DefiniteTopSpaceColumnValues.lean}{65}{Real dual coefficient at the definite endpoint}.

Integrate the resulting finite expansion of the corank-one
witness in \eqref{eq:principal-core-detection}. The core
forms are bounded and the lower row form is integrable.
At least one descending chain therefore gives a nonzero
initial pairing. Fix that chain and its last column.
The following finite support and ascent calculations
choose the new row tests with these data fixed.

\teamanchor{proof.principal-finite-coset-support}%
\subsubsection{Finite support for the fixed ascent tests}\label{sec:principal-ascent-finite}
In a transfer write \(V=Y_{i+1}\), \(W=Y_i\) and
\(\Lambda_V,\Lambda_W\) for their prescribed lattices.
The marked tuple \(x\) in \(W\) has rank \(j\) at
stage \(i=2j+1\) and rank \(j+1\) at stage \(i=2j+2\).
The even tuple has one final nonisotropic column.
Let \(R_x\) be the subgroup carrying \(x\) to \(x u\)
with \(u\) upper unitriangular. In rank zero \(R_x=U(W)\).
The point condition on a local component \(h\) is
\(h^{-1}x u=x+a\), with \(a\) in the first block
\(A_1\) of the actual mixed model. Write \(T_a\) for the
corresponding map \(V\to W\) obtained by extending \(a\)
along the input isotropic block. Its columns satisfy
\[
 a_j=T_a e_j^V ,
\]
where \(e_j^V\) is the paired isotropic frame vector in \(V\).
The shift condition is triviality of \(\psi_E\) on the
sum of the adjacent entries of strictly upper triangular
\(N\) whenever \(h^{-1}xN\in A_1\). Here \(\psi_E\) is
the trace character used in the standard upper triangular
entry character. At odd stages there is also
the isotropic-block character condition. The localization
of these conditions is by projection to each finite
component. For the odd case, right invariance of the
test gives the conditions on the full right level coset.
\releasecite{Development/AbelianVarietyTheory/Definitions/TwistedInnerProducts/IsAuxiliaryVectorWittDec.lean}{65}{The point, shift and isotropic-block conditions}
\releasecite{Development/AbelianVarietyTheory/Proofs/TwistedOperators/FlatMapNonvanishingConditionLocalise.lean}{67}{Localization}.

Here is the mixed-coordinate derivation. For
\(F_h=\omega(1,h)\Phi\), a nonzero flattening integral
has a nonzero first mixed coefficient at some \(x u\).
For \(a_2\in A_2\) the dual-block translation by \(h a_2\)
acts there by the phase of the pairing of \(x u\) with
\(h a_2\). By Weil covariance and the \(A_2\) eigencondition
on \(\Phi\), this is also the phase of the pairing of \(x\)
with \(a_2\). The difference \(h^{-1}x u-x\) thus pairs
trivially with \(A_2\). Since \(A_2\) has finite support,
duality gives that its finite part belongs to \(A_1\).

The flattening is the second mixed coefficient at the
fixed \(y\), integrated over \(u\) with weight
\(\overline{\psi_{\rm uni}(u)}\). Translation in the
first isotropic block by \(xN\) replaces \(u\) by
\(u+N\). Additive translation of the strictly upper
entries preserves the unitriangular Haar measure. Indeed
ordinary group translation is triangular in these entries
with all diagonal coefficients \(1\), so product measure
is Haar. The change \(u\mapsto u+N\) gives the multiplier
\(\psi_E(\sum_j N_{j,j+1})\). If \(h^{-1}xN\) lies in
\(A_1\), Weil covariance says this same translation
fixes \(F_h\). Nonzero flattening forces the multiplier
to be \(1\).

If \(h\,\mathrm{inc}_X(\xi)=\mathrm{inc}_X(\xi')\)
for \(\xi\in B_1\), the \(B_1\) eigencondition on
\(\Phi\) gives the phase at \(\xi\) for translation on
\(F_h\) by \(\mathrm{inc}_X(\xi')\). In the second
mixed coordinate this translation is multiplication
by the phase of the pairing of \(y\) with \(\xi'\).
The nonzero flattened value forces the two phases
to agree. This gives the third condition, including its
trivial form when \(y=0\). Right invariance under the
actual finite level gives the same statements for \(hK\).
\releasecite{Development/AbelianVarietyTheory/Definitions/TwistedInnerProducts/WittSmallGroupEmbeddingMemPointwiseStabilizerIso.lean}{253}{The two evaluations and the weighted Haar integral}
\releasecite{Development/AbelianVarietyTheory/Proofs/TwistedOperators/UnfoldingConditionsOfFlatMapWeilActionNeZero.lean}{80}{Mixed-coordinate phases}.

\paragraph{The selected split places.}
For \(v\in S_0\) label the half at \(w(v)\) by \(+\),
the other by \(-\), and put \(I=(\varpi^{d_v})\) there.
The staircase ideal of an auxiliary column with
zero-based index \(j\) in a block of \(c\) columns
is
\[
 \mathfrak a^{c-j}(\overline{\mathfrak a}^{\,j+1})^{-1},
 \qquad
 \mathfrak a=\mathrm{AA}\,\overline{\mathfrak m}^{-1},
 \quad \overline{\mathfrak a}
       =\overline{\mathrm{AA}}\,\mathfrak m^{-1}.
\]
The two completions of \(\mathrm{AA}\) here are unit
ideals. Their integrality follows from the excluded
primes, and the integral element
\(\lvert\mathrm N\mathfrak c_{\!K}\rvert\) of
\(\mathrm{AA}\) is a unit at \(v\). The upper block
has \(c+1\) columns when the lower one has \(c\).
At \(+\) the \(j\)-th raw auxiliary vector of either
lattice has coefficient ideal \((\varpi^{d_v(j+1)})\).
At \(-\) the lower and upper ideals are respectively
\((\varpi^{-d_v(c-j)})\) and
\((\varpi^{-d_v(c+1-j)})\). These are paired with
the opposite columns by the antidiagonal Hermitian
matrix of the auxiliary blocks. In the even lower block its
middle diagonal entry has sign \(-1\).
\releasecite{Development/AbelianVarietyTheory/Definitions/Divisors/CmHostImaginaryGeneratorNeZero.lean}{133}{Staircase ideals}
\releasecite{Development/AbelianVarietyTheory/Definitions/Divisors/BasePlaceNotMemChiGridBadPlaces.lean}{80}{The two local depth ideals}
\releasecite{Development/AbelianVarietyTheory/Proofs/TwistedOperators/SplitFrameOddLatticeRight.lean}{199}{Unit coefficients at the selected places}.

In both marked frames the scalar on the \(j\)-th
raw auxiliary column is
\(\eta_j=(-\xi_0)^{j+1}\), including the last even
column. Thus the marked columns in the \(+\) lattice
basis are normalized by
\(\varpi^{d_v(j+1)}/\eta_j^+\).
By the depth condition \(T_a^+\Lambda_V^+
\subset I\Lambda_W^+\), the same normalization sends
\(a_j^+\) into \(I\Lambda_W^+\).
At the other half \(T_a^-\) maps \(\Lambda_V^-\)
into \(\Lambda_W^-\) by the underlying homomorphism
lattice and the unit coefficient ideal. Consequently
\[
 \frac{\varpi^{-d_v(c-j)}}{\eta_j^-}\,a_j^-
   =\varpi^{d_v}\,
      T_a^-\bigl(\varpi^{-d_v(c+1-j)}a_j^{V,-}\bigr)
   \ \in I\Lambda_W^- ,
\]
where \(a_j^{V,-}\) is the raw upper auxiliary vector.
Pairing \(\Lambda_W^-\) with \(\Lambda_W^+\) and
absorbing the antidiagonal sign gives the normalized
covector congruences. The normalized base columns
and covectors are primitive coordinate frames in the
two halves. The Gram matrix of the transformed tuple
is unchanged, since \(h\) is unitary and \(u\) is
upper unitriangular. In the odd stage the Gram matrix
is zero. In the even stage the only possible nonzero
entry is the last diagonal entry and the last row of
\(u\) is fixed.
\releasecite{Development/AbelianVarietyTheory/Definitions/TwistedInnerProducts/OffQOfSizeRows.lean}{43}{The depth homomorphism lattice}
\releasecite{Development/AbelianVarietyTheory/Proofs/TwistedOperators/BesselStepBasePointScaledCoordinates.lean}{34}{The odd marked columns}
\releasecite{Development/AbelianVarietyTheory/Proofs/TwistedOperators/HermitianPairingFourierJacobiStepBasePointColumn.lean}{32}{The even marked columns}
\releasecite{Development/AbelianVarietyTheory/Proofs/TwistedOperators/DepthLatticeRowMulVecMemModulusScaledCompletion.lean}{78}{Application to the actual row lattice}.

Write \(A,B\) for the resulting base inclusion and
surjection and \(A',B'\) for the transformed pair.
They satisfy \(A'\equiv A\), \(B'\equiv B\pmod I\)
and \(B'A'=BA\). With \(L_A A=1\) and \(B'J'=1\)
over the local integer ring, set
\[
 T_0=1+(A'-A)L_A,\qquad
 T=T_0+J'(B-B'T_0).
\]
Then \(TA=A'\), \(B'T=B\) and \(T\equiv1\pmod I\).
The inverse transpose on the other half is its unitary
extension and preserves \(\Lambda_W^-\).
Thus one principal-level \(k\) satisfies \(kx=h^{-1}xu\).
It follows that \(h=(h k)k^{-1}\) is in \(R_x K_v\).
\releasecite{Development/AbelianVarietyTheory/Proofs/TwistedOperators/SplitPlaceTuplesOneOrbitModUniformizerPower.lean}{314}{The two-sided tuple matching}
\releasecite{Development/AbelianVarietyTheory/Proofs/TwistedOperators/FramedCongruenceElementMemPrincipalCongruence.lean}{49}{The one-sided modulus in the intrinsic level}.

\paragraph{The places away from the depth.}
Put \(\mathcal O=\mathcal O_{E,v}\), interpreted as
the product of the two local rings at a split place.
The depth modulus is a unit here. The completed
\(A_1\) condition means that \(T_a\) takes \(\Lambda_V\)
into \(\mathrm{AA}\Lambda_W\). For its \(j\)-th input
line put \(J_j=\{c:c e_j^V\in\Lambda_V\}\).
The same coefficient condition on \(x_j\) corresponds
to \(\mathrm{AA}\Lambda_W\), as the ratio of the upper
and lower staircase ideals in the same column is
\(\mathfrak a=\mathrm{AA}\) at \(v\). This is the
primitive-column property, with the scaled target
lattice. For adjacent columns,
\[
 J_j J_{j+1}^{-1}
   =(\xi_0)\,\mathrm{AA}\,\overline{\mathrm{AA}}
   =\operatorname{Ann}_{\psi_E}(\mathcal O).
\]
The first equality follows from the staircase formula
and \(\eta_{j+1}/\eta_j=-\xi_0\). The second is the
trace-character normalization at \(v\). These equalities
are of completed ideals.
\releasecite{Development/AbelianVarietyTheory/Proofs/TwistedOperators/DepthLatticeGridRowLocalEqCompletedHomLattice.lean}{49}{Depth condition off the selected places}
\releasecite{Development/AbelianVarietyTheory/Proofs/DecoratedLatticeProofs/MemHomLatticeCompletionIffSupportedMapsStatement.lean}{6}{The scaled target of the homomorphism lattice}
\releasecite{Development/AbelianVarietyTheory/Proofs/TwistedOperators/PsiPrimitiveStandardFramesBesselLevel.lean}{84}{The adjacent ideal ratio and trace annihilator}
\releasecite{Development/AbelianVarietyTheory/Proofs/TwistedOperators/PsiPrimitiveStandardFramesFourierJacobiLevel.lean}{53}{The same ratio for the other parity}.

Use the local Witt Iwasawa decomposition
\(h=n\,d(t)\gamma\kappa\) with \(t\) diagonal,
\(\gamma\) in the unitary group of the Witt core,
and \(\kappa\) in the full lattice stabilizer \(K_v\).
The unipotent \(n\) acts on the marked tuple by an
upper unitriangular matrix. In the even case this includes
the last column since \(n\) fixes it modulo the isotropic
space. For each hyperbolic column \(j\), multiply the
point condition on the left by \(\kappa\) and take
the projection to the \(j\)-th isotropic line in \(W\).
The projection preserves the lattice and annihilates
the earlier marked columns. Upper unitriangularity
therefore gives
\(t_j^{-1}x_j J_j\subset\mathrm{AA}\Lambda_W\).
The primitive-column property gives \(t_j^{-1}J_j
\subset J_j\), hence \(t_j^{-1}\in\mathcal O\).

For \(s\in J_jJ_{j+1}^{-1}\) use the adjacent matrix
\(t_j s E_{j,j+1}\), multiplied on the left by the
upper unitriangular matrix of \(n\). Its image under
\(h^{-1}x\) is \(\kappa^{-1}x\,s E_{j,j+1}\).
This extended map preserves the scaled target lattice
on the source by the definition of \(J_jJ_{j+1}^{-1}\).
The sum of adjacent entries of the test matrix is
\(t_j s\). The shift condition and the displayed
annihilator identity therefore give \(t_j\in\mathcal O\).
This treats every \(t_j\) in the even case and all
but the last in the odd case.
\releasecite{Development/AbelianVarietyTheory/Proofs/TwistedOperators/GoodPlaceOneCosetOfBesselLevel.lean}{119}{Odd Witt Iwasawa reduction}
\releasecite{Development/AbelianVarietyTheory/Proofs/TwistedOperators/FourierJacobiLevelOneCosetAtGoodPlace.lean}{158}{Projections for the even point and shift conditions}.

In the odd case with \(q>0\) isotropic columns, let
\(\mathfrak b\) be the last isotropic coefficient ideal
in \(W\) and \(\mathfrak c\) the coefficient ideal of the
distinguished line in the core of \(V\). The staircase
formula gives
\[
 \mathfrak b=(-\xi_0)^{-q}\mathrm{AA}^{q+1}
                       \overline{\mathrm{AA}}^{-q},
 \qquad
 \mathfrak c=(-\xi_0)^{-(q+1)}\mathrm{AA}^{q+1}
                       \overline{\mathrm{AA}}^{-(q+1)}.
\]
Thus \(\overline{\mathrm{AA}}\mathfrak b^{-1}
=(\xi_0^{-1})\mathfrak c^{-1}\).
The terminal covector of the isotropic-block condition,
multiplied before taking its adjoint by \(s\), has
image coefficient ideal
\((\xi_0^{-1})\overline{s}\,\overline{\mathfrak c}^{-1}\)
on the last dual isotropic line. That line has lattice ideal
\(\overline{\mathfrak b}^{-1}\) by self-duality. This
adjoint belongs to the completed homomorphism lattice
exactly when \(s\in\mathcal O\).
By right level saturation remove \(\kappa\) in the
isotropic-block condition. The last coordinate of the
unitriangular action of \(n\) is fixed and \(\gamma\)
acts trivially on the isotropic block. The condition
therefore says that the terminal covector multiplied
before adjunction by \(t_{\rm last}-1\) annihilates the
integral complementary isotropic block. It annihilates
the entire homomorphism lattice since it only pairs
with the projection to that block. Self-annihilation
and the terminal covector calculation give
\(t_{\rm last}-1\in\mathcal O\).
All \(t_j\) are now units, so \(d(t)\in K_v\).
The factor \(n\gamma\) belongs to \(R_x\) in the
odd case.
\releasecite{Development/AbelianVarietyTheory/Proofs/TwistedOperators/StaircaseIdealsCompletionEqOfModulusUnit.lean}{25}{The two terminal ideals}
\releasecite{Development/AbelianVarietyTheory/Proofs/TwistedOperators/LatticeMembershipAdjointIff.lean}{26}{The terminal adjoint map}
\releasecite{Development/AbelianVarietyTheory/Proofs/TwistedOperators/IsotropicConditionLastTorusDiagonalIntegral.lean}{17}{Using the terminal character and the self-annihilator}.

In the even case apply the same point projection to the
last column, projecting now to the Witt core \(W_0\).
If \(L_z\) is the completed lattice on its distinguished
line, this gives \(\gamma^{-1}(\mathrm{AA}L_z)
\subset\mathrm{AA}\Lambda_{W_0}\), and cancellation gives
\(L_z\subset\gamma\Lambda_{W_0}\).
The line lattice is self-dual. Orthogonal projection to it
preserves \(\gamma\Lambda_{W_0}\), since pairing with
\(L_z\) is integral. Its orthogonal complement in
\(\gamma\Lambda_{W_0}\) is self-dual. At these places the
étale extension is split or unramified. In the split case
a self-dual lattice is determined by either half and its
dual. In the unramified case an orthogonal unit basis is
obtained by splitting a vector with unit norm and using
surjectivity of the unit norm map. Consequently any two
such complement lattices are related by an isometry.
Extending this isometry by the identity on the line gives
\(\gamma=\gamma_0\kappa_0\) with \(\gamma_0\in R_x\) and
\(\kappa_0\) preserving the core lattice, hence also the
full lattice under the Witt embedding. This gives
\(h\in R_xK_v\) in the even case.
\releasecite{Development/AbelianVarietyTheory/Proofs/TwistedOperators/OrbitAny.lean}{19}{Splitting off the fixed line}
\releasecite{Development/AbelianVarietyTheory/Proofs/DecoratedLatticeProofs/SelfDualLatticesSingleOrbitTamePlace.lean}{124}{The self-dual complement orbit}.

At all but finitely many \(v\), the given \(h_v\) already
belongs to \(K_v\), so choose \(r_v=1\) there. The local
factorizations thus assemble in the restricted product
to \(h=r h_\infty k\) with \(r\in R_x(\mathbb A_F)\)
and \(k\) in the prescribed adelic level. This proves the
finite orbit-support condition used in the unfolding.
\releasecite{Development/AbelianVarietyTheory/Proofs/TwistedOperators/OneCosetDecompositionStatement.lean}{21}{Assembly at the selected places}
\releasecite{Development/AbelianVarietyTheory/Proofs/TwistedOperators/StandardAuxiliaryFramesGoodRowGrids.lean}{122}{Putting the support into the odd and even transfers}.

\subsubsection{The depth blocks in the two mixed coordinates}
The ordinary homomorphism lattice, before applying
the depth condition, is the lattice of maps
\(\Lambda_V\to\mathrm{AA}\Lambda_W\). It is self-dual
for the symplectic character. In integral dual bases the
trace character pairs the matrix entries, and the product
\((\xi_0)\mathrm{AA}\overline{\mathrm{AA}}\) is its
coefficient annihilator. The same conclusion holds for
the larger lattice used to define \(B_1,B_2\).
Its target lattice is shifted by
\(\mathfrak m\overline{\mathfrak m}^{-1}\), whose
product with its conjugate is the unit ideal. This
lattice contains the depth lattice \(\Gamma\), since it
has the same depth on the selected half and has the larger
inverse-depth lattice on the conjugate half.

The upper and lower Witt projections preserve the
completed integral lattices. In the auxiliary blocks
these are the projections to the indicated coordinate
summands with their staircase coefficients. Consequently
they also preserve the ordinary, shifted and depth
homomorphism lattices. For a self-dual homomorphism
lattice, the two projected pieces of each hyperbolic
pair are dual, since the pair forms a symplectic
orthogonal summand with complement given by the
other mixed blocks. The shifted lattice makes
\(B_1,B_2\) mutually dual in
the second block. In the first block \(A_2\) is defined as
the annihilator of \(A_1\), and finite local biduality gives
the converse. These are compact open lattices in their
finite coordinate spaces. Applying the projections
decomposes \(\Gamma\) into components in the five
prescribed blocks. The first dual component belongs to
\(A_2\) because the ordinary lattice is self-dual.
The middle components belong to the shifted lattice as
above, and the last is in the core depth lattice by its
definition as the pullback along the core embedding.

For the eigencharacters on components coming from
\(\Gamma\), extension of the marked \(x\) as a first
isotropic block map takes the integral source into the
target scaled by \(\mathfrak a\), the ratio of the two
same-column staircase ideals. This is the ordinary
homomorphism lattice off the selected places. At the
selected places the map has the ordinary lattice on
the \(+\) half and the inverse-depth lattice on the
\(-\) half. These halves constitute the annihilator
of \(\Gamma\), since its \(+\) half has depth and its
\(-\) half has the ordinary lattice. In the odd case
the terminal ideal identity before evaluating away
from the depth is
\[
 \overline{\mathrm{AA}}\mathfrak b^{-1}\mathfrak m^{-1}
       =(\xi_0^{-1})\mathfrak c^{-1}.
\]
Thus the adjoint of the marked \(y\) has the same
ordinary and inverse-depth pattern by the terminal
covector calculation. It too annihilates the relevant
projection of \(\Gamma\). This proves that the two
prescribed characters are trivial on the appropriate
components of the depth lattice. The even stage has
\(y=0\).
\releasecite{Development/AbelianVarietyTheory/Definitions/TwistedInnerProducts/IsAuxiliaryVectorWittDec.lean}{151}{The shifted target and exact mixed pieces}
\releasecite{Development/AbelianVarietyTheory/Proofs/TwistedOperators/BesselToFourierJacobiDepthLatticeSplits.lean}{19}{Odd projections and phase conditions}
\releasecite{Development/AbelianVarietyTheory/Proofs/TwistedOperators/FourierJacobiToBesselLatticeSplits.lean}{67}{Even projections and phase conditions}
\releasecite{Development/AbelianVarietyTheory/Proofs/TwistedOperators/BasePointCharacterTrivialOnDepthLattice.lean}{37}{The first marked character}
\releasecite{Development/AbelianVarietyTheory/Proofs/TwistedOperators/BesselSecondBasePointCharacterTrivialDepthLattic.lean}{330}{The second marked character}.

The core block has a nonzero vacuum because its depth
lattice is contained in the ordinary self-dual
homomorphism lattice in the residual symplectic space.
In a symplectic integral polarization of that lattice,
the characteristic function of one half in the
Schr\"odinger representation is fixed by translations
and by the paired character operators. At split places
one may take the two split halves as the polarization.
This also gives the required half-lattice invariance
at the dyadic split places, and the translations by
twice the lattice satisfy the same invariance. Taking
the restricted tensor product of these local vectors and
any nonzero real Schwartz vector gives an adelic vacuum.
\releasecite{Development/AbelianVarietyTheory/Proofs/TwistedOperators/BesselStepCorePieceFiniteNonzeroVacuum.lean}{75}{The odd residual vacuum}
\releasecite{Development/AbelianVarietyTheory/Proofs/TwistedOperators/JacobiToBesselCorePieceNonzeroVacuum.lean}{29}{The even residual vacuum}.

\paragraph{A nonzero flattened vector.}
The character in the flattening is trivial on the finite
shifts satisfying \(x(u-1)\in A_1\). In fact write
\(J_j^W,J_j^V\) for the completed normalized coefficient
ideals of the lower and upper frames. Before specializing
to the places off the depth, the ratio is
\[
 J_j^W(J_{j+1}^V)^{-1}=(\xi_0)\overline{\mathfrak a}.
\]
The \(j,j+1\)-entry of \(u-1\) belongs to this ideal
multiplied by the coefficient ideal of the scaled target.
On the selected half this multiplier is the depth ideal,
and the product is
\((\xi_0)\overline{\mathrm{AA}}\). Here
\(\mathrm{AA}\) is a unit. On the other finite halves
the multiplier is \(\mathrm{AA}\) and
\(\overline{\mathfrak a}=\overline{\mathrm{AA}}\).
In both cases the product is the trace annihilator
\((\xi_0)\mathrm{AA}\overline{\mathrm{AA}}\).
The adjacent entries therefore all have trivial local
characters. This proves the stated character property.
\releasecite{Development/AbelianVarietyTheory/Proofs/TwistedOperators/GenericCharacterTrivialBesselStepUnitriangular.lean}{25}{The odd finite shift character}
\releasecite{Development/AbelianVarietyTheory/Proofs/TwistedOperators/GenericCharacterTrivialFourierJacobiStep.lean}{103}{The even finite shift character}.

Given the mutually dual compact-open blocks
\(A_1,A_2\) and \(B_2,B_1\), and a nonzero core
vacuum \(\phi_{\mathrm{core}}\) in the fixed \(\Gamma_0\)-space,
take in the two mixed coordinates a product
\(b_1(z)b_2(z')\phi_{\mathrm{core}}(w)\). The finite factors of the
bumps are the indicators of \(x_f+A_1\) and
\(y_f+B_2\), and their real-place factors are smooth nonnegative
bumps positive at \(x_\infty,y_\infty\). Fourier inversion
in the two hyperbolic summands gives its theta vector,
with the positive Fourier normalization for the prescribed
measures. Finite translations by \(A_1,B_2\) fix the bumps,
and multiplication by their dual characters takes its
constant phase at \(x,y\) on their supports. Thus the
vector lies in the required mixed class.

The support of \(b_1(xu)\) is compact in the
unitriangular group since \(x\) has a rational left inverse.
On its finite part the weight \(\psi_{\rm uni}(u_f)\)
is trivial by the preceding ideal calculation.
Choose the real support sufficiently small so that the
real part of \(\overline{\psi_{\rm uni}(u_\infty)}\)
is positive wherever \(b_1(xu)\ne0\). Such a support
exists by the left inverse and continuity at \(u_\infty=1\).
Then the flattening is a nonzero scalar times
\(b_2(y)\phi_{\mathrm{core}}\), with nonzero scalar by positivity of
the real part of the integral on a neighborhood of
the identity. The mixed class is stable under real
Heisenberg translations because the prescribed phases
are all finite. It is also stable under translations by
the embedded core annihilator, which commutes with
the outer hyperbolic blocks and with the translations
that define the fixed core vacuum. This supplies the nonzero
flattened test and the stability used by the Fourier
and Moyal arguments, given the compact-open block
duality and the fixed core vacuum.
\releasecite{Development/AbelianVarietyTheory/Proofs/TwistedOperators/ExistsUnfoldingFlatMapNeZeroFullRank.lean}{42}{Positive real bump and the two Fourier inversions}
\releasecite{Development/AbelianVarietyTheory/Proofs/TwistedOperators/DoubleMixedModelClassStableFlatVacua.lean}{106}{Stability and induced core vacuum}.

\paragraph{The core and line vacuum at each transfer.}
Write \(\mathfrak j=\mathfrak m\overline{\mathfrak m}^{-1}\)
and \(\mathfrak l_\rho
=\mathfrak a^\rho\overline{\mathfrak a}^{-\rho}\).
In the raw line frame \(\zeta_0\) the central auxiliary
line at stage \(2\rho\), for \(\rho>0\), has completed
lattice \(\mathfrak l_\rho\mathcal O_E\zeta_0\).
The line vacuum space uses the target
\(\mathfrak j\mathfrak l_\rho\mathcal O_E\zeta_0\).
Both ideals are inverse to their conjugates, so
these are self-dual line lattices. At rank zero use
the projection of the bottom core depth lattice to the
unit line. It is contained in the ordinary line
homomorphism lattice. This gives a nonzero line vacuum
in every case by the same integral polarization as
for the residual core.

Projection of the residual row depth map to the two
core lattices preserves the ordinary map condition and
the depth condition. This follows because their
embeddings are the original orthogonal summands in the
row lattices, with their prescribed lattices. In an even
step the remaining target projection is onto the line
of index \(\rho+1\). The containment in the shifted
homomorphism lattice gives its target
\(\mathfrak j\mathfrak l_{\rho+1}\mathcal O_E\zeta_0\)
with the ordinary coefficient factor \(\mathrm{AA}\).
This is exactly the line vacuum condition in the
even transfer.
\releasecite{Development/AbelianVarietyTheory/Definitions/TwistedInnerProducts/IsAuxiliaryVectorWittDec.lean}{180}{The line depth lattice}
\releasecite{Development/AbelianVarietyTheory/Proofs/TwistedOperators/FourierJacobiDepthLatticeLineComponentsMem.lean}{35}{The even line target}
\releasecite{Development/AbelianVarietyTheory/Proofs/TwistedOperators/FourierJacobiDepthDatumMulLineVacuumIsVacuum.lean}{39}{The even core and line factors}.

In an odd step take the Hermitian adjoint of the
line component. The link of lines uses
\(\zeta_{\rho+1}=(-\xi_0)\zeta_\rho\), and its
inverse on the line map is the Hermitian adjoint
\(\xi_0^{-1}\mathrm{starAdj}\) in the normalization
of the Weil action. On adjunction the ordinary
coefficient \(\mathrm{AA}\) changes to
\(\overline{\mathrm{AA}}\), and
\[
 \overline{\mathrm{AA}}\,\mathfrak l_{\rho+1}
    =\mathrm{AA}\,\mathfrak j\mathfrak l_\rho .
\]
This is the target required for the line vector at
\(\rho>0\). For \(\rho=0\), using this equality with the
ordinary original target gives the required depth on
the selected half. Applying adjunction with the
simultaneously allowed shifted original target
\(\mathfrak j\Lambda_{C_1}\) cancels \(\mathfrak j\)
in the displayed equality and gives the ordinary line
map condition. The two conditions together are the
rank-zero line depth lattice.

The link reverses the symplectic form. Complex
conjugation of the line vector therefore transports
its Heisenberg invariance. It also exchanges the
two split halves in adjunction, so the one-sided
translation condition transfers along the link. For
the tensor of the core and line vectors, the Weil
translation operators factor under the orthogonal
decomposition. The two projection conditions just
proved make this tensor a vacuum in the flattened
space. In the bottom rank-zero seesaw, the depth
lattice on the target \(Y_0\perp L_0\) is the direct
sum of the two projected depth lattices. The
entrywise Heisenberg annihilator conditions thus
restrict to each summand, and the tensor is a vacuum
for the bottom core depth lattice.
\releasecite{Development/AbelianVarietyTheory/Proofs/TwistedOperators/BesselStepLineComponentEqKappaLatticePoint.lean}{51}{The adjoint line ideal}
\releasecite{Development/AbelianVarietyTheory/Proofs/TwistedOperators/HermitianSkewAdjointMemAdjointHomLattice.lean}{45}{Adjoint normalization}
\releasecite{Development/AbelianVarietyTheory/Proofs/TwistedOperators/BesselStepZeroLineComponentAdjointMemLattice.lean}{89}{The rank-zero depth restriction}
\releasecite{Development/AbelianVarietyTheory/Proofs/TwistedOperators/LinkedLineHomTransportLatticeVacua.lean}{116}{Heisenberg invariance under the linked adjoint}
\releasecite{Development/AbelianVarietyTheory/Proofs/TwistedOperators/LastSeesawLatticesNonzeroLineVacua.lean}{102}{The bottom split depth lattice}.

\teamanchor{proof.principal-fixed-ascent}
\subsubsection{Ascent with a prescribed descending chain}
Write \([H]=H(F)\backslash H(\mathbb A_F)\). Work in the
conjugated lattice realization of the row. Complex conjugation and
reversing the sign of the first Hermitian form each reverse the
symplectic phase. Their combination retains the prescribed standard
additive character \(\psi\), and \(\psi_E=\psi\circ
\operatorname{Tr}_{E/F}\). Write \(C_i\) for the fixed rapidly
decreasing core forms, obtained by the actual descending lifts
\(C_i=\Theta_{D_i}(C_{i+1})\) with the prescribed core tests
\(D_i\). Write \(G_0=u_w\), and let \(G_{i+1}\) be a lift of
\(G_i\) with one of the permitted tests of the \(i\)-th row arrow.
The desired implication is the following. If the rank-zero pairing
of \(G_0\) and \(C_0\) is nonzero, then the tests on the row
arrows can be chosen so that the final Bessel period with \(C_m\)
is nonzero. This assertion holds for each compatible grid, each
prescribed descending sequence and every fixed \(w\) giving the
initial pairing. In addition to the core forms, only the initial row
and its partial lifts are required to be rapidly decreasing. From the
first lift onward the row is fixed by its prescribed finite group.
Indeed the resulting theta integral is fixed in its target variable
by any element that fixes its Weil test on that side.

Put \(\alpha_i=(\chi^{\mathrm{grid}}_{m-i})^{-1}\), and let
\(\eta\) be the odd line splitting character. The ordered pair
on the row arrow, with the larger Hermitian space first in the
oriented symplectic realization \(\operatorname{Hom}(V,W)\), is
\((\alpha_i\eta,\alpha_i)\) for odd \(i\), and
\((\alpha_i,\alpha_i\eta)\) for even \(i\). On the core
arrow it is \((\alpha_i,\alpha_i)\). These are the inverse
characters of the un-conjugated data because complex conjugation
changes every unitary scalar in the Weil action to its inverse.
Orthogonal direct sum multiplies the character attached to the
summands and leaves the opposite character unchanged. Thus on
splitting off the one-dimensional summand its character is
\(\eta\), with the common opposite character on the same
oriented space. For the other orientation its realization is
transported by the prescribed Hermitian adjoint and complex
conjugation. In particular, the additive characters on the central
and side coordinates in the following computation come from the
same adapted action.
\releasecite{Development/AbelianVarietyTheory/Definitions/TwistedInnerProducts/PartFixed.lean}{87}{The conjugated Weil datum and the actual integrals}.

Here is a way to check the simple action laws used below.
With \(B\) the oriented symplectic form, use
\[
 T_w F(v)=\psi(B(w,v)/2)F(v+w),\qquad
 T_wT_t=\psi(-B(w,t)/2)T_{w+t}.
\]
At a rational translate the lattice relation gives
\(\psi(B(v,w)/2)\) with the opposite sign, so their product
has full symplectic phase after moving the translate past a Weyl
operator. The adjoint identification of the linked line has
\(B'(Lx,Ly)=-B(x,y)\). Consequently evaluation of a mixed
coefficient translated by that line equals the conjugate line
coefficient, with its full central phase conjugated. At split finite
places the same normalization, in matrix coordinates, is
\(\chi_W(\det k_V)^{-1}\chi_V(\det k_W)\) times pullback
of the Schwartz function. The two determinant factors are on the
two sides of the same action. At the other finite places the
vacuum-normalized splitting gives the identity on the self-dual
vacuum. The finite-coset proof gives the stronger depth, support and
annihilator properties for these particular adapted lattices.

\paragraph{The two unfoldings with their test functions.}
In either step let \(V\) be the upper row space and \(W\) the
lower. Write \(V=X\oplus V_0\oplus Y\) with \(X,Y\)
dual isotropic, and \(\ell=\dim_E X\). In the flag unipotent
write \(n=n_C(c)n_B(b)m(u)\), with \(c^*=-c\),
\(b\colon V_0\to X\), and \(u\in U_\ell\) upper
unitriangular. In the sign convention for \(B\), the translate of
the first mixed row by \(n_B(b)\) has components
\((x,xb,-xbb^*/2)\). Denote the first mixed transformation of a
lattice vector by \(F^X(x;t)\). Precisely, it is the evaluation
at zero of the product of the Weyl operators for \(j_Y(A)\),
\(j_X(x)\) and the middle vector \(j_0(t)\), integrated over
the rational quotient in \(A\). This definition fixes the half
phases. If the lower flag is nonempty, its analogous transformation
is denoted \((F^X)^Y(y;t)\); its dual argument is inserted using
the adjoint for \(B\). Give all additive quotients covolume one.
Give \(U_\ell(\mathbb A_F)\) the Haar measure for which
\([U_\ell]\) has volume one. These choices realize the
unipotent quotient Haar of the periods. With any prescribed
positive quotient Haar proportional to them, both sides acquire
the corresponding positive scale.

Fourier expansion of the evaluation at zero along the compact
dual block gives the sum of \(F^X(x;0)\) over rational frames
\(x\). On \(n_C(c)\) its character is
\[
 \psi\bigl(B(j_X(x),j_Y(xc))/2\bigr).
\]
The coefficient selected by the upper Bessel weight therefore has
Gram matrix zero. For the upper Fourier--Jacobi weight the product
of the conjugate line function and the conjugate adjacent-root
character has central weight equal to the conjugate of this phase
at \(z e_\ell^t\), where \(z\) is the lower distinguished
line vector. In that case the selected Gram is
\(\operatorname{diag}(0,\ldots,0,\langle z,z\rangle)\).
The trace pairing on the space of skew-Hermitian matrices is
nondegenerate, so integrating over \([c]\) imposes precisely the
indicated Gram condition.

Only independent frames contribute. To see this, take the first
dependent column \(j\) and a relation \(x d=0\) with
\(d_j=1\) and \(d_k=0\) for \(k>j\). For
\(j<\ell\), the stabilizer contains
\(1+d s e_{j+1}^t\), with \(s\in\mathbb A_E\). Its
adjacent-root weight is \(\psi_E(s)\), while it fixes \(x\).
It also fixes the Jacobi line row \(z e_\ell^t\), and in the
Bessel case it leaves the last-row side weight unchanged. With last
Gram entry nonzero necessarily \(d_\ell=0\). In the other case
with \(j=\ell\), the side translation
\(b=s d z^*/\langle z,z\rangle\), with \(z\) now in the
upper Bessel middle space, satisfies \(xb=0\) and has side
weight \(\psi_E(s)\). In either case the orbit integral
contains the integral of a nontrivial character on
\(E\backslash\mathbb A_E\), hence is zero. For an independent
frame the specified Gram gives a single rational orbit, by Witt's
extension theorem applied to the isometry of the two ordered
subspaces. Choose \(x_*\) to be the reference lower isotropic
frame in the first step, and that frame followed by \(z\) in
the second. Write
\[
 R=\{g\in U(W):g x_*=x_* u_g,\ u_g\in U_\ell\}.
\]
Its reductive part is the full lower middle group for the
isotropic frame; with the last line appended it is the stabilizer
of the distinguished line vector in that middle group. Its
unipotent part is the lower flag unipotent. These descriptions
follow by acting on the flag columns and on the orthogonal
quotient of the isotropic span. In the appended-line case, the
last adjacent entry of \(u_g\) is the lower side coordinate.

Let \(a\) be the upper core variable and embed it as
\(\iota(a)\) in \(U(V)\). Let \(\phi=C_{i+1}\), let
\(G=G_i\) with the conjugations of the upper period included in
its weight, and set \(F_{a,g}=\omega(\iota(a),g)\Phi\).
Unfolding the independent rational orbit and changing the
rational translates of \(u\) gives the upper period as
\[
 \begin{gathered}
 \int_{R(F)\backslash U(W)(\mathbb A_F)} f_\Phi(g)\,dg,\\
 f_\Phi(g)=G(g)\int_{[H]}\!\phi(a)
 \int_{U_\ell(\mathbb A_E)}\!\int_{[b]}
    \vartheta_a(b,u) F_{a,g}^X(x_*u;x_*b)\,db\,du\,da .
 \end{gathered}
 \tag{A}
\]
Here \(H\) is the upper core group and the unipotent group
in the \(u\)-integral has the stipulated rational covolume.
The two weights and the associated contractions are as follows.

\begin{enumerate}
\item In the Bessel-to-Jacobi unfolding the weight is
\(\overline{\psi_E(\sum_j u_{j,j+1})}
\overline{\psi_E(e_\ell^t b z)}\). Put
\(y_*=-z e_\ell^t\) in the adjoint block. For \(\ell=0\)
both frame blocks and their characters are trivial and these
arguments are zero. The equality
\[
 \overline{\psi_E(e_\ell^t b z)}F^X(x_*u;x_*b)
  =\bigl(T_{j_{W,X}(b)}T_{\operatorname{incAdj}(y_*)}
              F^X(x_*u;\cdot)\bigr)(0)
 \tag{B}
\]
is the Weyl law followed by the lattice relation at the rational
vector \(\operatorname{incAdj}(y_*)\). The two half phases
add to the side phase. Thus integration of (B) over \([b]\) is
the second mixed coefficient at \(y_*\). Define
\(\mathcal F_{BF}\Phi\) by integrating the resulting mixed
coefficient over \(u\) against the conjugate adjacent character;
retain the remaining core variable.

For \(n\) in the lower flag group, substitution of
\(n x_*=x_*u_n\) in \(u\) changes the adjacent factor by
\(\overline{\psi_E(\sum_j(u_n)_{j,j+1})}\). The second
adjoint argument changes from \(v=\operatorname{incAdj}(y_*)\)
to \(n^{-1}v\). Its projection to the remaining line block is
the line Weyl translate. Moving the remaining rational block
through the first evaluation gives the scalar
\(\psi(B(v,n^{-1}v)/2)\). These are respectively the
flag factor, the line side action and its central scalar in the
lower Jacobi datum. For \(h\) in the lower middle group and
\(a\) in the upper core, the remaining action is
\(\omega_0(a,h)\). This follows by commuting the two embedded
middle actions with the mixed transformations. Equivalently, the
new line action is the Weil action on the lower Jacobi line
brought across the Hermitian adjoint; conjugation on the line
and the sign in \(y_*\) are essential in this identification.
For a tested tensor \(D_i\otimes L\) on the remaining core
and line, the integral over \([R]\) in (A) is consequently the
lower Fourier--Jacobi period with core
\(\Theta_{D_i}(\phi)=C_i\) and the transported line test \(L\).
The product formula follows by inserting the tensor and
integrating \(\phi(a)\omega_0(a,h)D_i(0)\) in \(a\).

\item In the Jacobi-to-Bessel unfolding let
\(Y_a=\omega_{\mathrm{line}}(a,1)\Phi_0\) be the fixed
upper line test after the upper core action. The weight is
\[
 \vartheta_a(b,u)
 =\overline{\psi_E(\sum_j u_{j,j+1})}
                  \,\overline{Y_a(z e_\ell^t b)} .
\]
For the first
\(\ell-1\) rows of \(b\) the integral is the second
mixed coefficient at zero. The last row maps isomorphically
over \(E\) to the line block. Thus the remaining integral is
\[
 \int_{[\mathscr W_{\mathrm{line}}]}
   \bigl(\omega_0(a,(h,1))\Xi\bigr)(0,t)
                  \overline{Y_a(t)}\,dt,
 \qquad \Xi=\mathcal F_{FB}\Phi ,
 \tag{C}
\]
where \(\mathcal F_{FB}\) is the weighted double mixed
coefficient at \((x_*,0)\), and \(h\) fixes the appended
line. The rational line isomorphism carries the mass-one
measure on the first compact quotient to the mass-one measure
on the second. Substituting \(u_g\) as above produces the
lower Bessel character: its last adjacent coefficient is the
side term \(e^t b z\), and the preceding adjacent
coefficients are those of the lower flag. The second evaluation
at zero has no remaining central phase. This verifies the
lower mean as the Bessel integral with the handed form
\(h\mapsto\int_{[H]}\phi(a)\,\textup{(C)}\,da\).
For \(\Xi=D_i\otimes\Phi_0\), factor the Weil action
on this tensor. The line inner product in (C) is then
\(\|\Phi_0\|^2\) by unitarity, and the handed form is
\(\|\Phi_0\|^2 C_i(h)\).
\end{enumerate}
\releasecite{Development/AbelianVarietyTheory/Proofs/TwistedOperators/VanishingDescendsBesselToFourierJacobi.lean}{127}{First mean and normalized character comparison}
\releasecite{Development/AbelianVarietyTheory/Proofs/TwistedOperators/VanishingDescendsFourierJacobiToBessel.lean}{237}{Second mean with the actual line test}.

\paragraph{Absolute convergence and the topology of the tests.}
The lattice realization is obtained from the adelic Schwartz
realization by the Weil--Brezin transform: in a polarization it
is the sum of \(\phi(x+\lambda)\), \(\lambda\in X(F)\),
against the ordinary character of \(\lambda\) in the dual
coordinate, with the half phase of the Weyl law. Conversely
Fourier coefficients in that compact dual coordinate recover
\(\phi\). On fixed finite support and level, differentiation
on the real compact quotient gives Schwartz seminorms of the
inverse, and summing against rapid decay gives seminorms of
the forward transform. In particular the mixed transforms here
are restrictions and partial Fourier transforms on adelic
Schwartz space. For fixed finite datum the \(u\)-integral
defining a flattened test is continuous in these seminorms.
The map \(u\mapsto x_*u\) is a proper affine linear
embedding on upper-entry coordinates with a left inverse from
the full rank of \(x_*\); its real integral is bounded by
integrating a Schwartz majorant in those coordinates. The
finite part ranges over a compact set for a fixed finite test.

For absolute convergence of (A), first sum the absolute values
of the mixed evaluations on the original compact flag quotients
before unfolding. A Schwartz seminorm bounds the real sum over
rational \(x\) by summation of
\((1+\|x\|)^{-B}\) on a lattice for \(B\) larger than
its real dimension; a second power bounds the lattice-point
count after increasing denominators. In detail, choose a product
finite symplectic lattice \(L_0\) small enough that its Weyl
operators fix the original Schwartz vector. After acting by
\(g\) the vector is fixed under \(gL_0\). Locally
\(\varpi_v^{d_v}L_{0,v}\subset g_v L_{0,v}\) with
\(\prod_v q_v^{d_v}\ll H(g)^A\), by the entries of the
inverse action on this lattice. The dual split coordinates of
this smaller lattice force support in a denominator lattice
with index polynomial in that product, and its translation
coordinates give constancy on cells of inverse polynomial
volume. The finite supremum of the transformed vector is
therefore bounded by a power of \(H(g)\) times its fixed
\(L^2\) norm. At infinity use a Cartan decomposition in the
real symplectic group. The two compact factors act by a
bounded family on each Schwartz seminorm, with finitely many
input seminorms. The diagonal part acts by a determinant
factor and linear coordinate changes, which bound its
Schwartz seminorms by powers of matrix height. These arguments
give a polynomial bound in the upper
core and lower quotient heights, uniform for \(a\) or \(g\)
translated in a fixed compact set. For the Jacobi case the
absolute value of the upper line function is periodic in the
line translation and obeys the same bound uniformly on its
compact quotient. The estimates descend to quotient heights:
rational changes of the middle bases permute the frame sum
and preserve the mass-one measures on the additive quotients,
so one may take reduced representatives for \(a\) and \(g\).

The same estimate applies when restricting \(g\) to \(R\)
and when integrating the means for a compact family of right
translates. The heights of \(n\iota(b)\) and \(b\) are
comparable up to powers on these quotients. One can see the
inequality in the forward direction by reducing \(b\) in its
middle group and then choosing \(n\) in a compact fundamental
set for the rational unipotent. In the reverse direction use
the wedge \(e_X\) of the rational basis of the isotropic
flag. On \(e_X\wedge V_0\), the unipotent action is the
identity, and the embedded middle acts by
\(e_X\wedge v\mapsto e_X\wedge b v\). For every nonzero
\(v\in V_0(F)\) and every rational \(\gamma\) in the
ambient group, the vector
\((\gamma n\iota(b))(e_X\wedge b^{-1}v)
=\gamma(e_X\wedge v)\) is nonzero rational. Its product
norm is bounded below by the product formula on any nonzero
rational coordinate. Thus the product norm of \(b^{-1}v\) is
bounded below by a fixed inverse power of
\(H(\gamma n\iota(b))\). Take the infimum over rational
\(\gamma\). Mahler reduction for the middle lattice bounds
its quotient height by a power of the reciprocal of this
shortest-vector bound. The same argument with the rational
middle vector in the direct summand applies when the
distinguished line is fixed. This comparison is invariant under
changing the chosen compact set of right translates. Choosing
the rapid decay powers of \(\phi\) and \(G\) larger than
the polynomial powers thus bounds the absolute integrals by an
integrable majorant on the core and lower quotients, which have
finite volume; the unipotent quotients are compact. In (A)
one can equivalently unfold the nonnegative sum first, by
Tonelli, obtaining exactly the integral of the absolute
independent-orbit integrand. These arguments prove the two
unfoldings, their disintegrations over \([R]\), continuity
of the \([R]\)-mean in the real right translate, and
continuity as a linear functional of the remaining flattened
Schwartz vector. They also justify passage from finite sums of
real Fock tests to Schwartz tests by the convergence of the
Hermite-function expansion in the Schwartz topology.

\paragraph{From a global zero to the mean on the finite coset.}
Fix the preceding row, together with a nonzero upper line
test in the Jacobi-to-Bessel step. Assume all the permitted
new row tests, with the allowed upper line test in the Jacobi
case, give zero upper period. The finite Heisenberg Weyl
operators on the depth space permute its actual coset basis,
up to unit scalars. Explicitly, in a split polarization
write the depth lattice as \(I X\oplus X^\vee\). Its
annihilator has split coordinates \(X\oplus I^{-1}X^\vee\).
Translation by \(x\in X\) sends
\({\bf1}_{\kappa+IX}\) to
\({\bf1}_{\kappa-x+IX}\); the dual coordinate acts on
that coset by the character evaluated at \(\kappa\).
The Weyl half phase is a constant. The conjugated model
conjugates these scalars. The rational coset labels give all
finite residues by the Chinese remainder theorem. Equivalently
the allowed lattice vectors are fixed by \(T_{2u}\) for
every depth-lattice vector and by \(T_u\) on both split
halves. This is the same condition at dyadic split places,
where requiring only the first invariance would lose the
coset bound. Off the depth primes the space is the self-dual
vacuum line. Thus the coset tests span the full finite tested
space. The Hermite expansion extends the assumed vanishing
to all real Schwartz tests in that finite space. In particular
it holds on the family of tests with the finite support and
annihilator properties of the finite-coset proof, a family
closed under real Weyl translation.

Modulate \(\Phi\) by any Weyl operator on the first dual
block supported at infinity. The mixed-coordinate Weyl law
multiplies the integrand before the \(a,u,b\) integrations
in (A) by
\[
 \psi_\infty\bigl(B(j_X((g^{-1}x_*u)_\infty),j_Y(c))\bigr).
 \tag{D}
\]
All real frequencies occur by nondegeneracy. The absolutely
integrable unfolded integrand thus defines a finite complex
measure whose pushforward under
\((g,a,u,b)\mapsto (g^{-1}x_*u)_\infty\) has every Fourier
coefficient zero. For a compact set \(Q_\infty\) in the
real lower group, the subset
\(Q_\infty^{-1}x_*U_\ell(F\otimes\mathbb R)\) of the
real frame space is closed. Indeed, from a converging
sequence in it, take a convergent subsequence in
\(Q_\infty\) and apply a left inverse of \(x_*\) to
obtain a limit in the closed upper unitriangular group.
The inverse image of this set under the displayed pushforward
condition is exactly
\(g_\infty\in R_\infty Q_\infty\), independent of
\(u\). Fourier uniqueness gives zero integral over every
such strip. This argument uses the joint frame \(g^{-1}x_*u\)
and the joint side and central weights of the single period.

Use now the finite support result for the flattened evaluation
of the actual test: the integrand after the inner integrations
in (A) is supported on
\(R(\mathbb A_F)U(W)(F\otimes\mathbb R)J\) with \(J\)
the lower right finite group. It is right \(J\)-invariant by
the Weil action on the test and the preceding row invariance.
Disintegrate the unfolded quotient along
\(R(F)\backslash R(\mathbb A_F)\). Its quotient Haar
has the product description from the lower flag unipotent and
its middle group; the determinant of the latter on the
unipotent is one, as it is an \(F\)-rational character of
the middle unitary group. Over the support there is a single
finite orbit, whose positive measure is a constant multiple
of \(\operatorname{vol}(J)/\operatorname{vol}(J\cap
R(\mathbb A_{F,f}))\). The vanishing on strips therefore
implies that the real quotient integral of the \([R]\)-mean
on every compact set is zero. The mean is continuous, so its
value at the identity is zero. By the computed covariance this
is vanishing of the lower functional on the flattened vectors.

\paragraph{Using a prescribed lower test.}
The chosen compact support bump gives a nonzero flattened
depth vector, and every translate by its actual annihilator
is the flattening of another permissible Schwartz vector:
extend the translate in the core block before the two mixed
transforms, which commute with it. To see why these translates
test the entire depth space, let \(f_0\) be the nonzero
flattened test and use the self-dual identity
\[
 \|f_0\|^2 f
   =\int_{\mathscr W(\mathbb A_F)}
         \langle f,T_v f_0\rangle T_v f_0\,dv .
 \tag{E}
\]
In the Schr\"odinger realization (E) follows by Fourier
inversion in the modulation coordinates followed by
integration in the translation coordinates. Both sides have
the same normalization of Haar, so there is no extra scalar.
At fixed finite test data its coefficient has compact finite
support. At infinity its coefficient is Schwartz, by repeated
differentiation and integration by parts in the translation
and modulation coordinates. Since the seminorms of \(T_v f_0\)
grow at most polynomially, (E) converges in Schwartz space.
By applying the common lattice invariances to both vectors
in the coefficient and commuting their Weyl operators, the
coefficient vanishes off the annihilator, including the two
split-half conditions. Applying the continuous lower mean to
(E) now shows that it vanishes on every lower depth vector.

In the first unfolding insert the fixed tensor
\(D_i\otimes L\) using the given lower line test transported
by the linked line map. In the second insert
\(D_i\otimes\Phi_0\) with the very same upper line vector
in its contraction. The split depth block decomposition and
the conjugate adjoint matching ensure these are vectors of
the flattened depth space. By (B) and (C), respectively,
vanishing would give the lower period equal to zero, or
\(\|\Phi_0\|^2\) times that period equal to zero. The upper
line test in the second case can be any nonzero prescribed
depth test. This proves the contrapositive at each step. Since
finite cosets and real Hermites span densely in the topology
used in (A), at least one individual permitted coset and
polynomial succeed on the new arrow. Earlier choices and
the automorphic input remain fixed.

At rank zero \(C_0=Y_0\perp L_z\), and \(Y_1=C_1\).
The chosen line space of the Jacobi datum is \(-L_z\), so a
product test in the core descent has row factor \(D\) and
line factor \(\overline{L}\), with \(L\) on the Jacobi
line space. For this test, use the split Weil action in the two
blocks to write the bottom pairing as
\[
 \begin{split}
 &\int_{[U(Y_0)]}\int_{[U(Y_1)]}
      C_1(a)\,\theta_D(a,b)\,\overline{\theta_L(a,1)}
                       G_0(b)\,da\,db \\
 &\qquad=\int_{[U(Y_1)]} C_1(a)\,\overline{\theta_L(a,1)}
                                    G_1(a)\,da,
 \qquad G_1(a)=\int_{[U(Y_0)]}\theta_D(a,b)G_0(b)\,db .
 \end{split}
\]
This is the rank-zero Jacobi period, with the prescribed
upper core measure and lower row measure. If the bottom core
measure differs from the bottom row measure, their constant
ratio is positive. The equality is justified by the same
Schwartz and rapid decay bounds. Finite sums of these tensor
tests are dense in the depth space on the split block, by the
finite residue basis and real Hermite expansion. Consequently
some permitted first row test and linked line test give a
nonzero Fourier--Jacobi period. The alternating transfers
then retain each lower test. Since \(m=2\rho+2\) is even,
the last period is Bessel.
\releasecite{Development/AbelianVarietyTheory/Proofs/TwistedOperators/NonvanishingThetaTowerBottomRung.lean}{27}{Rank-zero splitting and descent input}.

At the last step the auxiliary isotropic frame is
\(e_j=(-\xi_0)^j e_j^{\mathrm{grid}}\), \(1\le j\le\rho\),
and the marked line vector is
\((-\xi_0)^{\rho+1}z^{\mathrm{grid}}\). The stabilizer and
flag subgroup are the original ones. For the simple-root
coordinates on all consecutive columns, including the final
line coordinate, this changes the matrix coefficient by
\((-\xi_0)^{j+1}/(-\xi_0)^j=-\xi_0\).
The grid Bessel character uses \(\psi_E(\xi_0\sum_j s_j)\)
in these unscaled coordinates. The auxiliary one is therefore
its conjugate. Applying conjugation to the terminal auxiliary
integral recovers \(B^+\) in the original model up to the
positive Haar ratio of the top definite quotient. The flag
quotients on both sides have volume one, because the rational
frame changes identify those quotients. This proves that the
terminal integral in the ascent detects the prescribed
\(B^+\).
\releasecite{Development/AbelianVarietyTheory/Proofs/TwistedOperators/StandardAuxiliaryFramesFitGridEnds.lean}{201}{Endpoint character}
\releasecite{Development/AbelianVarietyTheory/Proofs/TwistedOperators/StandardAuxiliaryFramesFitGridEnds.lean}{262}{Endpoint pairing and Haar}.

\teamanchor{proof.principal-real-bessel-line}
\subsubsection{The real Bessel line in the arithmetic tests}
Here is the real identification on the row tower. At a given
real place let \(P_i\) be the unshifted signed row list, put
\(s=r+6\) and
\(\delta_i=-1/2\) for odd \(i\) and \(0\) for even \(i\).
The actual row has the dual of \(P_i+\delta_i\).
For the arrow of index \(i\), write \(c=c_{m-i}\).
In the original action the two splitting exponents, in the
order larger, smaller, are \(c+\delta_i,c+\delta_{i+1}\).
They are opposite-parity because \(s\) is odd. The larger
space in this orientation has negated form. After matching
the two centers, reversing the signature signs below the
center, and adding the centered string
\[
 ((s-1)/2,(s-3)/2,\ldots,-(s-1)/2)
\]
in the negative list, one obtains precisely the parameter of
the larger negated row on the determinant cover, before
taking the dual. In fact every old entry obeys
\(x+(s+1)/2\le c\) in \(P_i\).
Thus every old entry reverses its sign label, which is exactly
the signature flip from negating the form, and the new block
agrees with \(P_{i+1}\) with its shift. The inequality is
the grid entry bound obtained by putting its new center above
all old blocks and the old lists. The rank and parity shifts
make each real infinitesimal character regular and in the
discrete-series lattice for the corresponding row group.
This is the signed pair in
\href{https://doi.org/10.4153/CJM-1999-029-6}{Paul, Theorem~2.7(b)}
with the displayed inserted string. For first occurrence, in
the notation of \href{https://arxiv.org/abs/2008.06174}{Ichino,
Theorem~6.1} the smaller centered parameter has no zero and
\(k_\pi=-1\), so we are in the odd difference case with
\(t=(s-1)/2\) and \(l=0\) in its tested tower. The old
entries are outside the range for \(C_\pi^\pm(t)\). Deleting
one hyperbolic plane decreases \(l\) below zero, so the
lower real lift vanishes. The global constant-term formula
then gives cuspidality for these row lifts, and real Howe
duality gives the specified target type. The same argument
works after taking the original dual convention on both
sides.
\releasecite{Development/AbelianVarietyTheory/Proofs/TwistedOperators/ContragredientThetaParamsAddBlockArrow.lean}{17}{Original row lists, shifts and centers}
\releasecite{Development/AutomorphicArithmetic/ThetaCorrespondence/ThetaPairLifts/LocalTypeDictionary/SeparatingFamilyExtendsByTwistedRow.lean}{16}{The corresponding real type for the placed data}.

For each real place take these fixed smooth irreducible modules
and the indicated top compact dual. The real
polynomial row parameters \(\rho\) and a smooth bottom vector
\(w\) map to the same vector \(v(\rho,w)\) of this row
module for all multiplicity vectors and all finite tests.
Indeed real local Howe intertwining at each arrow has one
line for the specified source and target types. Fix one
intertwiner on the Fock polynomials at each real place. With
all finite factors held fixed, write the kernel integral on
the real isotypic target in its multiplicity space. Each
component map of the source and real oscillator test is an
intertwiner and hence a scalar multiple of the fixed map.
The scalar can depend on the finite factors and the input
multiplicity, and is independent of the real polynomial and
input real vector. Use the smooth globalizations to extend
this equality in the source vector: continuity of theta
integration on smooth cuspidal vectors follows from uniform
rapid decay seminorms and the polynomial bounds above.
Iterating
gives \(\theta_{\rho,\Phi}(t w)=F_{t,\Phi}(v(\rho,w))\).
The finite multiplicity map \(F_{t,\Phi}\) carries the
remaining scalar and may be zero. The real column module
is the tensor product of the local dual highest-weight
modules, and each column map with fixed \(\Psi\) is
equivariant on this module. These factorizations use the
same real modules at every varying \(t,\Psi,\Phi\).

The iterated Bessel pairing for given maps is a continuous
bilinear functional on these real modules. In the terminal
period the top definite unitary group is anisotropic, so the
space of evaluation of the two automorphic functions is
covered by a compact set from that quotient and the
flag-unipotent quotient. The integrals converge absolutely
and are bounded by smooth seminorms on this compact set.
A simultaneous right translation of the two arguments by
the embedded top real group leaves the pairing invariant;
right translation by the flag changes it by the fixed Bessel
character, since replacing \(n\) by \(nn_0\) in the
weighted integral produces \(\psi_{\mathrm{Bes}}(n_0)\)
from the conjugate weight. For each real place it is generic in
the sense of
\cite[Theorem~A and \S2.2]{p2jiang-sun-zhu}. All adjacent-root characters are
nontrivial, and the side character is nontrivial on the
distinguished nonisotropic line, including its prescribed
nonzero scalar from the terminal frame. The theorem
applies to the irreducible smooth Fr\'echet
globalizations in question. Applying it with the other
local vectors held fixed, successively at every real place,
gives a single line of such bilinear functionals. More
explicitly, choose pure tensors at which a nonzero
functional is nonzero. At each place the local functional
obtained by fixing those other factors spans its
one-dimensional space. Proportionality then holds on pure
tensors by induction on the number of varying factors and
extends by continuity. Pulling back along \(v(\rho,w)\)
gives the same function \(a(\rho,\xi)\) up to a scalar
depending on \(t,\Psi,\Phi\), for fixed \(w\).

If the input space is nonzero, bottom separation and the
proved ascent give one triple with nonzero terminal period
at some real pair. Choose that pair \((\rho_*,\xi_*)\).
For every other nonzero input, its triple supplied by the
same separation has nonzero scalar on the common function,
and hence the same pair \((\rho_*,\xi_*)\) works. This
choice depends on the grid, volumes and \(w\) and precedes
the quantifier over the nonzero inputs. For an empty
separation problem any real pair suffices.

\teamanchor{proof.principal-test-column-span}
For the final arithmetic-to-period implication, fix a finite row
test \(\Phi\). The arithmetic list for this \(\Phi\) runs through
all pairs of a terminal double coset and a basis vector of the
prescribed column coefficient space. The terminal group is definite,
so the rational lattice stabilizers are finite. Its congruence
level is principal at places above the two distinct type primes.
A finite-order matrix acting trivially modulo one of those places
has order a power of that residue characteristic. Indeed, in the
filtration by deeper principal subgroups each successive quotient
has that characteristic, and the intersection of the filtration
is the identity. The two places therefore force every
stabilizer at this level to be trivial.

In the terminal column level the finite right character is
trivial. A tuple of values at the class representatives extends
uniquely to a column by the rational coefficient covariance and
right-level invariance. The class stabilizer calculation makes
this extension well defined for arbitrary coefficient values.
The arithmetic test attached to a representative and a basis
vector has precisely that value at the representative and zero
at the other classes. These tests thus span all columns of the
specified compact dual type on the same level. In particular the
column supplied by separation is in their span for this fixed
\(\Phi\). Contraction by \(\xi_*\) preserves this linear combination.
\releasecite{Development/AbelianVarietyTheory/Definitions/Divisors/GridTestLatticeLargestSlopeBound.lean}{116}{The arithmetic test index and coefficient read}
\releasecite{Development/AbelianVarietyTheory/Proofs/DivisorProofs/EqOneOfConjMemOfPow.lean}{170}{Triviality of the actual class stabilizers} and
\releasecite{Development/AbelianVarietyTheory/Proofs/TwistedOperators/EqAdFormsOfValues.lean}{102}{The span for fixed finite row test}.

At this same real
pair, absolute convergence gives linearity of both
integrals in the column. If the arithmetic functionals
vanish, their integral comparison gives vanishing on each
column and therefore on the whole span. Here the volume
and normalization factors in the comparison are nonzero.
If starting from a preselected arithmetic real pair,
apply the stipulated real-replacement property of that
choice to the vanishing family before using the newly
chosen pair. Equivalently one may arrange the arithmetic
choice to use \((\rho_*,\xi_*)\) from the outset. The
base-choice conditions concern the Weil data, measures,
nonzero \(w\), and square constants, so this replacement
preserves them. Injectivity at the good pair then yields the
real-replacement property for that choice, since a
simultaneous zero there forces the input to be zero.
The separating nonzero period then gives the contradiction
for every nonzero realizable input.
\releasecite{Development/AbelianVarietyTheory/Proofs/TwistedOperators/InjKOfSeparates.lean}{57}{Use of the arithmetic column span}
\releasecite{Development/AbelianVarietyTheory/Definitions/TwistedInnerProducts/ExistsMatchOver.lean}{255}{Real replacement from an injective choice}.

\subsection{From top theta periods to bottom Bessel periods}

\teamanchor{proof.relation-transport}
The next argument concerns the simultaneous comparison stated
at the end of the section. Its full base lift and nonzero reference
are used to derive that comparison. The principal-level tests above
and the arithmetic realizations below provide separate inputs to
the descended pairing.
At the selected translate, the period comparison uses the
geometric embeddings, real types, twisting characters, levels and
compatible local tests in its hypotheses. Here \(r_S\geq1\), so
\(m\geq4\). In the argument establishing the comparison there are
two surjectivity assertions. The projected rows descend onto
\(E_{\rm row}\), and the base chain on the larger top space reaches
every row in the full permitted bottom multiplicity space for the
fixed level, character and discrete series. The second assertion is
used for the nonzero reference in that argument. The first is also
needed for descent in the arithmetic application below.
For the theta norm, let \(A\) have rank \(n\) and let the target \(B\)
have rank \(M>n\) and global Witt index \(q\). In the norm formula
write \(\chi\) for the character inducing the normalized Siegel
section. It is \(\chi_V\) in Gan--Qiu--Takeda and is the inverse of
the character denoted \(\chi_B\) in the release's Weil datum. Its
restriction to \(\mathbb A_F^\times\) is \(\eta_{E/F}^{M}\).
The doubled source is
\(A^\square=A\oplus(-A)\). Its Siegel parabolic stabilizes the
diagonal isotropic subspace of dimension \(n\). For a Schwartz test
\(\Phi\) in the doubled Weil representation \(\omega^\square\) that
is a polynomial Fock vector at infinity, let \(\Phi_s\) be the
standard Siegel section with value
\(\Phi_{s_0}(g)=(\omega^\square(g)\Phi)(0)\) at \(s_0=(M-n)/2\),
in the normalized induction by \(\chi\circ\det\,|\det|^s\).
If \(M-q>n\), Yamana's convergent Siegel--Weil formula gives
\[
 E_{\rm Sieg}(g,s_0;\Phi_{s_0})
 = I(g,\Phi),\qquad
 I(g,\Phi)=\int_{[U(B)]}\theta^\square(g,h;\Phi)\,dh_1 .
\]
Here \(dh_1\) gives the quotient volume one. The theta integral
converges absolutely and the Eisenstein series is holomorphic at
\(s_0\). For \(M>2n\) the defining Eisenstein series itself converges
absolutely there.
\cite[Theorem~2 and the convergence statements preceding it]{tr3yamana}
The column tower has definite targets. For the upward adjoints
of the row maps, the rank increase is \(t\), which exceeds the
negative index at each real place of the target and hence its
global Witt index.
For an anisotropic target and \(s_0>0\), Gan--Qiu--Takeda's
Theorem~7.1 gives the same identity. The factor in Tamagawa
conventions is \(\tau(U(B))/[E:F]=1\) times \(I(g,\Phi)\).
\cite[\S3.1 and Theorem~7.1]{gqt}

Let \(f_1,f_2\) be cusp forms in an irreducible cuspidal representation
\(\pi\) of \(U(A)(\mathbb A_F)\). In the bilinear convention the lift
with test \(\varphi_i\) is
\(\theta_i(h)=\int_{[U(A)]}\theta(g,h;\varphi_i)f_i(g)\,dg\).
The doubling seesaw uses the doubled test
\(\Phi=\delta(\varphi_1\otimes\overline{\varphi_2})\),
where \(\delta\) intertwines the doubled Weil representation on the
two copies of \(U(A)\) with the tensor product of the original
representation and its dual twisted by \(\chi\circ\det\), and
\(\delta(\varphi_1\otimes\overline{\varphi_2})(0)
=\langle\varphi_2,\varphi_1\rangle\).
In Tamagawa measures, in the convergent range where the lifts have
Petersson pairings and the seesaw integrals are absolutely convergent,
\[
 \langle\theta_2,\theta_1\rangle
   =[E:F]\,Z(s_0),\qquad
 Z(s)=\int_{[U(A)]^2} E_{\rm Sieg}((g_1,g_2),s;\Phi_s)\,
 f_1(g_1)\overline{f_2(g_2)}\,
 \chi(\det g_2)^{-1}\,dg_1\,dg_2 .
\]
The factor \([E:F]\) converts the volume-one theta integral to
Tamagawa measure. In \(n<M\leq 2n\), \(M>n+q\), Gan--Qiu--Takeda's
Propositions~11.1(ii) and~11.2 give this identity when the lower
lifts of \(\pi\) in the target Witt tower vanish, so the lift under
consideration is cuspidal. The lower-lift condition is vacuous for
an anisotropic target.
\cite[\S\S11.3--11.5, Propositions~11.1(ii) and 11.2]{gqt}
For the selected chains the release uses the absolutely convergent
seesaw with the Siegel--Weil value above. Its doubling integrand
uses the conjugates of the bilinear inputs as its two arguments, so
with inducing character \(\chi_B^{-1}\) the scalar on the second
copy is \(\chi_B(\det g_2)\). In the prescribed measures the target
volume is the actual quotient volume. The pairing uses \(L^2\)
bounds on both theta kernels and rapid decay of the source cusp forms.
\releasecite{Development/AbelianVarietyTheory/Proofs/DecoratedLatticeProofs/GlobalFormFrameLattice.lean}{109}{square-integrable kernels and Fubini for the norm}

With factorized vectors, pairings and measures the doubling integral
unfolds, initially for \(\operatorname{Re}s\) sufficiently large, to
\[
 Z(s)=\prod_\nu Z_\nu(s),\qquad
 Z_\nu(s)=\int_{U(A)(F_\nu)}\Phi_{\nu,s}(g,1)\,
     \langle f_{2,\nu},\pi_\nu(g) f_{1,\nu}\rangle\,dg .
\]
The local factors have meromorphic continuation. For unramified data
the spherical pair has pairing one, the spherical section has value
one on the hyperspecial compact, and the compact has volume one.
In those normalizations
\(Z_\nu(s)=L_\nu(s+\tfrac12,\pi_\nu\times\chi_\nu)/
d_\nu(s,\chi_\nu)\).
\cite[\S\S11.5--11.6]{gqt}
For these spherical data at split and inert places outside the
exceptional primes of Eischen--Harris--Li--Skinner's calculation,
the denominator is
\[
 d_\nu(s,\chi)=\prod_{i=0}^{n-1}
      L_\nu(2s+n-i,\chi|_{F_\nu^\times}\eta_{E/F,\nu}^{\,i}).
\]
Here \(\eta_{E/F,\nu}\) is the local quadratic character, trivial
at split places, and the numerator is the standard local factor
of the base change of \(\pi_\nu\) twisted by \(\chi_\nu\).
\cite[Hypotheses~1.4.4 and \S4.2.1, (53)]{p2eischen-doubling}
Since \(\chi_\nu|_{F_\nu^\times}=\eta_{E/F,\nu}^{M}\), the characters
in the denominator take values of absolute value one on a
uniformizer. Every argument \(2s+n-i\) is positive at real \(s>0\),
so the local denominators are finite and nonzero there.
With \(S_0\) containing the archimedean places, the places over the
auxiliary prime and all other places where the data fall outside
this unramified calculation, the meromorphic expression for \(Z(s)\)
is
\[
 \frac{L^{S_0}(s+\tfrac12,\pi\times\chi)}
      {\prod_{\nu\notin S_0}d_\nu(s,\chi_\nu)}
 \prod_{\nu\in S_0}Z_\nu(s).
\]
The norm uses the value of the meromorphic expression as a whole
at the holomorphic point \(s_0\), with the prescribed measure conversion.

For the real local criterion, take \(A\) skew-Hermitian and \(B\)
Hermitian, with \(M\geq n\), and use the Weil representation of the
pair with additive character \(x\mapsto e^{-2\pi i x}\) and unitary
splitting characters of exponents congruent to the two ranks.
Given an irreducible tempered \(\tau_\infty\) of \(U(A)(\mathbb R)\),
Ichino's local form
\[
 (\varphi_1,\varphi_2,u_1,u_2)\ \longmapsto\
 \int_{U(A)(\mathbb R)}\!
  \langle\varphi_2,\omega(g)\varphi_1\rangle\,
  \overline{\langle u_2,\tau_\infty(g)u_1\rangle}\,dg
\]
is absolutely convergent on Schwartz tests and smooth vectors.
It is nonzero as a form exactly when the small theta lift of
\(\tau_\infty\) to \(U(B)(\mathbb R)\) is nonzero. One can test
nonvanishing on polynomial Fock vectors and \(K\)-finite vectors.
\cite[\S4.1, Proposition~7.1 and Remark~7.2]{p2ichino-real-theta}
For the bilinear global inputs \(f_i\), the local form uses
\(\tau_\infty=\pi_\infty^\vee\) and \(u_i\) dual to \(f_{i,\infty}\)
under the invariant Hermitian identification with the conjugate
representation. On the first copy of the source group the doubled
section at \(s_0\) has value
\(\langle\varphi_2,\omega(g)\varphi_1\rangle\) in the corresponding
local normalization. The release proves the Hecke norm formula
using the split unramified formula, its inert calculation and
the convergent seesaw.
\releasecite{Development/AutomorphicArithmetic/ThetaCorrespondence/ThetaPairLifts/LocalTypeDictionary/HeckeFormOfDoubling.lean}{19}{Assembly of the Hecke norm formula}.
The
\releasecite{Development/AutomorphicArithmetic/ThetaCorrespondence/ThetaPairLifts/LocalTypeDictionary/OntoAllOfPrintsAndHecke.lean}{28}{surjectivity along the prescribed chains}
uses the norm in the specified quotient measures, the unramified
factors and nonzero theta maps on the prescribed types and levels.
With that surjectivity and the selected real, twisting, level and
geometric hypotheses, the release uses rapid decay of cusp forms,
discrete series, Harish--Chandra modules and strong approximation
for the special unitary groups to
\releasecite{Development/AutomorphicArithmetic/ThetaCorrespondence/ThetaPairLifts/ThirdPlaceSeries/FamAllQOfAllTrue.lean}{38}{choose the compatible top and bottom models}.

In these models a relation between top periods is carried to the
bottom in two stages. First the seesaw comparison at the last two
levels uses the codimension-one top pairing. Then each two-level
step uses a Bessel or Fourier--Jacobi coefficient of the corresponding
lower row. Composing the two Weil kernels and unfolding the quotient
integrals gives the lower coefficient at the prescribed translate.
The ramified column conditions identify the needed coefficients and
rapid decay justifies the integrations. For prescribed real tests
and translates in a two-level step, one matching real test works for
all tuples. The kernel expansion uses finitely supported frame
coefficients chosen independently of the tuple, with the tuple
dependence in the normalization and the prescribed vectors. These
are the quantifiers in the
\releasecite{Development/AbelianVarietyTheory/Definitions/LatticeFrames/ValLastX.lean}{186}{test identities}.
Induction transports any finite vanishing linear combination of the
normalized top periods to the corresponding combination of bottom
coefficients at every permitted frame and real vector. This is the
\releasecite{Development/AbelianVarietyTheory/Proofs/DecoratedLatticeProofs/SeesawTransportRelationsOfFamilyLaws.lean}{25}{comparison by the last two levels and successive two-level steps}.

For each word \(w\), combine the invertible column change from the
theta comparisons with the isometry identifying the chosen bottom
columns with the columns of the Poincar\'e series. The composite
multiplies norms by a constant \(r_{\mathrm{col},w}>0\). Dividing by
\(r_{\mathrm{col},w}\) gives a unitary \(U_w\) on \(\mathcal M_w\) and
divides both bottom pairings by the same constant. For a permitted full
bottom row \(T\), choose a preimage under the base chain and let
\(p_{w,T}(\phi_0)\) be its normalized top period on the base column
\(\phi_0\). Let \(c_{\mathrm{top}}(w,\boldsymbol\mu)\) and
\(c_{\mathrm{bot}}(w,\boldsymbol\mu)\) be the products of the upper and
lower factors in that comparison. The upper product includes the two
column factors and the real comparison factor from each two-level step,
and the factor for the top pairing. The lower product contains one
factor for each downward step. The scalar in the period relation is
\[
 d_{w,\boldsymbol\mu}
   =s_w\frac{c_{\mathrm{top}}(w,\boldsymbol\mu)}
              {c_{\mathrm{bot}}(w,\boldsymbol\mu)}\ne0,
\]
where \(s_w\ne0\) comes from comparing the two column lifts and is
independent of the tuple. For a tested row \(F_m\) in the projected
top space, the relation
\[
 \alpha\,p_{w,T}(\phi_0)
   +\alpha'\,d_{w,\boldsymbol\mu} b_{w,\boldsymbol\mu}(\phi,F_m)=0
\]
gives a vanishing combination of the two bottom pairings with
coefficients \(\alpha,\alpha'\). The first uses \(T\) and \(U_w\phi_0\),
and the second uses the downward image of \(F_m\) and \(U_w\phi\).
The release obtains this implication first for the chosen
top coefficient coordinates, including all choices of the two
column inputs. The rule for \(d_{w,\boldsymbol\mu}\) follows by
\releasecite{Development/AutomorphicArithmetic/ThetaCorrespondence/ThetaPairLifts/ThirdPlaceSeries/RelOfLinks.lean}{22}{dividing by the nonzero lower factors}.

\teamanchor{proof.common-central-translate}
Let \(t_0\) be the fixed finite element matching the bottom row level
and character with those of \(\mathcal H\), and let \(J_{\rm row}\)
be the common bottom row level subgroup in these coordinates. The
construction starts with a reference placed and framed Bessel datum.
Comparing its translate with \(t_0\) and absorbing a central adjustment
in the embedded column group places it in \(zt_0J_{\rm row}\), for
one finite central element \(z\in U(Y_0)(\mathbb A_F)\).
The
\releasecite{Development/AutomorphicArithmetic/ThetaCorrespondence/ThetaPairLifts/ThirdPlaceSeries/MatchingCoset.lean}{66}{pairwise coset comparison}
then puts every word in that same coset. This gives the
\releasecite{Development/AutomorphicArithmetic/ThetaCorrespondence/ThetaPairLifts/ThirdPlaceSeries/OneCentralPin.lean}{41}{common placement}.
Matching at \(t_0\) also gives matching at \(zt_0\), since
conjugation by \(z\) fixes the level subgroup and its character.
Right translation by \(z\) and its inverse preserve the Petersson norm
and distinguished row conditions, giving a unitary \(V\) on
\(E_{\rm row}\).

Evaluation on the chosen real unit vector at \(zt_0\), followed by
complex conjugation, gives a conjugate-linear bijection from the full
permitted bottom multiplicity space onto \(\mathcal H\). This uses the
matched level and character, the real discrete-series identifications
and the real gap hypothesis. Write \(\operatorname{ev}_z\) for this
map with the bottom pairing normalization. For a full row \(T\)
and a ramified column \(\psi\) the pairing is
\[
 \langle\operatorname{ev}_z(T),\mathcal P_w\jmath_w\psi\rangle .
\]
For the downward image of a projected top row, the matching to the
distinguished space represents this pairing by
\(\mathcal E_w(\psi,V\iota_{\boldsymbol\mu}F_m)\), up to a fixed
nonzero conversion factor in the row normalization. We use the
same conversion on both terms of the bottom relation. The release
compares these pairings in the
\releasecite{Development/AutomorphicArithmetic/ThetaCorrespondence/ThetaPairLifts/ThirdPlaceSeries/MiStar2AtRealisation.lean}{84}{chosen row and column coordinates}.

\teamanchor{proof.nonzero-anchor}
Use the full-space diagonal estimate \eqref{eq:trace-single-diagonal}
and the ramified dimension bound \eqref{eq:trace-ramification} for
these same full operators and inclusions. On \(\ker\mathcal P_w\),
the operator \(K_{ww}-I\) acts by \(-I\). Extending an orthonormal
basis of this kernel to one of \(H_w^{\rm full}\) gives
\[
 \dim\ker\mathcal P_w
 \leq \|K_{ww}-I\|_{\rm HS}^2
 \leq e_0d_w
 <r_w .
\]
Indeed \(r_w\geq(1-2f/b)d_w\), \(d_w>0\) and
\(e_0<1-2f/b\) at sufficiently large \(x\), since \(f\leq x^5\),
\(b=\exp(x^\varepsilon)\) and
\(\sqrt{e_0}\leq\exp(-x^{\varepsilon/8})\).
Choose \(\psi_*\in\mathcal M_w\) with
\(y=\mathcal P_w\jmath_w\psi_*\ne0\). Take the full bottom row
\(T_*\) with \(\operatorname{ev}_z(T_*)=y\). Its pairing with
\(\psi_*\) is \(b_{*,w}=\langle y,y\rangle\ne0\). Both the choice of
\(T_*\) and its pairing use the full bottom space. See the
\releasecite{Development/AutomorphicArithmetic/ThetaCorrespondence/ThetaPairLifts/ThirdPlaceSeries/AnchorPairs.lean}{48}{construction of the reference pair}.

Use the input \(\phi_0=U_w^{-1}\psi_*\) and the base lift of this
same \(T_*\) to form the reference top period \(p_{*,w}=p_{w,T_*}(\phi_0)\).
If \(p_{*,w}=0\), the relation with \((\alpha,\alpha')=(1,0)\) would
give \(b_{*,w}=0\). Thus \(p_{*,w}\ne0\). In the two-term relation take
\(\alpha=d_{w,\boldsymbol\mu}b_{w,\boldsymbol\mu}(\phi,F_m)\)
and \(\alpha'=-p_{*,w}\). The resulting bottom equality is
\begin{equation}
\label{eq:period-reference}
 d_{w,\boldsymbol\mu}\,b_{w,\boldsymbol\mu}(\phi,F_m)\,b_{*,w}
 =p_{*,w}\,\mathcal E_w(U_w\phi,V\iota_{\boldsymbol\mu}F_m).
\end{equation}
Here the fixed, nonzero conversion from the bottom row normalization
is included in \(p_{*,w}\). The equality holds on the compatible
coefficient coordinates, which span the projected complex top row
spaces. It therefore extends to those spaces by bilinearity.

\teamanchor{proof.common-period-identity}
For the base tuple, \eqref{eq:period-reference} makes the top period
vanish on \(\ker\iota_{\boldsymbol\mu_0}\) in its projected domain.
Evaluate the base period on any preimage of \(F\in E_{\rm row}\)
to define \(B_w(\phi,F)\). Surjectivity gives a preimage for every
\(F\), and kernel vanishing makes its choice immaterial. With
\(a_w=p_{*,w}/(b_{*,w}d_{w,\boldsymbol\mu_0})\)
we obtain
\begin{equation}
\label{eq:period-bessel}
 B_w(\phi,F)
   = a_w\big\langle s(VF),\mathcal P_w\jmath_wU_w\phi
          \big\rangle_{\mathcal H},\qquad a_w\ne0 .
\end{equation}
The same \(V\) applies to every word, with \(a_w\) and \(U_w\) chosen
for each \(w\). At another tuple the identity on its projected
rows reads
\begin{equation}
\label{eq:period-tuple}
 b_{w,\boldsymbol\mu}(\phi,F_m)
 =\frac{d_{w,\boldsymbol\mu_0}}{d_{w,\boldsymbol\mu}}\,
     B_w(\phi,\iota_{\boldsymbol\mu}F_m).
\end{equation}

\subsection{The coefficient-field map and finite-place tests}

\teamanchor{proof.period-row-models}%
\subsubsection{The definite coefficient models}
Write \(K\) for the coefficient field, \(\mathcal O_K\) for its
integers and \(\iota_K\) for the chosen embedding into \(\mathbb C\).
Use the compact algebraic coefficient types at the top and on the
definite column tower. In the chosen matrix coordinates these are
polynomial representations with determinant twists, defined over \(K\)
using the specified embeddings of \(E\). At a definite vertex with
level \(J\), choose representatives \(c\) for
\(G(F)\backslash G(\mathbb A_{F,f})/J\). The stabilizer
\(\Gamma_c=G(F)\cap cJc^{-1}\) is finite. Evaluation at the representatives
identifies coefficient-valued forms with
\(\bigoplus_c W_c^{\Gamma_c}\), where \(W_c\) is the algebraic
coefficient space and the stabilizer action includes the level
character. In fact the value at any other point is determined by
covariance, and the stabilizer condition makes its definition
independent of the presentation in that double coset.
This description commutes with extension of scalars because the
invariants are kernels of finite systems of equations over \(K\).
The integral model consists of the vectors whose value at each \(c\)
belongs to its prescribed value lattice. The equations impose a
lattice in the invariant space: it is a submodule of the direct
sum of the value lattices and spans after clearing denominators.

For the principal-level space let \(\Gamma\) now denote the finite
group of level changes. For the chosen type characters \(\chi\) of
\(\Gamma\) and the character \(\alpha\) of a subgroup \(P\), the good
projector is the product of
\[
 \sum_{\chi}\frac{\chi(1)}{|\Gamma|}
                  \sum_{g\in\Gamma}\chi(g^{-1})\,g,
 \qquad
 \frac1{|P|}\sum_{g\in P}\alpha(g)^{-1}g .
\]
The first operator, by character orthogonality, selects the specified
isotypes. The second is the projector on the \(\alpha\)-equivariant
level subspace by averaging. They commute since the first is central.
With the class inner product the level actions are unitary, so the
projectors are orthogonal. The class weights are \(1/|\Gamma_c|\)
when the level and compact real group have volume one. For the
larger level use \(|P|^{-1}\) times the principal class inner product:
the orbit of a larger-level class contributes \(|P|\) times its
mass in the principal class sum by orbit--stabilizer.
The level and type matrices lie in \(K\) with the stated character
values. Their action can be read at every realigned class
representative. If a representative translates to \(cg=\gamma d k\),
covariance reads the result as the value at \(d\) acted on by the
coefficient representation of \(\gamma\) and the level character of
\(k\). The stabilizer condition makes this independent of the
decomposition. This is also the finite matrix action on the class
values. Thus for the prescribed transported frames and coefficient
lattices, the projector image \(A_{\boldsymbol\mu,K}\) complexifies
to the actual projected top space \(\mathcal V^0_{\boldsymbol\mu}\).
In this assertion the determinant twists, level characters and
coefficient lattices are those on the same definite groups as in
\eqref{eq:period-top}.
\releasecite{Formalism/LocalToGlobal/Divisibility/DivisionCompatibleToSub2Indef2.lean}{210}{Row model identities}

Pull back the value lattice to the image of the projector by
intersection. Equivalently, for a basis matrix \(E_{\boldsymbol\mu}\)
of that image and ambient class value lattice \(L_{\boldsymbol\mu}\),
its lattice in the basis coordinates is
\[
 M_{\boldsymbol\mu}=
     \{y:E_{\boldsymbol\mu}y\in L_{\boldsymbol\mu}\}.
\]
Indeed every integral vector in the image has unique coordinates
in \(E_{\boldsymbol\mu}\), and their integrality condition is
exactly the displayed one. This pullback uses the original value
lattice even at primes dividing projector denominators.

The same class model gives rational projections
onto the fitted auxiliary coordinates at the base level. Let \(E_J\)
be the inclusion of the level coordinates in the full class
coordinates and let \(E_J^\flat\) be its inverse on its image,
extended by the orthogonal projection there. Let \(V\) include
the good eigenspaces at that level, with inverse \(V^\flat\) on
its image, again extended by orthogonal projection. Let \(Z\)
include the selected complementary subspace of the good space
and let \(P_Z\) be the orthogonal coordinate map to this
complement, so \(P_Z Z=I\). The stage projection from the level
model is \(P_ZV^\flat E_J\). Its right inverse is
\(E_J^\flat VZ\), since \(E_J E_J^\flat VZ=VZ\) and
\(V^\flat V=I\). All matrices are over \(K\) for the class
Hermitian form. The rational surjection maps the level value
lattice onto a full lattice in these stage coordinates.
The arithmetic fitted subspace is the kernel complement
identified for the auxiliary theta map in
\S\ref{sec:period-fitted-stage}.
\releasecite{Formalism/LocalToGlobal/Divisibility/LiftLatLemmas.lean}{284}{Right inverse of the stage projection}.

\teamanchor{proof.period-selected-row-cuts}
\subsubsection{Descent on the selected row spaces}\label{sec:period-selected-row-cuts}
In this argument the ranks are \(n_j=r+jt\), with
\(t\) odd and \(t\ge r+5\). The final row group is
\(U(X')\), identified by the prescribed isometry with
\(U(Y_m)\). All preceding row groups are \(U(Y_j)\).
For a tuple \(\boldsymbol\mu\), write
\(\mu_{\boldsymbol\mu,i}=\mu_{0,i}\xi_{\boldsymbol\mu}\)
for its splitting characters before the parity adjustments.
The finite-order character \(\xi_{\boldsymbol\mu}\) is trivial
at infinity and on the base field. Write \(S\) for the
designated split places of the residue group and \(Q\)
for the additional split places of the tuple. At \(Q\) the
character \(\xi_{\boldsymbol\mu}\) is tamely ramified and
the base splitting characters are unramified. In the
orientation \(\operatorname{Hom}(Y_j,Y_{j+1})\) let
\((\chi_S,\chi_T)\) denote the two characters of this arrow.
At step \(i=j+1\) they are
\[
 (\mu_{\boldsymbol\mu,i}\eta^{-1},\mu_{\boldsymbol\mu,i})
 \quad(i\text{ odd}),\qquad
 (\mu_{\boldsymbol\mu,i},\mu_{\boldsymbol\mu,i}\eta^{-1})
 \quad(i\text{ even}).
\]
Here \(\eta\) is unramified at all finite places and supplies
the half-integral change at infinity. Use the selected
half of the split algebra to read all the matrices. The
specified local matrices below are in the transported
lattice frames. In particular the same description uses
the top frame on \(X'\) and the final datum on
\(\operatorname{Hom}(Y_{m-1},X')\).
\releasecite{Formalism/LocalToGlobal/IndefiniteReal/HasInfExpF.lean}{13}{Tuple twist and parity}
\releasecite{Formalism/LocalToGlobal/IndefiniteReal/LinksDefs.lean}{424}{Levels and ramified tuple characters}
\releasecite{Development/AbelianVarietyTheory/Definitions/LatticeFrames/NotBadPrimeBase.lean}{197}{Actual successive and final finite tests}.

\paragraph{Finite type and measure.}
Use the chosen subgroup \(P\) with unitary character \(\alpha\)
in the common residue group
\(\Gamma=\prod_{v\in S}\mathrm{GL}_r(k_v)\). On each stage the level at
\(S\) is the inverse image of the upper block parabolic
for \((r,t,\ldots,t)\), with first block in \(P\).
The character \(\lambda_j\) is \(\alpha\) on the first
block and the respective \(\chi_1,\ldots,\chi_j\) on
the determinants of the later blocks. Write
\(a=n_j\), \(b=a+t\). At \(S\) the finite test is
supported on integral matrices reducing to \(\binom{u}{0}\),
with \(u\) in this lower level. Its weight is
\[
 \lambda_j(u)\chi_{j+1}(\det u)^{-1}.
 \tag{F}
\]
Here \(\chi_{j+1}\) on the residue field is the inverse
of the restriction of \(\mu_{0,j+1}\) on the chosen
half. The tuple twist is unramified at these places.
The base test off \(S\cup Q\) is the unit self-dual
vacuum. At a nonsplit finite place the given characters
are unramified and this is the hyperspecial vacuum.

In the split Schr\"odinger model the source pullback is
\(X\mapsto Xg\), with unit factor
\(\chi_{j+1}(\det g)\) on integral units at \(S\).
Thus (F) transforms on the source by \(\lambda_j\).
For \(k\) on the target, the pullback is
\(X\mapsto k^{-1}X\), with unit factor
\(\chi_{j+1}(\det k)^{-1}\). Reducing
\(k=\left(\begin{smallmatrix}u_0&*\\0&d_0\end{smallmatrix}\right)\),
the change in the test is
\(\lambda_j(u_0)^{-1}\chi_{j+1}(\det u_0)\).
Multiplying the unit factor gives
\(\lambda_j(u_0)^{-1}\chi_{j+1}(\det d_0)^{-1}
=\lambda_{j+1}(k)^{-1}\).
To see exact coefficient support, suppose \(X,Xg\) are
both in the support. Write \(A,A'\) for their top
\(a\)-row square blocks. Both are invertible integral,
and \(g=A^{-1}A'\) reduces to the lower level.
Conversely on that level the coefficient is
\(\lambda_j(g)\|\Phi_v\|^2\).
The doubled height is one there. With the conjugate
Petersson input the two \(\lambda_j\) factors cancel.
The semilocal contribution at \(S\) is exactly the
level volume times the product of the norms of the
actual local tests.
\releasecite{Formalism/LocalToGlobal/IndefiniteReal/DatumVec.lean}{51}{The residue weight}
\releasecite{Development/AbelianVarietyTheory/Proofs/DecoratedLatticeProofs/ThetaChainStepVectorTwoLevelLaw.lean}{20}{Target character law on the placed test}.

Use the principal first-block level \(P=\{1\}\).
For each \(\gamma\in\Gamma\) lift it to the same
diagonal first block on the two stages. Acting
simultaneously on both legs conjugates the first block
in (F) and leaves all other diagonal determinants fixed.
The two splitting determinant factors cancel. The test
is therefore fixed by this simultaneous action. The
lift and its Petersson adjoint intertwine the two
actions of \(\Gamma\), by the change of variable in the
quotient integral. Both consequently commute with the
common good central projector. The \(P,\alpha\) test
is the sum of the principal tests translated on the
lower side by these lifts, with coefficient
\(\alpha(\gamma)^{-1}\), \(\gamma\in P\).
Indeed the term indexed by \(\gamma\) has first block
\(\gamma^{-1}\) and the source splitting multiplier
is \(\chi_{j+1}(\det\gamma)\), as in (F).
On the \(\alpha\)-isotype the summed integral is
\(|P|\) times the principal integral before Haar
conversion. The smaller-level Haar is \(|P|^{-1}\) times
the principal one, because the smaller level has
index \(|P|\) over the principal level. This cancellation
applies on every stage where this first-block level
changes, including the top restriction.

If \(d\nu_j\) is the underlying row measure in the
period definitions, the normalized stage measure is
\[
 d\nu_j^{\boldsymbol\mu}
   =m_Y(\boldsymbol\mu,j)\,d\nu_j,\qquad
 m_Y(\boldsymbol\mu,j)
   =\prod_{v\in Q}[\mathrm{GL}_{n_j}(\mathcal O_v):
                 J_{(r,jt),v}].
 \tag{G}
\]
At \(j=0\) the suffix is empty and this factor is one,
as it is for all stages of the base tuple.
Multiplying the product of the two form evaluations
by this row measure gives the Hilbert metric of the
corresponding stage. The downward map integrates using
\(d\nu_{j+1}^{\boldsymbol\mu}\). Thus the adjoint
integrates the conjugate kernel using
\(d\nu_j^{\boldsymbol\mu}\), by Fubini with the
product of these two measures. This gives precisely
the powers of \(m_Y\) in each direction, including
the measure after the final isometry to \(X'\).
\releasecite{Formalism/LocalRingCardinality/DataTheta.lean}{727}{Stage metrics and row maps}
\releasecite{Development/AutomorphicArithmetic/ThetaCorrespondence/ThetaPairLifts/ThirdPlaceSeries/InverseLevelCardinality.lean}{24}{Small-level Haar factor}.

\paragraph{Doubling at a tuple prime.}
Here the full-rank support argument at \(S\) must be
replaced by an integral on the suffix of the matrix.
Fix \(v\in Q\), omit \(v\) from the notation, and put
\(d=jt\). In the selected split coordinates the
finite vector is
\[
 f(X)={\bf1}_{\mathscr X}(X)\,
          \xi_{\boldsymbol\mu}(\det \overline{X}_{00}),
 \qquad
 \mathscr X=\left\{
 \begin{pmatrix} A & B \\ C & D\end{pmatrix}
 \text{ integral}: A\in\mathrm{GL}_r(\mathcal O),
                   C\equiv0\pmod\varpi\right\}.
 \tag{H}
\]
In (H) the first row and column blocks both have size \(r\).
The finite level is the parahoric of shape \((r,d)\)
with character \(\xi_{\boldsymbol\mu}^{-1}\) on the
suffix determinant. This follows by applying the
split Weil formula to (H). For instance the source
pullback contributes \(\xi(\det\bar g_{00})\), and the
scalar of the Weil action contributes
\(\xi(\det\bar g)^{-1}\).
The same calculation on the target gives the inverse
of its level character. This verifies level covariance
on both sides of the actual step, including its top version.
\releasecite{Formalism/LocalToGlobal/IndefiniteReal/DatumVec.lean}{68}{The additional-prime factor}
\releasecite{Formalism/LocalToGlobal/Divisibility/DefGroup.lean}{700}{The suffix determinant in the level character}.

Let \(\pi_j\) be the representation on the conjugate
automorphic input in the doubling integral. At \(j=0\)
it is spherical at \(v\). For a constituent obtained by
the first \(j\) steps, \(\pi_j\) is the normalized
parabolic induction of a spherical \(\mathrm{GL}_r\)
representation and a \(\xi\)-twist of a spherical
\(\mathrm{GL}_d\) representation. Indeed split type-II
Howe correspondence adds the determinant character
\(\chi_S\) on the new \(\mathrm{GL}_t\) block and
twists the old representation by \(\chi_S/\chi_T\).
The latter is unramified at \(v\). The successive
suffix factors before taking the common twist out are
unitary determinant characters. Normalized parabolic
induction of irreducible unitary representations of
split \(\mathrm{GL}\) is irreducible. The suffix
spherical representation at step \(j\) consequently
has the Satake strings from the earlier steps.
\href{https://www.numdam.org/item/ASENS_2008_4_41_5_717_0/}{M\'inguez, Corollary~6.3},
\href{https://www.numdam.org/articles/10.24033/asens.1510/}{Tadi\'c, \S0 and the unitary induction property}.

Include the scalar \(\chi_T(\det g)^{-1}\) from the Weil
action in the group coefficient, so now use
\(\pi'_j=\pi_j\otimes\chi_T^{-1}\). It is the
normalized parabolic induction of
\(\tau\boxtimes\sigma\) for the parabolic \(P_{r,d}\)
where \(\tau\) is a \(\xi^{-1}\)-twist of a spherical
representation and \(\sigma\) is spherical. Denote
by \(K=\mathrm{GL}_{r+d}(\mathcal O)\) and
\(J=J_{(r,d)}\). The \(J\)-line needed in the group
coefficient now has character
\(\xi^{-1}(\det\bar g_{00})\). Its dimension is one:
take a principal-series presentation surjecting onto
\(\pi'_j\) with the \(r\) characters of inertia type
\(\xi^{-1}\) first and the remaining \(d\) unramified.
Taking the \(J\)-character is exact by averaging.
In the induction, the double cosets for the upper
Borel and \(J\) are the subsets of \(r\) diagonal
positions assigned to the first block of \(J\).
On the intersection with the diagonal units the
character can match only if this subset is precisely
the \(r\) positions of type \(\xi^{-1}\), since
\(\xi\) is nontrivial on the residue units. That coset
contributes at most one dimension. In the parabolic
induced model a nonzero vector with this character
is supported on \(P_{r,d}J\), with its value at one
the product of the two spherical vectors after
their character twists. The equivariance prescriptions
agree on \(P_{r,d}\cap J\), so this defines the
actual line. Normalize this value to have norm one
in the inducing representation.

With \(H(g)\) the normalized doubling height, the local
integral, up to a positive scalar, is
\[
 \int_{\mathrm{GL}_{r+d}}
    |\det g|^{b/2}H(g)^{s-t/2}\,
          \langle\pi_j'(g)v,v\rangle
       \int f(Xg)\overline{f(X)}\,dX\,dg .
 \tag{I}
\]
When \(d=0\), the full-rank block in (H) forces the
support of the coefficient to be \(K\), so (I) is a
nonzero volume times two squared norms. Suppose \(d>0\).
In the compact induced realization, write the group
coefficient as
\(\int_K\langle v(kg),v(k)\rangle\,dk\).
The compact variable is supported on \(J\). Put
\(g'=kg\) and \(X'=Xk^{-1}\) in (I), so that
\(Xg=X'g'\). The covariances of \(f(X'k)\)
and \(v(k)\) cancel in the two conjugate slots.
This leaves a positive compact volume
times the same integral with
\(\langle v(g),v(1)\rangle\) in the integrand. It now
only integrates \(g\) in \(P_{r,d}J\).

For \(X\in\mathscr X\), eliminate its first-column
blocks by the target unit matrix
\[
 \begin{pmatrix} A^{-1}&0\\-CA^{-1}&I\end{pmatrix}.
\]
Then eliminate its first-row right block by the
integral upper unipotent matrix on its right. These
changes can be made in the inner \(g\)-integral: the
target change multiplies \(f(Xg)\) and \(f(X)\)
by the same unit scalar. The source change is in
the upper unipotent of \(P_{r,d}\) and is absorbed by
left translation in \(g\), using the left covariance
of the induced vector and the trivial character on
this upper unipotent. Doubling height is bi-invariant
under \(K\), by the lattice normalization on the
doubled source. The remaining rectangular block
\(D_X=D-CA^{-1}B\) ranges over all integral matrices
of size \((d+t)\times d\); the change from \(D\)
is an integral translation and preserves Haar.
In \(P_{r,d}J\) use \(g=m_0u k\), with
\(m_0=\operatorname{diag}(A_0,D_0)\), \(u\) upper
unipotent, and \(k\in J\). Its Haar integral is in
the order \(dm_0\,du\,dk\) up to a positive constant,
since writing the parabolic element as \(m_0u\)
absorbs the left modular density. In the test after
the eliminations, \(A_0\) is an integral unit and
\(u\) is integral, while \(D_XD_0\) must be integral.
The \(k\)-character cancels against the character
of \(f\), and the \(A_0\)-character cancels against
that of \(\tau\). The normalized inducing factor
is \(\delta_{P_{r,d}}(m_0)^{1/2}=|\det D_0|^{-r/2}\).
Multiplying it by \(|\det g|^{b/2}\) gives
\(|\det D_0|^{(d+t)/2}\). The doubled height of
\(\operatorname{diag}(I_r,D_0)\) is the doubled
height of \(D_0\), by the orthogonal product of
the doubled lattice and the identity block.

The integral that remains is exactly the spherical
split doubling integral for the pair
\(\mathrm{GL}_d,\mathrm{GL}_{d+t}\) with group coefficient
of \(\sigma\) and integral-matrix vacuum. The calculation
can first be made in the right half-plane of absolute
convergence. At \(s=t/2\) the local numerator is the
Euler factor of \(\sigma\) at \((t+1)/2\).
Every suffix Satake entry has modulus at most
\(q^{(t-1)/2}\), so the numerator has no pole or
zero there. The spherical denominator is a product
of nonzero finite Euler factors at positive arguments.
For absolute convergence use the Cartan index
\(\mathbf k=(k_1,\ldots,k_d)\). The oscillator
coefficient is \(q^{-(d+t)\sum_i|k_i|/2}\).
The volume of its double coset is
\(O(q^{2\langle\varrho,\mathbf k\rangle})\),
where \(\varrho\) is the half-sum of the positive
roots of \(\mathrm{GL}_d\). The spherical coefficient
is bounded by a fixed polynomial in \(\sum_i|k_i|\)
times
\(q^{-\langle\varrho,\mathbf k\rangle+
                 (t-1)\sum_i|k_i|/2}\).
This follows from the spherical-function formula,
including the limiting polynomial factors at
coincident Satake entries. \cite[Proposition~4.1 and Theorem~4.2]{p2casselman-spherical}
Since
\(\langle\varrho,\mathbf k\rangle\le
(d-1)\sum_i|k_i|/2\), their product has the
exponential bound \(q^{-\sum_i|k_i|}\)
at the critical point. The doubling height changes
this bound by at most a fixed exponential in
\(|\operatorname{Re}s-t/2|\sum_i|k_i|\), so it
still sums on a neighborhood of that point.
All omitted factors
in the reduction are the positive volumes of the
integral blocks. Hence the particular tested factor
at every tuple prime is nonzero.

\paragraph{Real factors, labels and passage to the cuts.}
For the actual \(Y_j\) parameters the center is
\(-e_j\), with the parity shifts of \(\chi_T\) and
\(\chi_S\) on the smaller and larger models. In the
convention for the local correspondence before
passing to the dual row, the centered smaller lists
on the determinant cover have
the form \((A,B;C,D)\), with the first and third
positive and the other two negative. Centering the
larger lists gives
\[
 (A,\ ((t-1)/2,\ldots,-(t-1)/2),D;\ C,B).
 \tag{J}
\]
The shift from \(-e_j\) to the two actual centers
is the half-step dictated by the corresponding rank
parities. The sign labeling flips the old entries
below the center and inserts the new block in the
positive list. By the prescribed row interlacing the
distance of any old entry from the center is at least
\((t+3)/2\). In computing the real gap, opposite-sign
entries lie in disjoint strips separated by two successive
\(b\)-endpoints. Subtracting the block length \(t\)
from the gap of these endpoints leaves a gap larger
than the cube of the larger row rank. These numeric
conditions are unchanged by the tuple twist. By Paul's
Theorem~2.7(b) the centered parameters in (J) are the
prescribed pair. The smaller centered parameter is
integral with no zero. It has first occurrence in this
tower by Ichino's Theorem~6.1 with
\(k_\pi=-1,t_{\mathrm{Ichino}}=(t-1)/2,l=0\);
the prior centered entries are outside the two near
sets. In the reverse direction one deletes the maximal
central string, as in Paul's Proposition~3.4(ii).
The global constant-term calculation gives cusp forms
in the designated real stages, since the lower real
members vanish. This applies with the definite final
signature on \(X'\).
\releasecite{Development/AbelianVarietyTheory/Proofs/TwistedOperators/ThetaChainStepCentredBlockParameters.lean}{65}{Centered signed lists and the gap}
\releasecite{Development/AbelianVarietyTheory/Proofs/DecoratedLatticeProofs/ThetaChainRealZetaNonvanishing.lean}{25}{Real exponents of the placed tests}.

Use Ichino's Proposition~7.1 on this pair, polarize
in the two oscillator slots and then the two source
slots, and use real Hermite density as in the core
norm argument. This gives a joint real polynomial and
compact-finite source vector at which the real norm
integral is nonzero. The actual finite vector already
fixed in (F) and (H) can be used with every real
polynomial. The same discrete-series gap and doubled
height estimate give holomorphy of this real integral
on \(\operatorname{Re}s>0\).

Off \(S\cup Q\), start from the unramified bottom label
\(B\) and form the usual theta transfer at each adjacent
step: multiply old entries by \(\chi_S/\chi_T\) and add
the fixed length-\(t\) string centered at \(\chi_S\).
These transfers are injective on labels, by cancellation
of the added multiset and reversal of the common twist.
At step \(i\), after twisting the label by
\(\chi_T^{-1}\) in the doubling numerator, the old
strings are centered at
\(\mu_{\boldsymbol\mu,h}/\mu_{\boldsymbol\mu,i}\)
for \(h<i\). This follows from the alternation of the
two parity shifts: the product of the intermediate
ratios is one or \(\eta^{\pm1}\) and cancels the
endpoint shift of the string. The initial rank-\(r\)
label is twisted by \(\mu_{\boldsymbol\mu,i}^{-1}\),
with unitary values. Thus the unramified numerator
at the critical point is the product of the initial
Euler function at \((t+1)/2\) and
\[
 \prod_{h<i}\prod_{b'=0}^{t-1}
 L_E^{S\cup Q}(t-b',
       \mu_{\boldsymbol\mu,h}/\mu_{\boldsymbol\mu,i}).
 \tag{K}
\]
The initial Euler product converges absolutely and is
nonzero by the elementary spherical bound
\(q^{r/2}\) and \(t\ge r+5\).
The factors of (K) at one are nonzero by Hecke
nonvanishing, since the quotient characters have
distinct real exponents. The other factors are in
absolute convergence. Omission of the finite sets
preserves these properties. At a tuple prime the
local transfer is also injective, either by the
multiset description of the split local correspondence
or by canceling the added suffix. The archimedean
types are fixed throughout.
\releasecite{Development/AbelianVarietyTheory/Proofs/DecoratedLatticeProofs/ThetaLiftChainHeckeParameterTwistNextTargetChar.lean}{20}{Cancellation of the parity factors}
\releasecite{Development/AbelianVarietyTheory/Proofs/DecoratedLatticeProofs/ThetaChainLValuesNonvanishing.lean}{62}{Continuation and the range of arguments}.

The target Witt index is at most its negative real
index, hence less than \(t\). The normalized
convergent-range norm formula therefore applies to
these pairs. With the chosen real witness, its
semilocal factor is a nonzero common scalar times
the norm square of a lower multiplicity vector of
one transferred label. At \(S\) this is the character
cancellation on the whole prescribed finite isotype.
At \(Q\) it follows from the one-dimensional line
and (I). Away from these places average between
spherical vectors and use that every spherical
operator acts by its scalar on a common label.
This verifies the product for arbitrary linear
combinations within the occurrence space, without
requiring the automorphic vector itself to be a
pure tensor. Together with (K) it proves
nonvanishing for every nonzero lower multiplicity
vector of a transferred label.

To pass from the tested functions to the adjoint on
multiplicities, write the real theta contraction as
one fixed Howe map on each placed pair. Its values
on varying polynomial and compact-finite upper
vector span the smaller compact-finite real module:
their span is nonzero and stable under its compact
group and Lie algebra, since differentiating acts
by quadratic oscillator operators on the polynomials.
Then Schur's lemma factors the test integral as this
real contraction and one map \(D_j\) on multiplicities.
With the measures (G), Fubini identifies the
upward test as the Petersson adjoint of this product.
A nonzero tested lift on any lower multiplicity
vector therefore implies its image under \(D_j^*\)
is nonzero. The normalized norm formula gives a
positive scalar norm square on each transferred
label for the fixed real witness. Dividing out the
norm square of its real image and polarizing shows
that \(D_jD_j^*\) is a positive scalar on the
corresponding lower label. The real witness can
depend on the arrow and label.

It remains necessary here to track the finite cuts.
An upward adjoint sends the full lower level character
to the actual upper level character by the covariance
just computed. At \(S\) it also intertwines the
common residue-group idempotent. For spherical
Hecke projections, its covariance is the Satake
transfer just computed, so it sends an eigenspace
of label \(B\) into the eigenspace of the transferred
label. These two families of projectors commute
because the spherical operators are at disjoint
places from the residue-group operators. The
factor calculation at \(Q\) applies on the actual
line of the corresponding transported representation.
In particular a cut formed by the good central
projector, the level-character projector and the
successive spectral labels obtained by transfer
contains each of these upward images. The adjoint
for the cut is the same map on these inputs.
A preimage of a vector with lower label is its
adjoint image divided by the positive scalar above.
Adding the finitely many labels gives surjectivity
on the cut. Iterating yields the common bottom
image for the normalized downward maps. Since
the bottom group and its level and cut are unchanged
as the tuple varies, this image is the same at
every tuple. For the full base image take
\(Q=\varnothing\) and retain every permitted
bottom label. The same induction applies before
the good projector.

For the analytic row descent write \(D_{0,j}=D_0\cdots D_{j-1}\).
Its kernel complement in the stage cut is the image of its adjoint,
and the inverse there is
\[
 D_{0,j}^*(D_{0,j}D_{0,j}^*)^{-1}.
\]
The middle operator acts on each selected bottom label by the
product of the positive stage scalars. This gives the analytic rank
and explicit preimages on the row chain.
\releasecite{Development/AutomorphicArithmetic/ThetaCorrespondence/ThetaPairLifts/ThirdPlaceSeries/LinkOntoOfOntoAll.lean}{47}{The full base chain}.

\teamanchor{proof.period-fitted-stage}
\subsubsection{The fitted auxiliary theta matrix at the base level}
\label{sec:period-fitted-stage}

For each \(0\le i<m\) and each selected tuple
\(\boldsymbol\mu\), consider the auxiliary pair
\(V_i,W_i\), with \(V_i\) definite of rank \(r+t\) and \(W_i\)
of rank \(r\). The record stage is \(j=i+1\). Write \(T^{\rm raw}_{i,\boldsymbol\mu}\)
for theta integration from the full definite class space at the
tuple level to the row multiplicity space, and \(T_{i,\boldsymbol\mu}\)
for its good projection. Both use the phase-one normalization.
The base level in this comparison is \((P_0,\psi_0)\).
Let \(\iota_K\) be the prescribed complex coefficient embedding.
The arithmetic realization used for taking adjoints and solving the
relative matrix is at the prescribed infinite place \(v_0\).
If its representative is conjugate to \(\iota_K\), the complex
source and the complex target, with their pairings and both group
actions, are simultaneously conjugated. The phase-one condition
makes the arithmetic lift in this realization exactly \(T_{i,\boldsymbol\mu}\)
with this scalar convention.
\releasecite{Formalism/LocalToGlobal/IndefiniteReal/ZeroPartEigenvalueConvention.lean}{120}{The distinguished-place realization and its phase}
\releasecite{Formalism/LocalToGlobal/IndefiniteReal/LinksDefs.lean}{684}{The two coefficient readings}
\releasecite{Development/AbelianVarietyTheory/Proofs/IwasawaDecomposition/LiftCuspFormSpaceThetaMapEqData.lean}{152}{Transport back to the complex model}.

\paragraph{The class evaluations in the two constructions.}
The definite source level \(J_{i,\boldsymbol\mu}\) is the principal
flag level for the tuple, including its character on the other
blocks. Its compact real factor has volume one, and the finite Haar
measure has volume one on \(J_{i,\boldsymbol\mu}\). At a refined
finite level \(J'\subset J_{i,\boldsymbol\mu}\) let \(h_c\) be the
prescribed double-coset representatives, with rational stabilizers
\(\Gamma_c=U(V_i)(F)\cap h_cJ'h_c^{-1}\), where the finite-level
condition is on the finite component. Lift the value of a class
form \(f\) at \(h_c\) to the coefficient polynomial \(x_c(f)\)
using the given semilinear coefficient map \(c_{i,\boldsymbol\mu}\)
in the placed frame. On a pure lowest-type tensor \(u\) and at
every row point \(g\), the evaluation is
\begin{equation}
 \label{eq:period-auxiliary-class-evaluation}
 T^{\rm raw}_{i,\boldsymbol\mu}(f)(u)(g)
 =\sum_c
   \frac{[J_{i,\boldsymbol\mu}:J']^{-1}}{|\Gamma_c|}
   \,\vartheta_{i,\boldsymbol\mu}
     \bigl(h_c,g;\,p_{i,\boldsymbol\mu}
            (u,c_{i,\boldsymbol\mu}(x_c(f)))\bigr).
\end{equation}
Here \(p_{i,\boldsymbol\mu}\) is the harmonic polynomial map
followed by its tensor with the specified finite oscillator
function. The quotient measure on a double coset contributes
the volume of \(J'\) divided by the stabilizer order. At
product levels this volume is the product
\(\prod_v[J_{i,\boldsymbol\mu,v}:J'_v]^{-1}\).
In the integrand the two level characters on \(J'\) are inverse, and
real equivariance pairs the coefficient with the same harmonic
vector at every point of the compact real factor. This proves the
class formula. In a coefficient extension used for the complex class
sum the polynomial \(x_c(f)\) is extended with \(\iota_K\); the
rational mass factor is unchanged. The sum uses the actual
\((i,\boldsymbol\mu)\) finite test and the same \(g\) on the
two sides.
\releasecite{Formalism/LocalToGlobal/IndefiniteReal/SubgroupCharExtensionOneApply.lean}{136}{Definite Haar convention}
\releasecite{Formalism/LocalToGlobal/IndefiniteReal/ThPinsR.lean}{73}{Refined level, placed coefficients and stabilizers}
\releasecite{Development/AbelianVarietyTheory/Proofs/DivisorProofs/ThetaLiftEqClassSumThetaKernelsLowestKType.lean}{21}{The lowest-type integral on the actual classes}.

The finite actions can be compared before any choice of fitted
basis. On the definite source at a residue-group element \(s\)
use its prescribed lift \(k_V\) and right translation normalized
by \(\chi_{V,J}(k_V)^{-1}\), where \(\chi_{V,J}\) denotes the
extended flag character. At the row lift \(k_W\) the principal
flag character is trivial. At a type place the common lifted
block has the same residue \(s\) on each side. In the rectangular
finite test, \(\bar k_V^{-1}\left(\begin{smallmatrix}I\\0\end{smallmatrix}\right)\bar k_W
=\left(\begin{smallmatrix}I\\0\end{smallmatrix}\right)\).
The remaining determinant character of the finite oscillator
action is \(\chi_{V,J}(k_V)^{-1}\) on the simultaneous lift.
The component lifts may be taken trivial at \(Q(\boldsymbol\mu)\).
Right translation in the definite integral now gives equivariance
with the row right-translation action. The quotient Haar measure
is right invariant. The same equality on the lowest type
extends to the whole real multiplicity map. With this source
action the input at level \((P_0,\psi_0)\) has residue
character \(\psi_0^{-1}\), since the source flag character
divided by the enlarged level character is \(\iota_K(\psi_0)\).
On the row the level character is \(\iota_K(\psi_0^{-1})\).
Thus the projectors on the two sides are the transports of
\[
 e_X=\sum_s e_X(s)s,\qquad
 p_0=\frac{1}{|P_0|}\sum_{s\in P_0}\psi_0(s)s
\]
under their respective normalized actions. Here \(e_X\)
is the good central selector of the chosen finite types.
These are unitary group actions in the class and Petersson
metrics. The projectors are self-adjoint and commute, so their
product is an orthogonal projection. Equivariance
intertwines each projector and its adjoint. The selector \(e_X\)
acts at the residue group and the \(c(\sigma)\) used to
label a global spherical eigenspace acts through the Hecke operators.
\releasecite{Development/AbelianVarietyTheory/Proofs/DivisorProofs/ThetaLiftIntertwinesLevelLiftsLowestKType.lean}{29}{Change of variable and the inverse flag character}
\releasecite{Development/AbelianVarietyTheory/Proofs/DivisorProofs/ThetaLiftMemMultiplicitySpaceBaseLevel.lean}{17}{The inverse character on the row level}
\releasecite{Development/AbelianVarietyTheory/Proofs/TwistedOperators/ThGEquiv.lean}{17}{Compatibility with the good summands}
\releasecite{Formalism/LocalToGlobal/Divisibility/RowLat.lean}{248}{The character in the auxiliary class projector}.

\paragraph{How the record is constructed.}
Let \(B_{i,\boldsymbol\mu}\) send coordinates to the class form
in the prescribed full basis.
For a level coordinate vector \(x\) the form used in the record
is \(f_x=B_{i,\boldsymbol\mu}E_J x\), where \(E_J\) is the
prescribed inclusion of the level coordinates. The source
identification \(J_{M_t}\) is obtained by sending the standard
vectors of the record to these class forms and extending
complex linearly. Its source Gram is the positive definite
class Gram pulled back by \(E_J\). In principal measures the
class pairing of these vectors is the principal definite
pairing. At the character level the source record pairing is
this pairing divided by \(|P_0|\). The division rescales the
Gram and does not change the measure \(\nu_{V_i,\boldsymbol\mu}\)
in \eqref{eq:period-auxiliary-class-evaluation}. On the row
the Petersson pairing rescales by the same \(|P_0|^{-1}\),
with the fixed product of the lowest compact-type dimensions
included in the multiplicity pairing at both levels.

The record's target space \(M_j\) is built by a Hilbert
identification \(\mathcal E_i\) to the good character-level
row space, with \(\mathcal E_i^*\mathcal E_i=|P_0|I\) for the
principal row metric. The orthonormal coordinates used to
construct \(M_j\) determine \(\mathcal E_i\) and the lattice
and basis pulled back along it together. Its lattice is the
inverse image of the prescribed good integral forms. The auxiliary
\(\Theta_{P_0,\boldsymbol\mu,j}\) is obtained as follows.
Evaluate \eqref{eq:period-auxiliary-class-evaluation} on
\(J_{M_t}(\,\cdot\,)\), extend from the lowest type to the
cuspidal row multiplicity space, apply \(e_X\) there and
read the result via \(\mathcal E_i^{-1}\).
Indeed the actual theta integral supplies an extension with
the required row character by the finite covariance above.
A row multiplicity map is determined by its value on the
lowest type since the tensor of these types generates the
product Harish-Chandra module and the multiplicity maps are
continuous on its smooth globalization. Consequently the
linear extension used in constructing \(M_j\) agrees
pointwise with the same theta integral, with the same
\(J'\), class coefficients and source Haar measure.
In particular, for \(I_{i,\boldsymbol\mu}
=j_{M_t}\circ e_{M_t}\), with \(e_{M_t}\) the fixed
identification of the level coordinates with the record's
\(K\)-space \(M_{t,K}\),
\begin{equation}
 \label{eq:period-auxiliary-record-read}
 \mathcal E_i\Theta_{P_0,\boldsymbol\mu,j}(I_{i,\boldsymbol\mu}x)
       = T_{i,\boldsymbol\mu}(f_x).
\end{equation}
In the final record the source and target identifications are
the same maps onto their images, and the construction inserts
this very \(\Theta_{P_0,\boldsymbol\mu,j}\) as its auxiliary
theta map.
\releasecite{Formalism/LocalToGlobal/IndefiniteReal/IsMultMapUnivZero.lean}{117}{The two Hilbert pairings and the unscaled integration Haar}
\releasecite{Development/AbelianVarietyTheory/Proofs/IwasawaDecomposition/LiftDefiniteFormsEqUnitaryFormsClassSumPairing.lean}{74}{Source standard vectors and their class pairing}
\releasecite{Development/AbelianVarietyTheory/Proofs/IwasawaDecomposition/ExistsLiftSpacesOverDefiniteForms.lean}{44}{Construction by the lowest-type extension and the good projection}
\releasecite{Development/AbelianVarietyTheory/Proofs/IwasawaDecomposition/ExistsLiftSpacesOverDefiniteForms.lean}{185}{Insertion of the lift spaces, lattice and maps}
\releasecite{Development/AbelianVarietyTheory/Proofs/IwasawaDecomposition/LiftIdentificationsOfLiftsComparison.lean}{15}{The original identifications viewed as maps onto their images}
\releasecite{Development/AbelianVarietyTheory/Definitions/AutomorphicMembership/RowsRestWordColumnUnitBaseTupleAlt.lean}{392}{Assembly of the auxiliary spaces and theta map}.

\teamanchor{proof.period-auxiliary-field}
\paragraph{The coefficient field of the auxiliary lift.}
Fix an auxiliary index \(i=j-1\) and the prescribed pair of
ranks \(r+t,r\), with \(t\ge r+5\). Use the same definite
class space, finite test and source measure as in
\eqref{eq:period-auxiliary-class-evaluation}. Write
\(A_{i,\boldsymbol\mu,K}\) for the full rational class space
of the auxiliary definite group at its prescribed tuple level.
For \(x\) in this space, denote the unprojected lift evaluated
at the chosen unit lowest-type vector by \(U_{\boldsymbol\mu}(x)\).
We show that the multiplicity vector of the projected lift
lies in the prescribed \(K\)-space \(F_{i,K}\) of auxiliary cusp forms for every such
\(x\) and every selected tuple. This includes the inputs before
forming the complement to the kernel.

The base PEL marking has integer modulus \(M=g_0g_1\) and
auxiliary marking parameter \(A=1\). Write \(Q_{\boldsymbol\mu}\)
for the actual marked twist places and put
\[
 p_{\boldsymbol\mu}=\prod_{q\in\{\# k(P):\,P\in Q_{\boldsymbol\mu}\}}q,
 \qquad M_{\boldsymbol\mu}=M p_{\boldsymbol\mu}.
\]
The residue degrees are one. Take the half ideal
\(I_{\boldsymbol\mu}\) from the product of the distinct chosen
primes below \(Q_{\boldsymbol\mu}\) in the host imaginary
quadratic field and extend it to \(E\). The one-sidedness
condition at each rational prime gives
\(I_{\boldsymbol\mu}\overline{I}_{\boldsymbol\mu}=(p_{\boldsymbol\mu})\).
The marked ideal on the cover is the base marked ideal times
\(I_{\boldsymbol\mu}\), and the cover still has \(A=1\).
The field selection adjoins the finite CM Weber values at these
moduli before choosing the tuple. It suffices to include the
moduli formed from the subsets of the union of ramified places
of the finite allowed twist set. Together with the selected
roots of unity and character values, these choices give the
marked models and their component fields inside the prescribed
coefficient field. On each component the determinant character
of the cover is well defined on the dictionary points because
it is trivial on determinants of rational group elements and on the principal
congruence subgroup at \(M_{\boldsymbol\mu}\). Its finite-order
values give the algebraic component function \(c_{\delta,\boldsymbol\mu}\)
over \(K\). Its reading at a marked frame above \(q\) is
\(\xi_{\boldsymbol\mu}(\det_{\rm root} q)\), with the rational
translate from the dictionary understood. The multiplicative
constant in this reading is one.
\releasecite{Development/AbelianVarietyTheory/Proofs/TwistedOperators/ExistsNormalisedTupleDataOverCoefficientField.lean}{44}{The half ideal and the exact cover modulus}
\releasecite{Development/AbelianVarietyTheory/Proofs/IwasawaDecomposition/CoverFieldOfWeberLev.lean}{44}{Finite field selection before the tuple}
\releasecite{Development/AbelianVarietyTheory/Proofs/IwasawaDecomposition/CoverFieldOfWeberLev.lean}{85}{Component fields at the selected moduli}
\releasecite{Development/AbelianVarietyTheory/Proofs/TwistedOperators/CoversAtOverField.lean}{56}{The algebraic determinant function with constant one}.

Here is the class section on this cover. Read \(x\) at the
placed class representatives in its given coefficient frame
and apply the fixed algebraic harmonic map to the resulting
values. For each class \(c\), use the split-supported
lattice \(\Lambda_c\), its allowed odd denominator \(N_c\)
prime to \(M_{\boldsymbol\mu}\), and the weight
\(w_c=\widehat\phi_{\boldsymbol\mu}\) on the actual marked
quotient. The function \(\phi_{\boldsymbol\mu}\) is the
product of the prescribed type test and the finite twist test
\(\omega^Q\) with inverse determinant character.
Its values and its finite Fourier coefficients belong to \(K\).
At a framed point on the cover, the class jet is the sum of
harmonic derivatives of the pulled-back torsion translates as
in \eqref{eq:primitive-phase-class}, with the enlarged marking
in the Fourier sum. Its denominator is
\(m_cN_c^{d_i}|\mathcal Q_{\boldsymbol\mu}|\), where
\(\mathcal Q_{\boldsymbol\mu}\) is the finite space of marked
labels for all source columns. The Haar class weight is
\(v'_{\boldsymbol\mu}/|\Gamma_c|\), and the algebraic class
scalar is the one attached to the same split-supported lattice.
Multiply the weighted class sum by \(c_{\delta,\boldsymbol\mu}\).
The change of a lift of the marking scales the character
function and the finite Fourier trace by inverse factors.
The jet transforms by the prescribed Hodge weight. Thus the
product descends through the marked cover to a section
\(s_{\boldsymbol\mu}(x)\) over \(K\). The harmonic
contractions remove the change of polarization gauge, and
the last positive-block weight kills the constant term on
each isotropic line as in the primitive calculation below.
Hence this is a cusp section of the fixed base row model.
\releasecite{Development/AbelianVarietyTheory/Proofs/DivisorProofs/ExistsNormalisedTupleDataOverRowFieldConditions.lean}{31}{Finite test in the actual class frame}
\releasecite{Formalism/LocalToGlobal/ProjectiveLattices/ThShort.lean}{244}{The normalized short jet}
\releasecite{Formalism/LocalToGlobal/ProjectiveLattices/ThClass.lean}{184}{The class weights and determinant function}
\releasecite{Development/AbelianVarietyTheory/Proofs/SiegelThetaTransformation/ExistsSectionOneClassSeriesDescends.lean}{22}{Descent of the product on framed charts}
\releasecite{Development/AbelianVarietyTheory/Proofs/DivisorProofs/ReferenceTupleLiftsSpanAllTuples.lean}{59}{The character and cuspidal class section for each tuple}.

The normalized reading of this section must use the component
factors already installed in \(F_{i,K}\). In that installation
the empty-tuple cover has \(p_{\boldsymbol\mu_0}=1\),
\(I_{\boldsymbol\mu_0}=\mathcal O_E\) and
\(c_{\delta,\boldsymbol\mu_0}=1\). The construction initially
uses an auxiliary finite extension to put the reference datum
and all selected marked data in one model. It multiplies every
other cover function of reading \(a'\xi_{\boldsymbol\mu}\) by
\(a/a'\), where \(a=1\) is the reference constant. Denote its
cusp section for \(x\) by \(s^{\rm g}_{\boldsymbol\mu}(x)\).
The installed data keep this normalization after passage to the
fixed space of algebraic forms. They read, at every adelic
quotient point \(q\),
\[
 U_{\boldsymbol\mu}(x)(q)
    = L_i c_i^{\rm fr}(q)\,\operatorname{rd}_{q}
             s^{\rm g}_{\boldsymbol\mu}(x),
 \qquad L_i=(2\pi\xi_R)^{-d_i/2}.
\]
Here \(\operatorname{rd}_{q}\) uses the intrinsic Hodge frame
from the common adelic dictionary and includes its prescribed
determinant factor. The scalar \(c_i^{\rm fr}(q)\) is the
installed phase-one component factor. The level move and root
outside that factor are both one. These choices are part of
the construction of the geometric half and its fixed cusp
basis.
\releasecite{Development/AbelianVarietyTheory/Proofs/IwasawaDecomposition/ExistsReferenceTupleThetaComparisonNormalisedDat.lean}{213}{The identity reference cover}
\releasecite{Development/AbelianVarietyTheory/Proofs/IwasawaDecomposition/ExistsReferenceTupleThetaComparisonNormalisedDat.lean}{261}{The unit reference constant}
\releasecite{Development/AbelianVarietyTheory/Proofs/IwasawaDecomposition/ExistsConstantsThetaComparisonAllTuples.lean}{132}{The common cover constant in the merged data}
\releasecite{Development/AbelianVarietyTheory/Proofs/IwasawaDecomposition/ExistsGeometricHalfPhaseOneOverLawfulCore.lean}{29}{Installation of the component constants and the algebraic basis}
\releasecite{Development/AbelianVarietyTheory/Proofs/DivisorProofs/OneMapRowOfSerOfCore.lean}{17}{The identity move and root for every tuple}.

To compare with the section over \(K\), extend the same
source coefficient isometry to each geometric coefficient
field. At each class representative its value is that of
\(x\) in the fixed class embedding. Pull the residual level
action back from the same reference marking. These choices
leave the finite test, adelic point and lowest-type vector in
\(U_{\boldsymbol\mu}(x)\) fixed. Each decomposition uses the
prescribed source measure and its class weights
\(v'_{\boldsymbol\mu}/|\Gamma_c|\) for its representatives.
The class calculation then gives for both geometric
readings \(s\) of this input
\[
 U_{\boldsymbol\mu}(x)(q)
   = \lambda_i(q)\operatorname{rd}_{q}s,
 \quad
 \lambda_i(q)=k_i(-\delta_i)^{d_i}L_i c_i^{\rm vac}(q)\ne0.
\]
The finite Fourier expansion gives the common sign
\(\delta_i\in\{1,-1\}\). Multiplication by \(N_c\) in the
harmonic derivative cancels \(N_c^{-d_i}\). The even
representatives contribute \(z_c/m_c\), and the class scalar
supplies the exact finite lattice index and the inverse
determinant factor of the moved source vacuum. This leaves
the displayed harmonic normalization. The unit harmonic
scalar \(k_i\) is the same for the two fields because both
coefficient maps send the same extremal tensor to the same
nonzero harmonic polynomial before the Fock comparison. The
vacuum factor \(c_i^{\rm vac}(q)\) uses the base component of
the same dictionary point, with the marked half and Hermitian
frame transported by the level map. The twist of the analytic
vacuum by \(\xi_{\boldsymbol\mu}\) matches the value of
\(c_{\delta,\boldsymbol\mu}\) on that frame. Both cover
constants are one. Cancelling \(\lambda_i(q)\) for each
\(q\) proves
\[ \operatorname{rd}_{q}s^{\rm g}_{\boldsymbol\mu}(x)
      =\operatorname{rd}_{q}s_{\boldsymbol\mu}(x). \]
This equality includes zero lifts. The analytic vacuum and
harmonic factors are used in the comparison inside the
prescribed reading of a cusp section.
\releasecite{Development/AbelianVarietyTheory/Proofs/DivisorProofs/ReferenceTupleLiftsSpanAllTuples.lean}{115}{Retie using the same coefficient isometry and level action}
\releasecite{Development/AbelianVarietyTheory/Proofs/DivisorProofs/ExistsCcIsEquivAtSerOfOneMapRowOfCore.lean}{401}{The class index and character}
\releasecite{Development/AbelianVarietyTheory/Proofs/DivisorProofs/ExistsAdelicLiftEqGeometricClassSumAllTuples.lean}{168}{The harmonic and positive normalization factors}
\releasecite{Development/AbelianVarietyTheory/Proofs/DivisorProofs/CovarianceAtConstantXi.lean}{22}{Matching the harmonic map on the underlying class coefficient}
\releasecite{Development/AbelianVarietyTheory/Proofs/DivisorProofs/CovarianceAtConstantXi.lean}{65}{The determinant at the same marked frame}
\releasecite{Development/AbelianVarietyTheory/Proofs/DivisorProofs/CovarianceAtConstantXi.lean}{94}{The same nonzero multiplier at every quotient point}.

The algebraic cusp sections of the base model commute with
field extension. On each affine chart of the reduced framed
scheme the closed points have finite residue fields over the
ground field. Their maximal ideals intersect in the
nilradical, so a regular function vanishing at all algebraic
geometric points is zero. The same reasoning over
an algebraic closure, followed by scalar extension, shows
that finitely many such framed evaluations separate the
complexification of the finite-dimensional cusp space. This
also proves uniqueness in the class section comparison.
Evaluation coordinates at those points may use their
residue fields. The coefficients of \(s_{\boldsymbol\mu}(x)\)
are in \(K\) by descent and scalar extension on the fixed
cusp space. Apply the installed reading to a base-cusp basis
expansion of this section. Its value is precisely
\(L_i c_i^{\rm fr}(q)\operatorname{rd}_{q}s_{\boldsymbol\mu}(x)\)
at the designated lowest-type vector. Equality on every
quotient representative extends to the whole automorphic
function by left rational invariance. The lowest-type
evaluation is injective on the multiplicity space by the
positive multiplicity norm on the unit tensor. Thus the
unprojected lift is the image of a vector in \(F_{i,K}\).
\releasecite{Formalism/LocalToGlobal/ProjectiveLattices/ReadAlg.lean}{13}{Separation by algebraic geometric points}
\releasecite{Formalism/LocalToGlobal/ProjectiveLattices/ReadAlg.lean}{57}{Independence of evaluations}
\releasecite{Development/AbelianVarietyTheory/Proofs/TwistedOperators/RatOfOneMapOverL.lean}{50}{Reading the fixed cusp basis}
\releasecite{Development/AbelianVarietyTheory/Proofs/TwistedOperators/IpMSelfEqZero.lean}{24}{Injectivity at the unit lowest-type vector}.

If \(e_X=\sum_{g\in G} e_g[g]\) is the good projector, its
operation in the coefficient model is
\(\sum_{g\in G}e_g\chi_i(g)^{-1}\rho_i(g)\).
Here \(\rho_i\) is the action on the fixed algebraic cusp basis and
 \(\chi_i\) is its finite determinant multiplier.
The inverse scalar cancels the scalar of the finite level
action in the adelic reading. The character and all the
projector coefficients are in \(K\). Hence the same conclusion
holds for the actual projected lift, including the scalar
conjugation at the fixed place when the complex model uses the
conjugate form. For either face \(b\), the target vectors
\(w_{0,i,b}(e_k)\) read the prescribed \(K\)-basis of the
good character subspace of \(F_{i,K}\). The factor
\(\sqrt{|P_b|}\) in its common-space inclusion cancels the
inverse factor in its Hilbert character-level basis. The
coordinate functionals of the fixed basis of \(F_{i,K}\) are now
the separating tests. They give \(K\)-values on both the
full source lifts and these fixed target vectors.
\releasecite{Development/AbelianVarietyTheory/BlockTheory/TwistedDiscreteSeries/RayFieldTheta.lean}{41}{Projecting the unprojected algebraic vector}
\releasecite{Development/AbelianVarietyTheory/BlockTheory/TwistedDiscreteSeries/GoodLiftPhase.lean}{56}{Reading the fixed target frame and cancelling the index}
\releasecite{Development/AbelianVarietyTheory/Proofs/TwistedOperators/LiftTestBOfRatG.lean}{25}{The full coefficient coordinates and the two required families of values}.

\paragraph{The fitted projection and the arithmetic columns.}
The coordinate functionals on the fixed rational cusp space give
a \(K\)-matrix \(A\) of separating tests of \(T_{i,\boldsymbol\mu}\)
on the whole definite class space. Let \(G\)
be the \(K\)-Hermitian Gram of the full class space, up to an
overall positive real factor. Then extension of scalars
identifies \(\ker_K A\) with the kernel of \(T_{i,\boldsymbol\mu}\)
and identifies \((\ker_K A)^{\perp_G}\) with
\(\operatorname{im}T_{i,\boldsymbol\mu}^*\). In fact, the
first assertion follows from row reduction and separation,
and the second follows by applying the same extension to the
linear equations for the orthogonal complement. This
argument uses the source Gram.

Intersect the adjoint image with the good and level-character
images. This is the specified Hermitian complement, and we
use the chosen record basis for its rational form.
Since the orthogonal projectors commute with \(T^*T\),
they preserve its kernel and complement, and the intersection
is the complexification of that same rational intersection.
In the fixed class coordinates the basis has matrix \(VZ\).
Write \(\Pi_Z=P_ZV^\flat E_J\). Orthogonal projection
of \(E_J x\) to the good level complement is \(VZ\Pi_Zx\).
The difference is killed by \(T_{i,\boldsymbol\mu}\):
both vectors are in the character level, their difference
pairs to zero with the good component of the adjoint image,
and the good projection already fixes the image of \(T\).
Therefore
\[
 T_{i,\boldsymbol\mu}(B_{i,\boldsymbol\mu}E_Jx)
   =T_{i,\boldsymbol\mu}(B_{i,\boldsymbol\mu}VZ\Pi_Zx).
\]
This identifies the complement in the very class model used
to construct \(M_t\).

The \(W_i\) coordinate basis is the prescribed basis
\(b_{i,k}\) of the inverse-character good part of the
rational cusp forms. Denote its complex realization by
\(J_{W,i}\). In the arithmetic definition of the relative
matrix its columns are first realized with the factor
\(|P_0|^{-1/2}\) used for the smaller-level pairing. The
embedding into the principal multiplicity space multiplies
by \(|P_0|^{1/2}\), so its \(k\)-th column is precisely
\(J_{W,i}b_{i,k}\). In the record \(j_{W,i}\) is the
inverse image of this very basis under \(\mathcal E_i\).
Thus their evaluations at every point agree, with no
relative phase from level conversion. The definition of
\(U_{i,\boldsymbol\mu}\) solves for the columns of
\(T_{i,\boldsymbol\mu}B_{i,\boldsymbol\mu}VZ\)
in this basis using the prescribed embedding. Its rationality
uses the absolute field normalization of the tests at the same
cusp forms. In particular, the matrix is fixed by
\[
 J_{W,i}\!\left(\sum_k (U_{i,\boldsymbol\mu})_{kh} b_{i,k}\right)
   =T_{i,\boldsymbol\mu}(B_{i,\boldsymbol\mu}VZ e_h)
\]
for every fitted basis index \(h\). Applying this equality
with coefficients \(\Pi_Z e_a\) gives, for every standard
level basis vector \(e_a\),
\begin{equation}
 \label{eq:period-fitted-stage-vector}
 \Theta_{P_0,\boldsymbol\mu,j}(I_{i,\boldsymbol\mu}e_a)
      =j_{W,i}\!\left(U_{i,\boldsymbol\mu}
                          (P_ZV^\flat E_J)e_a\right).
\end{equation}
Indeed both sides have the same image under the injective
\(\mathcal E_i\) by \eqref{eq:period-auxiliary-record-read}.
Semilinearity with the same \(\iota_K\) then gives the
identity for every level coordinate vector. This argument
keeps the prescribed integral lattices and all coefficient
readings in place.
\releasecite{Formalism/LocalToGlobal/IndefiniteReal/PartsOrth.lean}{125}{Definition of the rational adjoint-image coordinates}
\releasecite{Development/AbelianVarietyTheory/Proofs/IwasawaDecomposition/LiftLevelFormOrthogonalKilledByThetaMap.lean}{18}{The orthogonal complement killed by the same lift}
\releasecite{Formalism/LocalToGlobal/IndefiniteReal/PartsOrth.lean}{152}{Definition of the relative lift matrix}
\releasecite{Formalism/LocalToGlobal/IndefiniteReal/LinksDefs.lean}{669}{The prescribed good basis and its norm rescaling}
\releasecite{Formalism/LocalToGlobal/Divisibility/DivisionCompatibleToSub2Indef2.lean}{299}{The identity on the actual level coordinates}.

The positive adjoint norm for this auxiliary pair in
\S\ref{sec:stage-petersson} gives surjectivity onto its
good principal row space for the selected tuple.
Averaging a preimage by \(e_Xp_0\) therefore gives a
preimage at the good character level.
The restriction of \(T_{i,\boldsymbol\mu}\) to the fitted
complement is injective and has this good image. It follows
that \(U_{i,\boldsymbol\mu}\) is square and invertible, and
\(E_J^\flat VZ U_{i,\boldsymbol\mu}^{-1}\) gives the explicit
small-level preimages of the prescribed row coordinates.
In particular this is the projection attached to \(M_j\)
and \(j_{W,i}\). On the base level, the separate comparison of
finite types in \eqref{eq:record-individual-isot-rank} and
\eqref{eq:record-base-rank-sum} identifies the dimension with
\(D_0\), the rank of the bottom part of the row chain on that
level.
\releasecite{Formalism/LocalToGlobal/Divisibility/LiftLatLemmas.lean}{198}{The coordinate right inverse}
\releasecite{Formalism/LocalToGlobal/HilbertRowConditions/ColumnsFunctionAtSubspaces.lean}{129}{The type-block partition}
\releasecite{Formalism/LocalToGlobal/NormalizedHeckeFactors/BaseRowsEqRecordRows.lean}{55}{The rank on the recorded blocks}.


\teamanchor{proof.period-row-matching}%
\subsubsection{The common downward coordinates}
Return to the descent from \(X'\) through the selected \(Y_j\)
to \(Y_0\). Section~\ref{sec:period-selected-row-cuts} gives
the common image on the analytic cuts for these normalized
maps. To identify the arithmetic kernel complement in
the original class model, use the principal-level test
identities, separation, and the rationality of
\(V_{\boldsymbol\mu,a}\) and the square ratios. The base row space has positive
dimension in the application. Surjectivity there and separation
give a nonzero left side in \eqref{eq:period-principal-test} for a
rational input, since the rational inputs span. Thus \(\kappa\ne0\).
Multiply each \(\psi_a\) by
\(\iota_K(e_{\boldsymbol\mu_0})c^\square_{\boldsymbol\mu_0}/\kappa\).
The resulting fixed linear functionals still separate. Their values
on every rational downward image are over \(K\), since by
\eqref{eq:period-principal-test} they are
\[
\frac{\iota_K(e_{\boldsymbol\mu_0})c^\square_{\boldsymbol\mu_0}}
     {\iota_K(e_{\boldsymbol\mu})c^\square_{\boldsymbol\mu}}\,
       V_{\boldsymbol\mu,a}(x).
\]
This uses the square ratios before constructing a common field
model of the row space.

For each tuple the kernel in the complexified top space is
the kernel of the matrix of these tests. All its entries are
over \(K\), so this kernel is the complexification of its
rational kernel \(U_{\boldsymbol\mu,K}\).
The rational adjoint image is the orthogonal complement of
\(U_{\boldsymbol\mu,K}\) for the class Gram matrix. Indeed the Gram
is over \(K\) with the involution induced by conjugation at
\(\iota_K\), and the equations for orthogonality extend to the
complex orthogonal complement. Include the central good projector
and impose either of the two prescribed level characters. In the
principal model write the latter projector as
\[
 e_{b,\boldsymbol\mu}
    = \frac{1}{|P_b|}\sum_{g\in P_b}\psi_b(g)^{-1}\rho_{\boldsymbol\mu}(g),
 \qquad b\in\{0,s\},
\]
where \(\rho_{\boldsymbol\mu}\) is the determinant-corrected
finite right action. The two actions and their covariance were
computed in \hyperlink{proof.principal-level-projection}{the
level-change calculation}. Their complex bottom averages intertwine
the same downward maps. In particular their image on the common
principal row space equals the image of each corresponding top
cut. These averages on the bottom preserve its \(K\)-form: for a
rational base preimage, multiplication by an element of the residue
group can be performed on that preimage before descent. Unitarity
gives the analogous action on the adjoints. The adjoint complement
on the cut is thus the intersection of the principal adjoint
complement with the good character space. Its Gram in the
character-level normalization is a positive common multiple of the
restricted principal Gram, so its orthogonal complement is
unchanged. This applies both to the base level used for the
auxiliary stage matrices and to the level of the periods indexed
by words.

In the good coordinates of the level under consideration, let
\(H_{\boldsymbol\mu}\) be the class Gram and
\(Z_{\boldsymbol\mu}\) the prescribed \(K\)-basis matrix of the
adjoint complement. The orthogonal projection to the
complementary coordinates is
\[
 P_{\boldsymbol\mu}=
  (Z_{\boldsymbol\mu}^*H_{\boldsymbol\mu}Z_{\boldsymbol\mu})^{-1}
   Z_{\boldsymbol\mu}^*H_{\boldsymbol\mu}.
\]
The indicated Gram matrix is invertible by positivity at the
chosen embedding. Thus \(P_{\boldsymbol\mu}Z_{\boldsymbol\mu}=I\)
and the complex downward map factors through these coordinates.
Use the images of the base complementary basis as the basis of
\(E_{\rm row}\), and write \(R_K\) for their \(K\)-span in
these coordinates. The matrix of the fixed, normalized tests on
this base basis has full column rank by separation, so select a
square invertible minor over \(K\). For the downward image of
any other complementary basis vector, the coordinates in the
base basis solve the linear system for this minor with right
side in \(K\). Hence their transition \(D_{\boldsymbol\mu}\)
is over \(K\), and is invertible by the equal-image and
kernel-complement properties. Let \(j_{\rm row}\) be the base
complex realization and \(j_{\boldsymbol\mu}\) the realization
of the projected top coordinates. We have proved
\begin{equation}
\label{eq:period-row-match}
 \iota_{\boldsymbol\mu}(j_{\boldsymbol\mu}y)
      = j_{\rm row}(D_{\boldsymbol\mu}P_{\boldsymbol\mu}y).
\end{equation}
For each cut the base realization is the image of its prescribed
complementary basis. Thus the transition in
\eqref{eq:period-row-match} is the unique relative matrix between
these realized bases. This identifies the record on its own
coordinates, including at the sharp level, even if the two
levels were given different bases. Here and below \(j_{\rm row}\)
and \(j_{\boldsymbol\mu}\) on rational inputs are semilinear for
\(\iota_K\).
\releasecite{Formalism/LocalToGlobal/IndefiniteReal/RationalAdjoints/RatZOfTests.lean}{18}{Tests on all rational inputs}
\releasecite{Formalism/LocalToGlobal/IndefiniteReal/PrimitiveStatementsWithPassage.lean}{58}{Finite equivariance}, \releasecite{Formalism/LocalToGlobal/IndefiniteReal/PrimitiveStatementsWithPassage.lean}{85}{the principal row and its two level cuts} and
\releasecite{Formalism/LocalToGlobal/IndefiniteReal/PartsOrth.lean}{134}{the relative matrices in the prescribed bases}.

For the metric at
\(\iota_K\), use the bottom Gram matrix pulled back along
\(j_{\rm row}\). If the basis is realized first on the principal
level, divide its Gram by \(|P|\) to obtain the metric at
the character level. Indeed the Petersson inner product on the
larger level uses volume one there, whereas its volume in the
principal normalization is \(|P|\). The two realizations can
use conjugate representatives of the infinite place, which
conjugates the pairing and preserves its norm. On a summand with
\(\iota_{\boldsymbol\mu}\iota_{\boldsymbol\mu}^*=\lambda I\),
the downward map multiplies squared norms by \(\lambda\) on
the kernel complement. Indeed a preimage there is
\(\lambda^{-1}\iota_{\boldsymbol\mu}^*e\), whose squared norm
is \(\lambda^{-1}\|e\|^2\).
The same argument gives the cross-pairing. If a fitted vector
\(y_p\) has image in that summand, then
\(\iota_{\boldsymbol\mu}^*\iota_{\boldsymbol\mu}y_p=\lambda y_p\).
Pairing this identity against any fitted vector \(z\) gives
\[
 h_{v_0}(D_{\boldsymbol\mu}y_p,D_{\boldsymbol\mu}z)
      =\lambda\,h^Q_{\boldsymbol\mu,v_0}(y_p,z)
\]
in the identified metrics at the chosen place \(v_0\).
Here \(\lambda=\kappa_{v_0}\prod_j\lambda_{j,p,v_0}\)
once the normalized step factors have been established.

At any other infinite place \(v\), use the fixed base
complement metric \(h^Q_v\) on the rational coordinates, with
the prescribed positive scalars \(\kappa_v,\lambda_{j,p,v}\).
The underlying class Hecke operators are the commuting
double-coset operators with their adjoints for the class
mass. Integration of a right translate and change of variables
give the adjoint by inverting the double coset and conjugating
its character. The class-sum version is the same calculation
on each source and target class. The adjoint matrix identities
transport to \(v\) because the operator and level matrices
have entries in the CM subfield, on which the field involution
maps to complex conjugation at every embedding.
The fitted class subspace is invariant under this family,
including its adjoints. Thus its orthogonal projection
commutes with the operators, and the induced operators on
fitted quotient coordinates are commuting normal operators.
Different joint spectral labels have disjoint joint eigenspaces,
which are orthogonal for the quotient form \(h^Q_v\).
For a joint spectral label \(p\) let \(E_p\) be the projector
to that summand. The prescribed row metric is
\[
 h_v(y,z)=\sum_p
     \kappa_v\left(\prod_j\lambda_{j,p,v}\right)
     h_v^Q(E_p y,E_p z).
\]
Each term is positive semidefinite and the projections sum to
the identity, so this gives a positive metric. Moreover for
\(y_p\) in one summand and any \(z\), it gives
\(h_v(y_p,z)=\kappa_v(\prod_j\lambda_{j,p,v})h_v^Q(y_p,z)\).
In particular \(h_v\) dominates \(h^Q_v\) times the smallest
of the positive factors,
\(h_v(y,y)\ge\kappa_v\min_p(\prod_j\lambda_{j,p,v})h^Q_v(y,y)\).
These identities give the dual norms and comparisons used in
Section~\ref{sec:integral} once the selected theta norm
scalars on the actual summands are determined.
\releasecite{Formalism/LocalToGlobal/IndefiniteReal/PrincipalLevelLattice.lean}{204}{Principal-level and character-level metrics}
\releasecite{Formalism/LocalToGlobal/IndefiniteReal/PrincipalLevelLattice.lean}{250}{Metrics at the other embeddings}
\releasecite{Formalism/LocalToGlobal/IndefiniteReal/CxCutPO.lean}{37}{Transport of adjoints and spectral orthogonality}.

\teamanchor{proof.stage-petersson}
\subsubsection{Petersson eigenvalues of the auxiliary row stages}\label{sec:stage-petersson}
Fix an auxiliary fitted row stage and a selected tuple
\(\boldsymbol\mu\). Its pair consists of the rank-\(r\) row space
\(W\) and the definite space \(V\) of rank \(n=r+t\), with
\(t\ge r+5\). These ranks are fixed over the stages. Write
\(f=[F:\mathbb Q]\). First choose a coefficient embedding over
the host imaginary quadratic field. The relative twist is
finite order and trivial at infinity, so for this embedding
its real test and real type are those of the base tuple.
The positive norm calculation below applies to every selected
tuple at the principal level. The quantitative bound needed
for \(\lambda_{j,p,v}\) concerns the base tuple.
\releasecite{Formalism/LocalToGlobal/IndefiniteReal/ExistsChosenExtremalVector.lean}{253}{Fixed auxiliary row rank}
\releasecite{Development/AbelianVarietyTheory/BlockTheory/TwistedDiscreteSeries/HeckeTransferRankOne.lean}{78}{Fixed auxiliary definite rank}.
Let \(T\) be the full downward lift from the source class space to
the cuspidal row multiplicity space before the principal row
projection. Use the Hermitian class-sum metric on the source.
On the row use the Petersson multiplicity metric with finite
principal-level volume one and real measures \(dg_\rho\) as
specified below. For an auxiliary subgroup \(H_P\) of the row
residue group, the character-level adjoint in these measures is
\(T^{*_{\rm lev}}=|H_P|^{-1}T^*\) when the source retains its
principal normalization. At principal level \(H_P=\{1\}\).
In \eqref{eq:period-auxiliary-record-read} both pairings are
divided by \(|P_0|\), so their simultaneous rescaling leaves
the adjoint of the restricted map unchanged. The factors \(|P|\) in the passage from the later
principal row metric to the character metric are handled by the
prescribed simultaneous norm rescaling.
\releasecite{Development/AbelianVarietyTheory/BlockTheory/DivisorBase/TwistedLevelValues.lean}{148}{The level adjoint}
\releasecite{Development/AbelianVarietyTheory/Proofs/DivisorProofs/AdelCOfInnerProductFormula.lean}{23}{The multiplicity metric}
\releasecite{Formalism/LocalToGlobal/IndefiniteReal/PrincipalLevelLattice.lean}{204}{The row metric conversion}.

For a pure unit lowest-type vector \(e=\otimes_\rho e_\rho\)
and an orthonormal basis \((b_k)_k\) of the definite coefficient
space \(C_V\), let \(\varphi_k\) be the actual oscillator test
associated to \(e\otimes b_k\) by the joint-harmonic map. Its
finite factor \(f_{\rm fin}\) is the prescribed level test and its
real factor is the polynomial applied to the prescribed Cartan
vacuum. The global oscillator measure has rational-lattice covolume
one. For \(y\) in the multiplicity space write \(u_y=y(e)\) for
lowest-type evaluation, and \(\Theta^\uparrow(\varphi_k,u_y)\)
for the scalar reverse lift to \(U(V)\) with kernel
\(\theta(g,h;\varphi_k)\) paired against \(\overline{u_y(g)}\).
Put \(d_K=\prod_\rho d_{K,\rho}\), where \(d_{K,\rho}\) is
the dimension of the lowest compact type at \(\rho\), and put
\(d_V=\dim C_V\).
Lowest-type evaluation satisfies
\(\langle y,y'\rangle=d_K\langle u_y,u_{y'}\rangle_{W}\)
in compatible Hermitian conventions. Taking the adjoint inside
the defining class sum therefore first multiplies it by \(d_K\).
Each source term is weighted by the inverse of the class
stabilizer order. Expanding the coefficient pairing in \((b_k)_k\)
gives
\[
 \langle T^*y,T^*y'\rangle
    =d_K^2\sum_k
       \langle\Theta^\uparrow(\varphi_k,u_y),
              \Theta^\uparrow(\varphi_k,u_{y'})\rangle_V .
\]
Here the definite Haar measure gives the source level and the
real compact factors total volume one, so integrating the
coefficient pairing gives exactly the above class sum.
\releasecite{Formalism/LocalToGlobal/IndefiniteReal/SubgroupCharExtensionOneApply.lean}{136}{Definite Haar measure}
\releasecite{Development/AbelianVarietyTheory/Proofs/DivisorProofs/FormOfClassSums.lean}{79}{Evaluation and class-sum adjoint}
\releasecite{Development/AbelianVarietyTheory/Proofs/DivisorProofs/AdelCOfInnerProductFormula.lean}{160}{The tested adjoint form}.

Apply the convergent norm formula to these reverse lifts, with
doubling source \(W\oplus(-W)\) and target \(V\).
The forms \(u_y\) are cuspidal and \(K\)-finite, the actual
\(\varphi_k\) are Schwartz with polynomial Fock part, \(V\)
has Witt index zero, and \(n>2r\).
With \(s_0=t/2>r/2\), the ordinary Siegel Eisenstein series
converges absolutely and equals the doubled theta integral
in quotient-volume-one measures by the convergent Siegel--Weil
formula. The ratio \(\tau(U(V))/[E:F]=1\) converts the
formula of Gan--Qiu--Takeda to these measures. The tensor
doubling map evaluates at zero by the oscillator scalar
product in the covolume-one measure. The inducing character
in the release convention is \(\chi_V^{-1}\), where
\(\chi_V\) is the \(V\)-splitting character denoted \(c_\mu\).
The inverse character on the second copy in the doubling
integral thus cancels the character in the tensor identity.
Cuspidality kills the lower rational orbits. The open
orbit unfolds to the integral of the oscillator matrix
coefficient paired with the conjugate of the row matrix
coefficient. The condition \(s_0>r/2\) gives the absolute
open-orbit unfolding. Rapid decay of the row cusp forms
and compactness of \([U(V)]\) justify the kernel integrations
in the seesaw. Thus with the actual
measures the norm formula is
\begin{equation}
\label{eq:actual-stage-upward-norm}
\begin{split}
 \langle\Theta^\uparrow(\varphi_k,u_y),
        \Theta^\uparrow(\varphi_k,u_{y'})\rangle_V
 ={}&\operatorname{vol}_{\nu_V}[U(V)]
 \int_{U(W)(\mathbb A_F)}
 \langle\omega(g)\varphi_k,\varphi_k\rangle_{\rm osc}\\[-2pt]
 &\hspace{20mm}\cdot
     \overline{\langle R(g)u_y,u_{y'}\rangle_W}\,d\nu_W(g),
\end{split}
\end{equation}
up to simultaneously transposing the arguments when using a
conjugate-linear-first convention. Changing the Haar scale on
\(U(W)\) multiplies both sides by the square of that scale,
and the factor on \([U(V)]\) gives the exact definite Haar
convention. This is the \(n>2r\) absolutely convergent
application of Siegel--Weil and the doubling open-orbit
calculation.
\cite[\S\S3.1, 7.1, 11.3, 11.5, Theorem~7.1]{gqt}
\releasecite{Development/AbelianVarietyTheory/Proofs/DecoratedLatticeProofs/UnitaryThetaInnerProductFormulaConvergentRange.lean}{36}{Placed doubling integral}
\releasecite{Development/AbelianVarietyTheory/Proofs/DecoratedLatticeProofs/ThetaInnerProductFormulaDefiniteOfRallisConverge.lean}{41}{Actual definite and row measures}.

The real coefficient in \eqref{eq:actual-stage-upward-norm}
at \(\rho\) is \(H_\rho c_\rho(g_\rho)\), where
\[
 c_\rho(g)=\langle\pi_\rho(g)e_\rho,e_\rho\rangle,\qquad
 H_\rho=\operatorname{hook}(\mathbf a^+_{j,\sigma,\rho})\,
        \operatorname{hook}(\mathbf a^-_{j,\sigma,\rho}).
\]
The highest joint harmonic is a product of leading and trailing
minor powers in disjoint sets of variables. The mixed
Laplacians annihilate it because the used column sets are
disjoint. The corresponding lowering operators annihilate the
real highest vector in the Fock model. Its compact orbit maps
equivariantly from the tensor of the lowest row type and the
definite coefficient type. The lowest-weight module it generates
has the row discrete-series parameter: the \(K_\rho\)-type fixes
the holomorphic parameter, and the lowering relations give its
cyclic lowest-weight model. Equivalently, its invariant scalar
product is fixed by the lowest-type pairing using the adjoint
raising and lowering operators and the Poincar\'e--Birkhoff--Witt
ordering. Schur's lemma on this tensor of compact types fixes
the constant from a unit pair mapping to the specified minor
product. The Fock squared norm with \(\|z^\beta\|^2=\beta!\)
is the Gaussian average of the squared modulus of the
polynomial. For a partition \(\kappa\) of length at most \(p\),
the product of leading minors is the matrix coefficient on a
unit highest vector of the Schur representation \(S_\kappa\)
of \(\mathrm{GL}_p\). This follows by realizing the highest
vector in the tensor of exterior powers with exponents
\(\kappa_i-\kappa_{i+1}\). Factor the square complex Gaussian
matrix as \(X=UR\) with \(U\) unitary and \(R\) upper triangular
with positive diagonal. Conditional on \(R\), \(U\) is Haar.
The upper triangular factor acts on the highest line by
\(\prod_i R_{ii}^{\kappa_i}\), and compact Schur orthogonality
contributes \(1/\dim S_\kappa\). Successive orthogonal
projections in Gram--Schmidt give independent squared diagonal
entries with gamma distributions of shapes \(p-i+1\)
and scale one. Hence the squared norm is
\[
 \frac{1}{\dim S_\kappa}
       \prod_{i=1}^p\frac{(\kappa_i+p-i)!}{(p-i)!}
 =\prod_{b\in\kappa}h(b)
\]
by the Weyl dimension and hook product formulas. The trailing
minors follow by coordinate relabeling. Disjoint variables
give the product of the two hook products and each \(b_k\) has
norm one. The other
matrix coefficient in the integrand has the opposite conjugation.
The resulting real factor is
\(H_\rho J_\rho\), where
\(J_\rho=\int_{U(W)(F_\rho)}|c_\rho(g)|^2\,dg_\rho\).
\releasecite{Development/AbelianVarietyTheory/Proofs/DecoratedLatticeProofs/JointHarmonicsLowestKTypeCoefficientReal.lean}{41}{Lowest-type Fock coefficient}
\releasecite{Development/AbelianVarietyTheory/Proofs/DivisorProofs/FockFormMinorProductEqProdHookProducts.lean}{140}{Fock norm of the prescribed minors}.

Let \(S\) be the finite places below the type and twist primes,
and \(q_v\) the cardinality of the residue field of \(F_v\).
Each place of \(S\) is split, with exactly one chosen prime
of \(E\) above it. On the chosen half the support of the test
is integral, with reduction
\(\left(\begin{smallmatrix}A\\0\end{smallmatrix}\right)\)
in the \(n\)-by-\(r\) matrix coordinates, where \(A\) is in
the specified subgroup \(H_v\subseteq{\rm GL}_r(k_v)\) at a
type place and in \({\rm GL}_r(k_v)\) at a twist place. The
other half uses the conjugate-dual lattice Fourier model.
Here \(k_v\) is the residue field and the multipliers on the
support are unitary level characters. If the support and its
translate meet, take a full-rank top block \(A\) there.
Integrality of both \(A\) and \(Ag\), with unit determinant
reductions, gives \(g\in{\rm GL}_r(\mathcal O_v)\).
At type places \(A\) and \(Ag\) in \(H_v\) modulo the
uniformizer give \(\bar g\in H_v\). The split dual frame
gives the inverse conjugate at the other half, so the same
conclusion holds in the unitary group. Thus the finite
oscillator coefficient at \(v\) is supported on the actual row
level \(J_v\), there transforming by the same unitary character
as the row functions. In the convolution the two characters
cancel. Integration over \(J_v\) contributes the finite squared
norm and the index of the principal subgroup in \(J_v\). The
global index is \(|H_P|\). At a tuple prime \(v\mid Q(\boldsymbol\mu)\)
the determinant weight in the prescribed test is the inverse
of the relative splitting character on row units. They cancel
in the row action, so \(J_v\) is the full hyperspecial
stabilizer and its character is trivial. The relative
character has modulus one, so the squared norm of the test
still counts every invertible top block with weight one.
\releasecite{Formalism/LocalToGlobal/IndefiniteReal/DatumVec.lean}{51}{The finite theta test}
\releasecite{Development/AbelianVarietyTheory/Proofs/TwistedOperators/MemLevelOfLevelDatumNeZeroTranslate.lean}{31}{Support intersection and the row level}
\releasecite{Development/AbelianVarietyTheory/Proofs/TwistedOperators/FiniteZetaIntegralThetaLiftEqPartialLValue.lean}{179}{Finite convolution on the actual Hecke part}.

Outside \(S\) the self-dual vacuum is spherical and the
stabilizer has volume one. The extension \(E/F\) is unramified
at all finite places. For the tested joint character let
\(\beta_{w,1},\ldots,\beta_{w,r}\) denote the normalized
standard Satake parameters at \(w\mid v\). In the split
chart they are the \({\rm GL}_r\) parameters at one prime and
their inverses at the other. In the inert chart they are
\(\alpha_1,\ldots,\alpha_{\lfloor r/2\rfloor}\), their inverses,
and a \(1\) if \(r\) is odd. The unramified character of the
spherical section at \(w\) is \(\chi_V^{-1}\).
With \(q_w=\mathrm Nw\) and \(\varpi_w\) a uniformizer, put
\[
 E_S=\prod_{w\not\mid S}\prod_{a=1}^r(1-x_{w,a})^{-1},
\qquad
 x_{w,a}=\beta_{w,a}\chi_{V,w}^{-1}(\varpi_w)\,
                         q_w^{-(t+1)/2}.
\]
Let \(\epsilon_v=1\) at a split place and \(\epsilon_v=-1\)
at an inert place. The spherical coset sum at \(v\) is the
local factor of \(E_S\) multiplied by \(d_v^{-1}\), where
\[
 d_v=\prod_{a=t+1}^{n}(1-\epsilon_v^a q_v^{-a})^{-1},
 \qquad d_{\rm off}=\prod_{v\notin S}d_v,\qquad
 d_S=\prod_{v\in S}d_v.
\]
For the split computation the two determinant coset factors
have inverse Satake lists and inverse character values, and the
common density correction is \(\prod_{a=t+1}^{n}(1-q_v^{-a})\).
For the inert computation let \(z=\chi_V^{-1}(\varpi_v)\).
The section on a graph of length \(\ell\) in the doubled
lattice has value \((zq_v^{-n})^\ell\). The spherical sum in
this length variable \(X=zq_v^{-n}\) has numerator
\(\prod_{i=0}^{r-1}(1-(-q_v)^iX)\) and denominator the
Euler polynomial on the list above at \(q_v^{r-1}X\).
Indeed the sum over the integral hyperbolic-frame Iwasawa
cells gives the length series, with the \(i\)-th geometric
factor carrying the residue norm \((-q_v)^i\). In this
calculation the trace-zero rank-one term in the radical has
one residue coordinate and the paired coordinates have two.
This uses the unramified involution at the ordinary inert places of the family. Since \(z=(-1)^n\) by
\(\chi_V|_{\mathbb A_F^\times}=\eta_{E/F}^n\), setting
\(a=n-i\) identifies the numerator with \(d_v^{-1}\) and
\(q_v^{r-1}X=zq_w^{-(t+1)/2}\).
Equivalently the normalized section at \(s_0=t/2\) has
denominator the factors at \(2s_0+r-i\) with inducing
characters of sign \(\epsilon_v^{n+i}=\epsilon_v^{n-i}\).
\cite[\S11.6]{gqt}
\releasecite{Development/AbelianVarietyTheory/Proofs/TwistedOperators/UnramifiedZetaSumGLCosetsStatement.lean}{14}{Split spherical coset sum}
\releasecite{Development/AbelianVarietyTheory/Proofs/TwistedOperators/UnramifiedDoublingIntegralInertPlace.lean}{84}{Inert length series and the local factors}
\releasecite{Development/AbelianVarietyTheory/Proofs/TwistedOperators/HasSumLocalSphericalZetaUnramified.lean}{101}{Comparison with the normalized base-change parameters}.

In the finite summation one may replace the actual matrix
coefficient by the normalized spherical coefficient times its
value at the identity. Aggregate over each hyperspecial
double coset using its Hecke eigenvalue on \(y\); this works
for \(y'\) in the full cuspidal multiplicity space.
The principal coset support above and the good spherical
sum therefore give the exact finite multiplier
\[
 |H_P|\,\|f_{\rm fin}\|^2_{\rm osc}\,E_S/d_{\rm off}.
\]
The trivial unitary bound is
\(|\beta_{w,a}|\le q_w^{(r-1)/2}\).
At a split place one can sum the normalized coefficient over
the integral sublattice cosets of each determinant valuation
in Hermite normal form. The generating series has roots the
Satake parameters scaled by \(q_w^{(r-1)/2}\), and its
coefficients are bounded by the corresponding counts for the
trivial representation. Its disk of convergence contains
\(|X|<q_w^{-(r-1)}\), which excludes poles there and gives the
asserted bound. The inverse bound follows by applying the
same argument to the dual. In the inert length formula, using
the trivial spherical coefficient shows that the series of
absolute values converges on \(|X|<q_v^{-2(r-1)}\).
Every unitary spherical coefficient is bounded by one, while
the numerator of the rational expression has no zero in that
disk. The absence of poles gives \(|\alpha_i|\) and
\(|\alpha_i^{-1}|\) at most \(q_v^{r-1}\), as required.
Here the underlying unitary structure is the right-regular
Petersson structure on the cusp space. The estimates also
give absolute convergence of the placed coset sums and
their Euler product at \(s_0\).
\releasecite{Development/AbelianVarietyTheory/Proofs/TwistedOperators/PowDvdOfHermMatEq.lean}{242}{Split trivial Satake bound}
\releasecite{Development/AbelianVarietyTheory/Proofs/DecoratedLatticeProofs/SatakeTrivialBoundOfBoundedZonalSphericalUnitary.lean}{168}{Inert trivial Satake bound}.

The definite volume is
\([K_{\rm lat}:K_V]\,\mathfrak m(V)\), with
\(\mathfrak m(V)\) the self-dual definite class mass.
At \(v\in S\) the definite source level is the
preimage of the parabolic residue matrices with diagonal
blocks in \(H_v\) and \({\rm GL}_t(k_v)\), with a free
off-diagonal block of size \(rt\). On the chosen finite
polarization, normalized additive lattice volume one gives
the local squared norm \(|H_v|q_v^{-nr}\). Thus
\[
 [K_{{\rm lat},v}:K_{V,v}]\,\|\varphi_{{\rm fin},v}\|^2
  =\frac{|{\rm GL}_n(k_v)|}{|H_v|\,|{\rm GL}_t(k_v)|q_v^{rt}}\,
      \frac{|H_v|}{q_v^{nr}}
  =\prod_{a=t+1}^{n}(1-q_v^{-a}).
\]
At good places both local factors equal one. This proves the
exact cancellation
\([K_{\rm lat}:K_V]\,\|f_{\rm fin}\|^2_{\rm osc}=d_S^{-1}\).
Combining the real and finite integrals with the adjoint
identity and dividing by \(|H_P|\) gives, on every nonzero
joint Hecke summand,
\begin{equation}
\label{eq:actual-stage-petersson-scalar}
 TT^{*_{\rm lev}}=\lambda I,\qquad
 \lambda=\mathfrak m(V)d_V\,\frac{|E_S|}{d_{\rm off}d_S}\,
             \prod_\rho H_\rho d_{K,\rho} J_\rho .
\end{equation}
In fact, testing the adjoint identity on a nonzero diagonal
vector shows \(E_S\) is positive real, since the other factors
are positive and the absolutely convergent \(E_S\) is
nonzero. One of the two factors \(d_K\) in that identity
is absorbed in the multiplicity pairing, and summing the
orthonormal basis gives \(d_V\).
\releasecite{Development/AbelianVarietyTheory/Proofs/DivisorProofs/VolEqIndexMulMass.lean}{26}{Volume in the definite metric}
\releasecite{Development/AbelianVarietyTheory/Proofs/TwistedOperators/ThetaLevelIndexMulSchwartzNormFormula.lean}{24}{Level index times squared norm}
\releasecite{Development/AbelianVarietyTheory/Proofs/DivisorProofs/AdelCOfInnerProductFormula.lean}{214}{The scalar on the full Hecke part}.

For each selected tuple, the good principal row space is the
orthogonal sum of its joint Hecke eigenspaces away from
\(S\). The spherical operators there commute and taking
the Petersson adjoint is induced by inverse double cosets
and conjugation of the coefficients, so the sum exhausts
this finite-dimensional space. On each nonzero summand
\eqref{eq:actual-stage-petersson-scalar} is positive. All
finite places of \(Q(\boldsymbol\mu)\) contributed the
compact factor computed above, and the remaining Euler
product is in absolute convergence with nonzero factors by
the unitary bound. The full adjoint is thus injective on
the good row multiplicity space for every \(\boldsymbol\mu\).
Explicitly, for any \(y\) in a summand its preimage is the
adjoint image divided by that summand's scalar. This gives
the principal-level surjectivity used in
\S\ref{sec:period-fitted-stage}. The estimates that follow
use the base tuple.
\releasecite{Development/AbelianVarietyTheory/Proofs/TwistedOperators/PrimeToHostCore.lean}{239}{Auxiliary norm formula for every tuple}
\releasecite{Development/AbelianVarietyTheory/Proofs/DecoratedLatticeProofs/UpInjWOfAForm.lean}{20}{Positive adjoint on the whole good row space}.

Both coefficient dimensions in \eqref{eq:actual-stage-petersson-scalar}
are at least one. The standard orthonormal lattice in the genus
contributes \((|\mu(E)|^n n!)^{-1}\) to \(\mathfrak m(V)\).
For an integral unitary matrix, summing the squared moduli
of a column over all embeddings gives \([E:\mathbf Q]\).
Each nonzero entry contributes at least that much by the
arithmetic--geometric mean inequality and its nonzero integral
norm. Thus the stabilizer consists of monomial matrices with
root-of-unity entries. The inequality
\(\varphi(u)^2\ge u/2\), checked on each prime power, gives
\(|\mu(E)|\le 8f^2\). This also proves
\(-\log\mathfrak m(V)\le n\log n+5nf\).
For the finite factors, \(|x_{w,a}|\le q_w^{-7/2}\).
Each \(\log|1-x_{w,a}|\le|x_{w,a}|\), and
\(\sum_w q_w^{-2}\le [E:\mathbf Q]\sum_p p^{-2}<2f\).
Similarly \(-\log(1-\epsilon_v^a q_v^{-a})\le2q_v^{-a}\).
As there are \(r\) factors at each \(v\) and \(t+1\ge2\),
the original bounds give
\(\log(|E_S|/(d_{\rm off}d_S))\ge-10fr\).
\releasecite{Development/AbelianVarietyTheory/Proofs/DivisorProofs/UnitaryMassLowerBoundOfOrthonormalLatticeFrames.lean}{148}{The orthonormal class mass}
\releasecite{Development/AbelianVarietyTheory/Proofs/DivisorProofs/IntegralEqClassSumOfHom.lean}{181}{Uniform local factor bounds}.

At \(\rho\), \(dg_\rho\) is the Haar measure of identity
density \(2^{-r(r-1)/2}\) in the release's real matrix
tangent coordinates. On the compact dual this has total
mass \(\operatorname{vol}_{0}U(r)=
\prod_{a=0}^{r-1}2\pi^{a+1}/a!\).
The normalized compact-dual measure transfers the invariant
identity density from the compact group of volume one.
Consequently \(dg_\rho\) is
\(\operatorname{vol}_{0}U(r)\) times the normalized
compact-dual Haar. Schur orthogonality for the genuine row
discrete series gives \(J_\rho=(d_\rho^D)^{-1}\), and
Harish--Chandra's degree polynomial, in the measure just
specified, is
\[
 d_\rho^D=\frac{1}{\operatorname{vol}_{0}U(r)}
        \prod_{a<b}\frac{|\Lambda_{\rho,a}-\Lambda_{\rho,b}|}{b-a}.
\]
The central circle is compact and the central character
unitary. The measure on the central circle scales the
identity by the same factor on the group and on its compact
dual, so the formula is for the full group \(U(p,q)\).
The actual \(\Lambda_{\rho,a}\) are strictly decreasing
and in \((r-1)/2+\mathbf Z\). In particular the degree
polynomial uses this shifted parameter. The labels over the
host quadratic field enumerate the original decreasing
row parameters. The spread is at most \(d_\rho+n+3\)
with \(d_\rho=|\mathbf a^+_{j,\sigma,\rho}|+
|\mathbf a^-_{j,\sigma,\rho}|\).
Indeed for \(q>0\) the first positive and first negative
Fock row lengths bound the extreme parameters by their
defining affine offsets. In the zero-based stage notation,
with stage offset \(e_j\), the first Fock row lengths are the
nonnegative parts of the floors of
\(\Lambda_{\rho,1}+e_j-(p-1)-(t+1)/2\) and
\(-(\Lambda_{\rho,r}+e_j+q+(t-1)/2)\), respectively. Adding
the resulting upper bounds on the two affine expressions,
with one added to each length for its floor, gives the
parameter spread bound \(d_\rho+n+1\).
For \(q=0\) the last parameter is \(-b_j\) in the release
parameter data, and \(e_j-b_j\ge(t+3)/2\) gives the same
bound from the positive length alone. Relaxing the spread by
two yields the displayed size used in the degree estimate.
\cite[\S41, Theorem~16]{ar4-harish-chandra}
\cite[(3.10), (3.12), Theorem~(8.1)]{ar4-atiyah-schmid}
\releasecite{Development/AbelianVarietyTheory/Proofs/DivisorProofs/StandardMeasureUnitaryPQEqVolumeSmulCompactDual.lean}{29}{Transfer of the real density}
\releasecite{Development/AbelianVarietyTheory/Proofs/DivisorProofs/LiftModelParameterDecreasingHalfIntegralWidth.lean}{52}{Actual parameter spread}
\releasecite{Development/AbelianVarietyTheory/Proofs/DivisorProofs/FormalDegreeRealPlaceUpperBoundOfHarishChandra.lean}{17}{Shift and degree of the selected representation}.

The stage sizes give \(d_\rho\le d_{\max}=3rb\),
\(L'=\log(b\,\mathrm{rd}_E g_0g_1 n)\ge5.5\),
\(\log n\le L'\), and \(\log(d_{\max}+2)\le2L'\).
The compact volume is at least \(r^{-r^2}\), from the
factorial product. Since \(d_\rho+n+3\le n(d_\rho+2)\)
for \(n\ge3\), one obtains
\[
 \log d_\rho^D\le r^2\log r+
         \frac{r(r-1)}2\log(n(d_{\max}+2))
       \le \frac54 n^2L' .
\]
Here the denominators \(b-a\) are at least one and \(2r\le n\).
With \(n\ge 8,\ f\ge1\), the total loss from the mass, the
finite factors and the real degrees is at most
\(n\log n+5nf+10fr+\tfrac54 fn^2L'\le8fn^2L'\).
It follows that
\begin{equation}
\label{eq:stage-petersson-lower-bound}
 \log\lambda\ \ge\ \sum_\rho\log H_\rho-8fn^2L'.
\end{equation}
\releasecite{Development/AbelianVarietyTheory/Proofs/DivisorProofs/AdelCOfInnerProductFormula.lean}{52}{Quantitative stage bound}.

The positive Fock hooks give exactly the base hook sum
\(\sum_\rho\log H^+_{j,\rho}\) at every host-field embedding:
the place labels permute the tested weights and the row
signature is constant along the real places. The negative
hooks have positive integral values. Thus
\(\exp(\sum_\rho\log H^+_{j,\rho}-8fn^2L')\) is a common
lower bound on \(TT^*\) at these embeddings and all nonzero
Hecke summands. If \(P_{\rm row}\) is the orthogonal
principal projection, then
\(\|(P_{\rm row}T)^*y\|^2=\|T^*y\|^2\) on
\(\operatorname{im}P_{\rm row}\). The same lower bound holds
for the squared singular values of \(P_{\rm row}T\) on its
kernel complement.

\teamanchor{proof.stage-primitive-class}
\subsubsection{The primitive class section of the base stage}
\label{sec:stage-primitive-class}
The phase-one comparison uses the base tuple of the
auxiliary pair \(V,W\) above. Fix its prescribed PEL row
model (the moduli model with polarization, endomorphisms and
level structure) over \(K\), the automorphic bundle for its
weight, and the framed theta cover with the two marked
torsion levels. The algebraic class section is built on this
cover, using the polarization line, its normalized theta
structure, and the \(E\)-linear harmonic differentiation in
the prescribed Hodge frames. At the primitive tuple the
relative twist is trivial and its ramified set is empty.
The level cover has degree one. Its integer multiplier
is \(1\), its half-level ideal is \(\mathcal O_E\), and its
scalar and finite determinant correction sections are
both \(1\). Thus both readings use the same marked half
of the row lattice.
\releasecite{Development/AbelianVarietyTheory/Proofs/IwasawaDecomposition/ExistsGeometricHalfPhaseOneOverLawfulCore.lean}{29}{Geometric data used in the phase-one construction}
\releasecite{Development/AbelianVarietyTheory/Proofs/IwasawaDecomposition/ExistsReferenceTupleThetaComparisonNormalisedDat.lean}{213}{Identity cover and trivial relative determinant character}
\releasecite{Development/AbelianVarietyTheory/Proofs/UnitaryExistence/ExistsDeltaDatumHigherLevel.lean}{72}{Scalar-one level cover}
\releasecite{Formalism/LocalRingCardinality/NFLeTwoMulTSq.lean}{33}{The selected size of the auxiliary definite rank}.

Use the definite source class representatives at the
refined level. Let \(\Lambda_c\) be the placed lattice at
class \(c\), \(\Gamma_c\) its finite stabilizer and \(v'\)
the refined level volume, so that the mass of this term is
\(v'/|\Gamma_c|\). Write \(P_c\) for the joint harmonic
polynomial of the coefficient vector at this class, of
degree \(d_c\). Let \(N_c\) be the prescribed odd
denominator prime to the marked level. In the covering
marked quotient \(Q_2\) let \(w_c(q)\) be the Fourier
weight of the actual finite test and let \(Y_{c,q}\)
be its marked lift after multiplication by \(N_c\)
in the refined lattice coordinates. This lift is
read by inverting \(N_c\) on the marked torsion module.
Multiplication by \(N_c\) is invertible there because
the level and \(N_c\) are coprime. Let \(Z_c\) be the
finite set of even representatives used in the projection
of the standard vacuum to \(\Lambda_c\). Write
\(z_c=|Z_c|\) and \(m_c\) for the index used to normalize
the short value. These are the two indices in
Section~\ref{sec:stage-conjugate-transport}.
\releasecite{Formalism/LocalToGlobal/ProjectiveLattices/ThClass.lean}{19}{The marked Fourier weights and the class rule}
\releasecite{Formalism/LocalToGlobal/ProjectiveLattices/ThClass.lean}{97}{The class-scalar normalization}.

In the fibre over a framed point, pull the theta section
back along the isogeny \([N_c]\) and translate its \(n\)
factors by the theta-group lifts prescribed by
\((Y_{c,q},u)\), \(u\in Z_c\).
Give the product of translates the algebraic
differential \(D_{P_c}\) at the zero section, using
the prescribed invariant \(E\)-differentials. The maps
of the abelian scheme, the finite theta-group action
and the finite trace are algebraic. In characteristic
zero, taking the indicated derivative is extraction
of a jet coefficient in the formal group, multiplied
by its factorial, so it is algebraic in the given
frames. The harmonic projection removes the quadratic
change of gauge. In fact, changing the \(E\)-linear
Hodge splitting changes the polarization connection
in the fibre by a linear differential form from the
polarization pairing. In logarithmic coordinates,
its change of gauge is \(\exp(Q(z))\); \(Q\) is the
sum of quadratic terms pairing a \(\tau\)-direction
with its conjugate direction, traced over the \(n\)
source coordinates. In the prescribed frames these
terms act on the polynomial by the contractions
\[
 \sum_{b=1}^{n}
     \partial_{\tau,i,b}\partial_{\bar\tau,i',b}.
\]
The formula for a jet in the changed gauge is obtained
by expanding the exponential. For any polynomial \(P\)
the extra terms are \(D_{Q(\partial)^jP}G(0)/j!\),
with \(j\ge1\). They are zero for the allowed joint
harmonics. This proves that the differential reads the
same holomorphic derivative in the prescribed fibre
coordinates once the constant in the theta-line frame
has been matched.
\releasecite{Formalism/LocalToGlobal/ProjectiveLattices/RowDeg.lean}{20}{The actual mixed contractions}
\releasecite{Formalism/LocalToGlobal/ProjectiveLattices/ThGerm.lean}{208}{The algebraic fibre derivative}
\releasecite{Development/AbelianVarietyTheory/Proofs/SiegelThetaTransformation/HarmonicConnectionValueEqPolynomialDerivativeStatement.lean}{14}{Comparison in the normalized chart}.

Let \(F_{c,q,u}\) denote the resulting product of
pulled-back translated theta sections before
differentiating, including the specified theta-group
phase. Write \(e_c\) for the source determinant
coefficient. On the primitive level the class section is
given in framed coordinates by
\begin{equation}
 \label{eq:primitive-phase-class}
 \sum_c
   \frac{v'e_c}{|\Gamma_c|\,m_cN_c^{d_c}|Q_2|}
   \sum_{q\in Q_2}w_c(q)
   \sum_{u\in Z_c} D_{P_c}F_{c,q,u}(0).
\end{equation}
The sums and weights are those of the actual class input,
so the expression is linear over \(K\). The theta-isogeny
formula identifies the pullback and the finite translations
with the standard Heisenberg expansion in the compatible
theta structures. It determines the relative term
coefficients up to the common factor of the frame.
\cite[\S1, Theorem~4]{mumford-theta-isogeny}
Multiplication by \(N_c\) in the fibre coordinate
introduces \(N_c^{d_c}\), which cancels its inverse in
\eqref{eq:primitive-phase-class}. The sum of the even
representatives is \(z_c\) times the projection to the
placed lattice vacuum by character orthogonality on the
quotient defining \(Z_c\). The remaining class index is
therefore \(z_c/m_c\), and the other two weights are
exactly the inverse Fourier normalization and the class
mass. This gives the same normalized derivative sum as in
the analytic lift. The theta section of the row fibre is the
same in all terms at the point in question.
\releasecite{Development/AbelianVarietyTheory/Proofs/SiegelThetaTransformation/ShortThetaValueComplexPointEqClassicalTheta.lean}{194}{The short value and its lattice projection}
\releasecite{Development/AbelianVarietyTheory/Proofs/SiegelThetaTransformation/ClassSumEvalComplexPointEqClassical.lean}{18}{Summation over actual source classes}.

We describe the vacuum and the framing in the reading of
\eqref{eq:primitive-phase-class}. On an analytic component
write \(L=p_\Lambda(\Lambda)\) for the period lattice of
the row and \(H\) for its positive Appell--Humbert form.
The imaginary part of \(H\) is the trace of the given
skew form and is unimodular on \(L\). The integral
half-root \(w_0+\bar w_0=1\) determines the semicharacter
for both the marked theta-line generator and the
standard oscillator vacuum. In a symplectic basis the
quasi-periodicity determines all Fourier coefficients
of the principal theta function from one coefficient.
This gives the one-dimensional fibre vacuum. Identify
the generator at a unit frame with its holomorphic
theta function \(\vartheta\), normalized by
\[
 \int_{\mathbb C^g/L}|\vartheta(z)|^2
          e^{-\pi H(z,z)}\,d\bar z=1,
\]
with probability Haar measure. The product of \(n\)
such generators reads \(C(x)\) times the classical
oscillator vacuum \(v_{\Lambda,x}\), with the Gaussian
factor removed, at the unit frame \(x\). The weighted
probability norm of \(\vartheta^{\boxtimes n}\) is also
one. In the self-dual model the classical vacuum's
weighted mean is independent of the unit frame. In
symplectic coordinates, orthogonality in the real half
leaves squared Gaussian coefficients on the other
half. Their sum unfolds to a Gaussian integral whose
covariance determinant cancels the Jacobian of the
period coordinates. Consequently \(|C(x)|\) is constant.
\releasecite{Development/AbelianVarietyTheory/Proofs/SiegelThetaTransformation/AnalyticFibreUnitFrameUniformisation.lean}{37}{The marked generator and its probability mean}
\releasecite{Development/AbelianVarietyTheory/Proofs/DivisorProofs/NormAdelicClassicalVacuumFactorEq.lean}{223}{Comparison of weighted and oscillator probability means}.

The phase of \(C(x)\) is constant as well. The marked
half is represented by the lattice \(L'=p_\Lambda(\mathfrak a\Lambda)\)
for \(\mathfrak a\bar{\mathfrak a}=(M)\). Since the trace
dual of \(\Lambda\) is \(\Lambda\), the trace dual of
\(\mathfrak a\Lambda\) for the form divided by \(M\) is
\(M\bar{\mathfrak a}^{-1}\Lambda=\mathfrak a\Lambda\).
On \(\mathbb C^g/L'\) the same theta function is a
section of the \(M\)-th power of a principal theta
line, up to a flat unitary character. Such a character
has an \(M\)-th root, so it can be included in the
principal line. Its \(M\)-th power is generated by
sections for \(M\ge2\). Indeed, translates of a
nonzero principal section by \(a_1,\ldots,a_M\) with
\(\sum a_j=0\) multiply to a section of that power by
the theorem of the square. One can choose the \(a_j\)
so that each factor is nonzero at any prescribed
point, by avoiding finitely many proper divisors in
the parameter space of the first \(M-1\) translates.
The space of sections of the power is the irreducible
finite theta representation of dimension \(M^g\).
For a symplectic decomposition of \(M\)-torsion, one
half acts with a full orbit of \(M^g\) distinct
characters, permuted transitively by the other half.
This proves irreducibility using the dimension from
Riemann--Roch for an ample line on an abelian variety.
The subspace of sections vanishing on all
\(M\)-division points is invariant under the theta
group and proper by generation at the origin,
so it is zero. There is thus a marked torsion value
of \(\vartheta\) which is nonzero at every given
fibre, with the torsion label allowed to vary.
The product can use this same label in its \(n\)
factors.
\cite[\S1]{mumford-theta-isogeny}
\releasecite{Development/AbelianVarietyTheory/Proofs/DivisorProofs/ExistsLevelOneThetaNeZeroOnIdealHalfLattice.lean}{132}{The half-lattice form and division points}.

Around a frame with a nonzero marked value, compare
the two powers at that torsion label. Both evaluations
are holomorphic in the domain coordinates, with the
same frame weight. For the theta series this follows
by writing the lattice sum in the holomorphic
Appell--Humbert gauge, where it converges normally on
compact subsets. The frame weight from the Gaussian
vacuum is the determinant power \(n/2\) in the
prescribed convention. For the algebraic evaluation
it is the corresponding Hodge determinant power.
Taking the quotient cancels this determinant even
when evaluating it in unit frames. The quotient is
\(C(x)\). It is holomorphic where the marked value
is nonzero and has constant modulus, so it is locally
constant. The domain of each component is connected,
and these neighborhoods cover it. This gives a single
nonzero constant \(C_C\) on each component \(C\).
\releasecite{Development/AbelianVarietyTheory/Proofs/DivisorProofs/PowGeneratorEqConstMulVacuumTheta.lean}{229}{Generator comparison in the unit frames}.

For the adelic reading choose a finite representative
placing the row lattice in \(C\). The transported
spherical vacuum reads \(\kappa_C v_{\Lambda,x}\) times
the prescribed determinant character \(\mu(g)\).
Indeed the finite lattice stabilizer acts on its
one-dimensional vacuum by the splitting character.
In the self-dual lattice realization the action before
splitting is pullback, fixing the defining lattice
vector. A change of finite representative for the
same lattice therefore multiplies the reading by the
same factor as \(\mu(g)\). The real action moves the
frame and has the determinant factor already in
the analytic weight. The probability inner products
agree under the self-dual lattice realization.
One can integrate the same Weil vector on the
adelic Heisenberg quotient, using the unit finite
lattice and a fundamental parallelotope at infinity.
The Gaussian-weighted real integral expressed in
period coordinates has probability measure and
is the product over the \(n\) source columns.
In particular \(|\kappa_C|=|C_C|\), so
\(\kappa_C/C_C\) is a unit scalar common to all inputs
and all points of \(C\).
\releasecite{Development/AbelianVarietyTheory/Proofs/DivisorProofs/ExistsVacuumMultiplierConstantPerLattice.lean}{225}{Adelic vacuum, stabilizer and real frame}
\releasecite{Development/AbelianVarietyTheory/Proofs/DivisorProofs/OscillatorModelValuesFirst.lean}{136}{The quotient of the two component scalars}.

The harmonic coefficient map and the normalized Fock
map are maps of the same definite source type to the
joint harmonic row type. They differ by a scalar
\(u_j\) of modulus one because their distinguished
extremal vectors have the same norm. Multiplicity one
here can be read from the extremal polynomial, a
product of flag minors in the positive and negative
blocks. In the adapted normalized frame a homogeneous
polynomial of degree \(d_j\) contributes the factor
\[
 u_j(-\delta_j)^{d_j}
       \bigl(\sqrt{2\pi\xi_R}\bigr)^{-d_j},
 \qquad \delta_j\in\{1,-1\}.
\]
The sign is that of the additive variable in the
short Fourier series, and the positive factor comes
from the normalized Gaussian derivatives. This factor,
the character \(\mu(g)\), and \(\kappa_C/C_C\) turn
\eqref{eq:primitive-phase-class} into the analytic
stage evaluation. Absorb the two unit scalars
into its component frame constant. At the primitive
tuple both the group element and the root of unity
outside this component constant are \(1\), giving
\eqref{eq:stage-base-phase-one}.
\releasecite{Development/AbelianVarietyTheory/Proofs/IwasawaDecomposition/ExistsCuspIsThClsDOfPre.lean}{272}{Unit scalar for the harmonic coefficient map}
\releasecite{Development/AbelianVarietyTheory/Proofs/DivisorProofs/ExistsAdelicLiftEqGeometricClassSumAllTuples.lean}{37}{Positive Lie factor and the short-series sign}.

The class formula is a section over the coefficient
field of the marked model. To see the gluing, change
a marking in the framed model. The finite Fourier
weights transform by the inverse character to the
theta-group action on the torsion sum, and the
homogeneous derivatives transform by the specified
weight of the Hodge frame. Thus the expressions
on overlaps satisfy the transition rule of the
automorphic bundle. Equivalently, on the fibre product
of the frame charts the two readings are equal as
sections of the pulled-back bundle. The charts are
smooth in characteristic zero and equality may be
tested on complex points. Faithfully flat descent
gives the section \(S_x\) over \(K\).
\releasecite{Development/AbelianVarietyTheory/Proofs/SiegelThetaTransformation/ExistsSectionOneClassSeriesDescends.lean}{22}{Descent for one normalized class series}.

The sections are cuspidal. For an isotropic row line,
choose a symplectic basis adapted to it and use the
Fourier series in the central line-unipotent coordinate.
In a theta term its frequencies, one at each real
embedding, are sums of squares of the leading row
entries of the definite-source columns, including the
fixed shifts. They vanish simultaneously only when
all these entries are zero. In the positive row
block the variables then lie in a hyperplane. Every
polynomial in the specified highest-weight component
is a sum of products of flag minors, and each product
contains a maximal positive-block minor. The latter
follows from positivity of its last weight. Indeed
for the \(j\)-th stage the positive block has size
\(k+j\), and its last polynomial weight in the
given convention is
\(\lfloor-b_{\tau,j}+e_{\tau,j}-(t+1)/2\rfloor\ge1\)
by the prescribed gap at \(j\). The maximal minor
vanishes on the hyperplane. Therefore the zero
frequency vanishes in every normalized class sum.
The same Fourier calculation applies to every
rational isotropic line. Any proper rational
parabolic stabilizes an isotropic flag. Integrate its
constant term first over the centre of the line
unipotent to get zero. At an anisotropic row group
there are no proper rational parabolics.
\releasecite{Development/AbelianVarietyTheory/Proofs/DivisorProofs/KapTopLastPosRun.lean}{16}{Last positive weight at the actual stage}
\releasecite{Development/AbelianVarietyTheory/Proofs/SiegelThetaTransformation/ExistsExpansionOfRowHighest.lean}{46}{Adapted line expansion and the maximal minors}.

On a toroidal compactification of the neat PEL model the
cusp sections are global sections of the subcanonical
extension of the automorphic bundle. The space is
finite-dimensional over \(K\), and extension of
scalars commutes with its sections, by computing
sections as the kernel of the restriction maps of
a finite affine cover. Thus the class values give
the \(K\)-linear coordinate map on the cusp basis.
The lowest-type vector is cyclic, so equality of
its evaluations gives equality of the multiplicity maps.
\href{https://msp.org/ant/2014/8-8/ant-v8-n8-p01-s.pdf}{Lan--Stroh, subcanonical cusp bundles on PEL models.}
\releasecite{Development/AbelianVarietyTheory/Proofs/IwasawaDecomposition/ExistsCoreConstantsBaseSolutionEquivariant.lean}{57}{The component constants at phase one}
\releasecite{Development/AbelianVarietyTheory/Proofs/IwasawaDecomposition/ExistsConstantsThetaComparisonAllTuples.lean}{227}{Comparison on the base class sections}
\releasecite{Development/AbelianVarietyTheory/Proofs/IwasawaDecomposition/GeometricHalfCoreFormsUniformisationIntegral.lean}{60}{The cusp-space realization}.

The construction uses the PEL framing whose conjugate
of the lattice module is the placed row lattice,
with the prescribed polarization, marked module
and coefficient isometry. These identifications are
needed at every class representative of the selected
stage, with the same phase-one choices before the
conjugate comparison.

\paragraph{The base reading on a marked short section.}
\teamanchor{proof.stage-marked-read}
In the common construction of the base and enlarged tuple data,
retain the vacuum comparison \(c_i(C)\) before
multiplying by the harmonic phase. The base marking is the
identity cover, so its determinant section is constant. Write
\(c_{0,i}\) for its value in the comparison frame. The common
comparison multiplies the vacuum factor by
\(k_i(-\delta_i)^{d_i}/c_{0,i}\), using the harmonic map and
additive sign of \eqref{eq:primitive-phase-class}. The
coefficient-field extension agrees with the original harmonic
map on the base coefficients. The construction installs this
product as the base component factor \(c_{E,i}(C)\) of the
row model. In particular, at the distinguished embedding and
the chosen lowest-type vector, the reading of a section \(f\)
in that model is
\begin{equation}
 \mathcal R_i(f)(g)=L_i\,
   \frac{k_i(-\delta_i)^{d_i}}{c_{0,i}}\,
   c_i(C(g))\,\mu_i(g)\operatorname{rd}_{g}(f).
 \label{eq:stage-marked-read}
\end{equation}
Here \(L_i=(2\pi\xi_R)^{-d_i/2}\), \(C(g)\) is the component
of the framed dictionary point, and \(\operatorname{rd}_g\)
uses its Hodge coordinates. The formula is the definition
of the installed reading on the fixed principal cusp space.
The component factor together with \(\mu_i(g)\) is
\(c_i^{\rm fr}(g)\) above. The class calculation identifies
these factors before choosing a finite class input. For either
prescribed level cut, first apply the corresponding finite
character projector in this same cusp space.
\releasecite{Development/AbelianVarietyTheory/Proofs/IwasawaDecomposition/ExistsConstantsThetaComparisonAllTuples.lean}{275}{Base restriction of the common datum}
\releasecite{Development/AbelianVarietyTheory/Proofs/DivisorProofs/ExistsAdelicLiftEqGeometricClassSumAllTuples.lean}{185}{The phase multiplying the base vacuum comparison}
\releasecite{Development/AbelianVarietyTheory/Proofs/IwasawaDecomposition/ExistsCoreConstantsBaseSolutionEquivariant.lean}{65}{Installation on the base components}
\releasecite{Development/AbelianVarietyTheory/Proofs/DivisorProofs/ShortSeriesBaseDatumAdelicFormMultiple.lean}{28}{Reading in the installed cusp space}.

Use this same formula for the short section on the run's
marked cover at \(M p_o\) after descent to its base weight
space. The level map pulls back the framed Hodge coordinates;
the cover coefficient-field isomorphism carries the
coefficients of the finite harmonic polynomial in these same
coordinates. Its complex embedding restricts to the
distinguished embedding of the base field. Equivariance here
follows by applying the coefficient map to the polynomial,
the finite Fourier weights and the theta-group phases in the
same finite trace. The adapted embedding fixes the small
field containing the deletion character at \(p_o\), so the
reference character and its marking stay fixed while the
extra ray-character values are read through the adapted
embedding. The joint algebraic determinant coefficient on
this lattice is the one in \eqref{eq:stage-determinant-ideal-reading}.
\releasecite{Development/AbelianVarietyTheory/Proofs/DivisorProofs/ShortValueAtPoint.lean}{42}{The coefficient map and the marked fibre of the short rule}
\releasecite{Development/AbelianVarietyTheory/Proofs/DivisorProofs/AdaptedExistence.lean}{163}{The adapted embedding fixing the small field}.

\teamanchor{proof.stage-conjugate-transport}
\subsubsection{Transport of the phase-one base stage}\label{sec:stage-conjugate-transport}
Fix one of the primitive base-tuple stages. Use the same
definite source \(V\), row space \(W\), principal
projection \(P_\sigma\), and degree and coefficient fields
as in the stage norm calculation. Write \(\iota\) for the
chosen coefficient embedding, \(k_{\rm CM}\) for the host
imaginary quadratic subfield, and \(K\) for the row
coefficient field. The embeddings over the chosen inclusion
of \(k_{\rm CM}\) give the holomorphic framed models in the
released comparison. They permute the embeddings of \(E\)
in the chosen CM type, since that type consists of the
extensions of this inclusion. The signatures and hence
the positive hook sum are preserved under the permutation.
The embedding \(\sigma^+=\overline{\sigma}\) over the
chosen inclusion is the partner of any \(\sigma\) whose
restriction to \(k_{\rm CM}\) is the opposite one.
\releasecite{Development/AbelianVarietyTheory/BlockTheory/DivisorBase/WeilGramSignature.lean}{18}{Embeddings over the host field}
\releasecite{Development/AbelianVarietyTheory/Proofs/DivisorProofs/CoefMapNormEquivariantBasisOfFive.lean}{148}{The geometric embedding condition}.

The fixed-basis comparison uses the algebraic cusp
sections of the framed unitary moduli space of the row
in this stage. Use the chosen weight and geometric cusp
basis. For each rational class input \(x\), set \(S_x\)
to be its algebraic theta class section, a linear combination
\(\sum_\ell a_\ell s_\ell\) of the fixed cusp sections
after base change to \(K\). Both the action of the finite
level group \(G\) on sections and this decomposition are
over the field. The character of a level move in the
adelic row comparison is denoted \(\chi_W\). It is the
product of the prescribed residue-determinant characters,
so has finite image in \(K^\times\).
The geometric row model reads a section \(s_\ell\) at the
fixed unit lowest-type vector as a function \(J_\sigma(s_\ell)\),
including the specified adelic and weight factors, in the
cuspidal multiplicity model. The normalized base class
identity of the phase-one realization is
\begin{equation}
\label{eq:stage-base-phase-one}
 P_\iota T_\iota j_\iota(x)
   =\iota(\zeta_0)P_\iota j_{W,\iota}(A_{z_0}S_x).
\end{equation}
Here \(A\) is the geometric level action, \(z_0\) a fixed
central element of \(G\), and \(\zeta_0\in K\) a root of
unity. One may take both to be one in the phase-one base
choice which absorbs the common vacuum phase into its base
component factors. This comparison starts with the framed cusp sections, their
linear decomposition over \(K\), and the normalized phase-one
class-section formula on the actual source classes. In that
formula the theta generator, differential coordinates, cover
marking and term volumes are part of the realization. The equality at
the lowest-type vector determines the map because the
vector is cyclic in the real representations and the
multiplicity maps are equivariant.
\releasecite{Formalism/LocalToGlobal/ProjectiveLattices/ThClass.lean}{184}{The algebraic class section and its base comparison}
\releasecite{Development/AbelianVarietyTheory/Proofs/IwasawaDecomposition/ExistsConstantsThetaComparisonAllTuples.lean}{227}{The normalized base choice}
\releasecite{Development/AbelianVarietyTheory/Proofs/DivisorProofs/ConjugateLiftUniqueAndBaseDatum.lean}{111}{The phase-one base reading}.

We compare the same section with the actual stage lift
after an embedding \(\sigma\) over \(k_{\rm CM}\).
Let \(M\) be the row marking modulus, \(M_2\) the covering
marking modulus, with \(M\mid M_2\), and \(Q_2\) the finite
quotient of the marked half at level \(M_2\), taken with
one copy per source coordinate. Since the auxiliary base
twist is trivial, its twisting support is empty. The
finite base test is the characteristic function \(\phi\)
of the identity leading block and the zero complementary
block. In particular its values are \(0,1\). Let
\(w_c(q)\), \(q\in Q_2\), be its finite Fourier
coefficients on the \(c\)-th source lattice class, with
the inverse transform normalized by \(|Q_2|^{-1}\).
The pairings in the Fourier matrix have values in
\(\mu_{M_2}\). An embedding of the algebraic class
coefficient field has
\(\sigma(\zeta_{M_2})=\zeta_{M_2}^{m_\sigma}\) with
\((m_\sigma,M_2)=1\). Extending \(\sigma\) to the covering
coefficient field if necessary, the same formula holds
for the chosen Fourier matrix. Consequently
\[
 \sigma(w_c(q))=w_c(m_\sigma q)
\]
for every \(c,q\), with one multiplier \(m_\sigma\) for
all classes. The actual comparison of the finite
additive character with the dual marking has one sign
\(\epsilon\in\{\pm1\}\). With the sign included in the
first marking, replacing the labels \(q\) by \(m_\sigma q\)
is exactly the change to the marking \(\epsilon m_\sigma\)
times the original marking. Indeed, the rational roots
at label \(q\) are specified by a congruence modulo the
lattice. Taking a representative of \(m_\sigma q\) in
the marked half is equivalent to multiplying a
representative of \(q\) by \(m_\sigma\), modulo \(M_2\)
on that half. The difference yields an integral
lattice translation. The theta-group lifts use the
same normalized signs, so the resultant sum is the
reindexing of the finite translations. Multiplication
by \(m_\sigma\) is a bijection of \(Q_2\), with inverse
given by a B\'ezout coefficient.
\releasecite{Development/AbelianVarietyTheory/Proofs/DivisorProofs/CoefMapNormEquivariantBasisOfFive.lean}{2008}{Actual finite Fourier average and marking sign}
\releasecite{Development/AbelianVarietyTheory/Proofs/DivisorProofs/CoefMapNormEquivariantBasisOfFive.lean}{2428}{The common Galois Fourier dilation}
\releasecite{Development/AbelianVarietyTheory/Proofs/DivisorProofs/MarkingRootsDilationCoprimeLevel.lean}{56}{Dilation on the finite marking quotient}.

The chosen half of the row level has ideal
\(\mathfrak m_b=\prod_{w\in P_g}w\), with
\(\mathfrak m_b\overline{\mathfrak m_b}=M\mathcal O_E\).
The marked submodule before reduction is
\(\overline{\mathfrak m_b}\) times the row lattice.
At each chosen type prime send the residue of
\(\epsilon m_\sigma\) to the scalar matrix in
\({\rm GL}_r(k_w)\). At the conjugate prime the unitary
condition uses the inverse conjugate matrix.
These scalars define a central element \(z_\sigma\) in
the prescribed full product \(G\) of residue groups.
Its action on the marked submodule agrees with
multiplication by \(\epsilon m_\sigma\) modulo \(M\).
The extra marking at \(M_2\) can be lifted after this
move because \((m_\sigma,M_2)=1\), using the compatible
lift of the dilated marking just computed. In particular
the element of \(G\) is independent of the source class
and of the individual torsion label. The central element
lifts adelically to a scalar on the selected split
places, with inverse conjugate on the other half and
the identity elsewhere. Denote its right translation
on row forms by \(R_\sigma(z_\sigma)\).
\releasecite{Development/AbelianVarietyTheory/Proofs/DivisorProofs/CoefMapNormEquivariantBasisOfFive.lean}{2483}{The central level move}
\releasecite{Development/AbelianVarietyTheory/Proofs/DivisorProofs/CoefMapNormEquivariantBasisOfFive.lean}{2591}{The compatible lifted marking}
\releasecite{Development/AbelianVarietyTheory/Proofs/DivisorProofs/CentralShiftPhaseConjReadingsFiveClauses.lean}{19}{The central adelic lift}.

The analytic and geometric evaluations now use the same
finite class expression. Here are the scalar checks on
it. Take the same source representatives \(h_c\) for an
open subgroup \(J'\subset J_V\), with class stabilizers
\(\Gamma_c\). Passing to the finer representatives gives
\({\rm vol}=1/[J_V:J']\); the class weights
\({\rm vol}/|\Gamma_c|\) are rational. Let \(\Lambda_c\)
be the conjugate of the corresponding finite lattice
in the oscillator frames, \(I_c\) the index of its
intersection with the standard lattice in \(\Lambda_c\),
and \(N_c\) the denominator in the class construction.
Use \(m_c,z_c\) for the two integer cardinalities in the
class term \(\mathrm{mIdx}(N_c,r),\mathrm{zCard}(N_c,r)\).
In the finite lattice projection \(z_c=m_c I_c^r\).
This is the index formula on the inclusion
\(N_c\mathcal O_E^n\subset\Lambda_c\cap\mathcal O_E^n
 \subset\Lambda_c\), applied in the row variables.
The \(m_c\) normalizes the kernel average and \(z_c\)
counts the translated terms. The differentiation of a
test of total degree \(d\) on the rescaled fibre
introduces \(N_c^d\), which cancels the compensating
factor \(N_c^{-d}\) in the class definition. In particular
the holomorphic evaluation of \(S_x\) on the moved
framed point has the same polynomial derivatives of
the same \(|Q_2|^{-1}\)-weighted translates as the
analytic class sum. This is the theta pullback for
compatible theta structures, followed by the finite
trace and the harmonic derivative in the marked
coordinates. One can compute it by first expanding the
finite theta-group projection as a sum of translations
of the generator. The sum over the even representatives
is \(z_c\) times the normalized lattice projection.
Linearity of the derivative then leaves the above
cardinalities and a common factor \(\kappa'(g)\)
from the \(n\)-fold theta generator on the row fibre.
The generator is nonzero, and its conversion to the
oscillator vacuum is an equality of theta lines at the
same point. Thus \(\kappa'(g)\) is common to all classes,
polynomials and torsion terms at \(g\).
\cite[\S1, Theorem~4]{mumford-theta-isogeny}
\releasecite{Formalism/LocalToGlobal/IndefiniteReal/ThPinsR.lean}{73}{Finer source class volume}
\releasecite{Development/AbelianVarietyTheory/Definitions/SpanningFamilies/DatEqZero.lean}{19}{The exact finite derivatives and their coefficients}
\releasecite{Development/AbelianVarietyTheory/Proofs/DivisorProofs/OscillatorModelValuesThird.lean}{293}{The index ratio in the short theta value}
\releasecite{Development/AbelianVarietyTheory/Proofs/SiegelThetaTransformation/ShortThetaValueComplexPointEqClassicalTheta.lean}{194}{Trace and harmonic differential in the theta comparison}
\releasecite{Development/AbelianVarietyTheory/Proofs/SiegelThetaTransformation/ClassSumEvalComplexPointEqClassical.lean}{18}{Evaluation of the class section}.

The remaining class coefficient is the determinant
character of the source lattice move times a positive
square root. To transport it, let
\(\mathfrak d_c/\overline{\mathfrak d_c}\) be the determinant
ideal of the move, with \(\mathfrak d_c\) and its conjugate
coprime. This is obtained by taking the positive
valuations of the determinant; the opposite valuations
are their negatives by unitarity. Write \(y_c\) for the
specified half determinant id\`ele, and write
\(\mu_\sigma,\eta_\sigma\) for the placed splitting
characters. At the supporting primes outside the type
primes one transports
\(\eta(\varpi_w)(\mathrm Nw)^{-1/2}\), using the
fixed algebraic readings of this product. Indeed the
product over a principal ideal for the \(\eta\)-reading
is the inverse of the product of the embeddings of its
generator in the compatible CM type, by the half
exponent of \(\eta\). The conjugate embeddings transport
this principal-ideal identity. Restoring the positive
norm factors yields the unramified placed character
\(\eta_\sigma\). Its restriction to \(F\) is preserved,
since the norm-normalized restriction at a split pair
or an inert prime is rational. Thus the product that
transports at \(y_c\) is
\(\eta(y_c)^{-1}\sqrt{\mathrm N\mathfrak d_c}\), with
the character and positive square root taken together.
For \(\mu\), the prescribed readings on units and
uniformizers are algebraic. Conjugating them and
permuting the integral CM infinity exponents gives
the placed character \(\mu_\sigma\). The character is
trivial on principal id\`eles because their infinity
factor is an algebraic product of paired embeddings
with the prescribed exponents. The conjugated
finite readings on units are roots of unity, and
the values on uniformizers have absolute value one
by raising their ideals to principal ideals, with
the unit factors at the type primes included.

In the prescribed class coefficient \(B(\mathfrak d_c)\)
these readings give
\begin{equation}
 \sigma(B(\mathfrak d_c))
   =(\mu_\sigma\eta_\sigma^{-1})(y_c)\sqrt{\mathrm N\mathfrak d_c}.
 \label{eq:stage-determinant-ideal-reading}
\end{equation}
The coefficient is defined by multiplying the corresponding
readings at prime ideals. This applies to every nonzero
determinant ideal prime to the bad modulus of the base datum.
It therefore applies to the determinant ideal of an adapted
short-section lattice supported off that modulus. The
own-prime correction on its marked cover contributes the
separate value of the deletion character on that ideal.
\releasecite{Development/AbelianVarietyTheory/Proofs/IwasawaDecomposition/ExistsReferenceTupleThetaComparisonNormalisedDat.lean}{91}{The joint reading on ideals prime to the bad modulus}
\releasecite{Development/AbelianVarietyTheory/Proofs/DivisorProofs/TwistScalarMemRange.lean}{124}{The same base reading for the short rule}.
The scalar in the primitive class section is
\(e_c=q_c/B(\mathfrak d_c)\), with \(q_c\in\mathbf Q_{>0}\)
and \(q_c^2 I_c^r=\mathrm N\mathfrak d_c\).
Therefore
\[
 \sigma(e_c)\,\frac{z_c}{m_c}
   =(\mu_\sigma\eta_\sigma^{-1})(y_c)^{-1} I_c^{r/2}.
\]
The square root \(I_c^{r/2}\) is exactly the conversion
between the projection of a unit vacuum and the
unit vacuum of the moved tensor lattice. The formula
fixes the scalar of each source class, including when
the index has odd valuation.
\releasecite{Formalism/LocalToGlobal/ProjectiveLattices/ThClass.lean}{97}{The class-scalar normalization}
\releasecite{Development/AbelianVarietyTheory/Proofs/DivisorProofs/EmbeddingConjugateOscillatorSixth.lean}{185}{Determinant ideal and the joint square-root reading}
\releasecite{Development/AbelianVarietyTheory/Proofs/DivisorProofs/CoefMapNormEquivariantBasisOfFive.lean}{2202}{The transported class coefficient}.

There are three remaining unit factors. First,
in the Fock coefficient map the fixed rational highest
vector goes to the prescribed minor product. In the
comparison at \(\sigma\) both coefficient highest
vectors are unit for the standard Hermitian form, and
the same minor product is transported by relabelling
its variables. The Cartan-frame change on the definite
side is unitary. Thus the two maps have a scalar of
absolute value one on the highest line. Since the
coefficient type is irreducible and the polynomial
maps intertwine its action, this scalar \(\alpha_\sigma\)
applies to the whole coefficient space. Second, the
cover factor on the base tuple is a constant
\(\gamma_\sigma\), the embedding of the prescribed root
of unity. Indeed its values on the marked framed
points are the root of unity times the auxiliary twist
character on the determinant, which equals one on the
base tuple. These points uniformize the complex charts,
so the cover section is that constant after base
change. Third, on a component \(c\) of the row charts
the two normalized vacua differ by a factor
\(c_{\rm fr,\sigma}(c)\) times the specified determinant
automorphy factor. In both descriptions the Gaussian
squared norm on the lattice torus has mean one. The
self-dual lattice gives a common one-dimensional
theta space, hence \(|c_{\rm fr,\sigma}(c)|=1\) at one
frame of the component. For another frame in the same
component, move first by a real unitary operator and
then by the finite self-dual lattice stabilizer. Their
oscillator multipliers are the determinant characters
appearing in the adelic automorphy. The comparison
with the transported framing therefore uses the same
\(c_{\rm fr,\sigma}(c)\). In the geometric reading the
fixed phase is \(c_{E,\sigma}(c)\), also of norm one.
Their quotient gives a function \(u_\sigma(g)\) of
absolute value one which depends only on the chart
component. It is unchanged under the left rational
group, the right real group, and the right principal
level. For any unipotent radical \(N\) and any \(g\),
additive approximation along a unipotent filtration
gives density of \(N(F)\) in \(N(\mathbb A_{F,f})\).
Thus for \(n\in N(\mathbb A_F)\) one can choose
\(\gamma\in N(F)\) with \(g^{-1}\gamma^{-1}ng\) in the
right real group times the principal finite level.
It follows that \(u_\sigma(ng)=u_\sigma(g)\), so
multiplication by \(u_\sigma\) preserves the vanishing
of the cuspidal constant terms. This multiplication
is unitary, since Petersson integration is unchanged. The multiplier
also passes through the unipotent averages. For a rational
unipotent radical \(N\), given \(n\in N(\mathbb A_F)\)
and \(g\), choose \(\gamma\in N(F)\) so that
\(g_f^{-1}\gamma_f^{-1}n_f g_f\) belongs to the principal
level. Such a \(\gamma\) exists by additive approximation
along a split unipotent filtration. The left rational,
right real and right principal invariance of the multiplier
then gives \(u_\sigma(ng)=u_\sigma(g)\).
\releasecite{Development/AbelianVarietyTheory/Proofs/DivisorProofs/CoefMapNormEquivariantBasisOfFive.lean}{827}{Fixed highest vector and unit Fock factor}
\releasecite{Development/AbelianVarietyTheory/Proofs/IwasawaDecomposition/OneMapTypingsEquiv.lean}{180}{The root of unity on the base cover}
\releasecite{Development/AbelianVarietyTheory/Proofs/IwasawaDecomposition/OmegaFormScalarExtension.lean}{162}{Vacuum movement by the determinant multiplier}
\releasecite{Development/AbelianVarietyTheory/Proofs/IwasawaDecomposition/OmegaFormScalarExtension.lean}{246}{A unit vacuum phase on each component}.

In detail, let \(D_{\sigma,c,x}\) denote the harmonic
derivative with the transported polynomial of the class
value \(x_c\), and let \(v_c\) be the projection of the
normalized real vacuum to the class tensor lattice.
Write \(\mathcal H(Y)\) for the normalized Heisenberg
translation by the marked root \(Y\). At the framed
point of \(g\) the common expression is
\[
 C_\sigma(x,g)=\sum_c
   \sigma\!\left(\frac{{\rm vol}\,e_c}{|\Gamma_c|}\right)
   \frac{z_c}{m_c |Q_2|}\,
   D_{\sigma,c,x}\!
    \left[\sum_{q\in Q_2}w_c(q)\mathcal H(Y_c(q))v_c\right].
\]
On the adelic side, evaluate \(\alpha_\sigma T_\sigma
j_\sigma(x)\) at the unit vector. This gives
\[
 c_{\rm fr,\sigma}m_W(g)\ell\,\kappa'(g) C_\sigma(x,g),
\]
with \(\ell\) the prescribed real derivative conversion
and \(m_W\) the determinant automorphy factor at
\(\sigma\). On the geometric side the
value of \(A_{z_\sigma}S_x\) at the frame is
\(\gamma_\sigma\kappa'(g)C_\sigma(x,g)\).
The adelic reading \(j_{W,\sigma}\) multiplies the last
value by \(c_{E,\sigma}m_W(g)\ell\).
Using \(u_\sigma=c_{\rm fr,\sigma}/c_{E,\sigma}\)
therefore gives
\begin{equation}
\label{eq:stage-conjugate-fixed-basis}
 \alpha_\sigma\gamma_\sigma T_\sigma j_\sigma(x)
     =U_\sigma j_{W,\sigma}(A_{z_\sigma}S_x).
\end{equation}
Here the equality is promoted from evaluated functions
by the injectivity of lowest-type evaluation. The use
of the same coefficient vectors \(x_c\) on both sides
is through the fixed algebraic coefficient basis and
its \(\sigma\)-linear realization, so the arithmetic
input on the right remains \(S_x=\sum a_\ell s_\ell\).
The common section factors are \(\sigma(a_\ell)\)
after this realization. The level action is read with
the same \(U_\sigma\). In fact a level move \(a\in G\)
changes a section frame by the analytic theta-line
scalar at that move raised to \(n\). The vacuum
comparison yields exactly the ratio of component
constants before and after that frame change. The
theta-line scalar cancels in this ratio and the
remaining multiplier is the placed determinant value
\(\sigma(\chi_W(a))\). Equivalently, writing
\(B_\sigma=U_\sigma j_{W,\sigma}\), one has
\[
 B_\sigma A_a
   =\sigma(\chi_W(a)) R_\sigma(a) B_\sigma .
\]
Both identities use the same marked section
evaluations. The chosen full conjugate-character lift
may use a different orthonormal realization of its
coefficient space. The two realizations are related
by the isometry obtained from their common Gram matrix.
The normalized oscillator vacua are on the same line
and the joint-harmonic couplings agree up to a single
unit scalar by their normalized extremal vectors.
Thus this replacement affects
\eqref{eq:stage-conjugate-fixed-basis} only through a
fixed unit scalar.
\releasecite{Development/AbelianVarietyTheory/Proofs/DivisorProofs/CoefMapNormEquivariantBasisOfFive.lean}{1244}{Level action and cancellation of the theta-line scalar}
\releasecite{Development/AbelianVarietyTheory/Proofs/DivisorProofs/CoefMapNormEquivariantBasisOfFive.lean}{2651}{Equality of the two framed class sums}
\releasecite{Development/AbelianVarietyTheory/Proofs/DivisorProofs/CoefMapNormEquivariantBasisOfFive.lean}{2760}{Fixed-basis identity at the chosen lowest-type vector}
\releasecite{Development/AbelianVarietyTheory/Proofs/DivisorProofs/ConjugateLiftUniqueAndBaseDatum.lean}{160}{Unitary independence of the conjugate model}.

Let \(e=\sum_{a\in G}e_a a\) be the \(K\)-rational group
idempotent defining the principal row projection on
the fitted summands, so
\(P_\sigma=\sum_a\sigma(e_a)R_\sigma(a)\). On the
geometric space put
\(P_K=\sum_a e_a\chi_W(a)^{-1}A_a\).
The level law gives \(B_\sigma P_K=P_\sigma B_\sigma\).
Define the stage map over \(K\) by
\(\theta_K(x)=\zeta_0 P_K A_{z_0}S_x\).
By \eqref{eq:stage-base-phase-one} this realizes the
actual projected map at \(\iota\), and it is unique
because the row realizations preserve independence.
Write \(z=z_0 z_\sigma^{-1}\in Z(G)\). Then
\[
\begin{split}
 U_\sigma j_{W,\sigma}\theta_K(x)
  &= \sigma(\zeta_0)P_\sigma B_\sigma A_{z_0}S_x \\
  &=\sigma(\zeta_0\chi_W(z))\alpha_\sigma\gamma_\sigma\,
             R_\sigma(z)P_\sigma T_\sigma j_\sigma(x).
\end{split}
\]
The last coefficient has absolute value one.
Right translation by \(z\) is an isometry for the
fixed Petersson measure and commutes with the
projection because \(z\) is central. This proves
the required equality of target Gram forms on all
rational inputs. The conjugate coefficient source
realizations have their class Gram forms by the
common coefficient pairing, with the rational class
weights computed above. Scalar extension of
\(\theta_K\), followed by the indicated invertible
unitary transformations, shows that the kernel of
\(P_\sigma T_\sigma\) is the complex span of the
realized rational kernel. Thus the fitted quotient
is computed using the same kernel complement.
\releasecite{Development/AbelianVarietyTheory/Proofs/DivisorProofs/ConjugatesOfRationalThetaLiftFromCharacterisatio.lean}{80}{Rational projector and central-translated comparison}
\releasecite{Development/AbelianVarietyTheory/Proofs/DivisorProofs/ConjugatesOfRationalThetaLiftFromCharacterisatio.lean}{227}{All embeddings of the phase-one stage}
\releasecite{Development/AbelianVarietyTheory/Proofs/DivisorProofs/LowFormOfEmb.lean}{44}{The complexified kernel}.

At an embedding opposite on \(k_{\rm CM}\), complex
conjugation of the real unitary representation at
\(\sigma^+\) replaces the highest type by the opposite
type. It gives the antiunitary mate of the chosen
unit lowest-type vector. The principal finite character
is trivial, so taking complex conjugates of its
automorphic values preserves its level and cusp
conditions. For the fixed cusp basis the readings
at the two embeddings are conjugate under this
identification. This is the conjugate reading used
in forming the rational row model. The multiplicity
Gram is Petersson evaluation at the unit vector
times the product of compact type dimensions.
Those dimensions agree under complex conjugation,
so the Gram forms are conjugate on basis vectors
and hence on the entire rational row. The rational
class-sum source forms are conjugate as well.
Finally the comparison of principal and character
volumes divides both the source quotient squared
norm and the row squared norm by the same \(|P|\).
Thus the squared singular value on the fitted
summand is measured by the stage norm estimate at
\(\sigma^+\), including at the embedding opposite
on \(k_{\rm CM}\).
\releasecite{Development/AbelianVarietyTheory/Proofs/DivisorProofs/HilbertModelConjugationAntilinearBijection.lean}{39}{Opposite embedding and principal level}
\releasecite{Development/AbelianVarietyTheory/Proofs/DivisorProofs/CentralShiftPhaseConjReadingsFiveClauses.lean}{52}{Conjugate Gram on the fixed basis}
\releasecite{Development/AbelianVarietyTheory/Proofs/DivisorProofs/RationalLiftLowerBoundOffBasePlace.lean}{239}{Both principal-to-character rescalings}.

\paragraph{The least stage factor.}
The isometries in this transport identify the rational kernel
complement with that of the actual projected lift and
carry the class-sum and Petersson forms to the prescribed
metrics. Apply \eqref{eq:stage-petersson-lower-bound} in the
model over the chosen CM type and use the conjugate
metric at its partner. At the character-level comparison
the principal factor divides both squared norms.
On a nonzero fitted row summand insert a vector of nonzero
quotient norm into the stage Gram equality, and divide by
that norm. One obtains
\[
 \log\lambda_{j,p,\sigma}\ \ge\ L_{6,j},\qquad
 L_{6,j}=\sum_\rho\log H^+_{j,\rho}-C_{\rm sz}fn^2L',
 \qquad C_{\rm sz}\ge8 .
\]
Empty summands contribute no constraint to the least
eigenvalue of the actual quotient. The inequality is uniform
in the fitted label \(p\) and the coefficient embedding
\(\sigma\). Consequently taking the minimum of the sums of
stage logs over fitted labels and averaging over all
coefficient embeddings preserves the lower bound
\(\sum_jL_{6,j}\) for the period-map height.
\releasecite{Development/AbelianVarietyTheory/Proofs/DivisorProofs/LiftAdjointEigenvalueLowerBoundBaseEmbedding.lean}{160}{Base projected lift}
\releasecite{Development/AbelianVarietyTheory/Proofs/DivisorProofs/RationalLiftLowerBoundOffBasePlace.lean}{196}{Transport of the conjugate row norm}
\releasecite{Development/AbelianVarietyTheory/Counterexamples/RotationExamples/PhaseBreaksTargetParts.lean}{29}{Unitary comparison of the two Gram forms}
\releasecite{Development/AbelianVarietyTheory/Proofs/DivisorProofs/LogLiftEigenvalueLowerBoundOtherPlaces.lean}{46}{Fitted eigenvalues away from the base place}
\releasecite{Development/AbelianVarietyTheory/Proofs/DivisorProofs/HookLogSumSubSizeTermLeLowerBound.lean}{78}{Base hook term and the common loss}.

\teamanchor{proof.rational-period-map}
\subsubsection{The normalized theta germs}
\label{sec:period-germs}
Let \(k_{\rm CM}\) be the host imaginary quadratic field, with
discriminant \(-d_{\rm CM}\), and use its chosen complex embedding.
The host curve ideal in the chosen family is the unit ideal, so
the period lattice of the host elliptic curve is
\(\Lambda=\delta\mathcal O_{k_{\rm CM}}\), with
\(\delta^{12}=\Delta(\mathcal O_{k_{\rm CM}})\) in Weierstrass
normalization, where \(\Delta=g_2^3-27g_3^2\).
Scaling gives \(\Delta(\Lambda)=1\). The CM \(j\)-invariant
is algebraic, so \(g_2(\Lambda)^3=j/1728\) and
\(27g_3(\Lambda)^2=g_2(\Lambda)^3-1\) are algebraic.
The invariant differential \(dx/y\) on
\(y^2=4x^3-g_2(\Lambda)x-g_3(\Lambda)\) pulls back to \(dz\).
This realizes the actual lattice, differential and
complex coordinate on an algebraic CM curve.

Put \(A=\operatorname{area}(\Lambda)\).
The reduced odd theta function of the divisor \([0]\) is
normalized by \(\theta'(0)=1\). Its multiplier on \(\gamma\in\Lambda\)
is the sign \(1\) for \(\gamma\in 2\Lambda\) and \(-1\) otherwise,
times \(\exp(\pi(z\bar\gamma+|\gamma|^2/2)/A)\).
The entire function with this multiplier, divisor and derivative
is unique, since the quotient of two such functions is elliptic
and entire. This is the odd theta function used in the vacuum.
Bannai--Kobayashi prove the following special case. For a CM
elliptic curve over a number field with the period coordinate
of its algebraic invariant differential, the reduced theta for
\([0]\) normalized by derivative one has algebraic Taylor
coefficients. Its reduced translate
\[
 U_a\theta(z)=
     e^{-\pi(z\bar a+|a|^2/2)/A}\theta(z+a)
\]
has algebraic Taylor coefficients at zero for every
\(a\in\mathbb Q\Lambda\).
Indeed their translation theorem works for a symmetric divisor,
with algebraic germ at zero. Divide by a rational function defining
the translated divisor locally at zero. The normalized coefficients
lie in a torsion field, and multiplying back by the rational germ
retains algebraicity. The reduced and Mumford translations differ
by a root of unity on torsion.
\cite[Example~1.9, Corollary~2.2 and Theorem~2.9]{bannai-kobayashi-theta}
Thus the even theta \(\theta_{\rm ell}=U_{\omega_0/2}\theta/
(U_{\omega_0/2}\theta)(0)\) and all its reduced torsion translates
have algebraic Taylor coefficients. The denominator is nonzero
since the chosen half period is nonzero modulo \(\Lambda\).

Here is the passage to the actual vacuum of the host \(E\).
Use the CM type over the embedding of \(k_{\rm CM}\).
The host fractional ideal \(\mathfrak A_0=\mathfrak b^{-1}\)
satisfies \(\mathfrak A_0\overline{\mathfrak A_0}
=\mathfrak D_{E/k_{\rm CM}}^{-1}\).
The lattice \(\Lambda_M\) is the type embedding of
\(\delta\mathfrak A_0\). For the Hermitian form \(H\) obtained by
dividing the standard form in the complex coordinates by \(A\),
its alternating form \(E_H=\operatorname{Im}H\) on that lattice
is
\[
 E_H(\delta x,\delta y)
   = \frac{2}{\sqrt{d_{\rm CM}}}
      \operatorname{Im}\operatorname{Tr}_{E/k_{\rm CM}}(x\bar y),
\]
with the trace in the chosen embedding.
An element pairing integrally with the lattice also pairs
integrally against its \(\omega\)-multiples, for
\(\omega=(1+i\sqrt{d_{\rm CM}})/2\). Hence the trace pairs
in \(\mathcal O_{k_{\rm CM}}\). The trace dual calculation
gives \(\mathfrak D_{E/k_{\rm CM}}^{-1}
\overline{\mathfrak A_0}^{-1}=\mathfrak A_0\).
Testing real vectors against a rational basis of the lattice
shows first that the dual vectors lie in the rational span,
since the alternating Gram matrix is nondegenerate and rational
there. This proves self-duality for \(E_H\).
The same calculation applies to the Hom lattice between two
of the self-dual Hermitian class lattices: the Hom lattice is
Hermitian self-dual, as can be checked in local bases using the
two unimodular Gram matrices. The relative trace dual then
introduces the same \(\mathfrak A_0\).

The fourfold isometry \(g_C\) of the vacuum is induced by an
isometry over \(k_{\rm CM}\) from the fourfold trace form to the
standard form. For existence, diagonalize the positive Hermitian
trace form over \(k_{\rm CM}\). Each diagonal value is a positive
rational number \(a\), and the four-squares identity gives
a rational matrix with transpose product \(a I_4\). These
matrices realize the isometry on the four copies of each
diagonal summand. For the chosen isometry, its complex matrix
has algebraic entries by expressing it in the type embeddings
of a \(k_{\rm CM}\)-basis. It takes rational period vectors
to rational period vectors. Thus
\(L_1=\Lambda^{4[F:\mathbb Q]}\) and \(L_2=g_C\Lambda_M^4\)
are commensurable and self-dual for the same alternating form.
\releasecite{Formalism/LocalRingCardinality/DataHost.lean}{153}{The inverse relative different factor}
\releasecite{Formalism/LocalToGlobal/KernelLattices/LateralTheory.lean}{128}{The chosen fourfold isometry}.

The common semicharacter of the two lattices is
\(\beta(x)=\exp(\pi i\,\sqrt{d_{\rm CM}}\,H(x,x)/2)\).
Its exponent divided by \(\pi i\) is \(E_H(\omega x,x)\),
an integer on either lattice. Polarizing this expression
shows that its change at \(x+y\) is \(E_H(x,y)\) modulo two,
since \(2E_H(\omega x,y)=E_H(x,y)+
\sqrt{d_{\rm CM}}\operatorname{Re}H(x,y)\).
On one standard elliptic factor, multiplication by \(\omega\)
has distinct eigenvalues zero and one on \(\Lambda/2\Lambda\).
This follows from its trace one and its even norm, using
\(d_{\rm CM}\equiv7\pmod8\).
The chosen \(\omega_0\bmod 2\Lambda\) is the sum of the two
nonzero eigenvectors, since the images of \(\omega_0/2\) under
\(\omega\) and \(1-\omega\) are nonzero modulo \(\Lambda\).
The even theta sign at
\(\gamma\) is the odd sign times
\((-1)^{E_H(\omega_0,\gamma)}\).
This sign and \(\beta(\gamma)\) both equal \(1\) on zero and
on the eigenvectors, and \(-1\) on their sum. This identifies
the sign in the actual product \(F_{\rm ell}\) of the
\(\theta_{\rm ell}\) over the standard elliptic factors.

For two such self-dual lattices \(L_1,L_2\), put \(L_0=L_1\cap L_2\).
The finite quotients \(L_1/L_0,L_2/L_0\) have a perfect pairing
through \(e^{2\pi i E_H}\), because the annihilator in each
lattice is the intersection. In particular they have equal
order. The transport \(T_{12}\) is the sum of
\(\beta(x)U_x\) over \(L_2/L_0\), divided by
\(\sqrt{[L_2:L_0]}\). Here \(U_x\) is the reduced translation
for \(H\). The translation law
\(U_xU_y=e^{-\pi iE_H(x,y)}U_{x+y}\) makes this
sum independent of representatives and invariant under
the operators \(\beta(x)U_x\) for \(x\in L_2\).
To compose with \(T_{21}\), commute the \(L_1/L_0\)
operators past the \(L_2/L_0\) operators on the source
theta space. Summing the commutator characters over
\(L_1/L_0\) leaves only the identity coset of \(L_2/L_0\).
The two square roots of the index cancel the order of the
surviving sum, so the transports are inverse. This also
proves nonvanishing of the transported vacuum. The standard
theta space is one-dimensional by the one-variable elliptic
theta condition applied in each variable. After transport,
pull back by \(g_C\). The result belongs to the theta line
for \(\Lambda_M^4\). A nonzero slice shows that the theta
space for \(\Lambda_M\) is nonzero. Two independent vectors
there would give independent fourfold products. Thus the
transported product is a nonzero scalar times a product of
one vacuum at four independent arguments.

The current normalization specifies
\[
 \theta_M(z_1)\theta_M(z_2)\theta_M(z_3)\theta_M(z_4)
   =\kappa_M (T_{12}F_{\rm ell})(g_C(z_1,z_2,z_3,z_4)).
\]
It can be imposed by taking a fourth root of the scalar in the
one-dimensional space. Here \(\kappa_M\) is the prescribed
finite product of powers of roots of
\(\varphi(\pi)/\varphi(\bar\pi)\) for generators of powers of
split prime ideals of \(k_{\rm CM}\), with \(\varphi\) its
complex embedding. In particular \(\kappa_M\) is algebraic
and nonzero. The right-hand side is a finite sum of products
of reduced torsion translates of the elliptic theta after
algebraic linear substitution, with algebraic coefficients.
Composing simultaneously with the four reduced translations
by a rational period vector \(a\) preserves this property.
Indeed \(g_C(a,a,a,a)\) is rational in the standard periods,
and the translation law adds only roots of unity.
Taking the coefficient of a multi-index \(\alpha\) in each
of the four independent argument blocks gives the fourth
power of the Taylor coefficient of \(U_a\theta_M\).
Thus each such coefficient is algebraic, also when the
value of \(\theta_M\) at zero vanishes.
\releasecite{Formalism/LocalToGlobal/KernelLattices/LateralTheory.lean}{23}{Fourth-tensor normalization}
\releasecite{Formalism/LocalToGlobal/KernelLattices/LateralTheory.lean}{96}{Algebraic factor in the vacuum}.

For a source and target class lattice \(N,N'\), take the
product of \(\theta_M\) in the matrix entries, then move its
source lattice and target lattice by the same normalized
Heisenberg averages. The two relative indices are finite since
all lattices are in the same rational span. The average at each
step and its torsion translates preserve algebraicity, by the
translation law and the preceding result. The determinant
factor \(\chi_F(N,N')\) in the moved vacuum is
\(\chi(\det N)^{\dim N'/2}\chi(\det N')^{-\dim N/2}\),
where the determinant ideals are measured from the standard
lattices. In the normalized germ its reciprocal removes this
factor. Thus this argument uses the stated normalization,
including when changing the chosen ideal character in the
paired construction.
\releasecite{Formalism/LocalToGlobal/KernelLattices/LateralTheory.lean}{165}{Moved vacuum and determinant factor}
\releasecite{Development/AbelianVarietyTheory/Proofs/IwasawaDecomposition/ChKerCCongr.lean}{75}{The normalized germ under the change of character}.

\subsubsection{Finite coset jets at the chosen points}
In one step with ranks \(a,b\), let \(\mathcal L\) be the
semilocal Hom lattice of the source and target class lattices
at the tested places. The finite symbol \(\Xi\) is supported
on \(c_A^{-1}\mathcal L\) and constant modulo \(c_B\mathcal L\).
The arithmetic kernel sums its values over
\(c_A^{-1}\mathcal L/c_B\mathcal L\). To see the normalization
of a term, write \(s\) for the lift of its coset in the
period-coordinate lattice. For the translation construction
of the indicator germ, let \(A_1,A_2\) be the two finite
sets of representatives in the Fourier sum and its dual
translation sum, and let \(\mathbf 1_M\) be the subgroup
indicator in that Fourier sum. If \(\alpha_H\) is the quadratic
phase \(\exp(\pi i\,\sqrt{d_{\rm CM}}\,H(x,x)/2)\) and \(G\)
the moved vacuum before the determinant factor, the normalized
germ for this coset is
\[
 G_s=\sum_{u\in A_2}
    \alpha_H(s)\alpha_H(u)e^{\pi i E_H(s,u)}
       \left(\frac1{|A_1|}
            \sum_{v\in A_1}\mathbf1_M(v)e^{2\pi iE_H(v,u)}\right)
     U_{u-s}G.
\]
Indeed Fourier transformation of the indicator gives the
inner average times the square root of its volume, and the
germ normalization contributes the reciprocal square root.
The multiplier for \(U_{-s}U_u\) is \(e^{\pi iE_H(s,u)}\).
The determinant factor \(\chi_F(N,N')\) cancels in \(G_s\).
This proves the formula by finite character expansion with the
normalized measures. Its summands are algebraic germs since the
phases are roots of unity on rational period vectors and the
translates of \(G\) have algebraic Taylor coefficients.
\releasecite{Formalism/LocalToGlobal/KernelLattices/XiIndCClass.lean}{18}{Finite Fourier convention}
\releasecite{Development/AbelianVarietyTheory/Proofs/DivisorProofs/CosetIndicatorThetaTorsionTranslates.lean}{102}{Indicator germ expansion}.

\paragraph{Reading the normalized real coefficients.}
Use the orthonormal definite frames on the source and target.
Write \(z_\tau\) for the real-place complex matrix coordinate
in these frames and \(\xi_R=|\tau(\xi_0)|>0\) for the host
form scale, which is constant on the CM type. With the
prescribed additive sign the Gaussian is
\(g(z)=\exp(-\pi\xi_R\sum_\tau\|z_\tau\|^2)\).
In a single entry of \(z_\tau\), translation by \(t\alpha\),
\(t\in\mathbb R\), acts in the Weyl model with derivative of
its phase equal to \(2\pi i\xi_R\operatorname{Im}(\bar\alpha z)\).
The derivative of the Gaussian contributes
\(-2\pi\xi_R\operatorname{Re}(\bar\alpha z)\).
The sum is \(-2\pi\xi_R\alpha\bar z\).
Consequently the combination of the directions \(1,i\),
with coefficients \(1,-i\), acts on \(g(z)H(z)\), after
division by \(2\sqrt{2\pi\xi_R}\), as
\[
 g(z)(2\pi\xi_R)^{-1/2}
       \bigl(\partial_z-2\pi\xi_R\bar z\bigr)H(z).
\]
The holomorphic derivative commutes with multiplication by
each \(\bar z\). In a word of creation operators evaluated at
zero every term involving one of those multiplication
operators vanishes. This leaves the iterated holomorphic
derivative with one factor of \((2\pi\xi_R)^{-1/2}\) for each
operator. The value of the Gaussian itself at zero is \(1\).

With the zero finite-argument vacuum in the standard lattice
model, \(H(z)=c_{\rm vac}G(\delta z)\). Indeed rescaling by
\(\delta\) identifies the Gaussian-removed quasiperiodicity
with that of the reference theta line of
Section~\ref{sec:period-germs}. Both nonzero vacua belong
to that line. Finite Heisenberg translation and the normalized
finite averages defining \(G_s\) preserve the common scalar
\(c_{\rm vac}\). Thus the monomial of each multi-degree
\((d_\tau)_\tau\) reads the derivative of \(G_s\) times
\[
 c_{\rm vac}\prod_\tau
    \left(\frac{\delta}{\sqrt{2\pi\xi_R}}\right)^{d_\tau}.
\]
In the conjugate Fock convention this formula is conjugated.
Division by the vacuum scalar and by these period powers,
followed by multiplication by \(\prod_\tau c_{D,\tau}^{d_\tau}\),
is exactly the derivative normalization below. These factors
must be used on the same real polynomial spaces and with
the same finite vacuum identification as in the lifted
coefficients.
\releasecite{Development/AbelianVarietyTheory/Proofs/DivisorProofs/FockJetGaussianEntireOriginGridFrames.lean}{260}{Directional Weyl derivative}
\releasecite{Development/AbelianVarietyTheory/Proofs/DivisorProofs/FockJetGaussianEntireOriginGridFrames.lean}{355}{Normalized creation at zero}
\releasecite{Development/AbelianVarietyTheory/Proofs/DivisorProofs/FockJetsUnitVacuumEqScaledThetaJets.lean}{33}{CM coordinates and period power}
\releasecite{Development/AbelianVarietyTheory/Proofs/DivisorProofs/FockVacEqJets.lean}{76}{Positive and conjugate Fock conventions}.

For a multi-index \(\alpha\) on the variables \(Z_{\tau,i,j}\),
the derivative used in the step is
\[
 J_s(\alpha)=
    \left(\prod_\tau c_{D,\tau}^{\sum_{i,j}\alpha_{\tau,i,j}}\right)
    \partial^\alpha G_s(0)
\]
with the prescribed algebraic differential scalars \(c_{D,\tau}\).
The factorial converting a Taylor coefficient to a derivative
is rational. To pass from algebraicity to the specified \(K\),
fix the modulus and the maximal degree for each word and step
before choosing the final class data. Definite class finiteness
gives finitely many self-dual lattice representatives in each
rank. For each pair, every positive divisor of the modulus
gives a finite Hom quotient and finitely many derivatives
below that degree. Include these normalized values, the
differential scalars and the embeddings of \(E\) in a finite
extension of \(\mathbb Q\).
Transport to any representative in the genus is an algebraic
linear substitution on the germs. Under this transport the
determinant ideals change by the principal ideals of the
rational frame changes, whose character values are algebraic
products of their embeddings. The chain and product rules
therefore express the transported derivatives in the same
finite extension.
In the family selection the modulus for a tuple and word is
\(c_Ac_B\), where \(c_A\) is the product of the residue
cardinalities on the selected row support to the power \(m+1\)
and \(c_B=T_w c_A\). The local support set belongs to the finite
list of subsets of the primes over the type and word primes
and the conductors of the finite list of twists.
The prescribed coefficient field contains these finite generators
before the grid is realized. This gives membership of
\(J_s(\alpha)\) in \(\iota_K(K)\) for the actual point computations.
\releasecite{Development/AbelianVarietyTheory/Definitions/AutomorphicMembership/JetsInIffStepField.lean}{99}{Normalized derivatives and degrees}
\releasecite{Development/AbelianVarietyTheory/Proofs/IwasawaDecomposition/FinAlgJetSet.lean}{53}{Finite field through lattice transport}
\releasecite{Development/AbelianVarietyTheory/Proofs/IwasawaDecomposition/GensJAlgebraic.lean}{61}{Actual moduli and field}.

Here is the resulting coefficient calculation. Multiply the
source coefficient polynomial \(P(X)\) by the determinant
powers \(\prod_\tau\det(X_\tau)^{n_\tau}\), with \(n_\tau\)
the source determinant shift for the step. Substitute
\[
 X_{\tau,i,j}=\sum_{h=1}^b Z_{\tau,h,i}Y_{\tau,h,j}.
\]
In the expanded polynomial replace \(Z^\alpha\) by
\(\sum_s\Xi(s)J_s(\alpha)\). Multiply by the ratio of the
inverse determinant character values of the source and target
frames. The resulting polynomial in \(Y\) is evaluated with
the target determinant shift divided out. The substitution
comes from evaluating the source frame against the columns
of the target frame via the Hom variable. Each monomial in
the derivative variables has the prescribed degree at every
embedding because the source polynomial has its prescribed
homogeneous weight and the determinant shift adds \(a n_\tau\).
Consequently only the indicated jets enter. Their field of
definition and the framed character readings give an operator
over \(K\).
\releasecite{Formalism/LocalToGlobal/KernelLattices/ThDistC.lean}{20}{Polynomial substitution in the jet}
\releasecite{Development/AbelianVarietyTheory/Proofs/GlobalInterpolation/KmapOfOkOfVecP.lean}{50}{Source and target determinant shifts}
\releasecite{Development/AbelianVarietyTheory/Proofs/GlobalInterpolation/KmapOfOkOfVecP.lean}{238}{Normalization of the finite characters}.

The determinant character readings may be checked on units
and a uniformizer at each place of the support. At a type
place they are the prescribed type characters on residue
units and the fixed inverse reference and twist value on
the uniformizer. At a twist place the reference is unramified
and the twist is tame, so the residue-unit value is that
of the inverse twist. The two rules agree with the value
on every determinant by writing a local element as a
uniformizer power times a unit. Off the support both
characters are unramified, and the determinant contributes
the inverse unramified value raised to the valuation of
the local frame determinant. The factors at conjugate split
places use inverse values because the characters are trivial
on the ground-field ideles. For the spherical factors,
in a self-dual lattice model with the indicated additive
character, the Fourier transform of the lattice indicator is
that same indicator. This follows by character orthogonality
on the lattice and self-dual volume one. It also applies at
the unramified inert places when both tests are in that
self-dual model.
\releasecite{Development/AbelianVarietyTheory/Proofs/GlobalInterpolation/KmapOfOkOfVecP.lean}{167}{Tame local readings}
\releasecite{Development/AbelianVarietyTheory/Proofs/GlobalInterpolation/KmapOfOkOfVecP.lean}{208}{Determinant character and frame valuations}.

\subsubsection{The finite symbols at the coupled class points}

Write \(A=B=m+1\) for the two depth exponents of the chosen
tuple's local box and \(M=m\) for its marking depth.
At each selected split place \(P\), set \(q=q_P\), let
\(\mathcal O\) be the local ring and use \(q\) as uniformizer
in the selected component. These primes have absolute
residue degree one. The matrix of a Hom variable \(x\) in the
two chosen frames is \(U=(f')^{-1}xf\).
At the primes of the word \(T_w\) use the indicator
\(\phi_T(x)\) that \(x\) maps the source local lattice into the
target local lattice and that the target's marked
residue covector pulls back to the source marking.
These conditions are unchanged under a perturbation by
\(T_w\) times the integral Hom lattice at these primes.
\releasecite{Development/AbelianVarietyTheory/Proofs/SatakeParameters/AbsNormEqOfSplitFactor.lean}{83}{Split places and their residue cardinalities}
\releasecite{Formalism/LocalToGlobal/Columns/FiniteSumNonzeroLargeDepth.lean}{641}{The framed matrix and the marked lattice indicator}.

On an ordinary arrow \(i<m-2\), the symbol is
\(\phi_T(x)\prod_P\Xi_{i,P}(U)\).
In \(\Xi_{i,P}\), let \(a,b\) be the two definite ranks,
and \(c',c\) the numbers of added source and target coordinates
in the mixed model. Use zero-based indices.
The enlarged matrix \(X(U,\mathbf t)\) has \(U\) in its first
\(b\) rows and first \(a\) columns. A position \((\ell,s)\)
is free if
\[
 s\ge a,\qquad
 (s-a=c)\ \text{or}\quad
 (s>a\ \text{and}\ (\ell<b\ \text{or}\ \ell-b<s-a)).
\]
The case \(\ell-b<s-a\) is read for \(\ell\ge b\).
At a position in an added column which is fixed, put
one when \(\ell=b+s-a\) and zero elsewhere. The remaining
positions below \(U\) are zero. Write \(\mathcal F\) for the
set of free positions, and
\(\mathfrak u(X)=\sum_{j=0}^{c-1}X_{b+j,a+j+1}\) for
the entry sum in the unipotent character. With the free entries
\(\mathbf t\) ranging over \((q^{-A}\mathcal O/q^B\mathcal O)^{\mathcal F}\),
the normalized local value is
\[
 \Xi_{i,P}(U)=
 \operatorname{vol}(C_{\rm be}\cap N_1
                  \cap\operatorname{Stab}(\Lambda_Y))^{-1}
 q^{-B|\mathcal F|}
 \sum_{\mathbf t}(\omega\Phi_i)\bigl(X(U,\mathbf t)\bigr)
                \psi_P\bigl(-\mathfrak u(X(U,\mathbf t))\bigr).
\]
Here \(\omega\Phi_i\) is the finite matrix symbol transported
by the two prescribed extended-frame automorphisms using their
splitting characters and determinant powers. The local volume
is for the indicated target unipotent subgroup in the
same measure as the integral defining the coefficient.
If the supported mixed-model kernel in these coordinates is
this integrand and is constant on its free cosets modulo \(q^B\),
partitioning the integral into these cosets gives the displayed
factor \(q^{-B|\mathcal F|}\). This calculation isolates the
support and frame checks needed for the ordinary integrals.
\releasecite{Formalism/LocalToGlobal/Columns/FiniteSumNonzeroLargeDepth.lean}{551}{Extended frames and the free entries}
\releasecite{Formalism/LocalToGlobal/Columns/FiniteSumNonzeroLargeDepth.lean}{937}{Normalization of the ordinary local symbol}.

On the penultimate arrow the extended source has one added
column and the target has none. A coupled point includes a
marked coset \(Y\) on the target of this arrow. In its frame
at \(P\) the coset is \(y'+q^{M_P}\mathcal O^b\), where
\(M_P=0\) at a type prime and \(M_P=M\) at a twist prime.
At a type prime \(y'\in q^{-1}\mathcal O^b\), and at a twist
prime \(y'\in q^{-M}\mathcal O^b\).
The prescribed symbol is \(\phi_T(x)\) times the product of
\((\omega\Phi_{m-2})(U\mid y')\), with \(y'\) the added column.
At a type prime the integral extended source basis scales
that column by \(q\). Thus a change of \(y'\) by an integral
vector changes the matrix in that basis by \(q\) times an
integral matrix. The source frame automorphism preserves
the rescaled lattice, so the change after applying it is
still divisible by \(q\) in the rescaled entries. At a twist
prime the source frame matrix has denominator exponent at
most \(m-2\), so a change by \(q^M\) times an integral
vector changes the framed extended matrix by a matrix
divisible by \(q\). The symbol conditions read the
integrality and the residues of these entries, and the
frame multipliers are unchanged. This gives independence
of the representative of \(Y\). The same argument for
changes of \(U\) by \(q^B\) times an integral matrix gives
the modulus in the Hom variable.
\releasecite{Formalism/LocalToGlobal/Columns/FiniteSumNonzeroLargeDepth.lean}{978}{The additional column}
\releasecite{Formalism/LocalToGlobal/Columns/FiniteSumNonzeroLargeDepth.lean}{1032}{The framed automorphism bounds}
\releasecite{Formalism/LocalToGlobal/Columns/FiniteSumNonzeroLargeDepth.lean}{1128}{The marked cosets and penultimate symbol}
\releasecite{Formalism/LocalToGlobal/Columns/ChainTw.lean}{218}{The box and marking exponents}.

On the last arrow both extended column counts are zero.
Now \(a=k+(m-1)(t+1)=r+(m-1)t\) and \(b=k+m(t+1)\).
Write \(U_{\rm up}\) for the first \(b-1\) rows of the
framed matrix and \(\ell_U\) for its last row.
At the top point over a row point, the frame on the added
line uses the scale \(\pi_{T_w}\). Thus \(\ell_U\) in the
unscaled last-line coordinate includes the factor
\(\pi_{T_w}^{-1}\) when read in the coupled row frame.
The prefix at this same point is read using the row frame.
Let \(\Phi_{\rm top}(U_{\rm up})\) be the corresponding
ordinary or sharp type symbol, or the tame twist symbol
at a twist prime. The local factor after the last partial
Fourier transform is
\[
 \Phi_{\rm top}(U_{\rm up})\,
 q^{-M_P a}\,
 \mathbf 1_{\ell_U\in q^{-M_P}(\mathcal O^a)^\vee}
 \psi_P(-\ell_U(y')) ,
\]
where \((\mathcal O^a)^\vee\) is the integral linear dual and \(y'\)
in the prescribed marked coset of the source. To verify the
factor, integrate \(\psi_P(-\ell_U(z))\) over
\(y'+q^{M_P}\mathcal O^a\). Translate by \(y'\). With
\(\operatorname{vol}(\mathcal O)=1\) and additive character of
conductor \(\mathcal O\), character orthogonality gives zero
unless the dual-lattice condition holds, and otherwise gives
the indicated volume times the phase. It also shows that the
phase is independent of the representative \(y'\).

The prefix on this arrow has integral support. In the type
symbols all columns are among the first \(r+(m-1)t=a\),
so the defining entry conditions impose integrality on each
prefix column. With no added columns its entry weights are
unscaled. The tame twist symbol also imposes integrality
on each entry. The prefix is constant under adding
\(q^B\) times an integral matrix since its other
conditions and characters read only residues. For the last
row, \(M\le A,B\), the type bound \(qy'\in\mathcal O^a\)
and \(B\ge1\) show that adding \(q^B\) times an integral
covector to \(\ell_U\) changes neither the dual-lattice
condition nor the phase. On its support the phase argument
has denominator exponent at most \(1\) at a type place and
\(2M\) at a twist place. These are within the prescribed
character depths. Thus the terminal symbol has the required
box and gives values in \(K\). The volume factor is a unit at
places of \(K\) over \(T_w\) since the primes in this row
support lie away from \(T_w\).
The local formula agrees with the symbol at a top column
point lifted from a row point. To check this, compare its
top frame \(f_{\rm top}\) with the block frame whose prefix
is the row frame and whose added vector is
\(\pi_{T_w}\) times the prescribed last basis vector.
At a twist prime that block frame is the chosen top frame.
At a type prime both frames are integral bases of the
same lifted lattice and give the same prescribed residue
frame of its decoration. Their transition therefore lies
in \(\mathrm{GL}_b(\mathcal O)\) and is the identity
modulo \(q\). The integrality of the entire matrix, which
is necessary for nonzero symbol at a type prime, is invariant
under this transition. On an integral matrix its prefix
residues are unchanged. Its last covector changes by an
element of \(q(\mathcal O^a)^\vee\), whose pairing with
\(y'\in q^{-1}\mathcal O^a\) is integral. Thus the
terminal Fourier phase is unchanged too. The decorated top
lattice at the primes over \(T_w\) is the same lattice in
the two descriptions, so the factor \(\phi_T(x)\) agrees.
This proves the equality of the two terminal symbols
at the lifted points, including their marks.
\releasecite{Formalism/LocalToGlobal/Columns/FiniteSumNonzeroLargeDepth.lean}{711}{The tame prefix}
\releasecite{Formalism/LocalToGlobal/Columns/FiniteSumNonzeroLargeDepth.lean}{780}{The type prefixes}
\releasecite{Formalism/LocalToGlobal/Columns/FiniteSumNonzeroLargeDepth.lean}{1150}{The marked-coset transform}
\releasecite{Formalism/LocalToGlobal/Columns/XiTopTilRefHolds.lean}{18}{Same lattice and residue frame at the type places}
\releasecite{Formalism/LocalToGlobal/Columns/ChainTopTilQ.lean}{43}{The same symbol at a lifted row point}
\releasecite{Development/AbelianVarietyTheory/Proofs/GlobalInterpolation/BadPrimeNormPrimeNeLevelPrime.lean}{88}{Support and the depth bound at the last step}.

The character used to move the arithmetic vacuum can be
changed on the class kernels if its restriction to norm-one
fractional ideals agrees with the original character. In
fact, for two self-dual lattices in the same Hermitian space,
the relative determinant ideal \(\mathfrak d\) satisfies
\(\mathfrak d\overline{\mathfrak d}=\mathcal O_E\).
In local bases, duality conjugates and inverts the
determinant ratio at the two conjugate places, which proves
the ideal identity. Both relative determinant ideals
of the source and target self-dual lattices therefore have
the same character readings before and after the change.
The vacuum scalar and the determinant normalization agree,
so the germs \(G_s\), the derivatives and the three types of
finite kernels agree in the same frames.
\releasecite{Development/AbelianVarietyTheory/Proofs/IwasawaDecomposition/EvaluationDatumOkBaseTuple.lean}{45}{Norm-one determinant ideals}
\releasecite{Development/AbelianVarietyTheory/Proofs/IwasawaDecomposition/ChKerCCongr.lean}{34}{Compatibility of the three kernels with the character change}.

\teamanchor{proof.period-normalized-entry}
\subsubsection{The normalized step entry}\label{sec:period-normalized-entry}

Fix two class points \(p,p'\) in one of the selected
definite steps, of ranks \(a,b\). Let \(\Xi_{p,p'}\) be
the finite symbol just described in their frames. This
means the term-level average in an ordinary step, the
marked-column restriction in the penultimate step,
and the marked-covector Fourier transform on the
last step. At the lifted row points the last symbol
is the coupled row symbol by the frame calculation above.
All three use \(\phi_T\) on the decorated \(T_w\)-lattice
and the indicated tests at the row-support places.

Consider the normalized mixed-model theta kernel
\(\mathcal T_{p,p'}(z)\) before real differentiation.
The finite support gives its decomposition into
indicators of the classes of \(c_A^{-1}\mathcal L/c_B\mathcal L\).
The self-dual vacuum takes care of the spherical places
and the Chinese remainder theorem gives the representative
sum at the remaining places. In a local matrix space the
frame change \(U\mapsto f_{p'}U f_p^{-1}\) has Jacobian
\[
 |\det f_{p'}|^a\,|\det f_p|^{-b}.
\]
Thus its half-density factor is the inverse square root
of the local lattice-volume ratio. The character
expansion of each indicator carries the square root of
the same volume, so the two factors cancel. The remaining
average has the denominator \(|A_1|\) and the two
Heisenberg translations with quadratic phase displayed
in the formula for \(G_s\). This gives the normalized
finite expansion, with the vacuum determinant included,
\[
  \mathcal T_{p,p'}(z)
      =\chi_F(N,N')\sum_s
            \iota_K\bigl(\Xi_{p,p'}(s)\bigr)G_s(z).
\]
Here \(G_s\) uses the same normalization for all cosets,
and finite differentiation commutes with this sum. The
normalization can also be checked by applying the two
finite character projectors to a single coset indicator.
The first projector uses the annihilator of
\(c_B\mathcal L\), so its average selects precisely the
given coset. The second transports to the paired half
of the period lattice. Multiplication of its two
translations yields \(e^{\pi i E_H(s,u)}\). Linearity
then gives the expansion for \(\Xi_{p,p'}\).

Let \(M_j(p)\) denote the inverse value of the
prescribed reference-and-twist character \(\mu_j\xi\)
on the determinant id\`ele of the framed class \(p\).
The unramified valuation and the ramified frame values
in the preceding calculation give \(M_j(p)\) in \(K^\times\).
At the conjugate half of a split place the local
determinant valuation is opposite; triviality on
ground-field id\`eles gives the corresponding inverse
character. At an inert place the determinant has
valuation zero. For the oscillator character on the
relative determinant ideals use \(\chi_\psi\). The two
complex class multipliers are
\[
 \begin{aligned}
 \rho_j(p)
   &=\iota_K(M_j(p))\,
           \chi_\psi(\det N)^{-b/2},\\
 \rho'_j(p')
   &=\iota_K(M_j(p'))^{-1}\,
           \chi_\psi(\det N')^{a/2}.
 \end{aligned}
\]
Their product is \(\iota_K(M_j(p)/M_j(p'))\,
\chi_F(N,N')^{-1}\). This cancels the character of the
moved vacuum before taking the normalized derivatives.
\releasecite{Formalism/LocalToGlobal/KernelLattices/XiIndCClass.lean}{199}{Frame volume, half-density, and unit-volume kernel}
\releasecite{Development/AbelianVarietyTheory/Proofs/SatakeParameters/ThetaSupportInDivLattice.lean}{81}{Indicator decomposition on the test box}
\releasecite{Development/AbelianVarietyTheory/Definitions/AutomorphicMembership/ReferenceCharactersFactorTame.lean}{571}{The two actual class multipliers}
\releasecite{Development/AbelianVarietyTheory/Proofs/GlobalInterpolation/KmapOfOkOfVecP.lean}{271}{Cancellation with the moved-vacuum character}.

Write \(\mathscr J_{p,p'}\) for the \(K\)-linear
coefficient contraction taking
\[
  Z^\alpha\longmapsto
       \sum_s\Xi_{p,p'}(s)\,\iota_K^{-1}\bigl(J_s(\alpha)\bigr)
\]
on the degrees selected by the input weight.
If \(P\) represents the input, write
\(P^\sharp(X)=P(X)\prod_\tau\det(X_\tau)^{n_\tau}\).
In the fixed complex frames the normalized analytic
entry evaluated at \(Y\) is
\begin{equation}
 \label{eq:definite-normalized-entry}
 \iota_K\left[
  \frac{M_j(p)}{M_j(p')}\,
  \mathscr J_{p,p'}\!\left(P^\sharp(Z^{\mathsf T}Y)\right)
 \right](Y)
       \prod_\tau\det(Y_\tau)^{-s_{{\rm div},\tau}} .
\end{equation}
Here the embedding is applied to the coefficients
of the polynomial before evaluation at \(Y\), and
\(s_{{\rm div},\tau}\) is the prescribed output
determinant exponent. Indeed the real operator
replaces each derivative variable by the
Gaussian-normalized creation operator. At zero its
evaluation is \(J_s(\alpha)\), by the calculation
of the creation operator and the prescribed
differential scale above. The source sharp factor
increases the determinant exponent from the
input coefficient convention to \(n_\tau\).
The exponent of the determinant in the
target coefficient model is \(-s_{{\rm div},\tau}\),
so removing the output shift gives exactly
that displayed evaluation. The exponents and
degrees are integral on these ranks. In particular
the output polynomial has the required covariance
over \(K\): substitution in the equivariance
identity for the complex map gives polynomial
identities with coefficients in \(K\), and the
embedding is injective.
\releasecite{Development/AbelianVarietyTheory/Proofs/GlobalInterpolation/KmapOfOkOfVecP.lean}{50}{Input sharp and output determinant exponents}
\releasecite{Development/AbelianVarietyTheory/Proofs/GlobalInterpolation/KmapOfOkOfVecP.lean}{96}{Normalized derivatives in the coefficient contraction}
\releasecite{Development/AbelianVarietyTheory/Proofs/GlobalInterpolation/TrueChainKernelEqPointKernelAtRepresentatives.lean}{14}{Application to each actual stage representative}.

Formula~\eqref{eq:definite-normalized-entry} therefore
gives the entry on every selected pair with the
specified tests and frame placements. Applying it at
the penultimate coupled points and restricting it
at the lifted top points proves the corresponding
entry identities there. The local orbital
realization of the opposite path in
\eqref{eq:period-principal-unfolding} is a separate
calculation on its indicated term levels. For the
definite top pairing, only the normalized test symbols
and placements of the column chain enter the entry
calculation.

\subsubsection{Descent of the period}
Write \(N_w\) for the prescribed column lattice before division.
Its scalar extension \(N_{w,K}\) has complexification
\(\mathcal M_w\). To read the top pairing, evaluate the normalized
step kernel on each pair of the chosen definite class
representatives, with the framed coefficient vectors. Denote its
prescribed arithmetic value by \(K_{j,K}(c,c')\) in these
coordinates. Equation~\eqref{eq:definite-normalized-entry}
identifies it under \(\iota_K\) with the step in
\eqref{eq:period-top}, including the source and target
determinant factors. At the final two steps the entries
are read at the coupled points, and at the final step
only at the images of the top row points.

The corresponding class measures are those of the term levels.
\releasecite{Development/AbelianVarietyTheory/Proofs/IwasawaDecomposition/KernelsOverOfFamily.lean}{139}{The step entries at the class representatives}
\releasecite{Development/AbelianVarietyTheory/Proofs/IwasawaDecomposition/KernelsOverOfFamily.lean}{183}{The coupled and top entries}.

Given these entries, start with the prescribed values
\(C_0(c)=\phi(c)\) after transport to the tuple's stage zero.
Successively set
\[
 C_{j+1}(c')=\sum_c |\operatorname{Stab}_j(c)|^{-1}
                                     K_{j,K}(c,c')C_j(c).
\]
If one refines a class into orbits at an intermediate level, the
sum of the reciprocal stabilizer orders over a finite orbit
decomposition is the orbit size divided by the acting group
order. Hence a refinement by an index \(h\) multiplies the
weighted sum of an invariant value by \(h\). This accounts for
the level volumes when using the coupled intermediate points.
At a row class \(a\), let \(\widetilde a\) denote the top column
point over \(a\), with its prescribed frame transport and coefficient
projection included in \(\operatorname{res}_a\). Pairing against
the class value \(x(a)\) gives
\[
 m_Y(\boldsymbol\mu,m)
   \sum_a |\operatorname{Stab}_Y(a)|^{-1}
      \bigl(\operatorname{res}_a C_m(\widetilde a),x(a)\bigr)_{\rm alg}.
\]
The algebraic pairing is invariant under changes of representatives.
Indeed the column and row carry inverse finite characters and the
branching projection pairs dual coefficient actions. The scalar
\(m_Y(\boldsymbol\mu,m)\) is a product of indices.
The class unfolding at each step and the entry identity imply by
induction that the displayed sum maps to the actual top pairing
under \(\iota_K\). All its operations are over \(K\).
The same entry identity enables the coefficient valuation for
the actual period in Section~\ref{sec:integral}.

\releasecite{Development/AbelianVarietyTheory/Proofs/IwasawaDecomposition/KernelAtPointPenultimateStepFirstCoupledKernel.lean}{286}{Class sums at the coupled points}
\releasecite{Development/AbelianVarietyTheory/Proofs/IwasawaDecomposition/KernelAtPointPenultimateStepFirstCoupledKernel.lean}{309}{Invariant pairing and the top class sum}.
\releasecite{Development/AutomorphicArithmetic/ThetaCorrespondence/ThetaPairLifts/ThirdPlaceSeries/RowEDcPtLinked.lean}{17}{Top period entry at the chosen points}.

Using the entry identity and \eqref{eq:period-row-match}, take the
base \(K\)-valued period on any rational preimage of \(F\in R_K\)
under \(P_{\boldsymbol\mu_0}\). The projected base map is onto,
and the simultaneous theta--Bessel identity gives zero on its
kernel. The embedded value determines the value in \(K\), so the
same kernel vanishes over \(K\). This defines uniquely
\[
 B_{K,w}:N_{w,K}\longrightarrow R_K^\vee,\qquad
 B_w(\phi_{\mathbb C},F_{\mathbb C})
   =\iota_K\bigl(B_{K,w}(\phi)(F)\bigr).
\]
The equality extends by bilinearity to the full complexifications.

\teamanchor{proof.local-unit-comparison}
For each finite place \(\nu\) of \(K\), choose the tuple
\(\boldsymbol\mu(\nu)\) prescribed by the family. The choice is the
same at conjugate places. At an untwisted place the base tuple has
trivial \(\xi_{H,\boldsymbol\mu_0}\). At a twisted place of residue
characteristic \(p\), the character \(\xi_{H,\boldsymbol\mu(\nu)}\)
is the product of the specified twist unramified above \(p\) and the
auxiliary character supported at \(p_{\rm de}\), using \(p'_{\rm de}\)
if \(p=p_{\rm de}\). The rational primes \(p_{\rm de},p'_{\rm de}\)
are distinct. In both cases the character is unramified above \(p\).
Every prime \(Q\) in the tuple's \(Q(\boldsymbol\mu(\nu))\) is a
ramification place of this character. A prime above \(p\) cannot
also lie above the chosen auxiliary prime, since these two
rational primes generate the unit ideal. At the base the trivial
character gives \(Q(\boldsymbol\mu_0)=\varnothing\).
In particular
\begin{equation}
 \label{eq:period-selected-twist-avoidance}
 Q(\boldsymbol\mu(\nu))\cap\{Q:Q\mid p\}=\varnothing .
\end{equation}
\releasecite{Development/AutomorphicArithmetic/ThetaCorrespondence/ThetaPairLifts/ThirdPlaceSeries/NotMemQBase.lean}{44}{The selected finite-place tuple and the two distinct auxiliary primes}.
In the compatible rational top models, evaluate
\eqref{eq:period-tuple} on a pair of coordinates over \(K\) where the
period is nonzero. The quotient of the top evaluation by the evaluation
on the descended row gives \(\zeta_{w,\boldsymbol\mu}\in K^\times\) with
image \(d_{w,\boldsymbol\mu_0}/d_{w,\boldsymbol\mu}\) at \(\iota_K\).
If all evaluations on these coordinates vanish, both sides vanish
on the complex span and we may choose scalar \(1\). In all cases the
base scalar is \(1\) and
\[
 b_{w,\boldsymbol\mu}(\phi,F_m)
   =\iota_K(\zeta_{w,\boldsymbol\mu})\,
        B_w(\phi,\iota_{\boldsymbol\mu}F_m).
\]

Use the support conclusion of the simultaneous theta--Bessel
comparison at this particular tuple. Its relevant power is a
nonzero rational number supported on its twisting primes. Applying
\(v_\nu\) to the power identity gives
\(n\,v_\nu(\zeta_{w,\boldsymbol\mu(\nu)})=0\) for some \(n>0\).
The possible scalar \(1\) also has valuation zero. Both conjugate
places have the same residue characteristic, so the factor is a unit
at both. Consequently the valuation of the tested top pairing equals
the valuation of \(B_{K,w}\) on the descended row.

\teamanchor{proof.row-lattice-dual}
Use the sharp level of the word periods. All integral test rows
below lie in the good part \(\mathcal V^0_{\boldsymbol\mu}\) on
which the simultaneous identity is stated. For each tuple let
\(A_{\boldsymbol\mu,K}\) be the rational top class space at that
character level, \(\mathscr L_{\boldsymbol\mu}\) its prescribed
value lattice and \(E_{\boldsymbol\mu}\) the inclusion matrix
of its good coordinates. Thus \(E_{\boldsymbol\mu}\) reads the
same coefficients as \(j_{\boldsymbol\mu}\) in
\eqref{eq:period-row-match}. The lattice of integral top test
rows in these coordinates is
\[
 M_{\boldsymbol\mu}
   = \{x:E_{\boldsymbol\mu}x\in\mathscr L_{\boldsymbol\mu}\}.
\]
In the full principal class coordinates the identical condition
uses the inclusion of the character level followed by
\(E_{\boldsymbol\mu}\) in the displayed read. A vector passes the
condition precisely when its class values, in the prescribed
frames, are integral. Multiplying a rational basis by a common
denominator shows that the cut lattice is full.

Use the basis of the adjoint image and the projection
\(P_{\boldsymbol\mu}\) in \eqref{eq:period-row-match}. The lattice
of projected coordinates is the image of \(M_{\boldsymbol\mu}\)
under \(P_{\boldsymbol\mu}\). This claim uses the integral
vectors before projection. Requiring their projected representatives
to have integral class values would give a different condition
at primes dividing projection denominators. In the prescribed
base row coordinates \(D_{\boldsymbol\mu}\) sends this image to
\[
 L_{\boldsymbol\mu}^{\rm row}
    =D_{\boldsymbol\mu}P_{\boldsymbol\mu}M_{\boldsymbol\mu}
       \subset R_K .
\]
The kernel of the projection is killed by the actual descent
and by the word period. Surjectivity on the rational class
coordinates makes \(L_{\boldsymbol\mu}^{\rm row}\) full.
For each \(\nu\) the specified row gluing uses the image lattice
for \(\boldsymbol\mu=\boldsymbol\mu(\nu)\).
Localization commutes with this matrix image, since each
localized vector has the form \(x/d\) for a scalar \(d\) outside
\(\nu\). In particular the prescribed local generators are the
actual descended integral rows:
\begin{equation}
 \label{eq:period-selected-local-row-lattice}
 R_\nu
  = (L_{\boldsymbol\mu(\nu)}^{\rm row})_\nu
  =\mathcal O_{K,\nu}\,
       D_{\boldsymbol\mu(\nu)}P_{\boldsymbol\mu(\nu)}
                                      M_{\boldsymbol\mu(\nu)} .
\end{equation}
\releasecite{Formalism/LocalToGlobal/Divisibility/DivisionCompatibleToSub2Indef2.lean}{210}{The sharp value lattice and the actual downward map}
\releasecite{Formalism/LocalToGlobal/Divisibility/GoodPartHermitianSpec.lean}{134}{The cut in class coordinates}
\releasecite{Formalism/LocalToGlobal/Divisibility/RowLat.lean}{138}{The projected-coordinate image} and
\releasecite{Formalism/LocalToGlobal/IndefiniteReal/PrincipalLevelLattice.lean}{266}{the prescribed localizations}.

These localizations glue to one lattice \(R\). To prove this, fix
the base image lattice \(L_0\). The tuple list is finite, so
there is an integer \(h>0\) such that all the \(L_{\boldsymbol\mu}^{\rm row}\)
lie between \(hL_0\) and \(h^{-1}L_0\). Outside the primes dividing
\(h\), they have the same localization as \(L_0\).
The finite \(\mathcal O_K\)-module \(h^{-1}L_0/hL_0\)
is the direct sum of its localizations at the primes dividing
\(h\), by the Chinese remainder theorem. In each summand select
the image of the required local lattice. Their sum pulls back to
the full lattice \(R\) with exactly these localizations. Membership in
a full lattice is tested at all the localizations: if \(x\) does
not belong, the proper ideal of \(a\in\mathcal O_K\) with \(ax\)
in the lattice lies in a maximal ideal where localization still
excludes \(x\). This also proves uniqueness.

Let \(C=R^\vee=\operatorname{Hom}_{\mathcal O_K}(R,\mathcal O_K)\),
viewed in \(R_K^\vee\). Its localization is the dual of \(R_\nu\),
which can be checked in a basis over the discrete valuation ring.
A valuation bound on the top periods extends to evaluations on
\(R_\nu\) by the unit identity above and the ultrametric inequality,
since the actual descended integral rows generate \(R_\nu\).
When the local estimate of Section~\ref{sec:integral} puts the
localized image in \(\mathfrak p_\nu^{m_{w\nu}}C_\nu\), multiplication
by the inverse prime power cancels it. Intersecting over all
finite places, with
\(I_w=\prod_\nu\mathfrak p_\nu^{-m_{w\nu}}\) for the prescribed
rounded gains, gives
\begin{equation}
\label{eq:period-lattice}
 B_{K,w}(N'_w)\subseteq C,\qquad N'_w=I_wN_w\quad (w\in W).
\end{equation}

\teamanchor{proof.period-common-family-cover}
The row lattice in \eqref{eq:period-selected-local-row-lattice}
uses the sharp cut. At a twisted place the arithmetic covering
in \eqref{eq:record-local-rank-sum}, for \(1\le j\le m\), uses
the full auxiliary class space \(M^L_{\boldsymbol\mu(\nu),j}\)
of the base level above.
Its source is the one-step class model of
\eqref{eq:period-fitted-stage-vector}. The localized Hecke algebra in
Section~\ref{sec:zeros} uses the normalization of its actual split
operators. Its maximal ideals index the chosen local records.
Applying their primary projections to a spanning set of the
twisted class lattice gives the projected coordinate generators
of each local module. Their spans over the local fraction field
are the primary summands, whose projections add to the identity.
This proves the covering and the joint detection on the prescribed
auxiliary space. The finite-place row lattice, the avoidance in
\eqref{eq:period-selected-twist-avoidance}, and this covering can
therefore all be used with the same function
\(\nu\mapsto\boldsymbol\mu(\nu)\).
\releasecite{Formalism/LocalToGlobal/NumberFieldFamilies/NPartCNd.lean}{132}{Local record indices and projected lattices}
\releasecite{Formalism/LocalToGlobal/NumberFieldFamilies/NPartCNdTie.lean}{65}{All residual ideals} and
\releasecite{Formalism/LocalToGlobal/NormalizedHeckeFactors/CoverThm.lean}{94}{Covering on the normalized class space}.

\subsection{The matrices in common row coordinates}

\teamanchor{proof.common-row-coordinates}
Choose an orthonormal basis \((F_i)\) of \(E_{\rm row}\).
The vectors \(s(VF_i)\) form one orthonormal basis of the
distinguished space \(H_{\rm row}=s(VE_{\rm row})\). Complete it
to an orthonormal basis of \(\mathcal H\) and let \(\Pi\) be the
orthogonal projection to these distinguished coordinates. In each
orthonormal basis \((\phi_{w,j})\) of \(\mathcal M_w\), the matrix \(J_w\) of
\(\jmath_w U_w\) into an orthonormal basis of \(H_w^{\rm full}\)
has orthonormal columns. Write \(P_w\) for the corresponding
matrix of \(\mathcal P_w\). By \eqref{eq:period-bessel} the
\(i,j\) entry of \(\Pi P_wJ_w\) is
\(a_w^{-1}B_w(\phi_{w,j},F_i)\). If \(Z\) is the coordinate
matrix of \((F_i)\) in the complexified rational row basis and
\(X_w\) is that of \((\phi_{w,j})\) in the complexified rational
column basis, expansion gives
\begin{equation}
\label{eq:period-block}
 \Pi P_wJ_w
 = Z^{\mathsf T}\,\iota_K([B_{K,w}])\,(a_w^{-1}X_w).
\end{equation}
The rational matrix uses the dual of the rational row basis on its
target. The transpose is ordinary because the pairing is linear in
the row vector. The invertible row change \(Z\) is common to the
words, and \(X_w\) is invertible for each word. Choosing the
conjugate representative of the infinite place conjugates the
whole identity.

The common row change is needed when taking sums of images.
For example, take \(E_{\rm row}=\mathcal H=\mathbb C^2\) with
\(s\) coordinatewise conjugation and two one-dimensional columns
given by \(z\mapsto(z,0)\) and \(z\mapsto(0,z)\). With common row
coordinates the resulting functionals span the dual. Using the
identity on the first row input and the coordinate swap on the
second makes both functionals multiples of \(F\mapsto F_1\).
With \eqref{eq:period-block}, all blocks use the same rows.
Section~\ref{sec:trace} estimates the full column pairings in
\(\mathcal H\) and then restricts to these common row coordinates.
Its Gram normalization is the positive volume factor \(c_0(w)\)
of \eqref{eq:period-gram}. The nonzero scalar \(a_w\) comes from the
period comparison.

\noindent\leavevmode\teamanchor{candidate.periods}%
\textbf{AI-proposed candidate (simultaneous matching of theta and Bessel periods).}
Fix a selected translate with the towers, levels and tests above
and with bottom rank difference \(2r_S+1\), where \(r_S\ge1\).
There is one finite central placement \(zt_0\) and one
conjugate-linear row isometry \(s\circ V\) for all words. For each \(w\)
there are a unitary \(U_w\) and \(a_w\in\mathbb C^\times\), independent
of the tuple and the vectors, such that
\[
 b_{w,\boldsymbol\mu}(\phi,F_m)
  = \frac{d_{w,\boldsymbol\mu_0}}{d_{w,\boldsymbol\mu}}\,
       a_w\big\langle s(V\iota_{\boldsymbol\mu}F_m),
                    \mathcal P_w\jmath_w U_w\phi\big\rangle_{\mathcal H}
\]
for every allowed tuple \(\boldsymbol\mu\), ramified column
\(\phi\in\mathcal M_w\), and projected top row
\(F_m\in\mathcal V^0_{\boldsymbol\mu}\). A positive power of the displayed
ratio is rational with prime support contained in the twisting primes
of \(\boldsymbol\mu\). In particular the identity matches the
prescribed coefficient-field tests and their actual downward images
at the base tuple and every finite-place tuple \(\boldsymbol\mu(\nu)\).

The stated identity gives kernel vanishing on the projected base
rows. Together with the arithmetic model and normalized comparisons
in the preceding subsection it supplies \(B_{K,w}\), the common block
identity \eqref{eq:period-block}, and the local units used for
\eqref{eq:period-lattice}.

For a definite unitary group and its codimension-one unitary
subgroup, Michael Harris constructs a rational bilinear period on
the coherent-cohomology models with compatible coefficient
projection. The pairing takes values in the compositum of the
two coefficient fields and is Galois-compatible with Tamagawa
measure. \cite[\S2.5, Lemma~2.5.4]{p2harris-adjoint}
This gives the rationality principle for the top pairing once
the lifts, coefficient models and measures are identified. The
finite placement, identification with the prescribed Bessel
vectors, and support of the twisting ratios require the
comparisons across the two towers.

For comparison, fix \(V_0\subset V_1\) over \(E/F\) with complement
given by hyperbolic planes and the norm line of \(E/F\). Write \(P\)
for the absolutely convergent integral of a cusp form over the
Bessel subgroup quotient against its generic flag character,
trivial on \(U(V_0)\). For a cuspidal
automorphic representation \(\sigma\) of \(U(V_0)\times U(V_1)\)
tempered at every place with discrete Hermitian weak base change
\(\Pi\), Beuzart-Plessis and Chaudouard give, for each nonzero
factorizable \(f=\bigotimes'_v f_v\in\sigma\),
\[
 \frac{|P(f)|^2}{\|f\|_{\mathrm{Pet}}^2}
 = |S_\Pi|^{-1}\mathcal L(\tfrac12,\sigma)
   \prod_v\frac{P_v^\natural(f_v,f_v)}{\|f_v\|_v^2}.
\]
Here \(S_\Pi\) is the component group of \(\Pi\), the Petersson norm
uses Tamagawa measure, and the invariant local norms are used in the
matrix coefficients as well. The local subgroup measures multiply
to Tamagawa measure. The form \(P_v^\natural\) is the local Bessel
functional on matrix coefficients divided by \(\mathcal L_v(\tfrac12)\).
Writing \(d_i=\dim V_i\), the local factor is
\(\mathcal L_v(s)=L_v(s,\Pi_v)\prod_{j=1}^{d_1}
L_v(s+j-\tfrac12,\eta_{E/F,v}^j)/L_v(s+\tfrac12,\Pi_v,\mathrm{As}')\),
with the Rankin--Selberg numerator and Asai signs
\((-1)^{d_0},(-1)^{d_1}\) in \(\mathrm{As}'\).
The global factor is continued from the product of the local factors.
\cite[\S\S1.3.1--1.3.2, \S1.3.5 and Theorem~1.3.6.1]{p2bp-chaudouard-bessel}

The seesaw identity supplies the interchange of compatible theta
kernels. The squared-period formula gives absolute squares on
factorizable vectors of a fixed \(\sigma\). The displayed matching
identity gives a bilinear comparison in common coordinates across
constituents and tuples, with the relative normalizations on the
prescribed tests.
